% Organic chemistry 2 : Study of other organic compounds. — Chemistry Upper Sixth — Cours signé OnBuch+
% Official MINESEC syllabus. Compiler avec : tectonic CHE02-organic-chemistry-2-study-of-other-organic-compounds.tex
\newcommand{\DOCMATIERE}{Chemistry}
\newcommand{\DOCNIVEAU}{Upper Sixth}
\newcommand{\DOCMODULE}{Upper Sixth — Chemistry}
\newcommand{\DOCLECON}{2}
\newcommand{\DOCTITRE}{Organic chemistry 2 : Study of other organic compounds.}
% OnBuch+ — préambule « cours signés » ENRICHI (compilé avec tectonic / XeLaTeX)
% Système visuel « Pop » d'OnBuch+ : fond crème (jamais blanc), bordures encre
% épaisses, ombres « dures » décalées, titres Archivo Black en capitales.
% Version MATHÉMATIQUES : graphes (pgfplots/tikz), boîtes pédagogiques
% (définition, propriété, méthode, attention, le savais-tu, exercice résolu,
% exercice, TP, bilan), page de garde et signature OnBuch+ ancrée en bas.
%
% Macros attendues dans main.tex AVANT \input{preamble} :
%   \DOCMATIERE  \DOCNIVEAU  \DOCMODULE  \DOCLECON (numéro)  \DOCTITRE
\documentclass[11pt]{article}

\usepackage{fontspec}
\usepackage[english]{babel}
\usepackage[a4paper,margin=2cm,top=3.4cm,bottom=2.8cm,headheight=1.4cm,footskip=1.3cm]{geometry}
\usepackage{amsmath,amssymb,amsfonts}
\usepackage{mathtools} % \xmapsto, \coloneqq... souvent utilisés par les modèles
\usepackage{cancel} % simplifications barrées (bilans de forces)
% \abs{z} (module, valeur absolue) et \norm{u} (norme), utilisés par les modèles
\providecommand{\abs}{}\let\abs\relax\DeclarePairedDelimiter{\abs}{\lvert}{\rvert}
\providecommand{\norm}{}\let\norm\relax\DeclarePairedDelimiter{\norm}{\lVert}{\rVert}
\usepackage[table]{xcolor} % [table] : fournit aussi \rowcolors (alternance de lignes)
\usepackage{pagecolor}
\usepackage{tikz}
\usetikzlibrary{arrows.meta,positioning,calc,shapes.geometric,shapes.misc,shapes.arrows,decorations.pathmorphing,decorations.pathreplacing,decorations.markings,patterns,shadows,mindmap,trees,fit,backgrounds,matrix,intersections,angles,quotes}
\usepackage{pgfplots}
\usepackage[version=4]{mhchem} % équations nucléaires \ce{^{235}_{92}U}
\usepackage[european,siunitx]{circuitikz} % schémas électriques
\pgfplotsset{compat=1.18}
\usepgfplotslibrary{fillbetween,groupplots} % aires entre courbes (calcul intégral)
\usepackage{enumitem}
\usepackage{fancyhdr}
% [most] charge toutes les bibliothèques tcolorbox (listings, minted, raster,
% documentation...) et gonfle inutilement la mémoire TeX sur un document long
% (~300 boîtes) : on ne charge que ce qui est réellement utilisé.
\usepackage[skins,breakable]{tcolorbox}
\usepackage{graphicx}
\usepackage{colortbl}
\usepackage{booktabs}
\usepackage{tabularx}
\usepackage{diagbox} % case d'en-tête diagonale des tableaux à double entrée
% Colonnes extensibles centrée / gauche / droite (C, L, R), souvent utilisées par les modèles
\newcolumntype{C}{>{\centering\arraybackslash}X}
\newcolumntype{L}{>{\raggedright\arraybackslash}X}
\newcolumntype{R}{>{\raggedleft\arraybackslash}X}
\usepackage{array}
\usepackage{multicol}
\usepackage{siunitx}
% parse-numbers=false : accepte aussi \SI{2,5 \times 10^{-3}}{...}
\sisetup{locale=UK,output-decimal-marker={.},group-minimum-digits=4,parse-numbers=false}
\DeclareSIUnit\ohm{\ensuremath{\symup{\Omega}}} % U+2126 absent de la police
\DeclareSIUnit\division{div} % réglages d'oscilloscope (ms/div, V/div)
\DeclareSIUnit\tr{tr} % tours (vitesse de rotation, ex. \SI{1500}{\tr\per\minute})
\usepackage{qrcode}
% \degree : commande fréquemment utilisée par les modèles pour un angle ou une
% température en dehors de \SI{}{} (non fournie par les paquets chargés).
\providecommand{\degree}{^{\circ}}
% \qed (fin de démonstration), utilisé par les modèles sans amsthm
\providecommand{\qed}{\hfill$\square$}
% \barre : double barre des valeurs interdites dans un tableau de variations (tkz-tab)
\providecommand{\barre}{\ensuremath{\|}}
% environnement proof (amsthm non chargé) utilisé par les modèles
\ifdefined\proof\else\newenvironment{proof}[1][Proof]{\par\noindent\textit{#1.}\ }{\qed\par}\fi
\usepackage{lastpage}
\usepackage{hyperref}
\hypersetup{hidelinks}

% ============================================================
% Palette « Pop » OnBuch+ — mêmes hex que la classe P (plus_ui.dart)
% ============================================================
\definecolor{poporange}{HTML}{F59321}
\definecolor{popdark}{HTML}{E07A0C}
\definecolor{popcream}{HTML}{FFF6E8}
\definecolor{popink}{HTML}{1C1714}
\definecolor{poppurple}{HTML}{7A5AE0}
\definecolor{popgreen}{HTML}{2E9E6A}
\definecolor{poppink}{HTML}{E0457B}
\definecolor{popblue}{HTML}{2F7DE1}
\definecolor{popgold}{HTML}{C98A12}
\definecolor{popmuted}{HTML}{8A7F76}
\definecolor{popline}{HTML}{EADFD2}
\definecolor{popgray}{HTML}{8A7F76}
\colorlet{popgrey}{popgray}
% Alias tolérants (le modèle nomme parfois les couleurs différemment)
\colorlet{popwhite}{white}
\colorlet{popblack}{popink}
\colorlet{popred}{poppink}
\colorlet{popyellow}{popgold}
% Teintes claires pour fonds de tableaux / zones
\colorlet{popblueL}{popblue!10!white}
\colorlet{poporangeL}{poporange!14!white}
\colorlet{popgreenL}{popgreen!12!white}
\colorlet{poppurpleL}{poppurple!12!white}
\colorlet{poppinkL}{poppink!12!white}
\colorlet{popgoldL}{popgold!14!white}

\pagecolor{popcream}

% ============================================================
% Typographie
% ============================================================
\defaultfontfeatures{Ligatures=TeX}
\setmainfont{PlusJakartaSans-Medium.ttf}[Path=../../../fonts/,BoldFont=PlusJakartaSans-Bold.ttf]
% Maths en sans-serif (Fira Math) avec chiffres et lettres droites de Plus
% Jakarta, pour que les formules se fondent dans le texte.
\usepackage[math-style=ISO,bold-style=ISO]{unicode-math}
\setmathfont{FiraMath-Regular.otf}[Path=../../../fonts/]
\setmathfont{PlusJakartaSans-Medium.ttf}[Path=../../../fonts/,range={up/num,up/{Latin,latin},\mathup}]
% (icomma retiré : décimales avec point en anglais)
% Caractères mathématiques Unicode absents des polices de texte (Plus Jakarta,
% Archivo Black) : rendus par leur équivalent mathématique, en texte comme en
% formule — sinon ils s'affichent en carré vide (ex. « Arithmétique dans ℤ »).
\usepackage{newunicodechar}
\newunicodechar{ℝ}{\ensuremath{\mathbb{R}}}
\newunicodechar{ℤ}{\ensuremath{\mathbb{Z}}}
\newunicodechar{ℂ}{\ensuremath{\mathbb{C}}}
\newunicodechar{ℕ}{\ensuremath{\mathbb{N}}}
\newunicodechar{ℚ}{\ensuremath{\mathbb{Q}}}
\newunicodechar{𝔻}{\ensuremath{\mathbb{D}}}
\newunicodechar{α}{\ensuremath{\alpha}}
\newunicodechar{β}{\ensuremath{\beta}}
\newunicodechar{γ}{\ensuremath{\gamma}}
\newunicodechar{δ}{\ensuremath{\delta}}
\newunicodechar{ε}{\ensuremath{\varepsilon}}
\newunicodechar{θ}{\ensuremath{\theta}}
\newunicodechar{λ}{\ensuremath{\lambda}}
\newunicodechar{μ}{\ensuremath{\mu}}
\newunicodechar{σ}{\ensuremath{\sigma}}
\newunicodechar{τ}{\ensuremath{\tau}}
\newunicodechar{φ}{\ensuremath{\varphi}}
\newunicodechar{ψ}{\ensuremath{\psi}}
\newunicodechar{ω}{\ensuremath{\omega}}
\newunicodechar{Σ}{\ensuremath{\Sigma}}
% majuscules produites par \MakeUppercase dans les titres (α-aminés -> Α)
\newunicodechar{Α}{\ensuremath{\alpha}}
\newunicodechar{Β}{\ensuremath{\beta}}
\newunicodechar{Γ}{\ensuremath{\Gamma}}
\newunicodechar{Δ}{\ensuremath{\Delta}}
\newunicodechar{Ε}{\ensuremath{\varepsilon}}
\newunicodechar{Θ}{\ensuremath{\Theta}}
\newunicodechar{Λ}{\ensuremath{\Lambda}}
\newunicodechar{Μ}{\ensuremath{\mu}}
\newunicodechar{Τ}{\ensuremath{\tau}}
\newunicodechar{Φ}{\ensuremath{\Phi}}
\newunicodechar{Ψ}{\ensuremath{\Psi}}
\newunicodechar{Ω}{\ensuremath{\Omega}}
\newunicodechar{↦}{\ensuremath{\mapsto}}
\newunicodechar{∈}{\ensuremath{\in}}
\newunicodechar{∉}{\ensuremath{\notin}}
\newunicodechar{∧}{\ensuremath{\wedge}}
\newunicodechar{⇔}{\ensuremath{\Leftrightarrow}}
\newunicodechar{⟺}{\ensuremath{\Longleftrightarrow}}
\newunicodechar{⇒}{\ensuremath{\Rightarrow}}
\newunicodechar{⟹}{\ensuremath{\Longrightarrow}}
\newunicodechar{∘}{\ensuremath{\circ}}
\newunicodechar{∪}{\ensuremath{\cup}}
\newunicodechar{∩}{\ensuremath{\cap}}
\newunicodechar{≡}{\ensuremath{\equiv}}
\newunicodechar{⊂}{\ensuremath{\subset}}
\newunicodechar{⊥}{\ensuremath{\perp}}
\newunicodechar{∃}{\ensuremath{\exists}}
\newunicodechar{∀}{\ensuremath{\forall}}
\newunicodechar{∛}{\ensuremath{\sqrt[3]{\,}}}
\newunicodechar{∜}{\ensuremath{\sqrt[4]{\,}}}
\newunicodechar{⃗}{\ensuremath{{}^{\to}}}
\newunicodechar{ⁿ}{\ifmmode^{n}\else\textsuperscript{\ensuremath{n}}\fi}
\newunicodechar{ˣ}{\ifmmode^{x}\else\textsuperscript{\ensuremath{x}}\fi}
\newunicodechar{ᵇ}{\ifmmode^{b}\else\textsuperscript{\ensuremath{b}}\fi}
\newunicodechar{ʳ}{\ifmmode^{r}\else\textsuperscript{\ensuremath{r}}\fi}
\newunicodechar{ᵘ}{\ifmmode^{u}\else\textsuperscript{\ensuremath{u}}\fi}
\newunicodechar{ʸ}{\ifmmode^{y}\else\textsuperscript{\ensuremath{y}}\fi}
\newunicodechar{ᵃ}{\ifmmode^{a}\else\textsuperscript{\ensuremath{a}}\fi}
\newunicodechar{ᵉ}{\ifmmode^{e}\else\textsuperscript{\ensuremath{e}}\fi}
\newunicodechar{ᵏ}{\ifmmode^{k}\else\textsuperscript{\ensuremath{k}}\fi}
\newunicodechar{ᵖ}{\ifmmode^{p}\else\textsuperscript{\ensuremath{p}}\fi}
\newunicodechar{ᴺ}{\ifmmode^{N}\else\textsuperscript{\ensuremath{N}}\fi}
\newunicodechar{ᶜ}{\ifmmode^{c}\else\textsuperscript{\ensuremath{c}}\fi}
\newunicodechar{ʰ}{\ifmmode^{h}\else\textsuperscript{\ensuremath{h}}\fi}
\newunicodechar{ᵠ}{\ifmmode^{q}\else\textsuperscript{\ensuremath{q}}\fi}
\newunicodechar{ₙ}{\ifmmode_{n}\else\textsubscript{\ensuremath{n}}\fi}
\newunicodechar{ₐ}{\ifmmode_{a}\else\textsubscript{\ensuremath{a}}\fi}
\newunicodechar{ᵢ}{\ifmmode_{i}\else\textsubscript{\ensuremath{i}}\fi}
\newunicodechar{ᵣ}{\ifmmode_{r}\else\textsubscript{\ensuremath{r}}\fi}
\newunicodechar{ₖ}{\ifmmode_{k}\else\textsubscript{\ensuremath{k}}\fi}
\newunicodechar{ᵦ}{\ifmmode_{\beta}\else\textsubscript{\ensuremath{\beta}}\fi}
\newunicodechar{ₓ}{\ifmmode_{x}\else\textsubscript{\ensuremath{x}}\fi}
\newfontfamily\popdisplay{ArchivoBlack-Regular.ttf}[Path=../../../fonts/]
\newfontfamily\poplogo{SpaceGrotesk-Bold.ttf}[Path=../../../fonts/]
\newfontfamily\popsemi{PlusJakartaSans-SemiBold.ttf}[Path=../../../fonts/]
\newfontfamily\popxbold{PlusJakartaSans-ExtraBold.ttf}[Path=../../../fonts/]
\newfontfamily\popxboldls{PlusJakartaSans-ExtraBold.ttf}[Path=../../../fonts/,LetterSpace=3.0]

\setlist[enumerate]{leftmargin=1.7em,itemsep=5pt,topsep=4pt}
\setlist[itemize]{leftmargin=1.7em,itemsep=4pt,topsep=4pt,label={\color{poporange}\textbullet}}
\setlength{\parindent}{0pt}
\setlength{\parskip}{5pt}
\renewcommand{\arraystretch}{1.25}

% pgfplots : style Pop par défaut
\pgfplotsset{
  every axis/.append style={axis line style={popink,line width=1pt},tick style={popink},
    grid style={popline,line width=0.5pt},label style={font=\small},tick label style={font=\footnotesize}},
  popcurve/.style={poporange,line width=1.8pt,no markers,smooth},
}
% Même style disponible dans un \draw TikZ simple (hors pgfplots)
\tikzset{popcurve/.style={poporange,line width=1.8pt}}
\tikzset{popgrid/.style={popline,very thin}}
\colorlet{popcurve}{poporange} % le nom du style est parfois utilisé comme couleur

% ============================================================
% Pill / squiggle
% ============================================================
\newcommand{\poppill}[3]{%
  \tikz[baseline=-0.55ex]\node[draw=popink,line width=1.2pt,rounded corners=2.6pt,%
    fill=#2,inner xsep=6.5pt,inner ysep=3pt]{%
    \popxboldls\fontsize{7}{7}\selectfont\color{#3}\MakeUppercase{#1}};%
}
\newcommand{\popsquiggle}[1]{%
  \tikz[baseline=-0.4ex]{
    \draw[#1,line width=2.6pt,line cap=round]
      plot [smooth] coordinates {(0,0.09) (0.34,0.01) (0.68,0.21) (1.02,0.01) (1.36,0.10)};
  }%
}
% Mot-clé surligné (marqueur orange sous le texte)
\newcommand{\cle}[1]{{\popxbold\color{popdark}#1}}

% ============================================================
% En-tête / pied
% ============================================================
\pagestyle{fancy}
\fancyhf{}
\renewcommand{\headrulewidth}{0pt}
\renewcommand{\footrulewidth}{0pt}
\lhead{\raisebox{-0.15ex}{\poplogo\fontsize{13}{13}\selectfont\color{popink}OnBuch\color{poporange}+}}
\rhead{\poppill{\DOCMATIERE\ \textbullet\ \DOCNIVEAU\ \textbullet\ Chapter \DOCLECON}{white}{popink}}
\lfoot{\popxbold\fontsize{7.6}{7.6}\selectfont\color{popmuted}\MakeUppercase{Signed course OnBuch+}\ \textbullet\ {\popsemi\DOCTITRE}}
\rfoot{\tikz[baseline=-0.6ex]\node[rounded corners=8.5pt,fill=popink,minimum height=17pt,minimum width=17pt,inner xsep=5pt,inner sep=0]{\popxbold\fontsize{8}{8}\selectfont\color{popcream}\thepage\,/\,\pageref*{LastPage}};}

% ============================================================
% Titres
% ============================================================
\newcommand{\chaptitre}[2]{%
  \begin{tcolorbox}[enhanced,colback=poporange,colframe=popink,boxrule=2.6pt,arc=11pt,
      drop shadow=popink,left=15pt,right=15pt,top=12pt,bottom=11pt,before skip=3pt,after skip=18pt]
    \poppill{#2}{popink}{popcream}\par\vspace{9pt}
    {\raggedright\hyphenpenalty=10000\exhyphenpenalty=10000\popdisplay\fontsize{22}{24}\selectfont\color{popcream}\MakeUppercase{#1}\par}
  \end{tcolorbox}
}
% Section numérotée : pastille orange avec le numéro + titre Archivo
\newcounter{popsec}
\newcommand{\coursec}[1]{%
  \stepcounter{popsec}\setcounter{popsub}{0}%
  \par\vspace{16pt}\noindent
  \begin{minipage}{\linewidth}
  \tikz[baseline=-0.7ex]\node[draw=popink,line width=1.6pt,fill=poporange,rounded corners=4pt,minimum size=22pt,inner sep=0]{\popdisplay\fontsize{12}{12}\selectfont\color{popcream}\thepopsec};%
  \hspace{8pt}{\popdisplay\fontsize{15}{17}\selectfont\color{popink}\MakeUppercase{#1}}\par
  \vspace{4pt}\noindent{\color{poporange}\rule{36pt}{3.6pt}}
  \end{minipage}\par\nopagebreak\vspace{8pt}
}
\newcounter{popsub}
\newcommand{\courssub}[1]{%
  \stepcounter{popsub}%
  \par\vspace{9pt}\noindent{\popxbold\fontsize{11.5}{13}\selectfont\color{popink}{\color{poporange}\thepopsec.\thepopsub}\hspace{6pt}#1}\par\nopagebreak\vspace{4pt}
}

% ============================================================
% Boîtes pédagogiques — même recette : fond blanc, bordure épaisse colorée,
% ombre dure encre, badge plein. Titre optionnel [#1].
% ============================================================
\tcbset{popbase/.style={breakable,enhanced,colback=white,boxrule=2.3pt,arc=9pt,
  shadow={3pt}{-3pt}{0pt}{popink},left=14pt,right=14pt,top=11pt,bottom=11pt,before skip=11pt,after skip=15pt,
  fontupper=\popsemi\fontsize{10.6}{14.5}\selectfont\color{popink},
  fontlower=\popsemi\fontsize{10.6}{14.5}\selectfont\color{popink}}}
% Coche dessinée (le glyphe ✓ n'existe pas dans les polices du document)
\newcommand{\popcheck}[1]{\tikz[baseline=-0.5ex]\draw[#1,line width=1.6pt,line cap=round,line join=round](-0.13,0.0)--(-0.04,-0.09)--(0.13,0.1);}
\providecommand{\coche}{\popcheck{popgreen}}
\providecommand{\resultat}[1]{\par\smallskip\noindent\fcolorbox{popgreen}{popgreenL}{\parbox{\dimexpr\linewidth-2\fboxsep-2\fboxrule\relax}{\textbf{Result.} #1}}\par\smallskip}
\newcommand{\popboxhead}[3]{\poppill{#1}{#2}{white}\ifx\relax#3\relax\else\hspace{6pt}{\popxbold\fontsize{10.6}{12}\selectfont\color{popink}#3}\fi\par\vspace{7pt}}

\newtcolorbox{aretenir}[1][]{popbase,colframe=popgreen,before upper={\popboxhead{Key points}{popgreen}{#1}}}
\newtcolorbox{exemplebox}[1][]{popbase,colframe=poporange,before upper={\popboxhead{Exemple}{poporange}{#1}}}
\newtcolorbox{definition}[1][]{popbase,colframe=popblue,before upper={\popboxhead{Definition}{popblue}{#1}}}
\newtcolorbox{propriete}[1][]{popbase,colframe=poppurple,before upper={\popboxhead{Property}{poppurple}{#1}}}
\newtcolorbox{methode}[1][]{popbase,colframe=popgold,before upper={\popboxhead{Method}{popgold}{#1}}}
\newtcolorbox{attention}[1][]{popbase,colframe=poppink,before upper={\popboxhead{Warning}{poppink}{#1}}}
\newtcolorbox{experience}[1][]{popbase,colframe=popblue,colback=popblueL,before upper={\popboxhead{Experiment}{popblue}{#1}}}
% « Le savais-tu ? » : carte violette pleine (contexte, culture, Cameroun)
\newtcolorbox{savaistu}[1][]{popbase,colback=poppurple,colframe=popink,
  fontupper=\popsemi\fontsize{10.4}{14.2}\selectfont\color{white},
  before upper={\poppill{Did you know?}{white}{poppurple}\ifx\relax#1\relax\else\hspace{6pt}{\popxbold\color{white}#1}\fi\par\vspace{7pt}}}
% Exercice résolu : énoncé en haut, \tcblower puis la solution sur fond crème
\newcounter{popexo}\newcounter{popexr}
\newtcolorbox{exoresolu}[1][]{popbase,colframe=poporange,colbacklower=popcream,
  segmentation style={popink,line width=1.2pt,dashed},
  before upper={\stepcounter{popexr}\popboxhead{Worked example \thepopexr}{poporange}{#1}},
  before lower={\poppill{Solution}{popink}{popcream}\par\vspace{6pt}}}
% Exercice (non résolu en place) — la correction est dans la partie Corrigés
\newtcolorbox{exercice}[1][]{popbase,colframe=popink,
  before upper={\stepcounter{popexo}\popboxhead{Exercise \thepopexo}{popink}{#1}}}
% Corrigé
\newtcolorbox{corrige}[1][]{popbase,colframe=popgreen,colback=popgreenL,
  before upper={\popboxhead{Solution}{popgreen}{#1}}}
% Objectifs / prérequis / bilan
\newtcolorbox{objectifs}{popbase,colframe=popink,colback=poporangeL,
  before upper={\popboxhead{Chapter objectives}{popink}{}}}
\newtcolorbox{prerequis}{popbase,colframe=popink,colback=popblueL,
  before upper={\popboxhead{Prerequisites}{popblue}{}}}
\newtcolorbox{bilan}{popbase,colframe=popink,colback=white,boxrule=2.8pt,
  before upper={\popboxhead{Fiche bilan}{poporange}{L'essentiel à connaître pour l'examen}}}
% Figure encadrée avec légende
\newtcolorbox{popfigure}[1][]{enhanced,breakable=false,colback=white,colframe=popline,boxrule=1.4pt,arc=8pt,
  left=10pt,right=10pt,top=10pt,bottom=8pt,before skip=10pt,after skip=12pt,halign=center,
  lower separated=false,halign lower=center,
  fontlower=\popsemi\fontsize{9}{11.5}\selectfont\color{popmuted},
  #1}
\newcommand{\legende}[1]{\par\vspace{6pt}{\popsemi\fontsize{9}{11.5}\selectfont\color{popmuted}#1}}

% ============================================================
% Signature OnBuch+
% ============================================================
\newcommand{\poplogoinline}[1]{{\poplogo\fontsize{#1}{#1}\selectfont\color{popink}OnBuch\color{poporange}+}}
\newcommand{\signatureonbuch}{%
  \begin{tcolorbox}[enhanced,colback=popink,colframe=popink,boxrule=2.2pt,arc=10pt,
      drop shadow=poporange,left=14pt,right=14pt,top=10pt,bottom=10pt,before skip=0pt,after skip=0pt]
    \tikz[baseline=-0.7ex]\node[circle,fill=popgreen,draw=popcream,line width=1.3pt,minimum size=22pt,inner sep=0]{\popcheck{white}};%
    \hspace{9pt}\begin{minipage}[c]{0.62\linewidth}
      {\popxbold\fontsize{12}{13}\selectfont\color{popcream}Signed\ \textbullet\ {\poplogo OnBuch\color{poporange}+}}\par\vspace{2pt}
      {\raggedright\popsemi\fontsize{8}{10.5}\selectfont\color{popline}\MakeUppercase{OnBuch+ teaching team}\\\MakeUppercase{Follows the official MINESEC syllabus}\par}
    \end{minipage}\hfill
    {\poplogo\fontsize{22}{22}\selectfont\color{popcream}OnBuch\color{poporange}+}
  \end{tcolorbox}%
}

% ============================================================
% Page de garde
% #1 titre · #2 matière · #3 niveau · #4 module · #5 n° leçon · #6 durée ·
% #7 description / objectifs (texte ou itemize)
% ============================================================
\newcommand{\pagegarde}[7]{%
  \thispagestyle{empty}
  \newgeometry{a4paper,margin=1.7cm,top=1.6cm,bottom=1.4cm}%
  \begin{tikzpicture}[remember picture,overlay]
    \fill[poporange!18] ([xshift=-3.2cm,yshift=-3.6cm]current page.north east) circle (4.6cm);
    \fill[poppurple!14] ([xshift=2.4cm,yshift=7.6cm]current page.south west) circle (3.8cm);
  \end{tikzpicture}%
  {\poplogo\fontsize{30}{30}\selectfont\color{popink}OnBuch\color{poporange}+}\hfill
  \raisebox{0.6ex}{\poppill{Signed course \textbullet\ Premium}{popink}{popcream}}\par
  \vspace{1pt}\popsquiggle{poppurple}\par
  \vspace{8pt}
  \begin{tcolorbox}[enhanced,colback=poporange,colframe=popink,boxrule=2.8pt,arc=13pt,
      drop shadow=popink,left=18pt,right=18pt,top=12pt,bottom=13pt,after skip=10pt]
    \poppill{#2 \textbullet\ #3}{popink}{popcream}\hspace{5pt}\poppill{#4}{white}{popink}\par\vspace{8pt}
    {\popdisplay\fontsize{14}{14}\selectfont\color{popink}CHAPTER #5}\par\vspace{4pt}
    {\raggedright\hyphenpenalty=10000\exhyphenpenalty=10000\popdisplay\fontsize{25}{27}\selectfont\color{popcream}\MakeUppercase{#1}\par}
    \vspace{8pt}
    {\popxbold\fontsize{9.5}{9.5}\selectfont\color{popink}\MakeUppercase{Suggested time: #6}}
  \end{tcolorbox}
  \begin{tcolorbox}[enhanced,colback=white,colframe=popink,boxrule=2.2pt,arc=10pt,
      drop shadow=popink,left=16pt,right=16pt,top=10pt,bottom=10pt,after skip=8pt]
    \poppill{In this chapter}{poppurple}{white}\par\vspace{5pt}
    \popsemi\fontsize{9.3}{12.6}\selectfont\color{popink}#7
  \end{tcolorbox}
  \noindent\begin{minipage}[t]{0.487\linewidth}
    \begin{tcolorbox}[enhanced,colback=popgreen,colframe=popink,boxrule=2.2pt,arc=10pt,
        drop shadow=popink,left=11pt,right=11pt,top=10pt,bottom=10pt,height=3.1cm,valign=center]
      \tcbox[colback=white,colframe=popink,boxrule=1.3pt,arc=5pt,boxsep=0pt,nobeforeafter,
        left=3pt,right=3pt,top=3pt,bottom=3pt,valign=center]{\qrcode[height=1.9cm]{https://play.google.com/store/apps/details?id=com.luvvix.onbuch&hl=en}}%
      \hspace{8pt}\begin{minipage}[c]{0.5\linewidth}
        {\popdisplay\fontsize{11}{12.5}\selectfont\color{white}\MakeUppercase{Download OnBuch}}\par\vspace{4pt}
        {\raggedright\popsemi\fontsize{8.4}{11}\selectfont\color{white}Free \textbullet\ Google Play\\AI tutor, past papers, results\par}
      \end{minipage}
    \end{tcolorbox}
  \end{minipage}\hfill
  \begin{minipage}[t]{0.487\linewidth}
    \begin{tcolorbox}[enhanced,colback=poppurple,colframe=popink,boxrule=2.2pt,arc=10pt,
        drop shadow=popink,left=11pt,right=11pt,top=10pt,bottom=10pt,height=3.1cm,valign=center]
      \tikz[baseline=-0.7ex]\node[circle,fill=white,draw=popink,line width=1.3pt,minimum size=1.6cm,inner sep=0]{\popdisplay\fontsize{24}{24}\selectfont\color{poppurple}+};%
      \hspace{8pt}\begin{minipage}[c]{0.6\linewidth}
        {\popdisplay\fontsize{11}{12.5}\selectfont\color{white}\MakeUppercase{Go OnBuch+}}\par\vspace{4pt}
        {\raggedright\popsemi\fontsize{8.4}{11}\selectfont\color{white}All signed courses, unlimited AI tutor, PDF sheets — OnBuch+ tab in the app\par}
      \end{minipage}
    \end{tcolorbox}
  \end{minipage}
  \vspace{8pt}
  \popcontenu
  \vfill
  \signatureonbuch
  \restoregeometry
  \newpage
}

% Rangée « Ce cours contient » (page de garde)
\newcommand{\poptile}[3]{% #1 couleur #2 gros texte #3 légende
  \begin{tcolorbox}[enhanced,colback=#1,colframe=popink,boxrule=2pt,arc=9pt,shadow={2.5pt}{-2.5pt}{0pt}{popink},
      left=6pt,right=6pt,top=9pt,bottom=9pt,halign=center,valign=center,height=1.75cm,width=\linewidth,nobeforeafter]
    {\popdisplay\fontsize{12.5}{13}\selectfont\color{white}\MakeUppercase{#2}}\par\vspace{3pt}
    {\centering\popsemi\fontsize{7.8}{9.5}\selectfont\color{white}#3\par}
  \end{tcolorbox}}
\newcommand{\popcontenu}{%
  \par\noindent{\popxboldls\fontsize{7.5}{7.5}\selectfont\color{popmuted}THIS SIGNED COURSE CONTAINS}\par\vspace{7pt}
  \noindent\begin{minipage}[t]{0.235\linewidth}\poptile{popblue}{Course}{complete, illustrated, step by step}\end{minipage}\hfill
  \begin{minipage}[t]{0.235\linewidth}\poptile{poporange}{Methods}{and worked examples}\end{minipage}\hfill
  \begin{minipage}[t]{0.235\linewidth}\poptile{poppink}{Exercises}{exam style + solutions}\end{minipage}\hfill
  \begin{minipage}[t]{0.235\linewidth}\poptile{popgreen}{Summary sheet}{the essentials to remember}\end{minipage}\par}

% Sommaire
\newtcolorbox{sommaire}{enhanced,colback=white,colframe=popink,boxrule=2.2pt,arc=10pt,
  drop shadow=popink,left=18pt,right=18pt,top=13pt,bottom=13pt,before skip=6pt,after skip=20pt,
  before upper={\poppill{Contents}{poppurple}{white}\par\vspace{9pt}\popsemi\fontsize{10.6}{14.5}\selectfont\color{popink}}}

% Signature de fin de cours — poussée tout en bas de la dernière page
\newcommand{\findecours}{%
  \par\vfill\nopagebreak
  \begin{center}\popsquiggle{poporange}\end{center}\vspace{4pt}
  \signatureonbuch
}
\AtBeginDocument{\def\llbracket{\mathopen{[\mkern-3.5mu[}}\def\rrbracket{\mathclose{]\mkern-3.5mu]}}} % intervalles d'entiers (glyphe ⟦ absent de la police)
% Glyphes absents de Fira Math : redéfinis à partir de glyphes disponibles
\AtBeginDocument{%
  \def\setminus{\mathbin{\backslash}}\let\smallsetminus\setminus
  \def\vdots{\mathord{\vbox{\baselineskip3.2pt\lineskiplimit0pt\kern4pt\hbox{.}\hbox{.}\hbox{.}}}}%
  \def\ddots{\mathinner{\mkern1mu\raise6.5pt\hbox{.}\mkern2mu\raise3.6pt\hbox{.}\mkern2mu\raise0.7pt\hbox{.}\mkern1mu}}%
  \def\triangle{\mathord{\scalebox{0.9}{$\symup{\Delta}$}}}%
  \def\nabla{\mathord{\raisebox{\depth}{\scalebox{1}[-1]{$\symup{\Delta}$}}}}%
  \def\checkmark{\popcheck{popdark}}%
  \def\star{\mathbin{\ast}}\def\bigstar{\mathord{\ast}}%
  \def\diamond{\mathbin{\scalebox{0.6}{$\Diamond$}}}%
  \def\bigcup{\mathop{\mathchoice{\raisebox{-0.3em}{\scalebox{1.7}{$\cup$}}}{\scalebox{1.25}{$\cup$}}{\cup}{\cup}}\displaylimits}%
  \def\bigcap{\mathop{\mathchoice{\raisebox{-0.3em}{\scalebox{1.7}{$\cap$}}}{\scalebox{1.25}{$\cap$}}{\cap}{\cap}}\displaylimits}%
  \def\lesseqqgtr{\mathrel{\lessgtr}}\def\bigoplus{\mathop{\oplus}\displaylimits}%
}
\providecommand{\ssi}{if and only if} % abréviation fréquente
\AtBeginDocument{\providecommand{\wideparen}{\overparen}} % arc de cercle

\begin{document}

\pagegarde{Organic chemistry 2 : Study of other organic compounds.}{Chemistry}{Upper Sixth}{Upper Sixth — Chemistry}{2}{35 periods}{This chapter extends organic chemistry to the major oxygen-, halogen- and nitrogen-containing families: haloalkanes, alcohols, phenols, aldehydes, ketones, carboxylic acids and their derivatives, nitro compounds, nitriles, amines and amino acids. For each family, students master nomenclature, preparation, characteristic reactions and distinguishing tests, with emphasis on reaction mechanisms required at GCE Advanced Level.
\vspace{4pt}
\begin{multicols}{2}\small
\begin{itemize}
\item By the end of this chapter I can name and classify halogeno compounds, alcohols, phenols, carbonyl compounds, carboxylic acids, their derivatives and amines using IUPAC rules.
\item By the end of this chapter I can describe the physical properties of each family and explain them using polarity and hydrogen bonding.
\item By the end of this chapter I can write balanced equations for the preparation of each family from appropriate starting materials.
\item By the end of this chapter I can explain and draw SN1, SN2 and elimination mechanisms for haloalkanes and predict the major product.
\item By the end of this chapter I can compare the reactivity of the -OH group in alcohols and phenols and account for the differences.
\item By the end of this chapter I can describe nucleophilic addition, oxidation, condensation and iodoform reactions of aldehydes and ketones.
\end{itemize}\end{multicols}}

\begin{sommaire}
\begin{multicols}{2}
\begin{enumerate}
\item Halogeno Compounds: Nomenclature, Properties and Preparation
\item Nucleophilic Substitution Reactions of Haloalkanes
\item Further Reactions of Haloalkanes: Substitution, Elimination and Aryl Halides
\item Comparative Reactivity and Uses of Halogeno Compounds
\item Alcohols: Classification, Nomenclature, Isomerism and Physical Properties
\item Preparation and Reactions of Alcohols
\item Distinguishing Tests for Alcohols and Introduction to Phenols
\item Ring Reactions of Phenol and Comparison with Alcohols
\item Aldehydes and Ketones: Structure, Properties and Preparation
\item Reactions of Aldehydes and Ketones
\item Carboxylic Acids: Structure, Properties, Preparation and Reactions
\item Carboxylic Acid Derivatives: Acyl Chlorides, Anhydrides, Esters and Amides
\item Nitrogen Compounds: Nitro Compounds, Nitriles, Amines, Diazonium Salts and Amino Acids
\item Integration activity
\item Methods and worked examples
\item Exercises
\item Solutions to the exercises
\item Summary sheet
\end{enumerate}
\end{multicols}
\end{sommaire}

\begin{objectifs}
\begin{itemize}
\item By the end of this chapter I can name and classify halogeno compounds, alcohols, phenols, carbonyl compounds, carboxylic acids, their derivatives and amines using IUPAC rules.
\item By the end of this chapter I can describe the physical properties of each family and explain them using polarity and hydrogen bonding.
\item By the end of this chapter I can write balanced equations for the preparation of each family from appropriate starting materials.
\item By the end of this chapter I can explain and draw SN1, SN2 and elimination mechanisms for haloalkanes and predict the major product.
\item By the end of this chapter I can compare the reactivity of the -OH group in alcohols and phenols and account for the differences.
\item By the end of this chapter I can describe nucleophilic addition, oxidation, condensation and iodoform reactions of aldehydes and ketones.
\item By the end of this chapter I can explain the acidity of carboxylic acids and the effect of substituents on acid strength.
\item By the end of this chapter I can interconvert carboxylic acids and their derivatives (acyl chlorides, anhydrides, esters, amides).
\item By the end of this chapter I can describe the basicity of amines, the formation and reactions of benzene diazonium chloride.
\item By the end of this chapter I can select and interpret distinguishing tests (Lucas, iodoform, Fehling's/Tollens', neutral FeCl3) to identify unknown organic compounds in GCE practical-style questions.
\end{itemize}
\end{objectifs}

\begin{prerequis}
\begin{itemize}
\item General organic chemistry from Lower Sixth: homologous series, functional groups, IUPAC nomenclature, structural and stereoisomerism
\item Structure and reactions of alkanes, alkenes and benzene (electrophilic substitution, free-radical substitution)
\item Bond polarity, electronegativity and types of intermolecular forces (hydrogen bonding, van der Waals)
\item Basic reaction kinetics and energetics: activation energy, rate equations, reaction profiles
\item Acids and bases: Brønsted-Lowry theory, pKa and equilibrium concepts
\end{itemize}
\end{prerequis}

\newpage

\coursec{Halogeno Compounds: Nomenclature, Properties and Preparation}

In Douala and Bamenda markets, traders sell palm wine that slowly turns sour, while pharmacists in Yaoundé dispense aspirin, paracetamol and chloroquine --- all organic molecules built from the functional families of this chapter. Why does the local anaesthetic used by the dentist contain a halogen? Why does the ``odour'' of over-ripe plantain reveal an ester? And why does palm wine left overnight develop the sharp smell of ethanoic acid? From the DDT once sprayed against malaria mosquitoes to the amino acids in the egusi soup, halogeno compounds, alcohols, phenols, carbonyls, acids and amines shape daily Cameroonian life. This chapter decodes the structures, preparations and reactions behind these familiar substances.

We begin with the \cle{halogeno compounds} --- molecules in which one or more hydrogen atoms of a hydrocarbon have been replaced by halogen atoms. They are the gateway family of organic synthesis: almost every other functional group in this chapter can be made \emph{from} a halogeno compound, so mastering their naming, properties and preparation is the foundation of everything that follows.

\courssub{Definition and Classification of Halogeno Compounds}

\begin{definition}[Halogeno compound]
A \cle{halogeno compound} (or organic halide) is an organic compound in which one or more hydrogen atoms of a hydrocarbon have been replaced by halogen atoms (\ce{F}, \ce{Cl}, \ce{Br}, \ce{I}). The functional group is the carbon--halogen bond, $\ce{C}-\ce{X}$, and the general formula of a saturated monohalogeno compound is \ensuremath{\mathrm{C_nH_{2n+1}X}}, often written \ce{R-X} where \ce{R} is an alkyl group.
\end{definition}

Halogeno compounds are classified in \textbf{two independent ways}.

\textbf{1. According to the carbon skeleton: alkyl, vinyl and aryl halides.}
\begin{itemize}
\item \cle{Alkyl halides} (haloalkanes): the halogen is attached to an $sp^3$ carbon of an alkyl group, e.g.\ \ce{CH3CH2Cl} (chloroethane).
\item \cle{Vinyl halides}: the halogen is attached directly to a doubly bonded carbon ($sp^2$), e.g.\ \ce{CH2=CHCl} (chloroethene, the monomer of PVC).
\item \cle{Aryl halides}: the halogen is attached directly to a benzene ring ($sp^2$), e.g.\ \ce{C6H5Cl} (chlorobenzene).
\end{itemize}

\begin{definition}[Primary, secondary and tertiary haloalkanes]
A haloalkane is \cle{primary} ($1^\circ$) if the carbon bearing the halogen is attached to \textbf{one} other carbon atom (or none); \cle{secondary} ($2^\circ$) if that carbon is attached to \textbf{two} other carbons; and \cle{tertiary} ($3^\circ$) if it is attached to \textbf{three} other carbons. Count the \emph{carbon} neighbours of the $\ce{C}-\ce{X}$ carbon, not the hydrogens.
\end{definition}

\begin{exemplebox}[Classifying haloalkanes]
\ce{CH3CH2CH2Br} (1-bromopropane) is \textbf{primary}; \ce{CH3CHBrCH3} (2-bromopropane) is \textbf{secondary}; \ce{(CH3)3CBr} (2-bromo-2-methylpropane) is \textbf{tertiary}. This classification matters enormously: it controls whether a haloalkane reacts by an $S_N1$ or $S_N2$ mechanism (Section 2).
\end{exemplebox}

\courssub{Nomenclature of Halogeno Compounds}

In IUPAC nomenclature the halogen is treated as a \emph{substituent}, exactly like an alkyl group, and is named by the prefixes \cle{fluoro-}, \cle{chloro-}, \cle{bromo-} and \cle{iodo-}.

\begin{methode}[IUPAC naming of a halogeno compound]
\begin{enumerate}
\item Identify the \textbf{longest carbon chain} containing the carbon(s) bearing the halogen(s); this gives the parent alkane name.
\item \textbf{Number} the chain from the end that gives the \emph{lowest set of locants} to all substituents (halogens and alkyl groups together).
\item Name the substituents with their locants, listed in \textbf{alphabetical order} (bromo before chloro before fluoro before iodo before methyl\ldots), ignoring multiplying prefixes (di-, tri-) when alphabetising.
\item Write locant--substituent pairs joined by hyphens, substituents fused to the parent name: e.g.\ 2-bromo-3-chloropentane.
\end{enumerate}
\end{methode}

\begin{exemplebox}[Worked example: naming \ce{CH3-CHBr-CHCl-CH2-CH3}]
The longest chain has 5 carbons $\Rightarrow$ parent \textbf{pentane}. Numbering from the left gives locants 2 (Br) and 3 (Cl); from the right it would give 3 (Br) and 2 (Cl). The set of locants is $\{2,3\}$ in both directions. To break the tie, the substituent that comes first alphabetically (\emph{bromo}) receives the lower number. Hence we number from the left. The name is \textbf{2-bromo-3-chloropentane}.
\end{exemplebox}

\begin{attention}[Common naming mistakes]
Do \emph{not} write ``chloroethane'' as ``ethyl chloride'' in an IUPAC answer --- \ce{CH3CH2Cl} is \textbf{chloroethane} (common name: ethyl chloride). Do not alphabetise by the multiplying prefix: in 4-bromo-2,2-dichlorobutane, \emph{bromo} still comes first. Finally, check numbering: the lowest \emph{set} of locants wins, so 2,3 beats 3,4 even though both contain a 3.
\end{attention}

Common (trivial) names are still widely used in industry and in the pharmacy: \ce{CHCl3} is \emph{chloroform} (trichloromethane), \ce{CCl4} is \emph{carbon tetrachloride} (tetrachloromethane), \ce{CHI3} is \emph{iodoform} (tri-iodomethane) and \ce{CH2=CHCl} is \emph{vinyl chloride} (chloroethene).

\courssub{Physical Properties of Halogeno Compounds}

The $\ce{C}-\ce{X}$ bond is \cle{polar}: the halogen is more electronegative than carbon, so it carries a partial negative charge ($\delta^-$) and carbon a partial positive charge ($\delta^+$). This polarity, and the large mass of the halogen, explain the physical properties.

\begin{propriete}[Trends in boiling point and density]
For a fixed alkyl group, boiling point and density \textbf{increase down the halogen group}:
\[ \ce{RF} < \ce{RCl} < \ce{RBr} < \ce{RI} \quad \text{(boiling point and density)} \]
because molecular mass and the strength of van der Waals forces increase. For a fixed halogen, boiling point \textbf{increases with chain length} but \textbf{decreases with chain branching} (branched isomers are more compact, with weaker intermolecular contact). All halogeno compounds are \textbf{insoluble in water} (they cannot form hydrogen bonds with water) but soluble in organic solvents. Chloro- and bromoalkanes are \textbf{denser than water}; many fluoro- and chloromethanes are gases at room temperature.
\end{propriete}

\begin{popfigure}
\begin{center}
\begin{tikzpicture}
\begin{axis}[
  ybar, bar width=0.45cm,
  width=11.5cm, height=6.2cm,
  ymin=0, ymax=120,
  ylabel={Boiling point / $^\circ$C},
  symbolic x coords={CH3Cl, CH3Br, CH3I, 1-Cl-butane, 2-Cl-butane, t-butyl-Cl},
  xtick=data,
  x tick label style={rotate=25, anchor=east, font=\small},
  nodes near coords, nodes near coords style={font=\small},
  ymajorgrids, grid style={popline!40},
]
\addplot[fill=popblue] coordinates {(CH3Cl,-24) (CH3Br,3.6) (CH3I,42.4)};
\addplot[fill=poporange] coordinates {(1-Cl-butane,78) (2-Cl-butane,68) (t-butyl-Cl,51)};
\end{axis}
\end{tikzpicture}
\end{center}
\legende{Boiling points of methyl halides (blue: rises down the halogen group) and of the isomeric butyl chlorides \ce{C4H9Cl} (orange: branching lowers the boiling point).}
\end{popfigure}

\begin{savaistu}[Halogeno compounds in daily life and in Cameroon]
Chloroquine, long used against malaria, is a chlorinated aromatic compound; the anaesthetic halothane (\ce{CF3CHBrCl}) contains three different halogens; and the refrigerant gases once used in fridges across Douala (CFCs such as \ce{CCl2F2}) were banned internationally because their chlorine atoms destroy the stratospheric ozone layer.
\end{savaistu}

\courssub{Preparation from Alkanes: Free-Radical Halogenation}

Alkanes react with chlorine or bromine in the presence of \textbf{ultraviolet light} (or strong heating, $>300^\circ\mathrm{C}$) by a \cle{free-radical substitution} mechanism. For methane:

\[ \ce{CH4(g) + Cl2(g) ->[UV light] CH3Cl(g) + HCl(g)} \]

The mechanism proceeds in three stages (details were met in Organic Chemistry 1): 
\textbf{Initiation:} $\ce{Cl2 ->[UV] 2Cl.}$ 
\textbf{Propagation:} $\ce{Cl. + CH4 -> HCl + CH3.}$ then $\ce{CH3. + Cl2 -> CH3Cl + Cl.}$ 
\textbf{Termination:} radicals combine (e.g.\ $\ce{Cl. + Cl. -> Cl2}$, $\ce{CH3. + Cl. -> CH3Cl}$, $\ce{CH3. + CH3. -> C2H6}$).

\begin{attention}[Free-radical halogenation gives mixtures]
This reaction is \textbf{rarely a clean preparative method}. The product \ce{CH3Cl} still contains \ce{C-H} bonds and is chlorinated further, giving a mixture of \ce{CH3Cl}, \ce{CH2Cl2}, \ce{CHCl3} and \ce{CCl4}. With higher alkanes the situation is worse: chlorination of propane gives a mixture of 1-chloropropane and 2-chloropropane. Fluorination is dangerously explosive; iodination does not proceed (thermodynamically unfavourable). Use this route only when a mixture is acceptable or when the alkane has only one type of hydrogen (e.g.\ methane under controlled conditions).
\end{attention}

\courssub{Preparation from Alkenes: Electrophilic Addition}

Alkenes react readily with halogens and hydrogen halides across the $\ce{C=C}$ double bond --- a much cleaner route than free-radical substitution.

\textbf{Addition of halogens, \ce{X2}} (in an inert solvent such as \ce{CCl4} or \ce{CH2Cl2}, room temperature), giving \emph{vicinal} dihalides:
\[ \ce{CH2=CH2(g) + Br2(l) -> BrCH2CH2Br(l)} \]
This is the standard test for unsaturation: the red-brown bromine is decolourised.

\textbf{Addition of hydrogen halides, \ce{HX}} (gaseous or in aqueous/ethanolic solution):
\[ \ce{CH2=CH2(g) + HBr(g) -> CH3CH2Br(l)} \]

With an \emph{unsymmetrical} alkene, \cle{Markovnikov's rule} applies: the hydrogen adds to the carbon of the double bond that already carries \textbf{more hydrogen atoms} (equivalently, the halogen attaches to the more substituted carbon), because the reaction passes through the more stable carbocation intermediate.

\begin{exemplebox}[Markovnikov addition to propene]
\[ \ce{CH3CH=CH2(g) + HBr(g) -> CH3CHBrCH3(l)} \]
The product is 2-bromopropane (secondary), not 1-bromopropane. The intermediate is the secondary carbocation \ce{CH3CH+CH3}, more stable than the primary alternative \ce{CH3CH2CH2+}.
\end{exemplebox}

\courssub{Preparation from Alcohols: Replacing $-$OH by $-$X}

This is the most important laboratory route to haloalkanes, because the position of the halogen is fixed by the position of the original $-\ce{OH}$ group. Three reagents are used.

\textbf{1. Hydrogen halides, HX.} The alcohol is heated with the concentrated acid (for \ce{HCl}, anhydrous \ce{ZnCl2} is added as catalyst --- the Lucas reagent):
\[ \ce{CH3CH2OH(l) + HBr(g) ->[heat] CH3CH2Br(l) + H2O(l)} \]
Reactivity order of HX: $\ce{HI} > \ce{HBr} > \ce{HCl} > \ce{HF}$; reactivity of alcohols: tertiary $>$ secondary $>$ primary.

\textbf{2. Phosphorus chlorides, \ce{PCl5} or \ce{PCl3}.} With phosphorus pentachloride at room temperature:
\[ \ce{CH3CH2OH(l) + PCl5(s) -> CH3CH2Cl(l) + POCl3(l) + HCl(g)} \]
The steamy fumes of \ce{HCl} (which turn damp blue litmus red and give white fumes with ammonia) are the standard test for an $-\ce{OH}$ group. With phosphorus trichloride (gentler, often used for higher alcohols):
\[ \ce{3CH3CH2OH(l) + PCl3(l) -> 3CH3CH2Cl(l) + H3PO3(l)} \]

\textbf{3. Thionyl chloride, \ce{SOCl2}.} The alcohol is refluxed with thionyl chloride (often with a trace of pyridine as catalyst):
\[ \ce{CH3CH2OH(l) + SOCl2(l) ->[reflux] CH3CH2Cl(l) + SO2(g) + HCl(g)} \]

\begin{attention}[Exam tip: why \ce{SOCl2} is preferred]
\ce{SOCl2} is the reagent of choice for converting an alcohol into a chloroalkane because \textbf{both by-products, \ce{SO2} and \ce{HCl}, are gases} that escape from the reaction mixture. The chloroalkane is therefore left almost pure and is easy to isolate --- unlike the \ce{PCl5} route, which leaves involatile \ce{POCl3} mixed with the product. Remember this justification; it is a favourite GCE question.
\end{attention}

\courssub{Preparation of Aryl Halides}

\textbf{Direct halogenation of benzene.} Benzene does not decolourise bromine under ordinary conditions (its delocalised ring resists addition), but in the presence of a \textbf{halogen carrier} catalyst such as \ce{FeBr3} (or iron filings, which form \ce{FeBr3} in situ), electrophilic \emph{substitution} occurs at room temperature:
\[ \ce{C6H6(l) + Br2(l) ->[FeBr3] C6H5Br(l) + HBr(g)} \]
Chlorination similarly uses \ce{Cl2}/\ce{FeCl3} (or \ce{AlCl3}) to give chlorobenzene. Iodination requires an oxidising agent; direct fluorination is too violent.

\textbf{From diazonium salts.} Aryl halides --- especially the iodo and fluoro compounds that cannot be made by direct halogenation --- are prepared by warming a benzene-diazonium salt with the appropriate reagent (e.g.\ \ce{CuCl} for chlorobenzene, \ce{CuBr} for bromobenzene, \ce{KI} for iodobenzene, \ce{HBF4} then heat for fluorobenzene):
\[ \ce{C6H5N2+Cl- + CuCl -> C6H5Cl + N2 ^} \]
This route will be developed fully in Section 13 of this chapter.

\begin{popfigure}
\begin{center}
\begin{tikzpicture}[
  box/.style={draw=popblue, fill=popblueL, rounded corners, align=center, font=\small, text width=2.6cm, minimum height=1.1cm},
  prod/.style={draw=popdark, fill=popgreenL, rounded corners, align=center, font=\small\bfseries, text width=2.8cm, minimum height=1.1cm},
  arr/.style={-{Stealth[length=3mm]}, thick, popdark},
  lbl/.style={font=\scriptsize, align=center, text width=2.6cm}
]
\node[box] (alkane) at (0,3) {Alkane \ce{R-H}};
\node[box] (alkene) at (0,0) {Alkene \ce{RCH=CH2}};
\node[box] (alcohol) at (0,-3) {Alcohol \ce{R-OH}};
\node[box] (benzene) at (8,1.5) {Benzene \ce{C6H6}};
\node[box] (diazo) at (8,-1.5) {Diazonium salt \ce{C6H5N2+}};
\node[prod] (halo) at (4,0) {Halogeno compound \ce{R-X} / \ce{C6H5X}};

\draw[arr] (alkane) -- node[lbl, above left]{\ce{X2}, UV light (mixture!)} (halo);
\draw[arr] (alkene) -- node[lbl, above]{\ce{X2} or \ce{HX} (Markovnikov)} (halo);
\draw[arr] (alcohol) -- node[lbl, below left]{\ce{HX}, \ce{PCl5}, \ce{SOCl2}} (halo);
\draw[arr] (benzene) -- node[lbl, above right]{\ce{X2} / \ce{FeX3}} (halo);
\draw[arr] (diazo) -- node[lbl, below right]{\ce{CuCl}, \ce{KI}, \ldots} (halo);
\end{tikzpicture}
\end{center}
\legende{Summary map of the five preparation routes to halogeno compounds studied in this section.}
\end{popfigure}

\begin{exoresolu}[Worked exercise]
\textbf{Statement.} A compound \textbf{A}, \ce{C4H10O}, reacts with \ce{PCl5} to give steamy fumes of \ce{HCl} and a compound \textbf{B}, \ce{C4H9Cl}, which is a \emph{secondary} haloalkane. (a) Identify \textbf{A} and \textbf{B}, giving their IUPAC names. (b) Write the equation for the reaction. (c) State one other reagent that converts \textbf{A} into \textbf{B}, and explain why \ce{SOCl2} would be the preferred choice.
\tcblower
\textbf{Solution.} (a) The steamy fumes with \ce{PCl5} confirm an $-\ce{OH}$ group, so \textbf{A} is an alcohol of formula \ce{C4H10O}. Since \textbf{B} is secondary, the $-\ce{OH}$ must be on a carbon attached to two other carbons: \textbf{A} is \textbf{butan-2-ol}, \ce{CH3CH2CH(OH)CH3}, and \textbf{B} is \textbf{2-chlorobutane}, \ce{CH3CH2CHClCH3}.

(b) Equation:
\[ \ce{CH3CH2CH(OH)CH3(l) + PCl5(s) -> CH3CH2CHClCH3(l) + POCl3(l) + HCl(g)} \]

(c) Concentrated \ce{HCl} with anhydrous \ce{ZnCl2} (Lucas reagent) would also work. \ce{SOCl2} is preferred because its by-products \ce{SO2} and \ce{HCl} are gases that leave the mixture, so 2-chlorobutane is obtained nearly pure.
\end{exoresolu}

\begin{aretenir}[Key points of this section]
\begin{itemize}
\item Halogeno compounds \ce{R-X} are classified as alkyl, vinyl or aryl halides, and haloalkanes as primary, secondary or tertiary according to the carbon bearing the halogen.
\item IUPAC names use the prefixes fluoro-, chloro-, bromo-, iodo-, with the lowest set of locants and alphabetical order of substituents.
\item Boiling point and density increase down the halogen group and with chain length, but decrease with branching; all are insoluble in water.
\item Preparations: alkanes + \ce{X2}/UV (radical, gives mixtures); alkenes + \ce{X2} or \ce{HX} (Markovnikov); alcohols + \ce{HX}, \ce{PCl5}, \ce{PCl3} or \ce{SOCl2} (cleanest, by-products gaseous); benzene + \ce{X2}/\ce{FeX3}; diazonium salts for aryl halides.
\end{itemize}
\end{aretenir}

\begin{exercice}[Practice]
(a) Give the IUPAC names of \ce{(CH3)3CBr} and \ce{CH3CHClCH2CH(Br)CH3}, and classify each as primary, secondary or tertiary.
(b) Arrange \ce{CH3CH2Cl}, \ce{CH3CH2Br} and \ce{CH3CH2I} in order of increasing boiling point and justify.
(c) Write balanced equations with state symbols for the preparation of bromoethane from: (i) ethene, (ii) ethanol (two different reagents).
(d) Explain why chlorination of propane is not a good method of preparing pure 2-chloropropane, and suggest a better route starting from propene.
\end{exercice}

\begin{corrige}[Exercise 1]
(a) \ce{(CH3)3CBr}: 2-bromo-2-methylpropane, tertiary (central C bonded to three carbons). \ce{CH3CHClCH2CH(Br)CH3}: 4-bromo-2-chloropentane; the C bearing Cl is secondary and the C bearing Br is also secondary.
(b) $\ce{CH3CH2Cl} < \ce{CH3CH2Br} < \ce{CH3CH2I}$: molecular mass and van der Waals forces increase down the group.
(c) (i) $\ce{CH2=CH2(g) + HBr(g) -> CH3CH2Br(l)}$; (ii) $\ce{CH3CH2OH(l) + HBr(g) ->[heat] CH3CH2Br(l) + H2O(l)}$ or $\ce{3CH3CH2OH(l) + PBr3(l) -> 3CH3CH2Br(l) + H3PO3(l)}$.
(d) Free-radical chlorination of propane substitutes at any of the six primary or two secondary hydrogens, giving a mixture of 1-chloropropane and 2-chloropropane plus polychlorinated products. Better: add \ce{HCl} to propene; Markovnikov addition gives 2-chloropropane as the single major product.
\end{corrige}

\coursec{Nucleophilic Substitution Reactions of Haloalkanes}

In the previous section we learnt to name halogeno compounds and to prepare them from alcohols and alkenes. We also noticed that the carbon--halogen bond is \emph{polar}: the halogen atom pulls the bonding electrons towards itself, leaving the carbon atom slightly electron-poor. That single observation is the key to almost everything a haloalkane does chemically. An electron-poor carbon is a natural target for any species that carries a pair of electrons to offer. In this section we study the most important family of reactions of haloalkanes --- \cle{nucleophilic substitution} --- and we use it to convert haloalkanes into alcohols, the first step of countless syntheses.

\courssub{The Polar C--X Bond and the Idea of Nucleophilic Substitution}

Halogens are more electronegative than carbon, so the \ce{C-X} bond is polarised:

\[ \ensuremath{\mathrm{C^{\delta^{+}}-X^{\delta^{-}}}} \]

The $\delta^{+}$ carbon is said to be \emph{electrophilic} (electron-loving): it can be attacked by any species that possesses a lone pair or a negative charge. Such a species is called a \cle{nucleophile} (``nucleus-loving'').

\begin{definition}[Nucleophile and nucleophilic substitution]
A \cle{nucleophile} is a species (anion or neutral molecule) that has a lone pair of electrons and uses it to form a new covalent bond with an electron-deficient atom. Common nucleophiles: \ce{OH-}, \ce{CN-}, \ce{NH3}, \ce{H2O}, \ce{RO-}.\\
A \cle{nucleophilic substitution} is a reaction in which a nucleophile replaces an atom or group (the \cle{leaving group}) attached to a saturated carbon atom. For a haloalkane:
\[ \ce{R-X + Nu^{-} -> R-Nu + X^{-}} \]
\end{definition}

The halogen leaves as a halide ion \ce{X-}, taking the bonding pair with it: it is the \cle{leaving group}. The whole reaction is a substitution because one group (\ce{Nu}) has taken the place of another (\ce{X}).

\courssub{Hydrolysis of Haloalkanes: Formation of Alcohols}

When a haloalkane is heated under \emph{reflux} with \textbf{aqueous} potassium hydroxide (or sodium hydroxide), the hydroxide ion replaces the halogen and an \textbf{alcohol} is formed. This reaction is called \cle{hydrolysis}.

\begin{exemplebox}[Hydrolysis of bromoethane]
\[ \ce{CH3CH2Br(l) + KOH(aq) ->[reflux] CH3CH2OH(aq) + KBr(aq)} \]
or, showing the actual nucleophile:
\[ \ce{CH3CH2Br + OH^{-} -> CH3CH2OH + Br^{-}} \]
Conditions: aqueous \ce{KOH}, heat under reflux. Product: ethanol.
\end{exemplebox}

\begin{attention}[Aqueous or alcoholic KOH? A classic GCE trap!]
The solvent decides the product:
\begin{itemize}
\item \textbf{Aqueous} \ce{KOH}, reflux $\Rightarrow$ \textbf{substitution} $\Rightarrow$ \textbf{alcohol}.
\item \textbf{Alcoholic} (ethanolic) \ce{KOH}, reflux $\Rightarrow$ \textbf{elimination} $\Rightarrow$ \textbf{alkene}.
\end{itemize}
For example, bromoethane with aqueous \ce{KOH} gives ethanol, but with alcoholic \ce{KOH} gives ethene. Writing the wrong solvent in an equation costs marks every year at the GCE.
\end{attention}

\courssub{Two Mechanisms: \texorpdfstring{$S_N2$}{SN2} and \texorpdfstring{$S_N1$}{SN1}}

Experiments show that hydrolysis does not proceed the same way for all haloalkanes. Depending on the structure of the substrate, the reaction follows one of two mechanisms.

\textbf{The \texorpdfstring{$S_N2$}{SN2} mechanism (Substitution, Nucleophilic, bimolecular).}\quad
For a \emph{primary} haloalkane such as bromoethane, the rate of hydrolysis depends on \emph{both} concentrations:
\[ \text{rate} = k\,[\ce{CH3CH2Br}][\ce{OH-}] \]
Two species are involved in the rate-determining step: the mechanism is \emph{bimolecular}, hence $S_N2$. The nucleophile attacks the $\delta^{+}$ carbon from the side \emph{opposite} to the leaving group (\cle{backside attack}); the new \ce{C-O} bond forms while the old \ce{C-Br} bond breaks, in \textbf{one single step} through a \textbf{transition state} in which carbon is partially bonded to five groups.

\begin{popfigure}
\begin{center}
\begin{tikzpicture}[scale=1.05, every node/.style={font=\small}]
% SN2 mechanism: OH- + CH3CH2Br
\node[font=\small\bfseries, popblue] at (2.2,3.1) {$S_N2$: one step, one transition state};
% reactants
\node at (0.4,1.9) {\ce{HO^{-}}};
\draw[->, thick, poporange] (0.75,1.75) .. controls (1.2,1.2) .. (1.75,1.05);
\node[font=\scriptsize, poporange] at (1.05,1.35) {lone pair};
\node at (2.6,1.0) {\ce{CH3CH2-Br}};
\draw[->, thick, poporange] (3.35,1.05) .. controls (3.6,1.5) .. (3.5,1.85);
\node[font=\scriptsize, poporange] at (3.95,1.5) {\ce{Br} leaves};
% arrow to TS
\draw[->, very thick, popink] (4.4,1.0) -- (5.3,1.0);
% transition state
\node[draw=popline, dashed, rounded corners, inner sep=3pt] at (6.9,1.0) {\ensuremath{\mathrm{HO^{\delta^{-}}--CH_{2}CH_{3}--Br^{\delta^{-}}}}};
\node[font=\scriptsize, popmuted] at (6.9,0.35) {transition state};
% arrow to products
\draw[->, very thick, popink] (8.4,1.0) -- (9.3,1.0);
\node at (10.6,1.0) {\ce{HO-CH2CH3 + Br^{-}}};
\end{tikzpicture}
\end{center}
\legende{Curly-arrow picture of the $S_N2$ attack of \ce{OH-} on bromoethane: the nucleophile approaches from the back as \ce{Br-} departs; both events occur in a single transition state.}
\end{popfigure}

\textbf{The \texorpdfstring{$S_N1$}{SN1} mechanism (Substitution, Nucleophilic, unimolecular).}\quad
For a \emph{tertiary} haloalkane such as 2-bromo-2-methylpropane, the rate depends \emph{only} on the concentration of the haloalkane:
\[ \text{rate} = k\,[\ce{(CH3)3CBr}] \]
Only one species appears in the rate-determining step: the mechanism is \emph{unimolecular}, $S_N1$. It occurs in \textbf{two steps}:
\begin{enumerate}
\item \textbf{Slow ionisation} (rate-determining): the \ce{C-Br} bond breaks on its own, giving a \cle{carbocation} and \ce{Br-}.
\item \textbf{Fast attack}: the nucleophile \ce{OH-} combines rapidly with the carbocation.
\end{enumerate}

\begin{popfigure}
\begin{center}
\begin{tikzpicture}[scale=1.0, every node/.style={font=\small}]
\node[font=\small\bfseries, poppurple] at (5.2,3.4) {$S_N1$: two steps, carbocation intermediate};
\node at (0.9,1.6) {\ce{(CH3)3C-Br}};
\draw[->, thick, poporange] (1.75,1.55) .. controls (2.1,1.1) .. (2.15,0.75);
\node[font=\scriptsize, poporange] at (2.6,1.15) {bond pair to \ce{Br}};
\draw[->, very thick, popink] (2.5,1.9) -- (3.3,1.9);
\node[font=\scriptsize] at (2.9,2.2) {slow};
\node at (4.6,1.6) {\ce{(CH3)3C^{+} + Br^{-}}};
\node[font=\scriptsize, popmuted] at (4.6,1.05) {carbocation};
\draw[->, very thick, popink] (5.7,1.9) -- (6.5,1.9);
\node[font=\scriptsize] at (6.1,2.2) {fast};
\node at (7.4,1.6) {\ce{OH^{-}}};
\draw[->, thick, poporange] (7.7,1.45) .. controls (8.1,1.2) .. (8.5,1.5);
\node at (9.5,1.6) {\ce{(CH3)3C-OH}};
\end{tikzpicture}
\end{center}
\legende{The $S_N1$ ionisation of 2-bromo-2-methylpropane: slow loss of \ce{Br-} gives a planar tertiary carbocation, which is then captured rapidly by \ce{OH-}.}
\end{popfigure}

\begin{methode}[The two mechanisms at a glance]
\begin{enumerate}
\item $S_N2$: \textbf{one step}; \ce{Nu-} attacks from the back as \ce{X-} leaves; transition state with 5 groups around C; rate $= k[\ce{RX}][\ce{Nu-}]$; typical of \textbf{primary} halides.
\item $S_N1$: \textbf{two steps}; slow ionisation to a \textbf{carbocation}, then fast capture by \ce{Nu-}; rate $= k[\ce{RX}]$; typical of \textbf{tertiary} halides.
\end{enumerate}
In curly-arrow conventions: a full-headed arrow shows the movement of an \emph{electron pair}, from the electron-rich end (lone pair or bond) to the electron-poor end.
\end{methode}

\courssub{Energy Profiles}

The difference between the two mechanisms appears clearly on an energy profile. The $S_N2$ profile has \textbf{one hump} (one transition state, no intermediate). The $S_N1$ profile has \textbf{two humps} separated by a valley: the valley is the carbocation \emph{intermediate}, and the first, higher hump corresponds to the slow ionisation step.

\begin{popfigure}
\begin{center}
\begin{tikzpicture}[scale=0.92]
% SN2 profile
\begin{scope}
\draw[->, thick] (0,0) -- (5.4,0) node[below, font=\scriptsize] {reaction coordinate};
\draw[->, thick] (0,0) -- (0,3.4) node[above, font=\scriptsize] {$E$};
\draw[very thick, popblue] (0.3,1.0) .. controls (1.6,1.0) and (1.9,2.9) .. (2.7,2.9) .. controls (3.5,2.9) and (3.8,0.7) .. (5.1,0.7);
\fill[popblue] (0.3,1.0) circle(1.5pt); \fill[popblue] (5.1,0.7) circle(1.5pt);
\node[font=\scriptsize, popblue] at (0.9,0.65) {\ce{RX + OH-}};
\node[font=\scriptsize, popblue] at (4.6,0.35) {\ce{ROH + X-}};
\node[font=\scriptsize, popblue] at (2.7,3.15) {TS};
\node[font=\small\bfseries, popblue] at (2.7,-0.7) {$S_N2$};
\end{scope}
% SN1 profile
\begin{scope}[xshift=7.2cm]
\draw[->, thick] (0,0) -- (5.4,0) node[below, font=\scriptsize] {reaction coordinate};
\draw[->, thick] (0,0) -- (0,3.4) node[above, font=\scriptsize] {$E$};
\draw[very thick, poppurple] (0.3,0.9) .. controls (1.0,0.9) and (1.2,3.0) .. (1.8,3.0) .. controls (2.3,3.0) and (2.4,1.9) .. (2.8,1.9) .. controls (3.2,1.9) and (3.3,2.5) .. (3.8,2.5) .. controls (4.4,2.5) and (4.6,0.6) .. (5.1,0.6);
\fill[poppurple] (0.3,0.9) circle(1.5pt); \fill[poppurple] (5.1,0.6) circle(1.5pt);
\node[font=\scriptsize, poppurple] at (0.95,0.55) {\ce{RX}};
\node[font=\scriptsize, poppurple] at (4.65,0.3) {\ce{ROH + X-}};
\node[font=\scriptsize, poppurple] at (2.8,1.55) {\ce{R+}};
\node[font=\scriptsize, poppurple] at (1.8,3.25) {TS1};
\node[font=\scriptsize, poppurple] at (3.8,2.75) {TS2};
\node[font=\small\bfseries, poppurple] at (2.7,-0.7) {$S_N1$};
\end{scope}
\end{tikzpicture}
\end{center}
\legende{Energy profiles. $S_N2$: a single transition state. $S_N1$: two transition states with a carbocation intermediate in the valley; the first step (ionisation) has the higher activation energy and is rate-determining.}
\end{popfigure}

\courssub{Why the Structure of the Substrate Decides the Mechanism}

The $S_N1$ route is only fast if the carbocation formed in step 1 is reasonably stable. Alkyl groups push electron density towards the positively charged carbon (positive inductive effect) and spread out the charge, so:

\begin{propriete}[Carbocation stability and choice of mechanism]
Carbocation stability increases with the number of alkyl groups on the charged carbon:
\[ \text{tertiary } (3^{\circ}) > \text{secondary } (2^{\circ}) > \text{primary } (1^{\circ}) > \ce{CH3+} \]
Consequences (examinable):
\begin{itemize}
\item \textbf{Tertiary} haloalkanes react by $S_N1$ (stable carbocation; also too crowded for backside attack).
\item \textbf{Primary} haloalkanes react by $S_N2$ (unstable carbocation; little steric hindrance, easy backside attack).
\item \textbf{Secondary} haloalkanes can follow either route, depending on conditions.
\end{itemize}
No proof is required at the GCE, but you must be able to \emph{use} this order to justify a mechanism.
\end{propriete}

Steric hindrance reinforces the trend: a tertiary carbon is shielded by three bulky alkyl groups, so a nucleophile simply cannot reach it from behind; a primary carbon is exposed and easily attacked.

\courssub{Other Factors Influencing \texorpdfstring{$S_N1$}{SN1} versus \texorpdfstring{$S_N2$}{SN2}}

\begin{methode}[Predicting the mechanism]
Ask, in order:
\begin{enumerate}
\item \textbf{Substrate:} $3^{\circ} \Rightarrow S_N1$; $1^{\circ} \Rightarrow S_N2$; $2^{\circ} \Rightarrow$ look at the conditions.
\item \textbf{Nucleophile:} a strong, concentrated nucleophile (\ce{OH-}, \ce{CN-}) favours $S_N2$ (it appears in the rate law); a weak nucleophile (\ce{H2O}) leaves time for ionisation, so $S_N1$.
\item \textbf{Solvent:} polar protic solvents (water, ethanol) stabilise ions and favour $S_N1$; less polar, aprotic solvents favour $S_N2$.
\item \textbf{Leaving group:} the better the leaving group (\ce{I-} $>$ \ce{Br-} $>$ \ce{Cl-}), the faster \emph{both} mechanisms.
\end{enumerate}
\end{methode}

\courssub{Stereochemical Consequences (Enrichment)}

Because the nucleophile attacks from the back in $S_N2$, the three groups left on carbon ``flip over'' like an umbrella in the wind: this is \cle{inversion of configuration} (Walden inversion). In $S_N1$, the carbocation intermediate is \emph{planar}, so the nucleophile can attack either face with equal probability, giving a nearly racemic mixture (\cle{racemisation}). This is enrichment beyond the strict syllabus, but it explains why examiners say $S_N2$ proceeds ``with inversion'' and $S_N1$ ``with racemisation''.

\courssub{Rates of Hydrolysis: Chloro-, Bromo- and Iodoalkanes}

\begin{experience}[Comparing rates of hydrolysis]
Three test tubes each contain a few drops of a different halogenoalkane (\ce{R-Cl}, \ce{R-Br}, \ce{R-I}) with ethanol and aqueous silver nitrate, warmed in a water bath. The halide ion released by hydrolysis precipitates as silver halide (\ce{AgCl} white, \ce{AgBr} cream, \ce{AgI} yellow). The precipitate appears \textbf{fastest with the iodoalkane} and slowest with the chloroalkane.
\end{experience}

The explanation is bond strength, not bond polarity. The \ce{C-I} bond is the longest and weakest (bond energies: \ce{C-I} $\approx$ 240 kJ mol$^{-1}$, \ce{C-Br} $\approx$ 280 kJ mol$^{-1}$, \ce{C-Cl} $\approx$ 340 kJ mol$^{-1}$), so it breaks most easily:

\[ \text{rate of hydrolysis: } \ce{R-I} > \ce{R-Br} > \ce{R-Cl} \]

\begin{attention}[Polarity vs bond strength]
Although \ce{C-Cl} is the \emph{most polar} bond, chloroalkanes hydrolyse \emph{slowest}. The rate of substitution is controlled by the strength of the \ce{C-X} bond that must break, not by its polarity. Do not write ``\ce{C-Cl} is most reactive because it is most polar'' --- this is a frequent error.
\end{attention}

\begin{exoresolu}[Predicting mechanism and product]
\textbf{Statement.} (a) Write the equation for the hydrolysis of 1-bromopropane by aqueous potassium hydroxide, state the conditions, and give the mechanism type. (b) Repeat for 2-bromo-2-methylpropane. (c) Which of the two reactions is faster if each halogen is replaced by iodine? Explain briefly.
\tcblower
\textbf{Solution.} (a) 1-Bromopropane is primary, so substitution occurs by $S_N2$:
\[ \ce{CH3CH2CH2Br + KOH(aq) ->[reflux] CH3CH2CH2OH + KBr} \]
Product: propan-1-ol. Conditions: aqueous \ce{KOH}, heat under reflux.\\
(b) 2-Bromo-2-methylpropane is tertiary, so $S_N1$ via the stable carbocation \ce{(CH3)3C+}:
\[ \ce{(CH3)3CBr + KOH(aq) ->[reflux] (CH3)3COH + KBr} \]
Product: 2-methylpropan-2-ol.\\
(c) In both cases the iodo-compound reacts faster, because the \ce{C-I} bond is weaker than the \ce{C-Br} bond and \ce{I-} is a better leaving group.
\end{exoresolu}

\begin{exercice}[Mechanisms and rates]
\begin{enumerate}
\item For each haloalkane, state the likely mechanism of hydrolysis by aqueous \ce{NaOH} and justify: (i) bromomethane; (ii) 2-bromobutane; (iii) 2-iodo-2-methylbutane.
\item Explain why the hydrolysis of \ce{(CH3)3CCl} is first order while that of \ce{CH3CH2Cl} is second order.
\item Arrange in order of increasing rate of hydrolysis: 1-chlorobutane, 1-iodobutane, 1-bromobutane. Justify.
\end{enumerate}
\end{exercice}

\begin{corrige}[Exercise 1]
\begin{enumerate}
\item (i) Bromomethane: primary (methyl), $S_N2$ --- no stable carbocation possible, no steric hindrance. (ii) 2-Bromobutane: secondary --- borderline; with strong \ce{OH-} in aqueous solution, predominantly $S_N2$ character, but $S_N1$ competes. (iii) 2-Iodo-2-methylbutane: tertiary, $S_N1$ --- stable $3^{\circ}$ carbocation and crowded carbon.
\item $S_N1$: the slow step is the ionisation of \ce{(CH3)3CCl} alone, so rate $= k[\ce{RX}]$ (first order). $S_N2$: the single step involves both \ce{CH3CH2Cl} and \ce{OH-}, so rate $= k[\ce{RX}][\ce{OH-}]$ (second order overall).
\item 1-chlorobutane $<$ 1-bromobutane $<$ 1-iodobutane. All are primary ($S_N2$); the rate increases as the \ce{C-X} bond weakens: \ce{C-Cl} $>$ \ce{C-Br} $>$ \ce{C-I} in strength.
\end{enumerate}
\end{corrige}

\begin{aretenir}[Key points]
\begin{itemize}
\item The polar \ensuremath{\mathrm{C^{\delta^{+}}-X^{\delta^{-}}}} bond makes haloalkanes react with nucleophiles by substitution: \ce{R-X + Nu- -> R-Nu + X-}.
\item Aqueous \ce{KOH}, reflux $\Rightarrow$ alcohol (substitution). Alcoholic \ce{KOH}, reflux $\Rightarrow$ alkene (elimination).
\item $S_N2$: one step, backside attack, transition state, rate $=k[\ce{RX}][\ce{Nu-}]$, favoured by $1^{\circ}$ halides; proceeds with inversion.
\item $S_N1$: two steps via a carbocation, rate $=k[\ce{RX}]$, favoured by $3^{\circ}$ halides; leads to racemisation.
\item Carbocation stability: $3^{\circ} > 2^{\circ} > 1^{\circ}$; polar solvents and weak nucleophiles favour $S_N1$.
\item Rate of hydrolysis: \ce{R-I} $>$ \ce{R-Br} $>$ \ce{R-Cl}, because the \ce{C-X} bond strength decreases down the group.
\end{itemize}
\end{aretenir}

\coursec{Further Reactions of Haloalkanes: Substitution, Elimination and Aryl Halides}

In the previous section we met the two great mechanisms of \cle{nucleophilic substitution}, $S_{\mathrm{N}}1$ and $S_{\mathrm{N}}2$, and we saw how a haloalkane can be converted into an alcohol using aqueous alkali. The C--X bond of a haloalkane is, however, far more versatile than that. By simply changing the nucleophile, the solvent or the temperature, the same molecule of 2-bromopropane can be steered towards an ether, an ester, a nitrile, an amine --- or towards an alkene by a completely different pathway, \cle{elimination}. In this section we complete our survey of haloalkane reactivity, and we end by asking why \emph{aryl} halides such as chlorobenzene refuse to react under the same conditions.

\courssub{More Nucleophilic Substitutions: Ethers, Esters, Nitriles and Amines}

All the reactions in this part follow the same logic as hydrolysis: a nucleophile attacks the electron-poor carbon of the C--X bond and displaces the halide ion. Only the identity of the nucleophile changes.

\textbf{Williamson ether synthesis.} An \cle{alkoxide} ion, $\ce{RO^-}$ (prepared by dissolving sodium in the corresponding alcohol), is a powerful nucleophile. Heated with a haloalkane in ethanol, it gives an \cle{ether}:

\[ \ce{CH3CH2O^-Na+(s) + CH3CH2Br(l) ->[ethanol, reflux] CH3CH2-O-CH2CH3(l) + NaBr(s)} \]

The mechanism is $S_{\mathrm{N}}2$: the alkoxide oxygen attacks the $\alpha$-carbon from the back side as bromide leaves. This is the standard laboratory route to \emph{unsymmetrical} ethers, $\ensuremath{\mathrm{R-O-R^{\prime}}}$.

\begin{attention}[Choosing the right pair for Williamson]
Because the alkoxide is also a strong base, the haloalkane should be \textbf{primary} (or methyl). With a tertiary haloalkane, the alkoxide acts as a base and \emph{elimination} wins, giving an alkene instead of the ether. To make $\ce{(CH3)3C-O-CH3}$, use \emph{tert}-butoxide $+$ $\ce{CH3Br}$, never $\ce{(CH3)3CBr}$ $+$ methoxide.
\end{attention}

\textbf{Formation of esters.} Warming a haloalkane with the \emph{silver salt} of a carboxylic acid gives an ester:

\[ \ce{CH3COO^-Ag+(s) + CH3CH2Br(l) ->[heat] CH3COOCH2CH3(l) + AgBr(s) v} \]

The insoluble silver halide precipitates, pulling the equilibrium forward. The carboxylate anion $\ce{RCOO^-}$ is the nucleophile, and the product is an ester, $\ensuremath{\mathrm{RCOOR^{\prime}}}$.

\textbf{Formation of nitriles: lengthening the carbon chain.} Refluxing a haloalkane with \emph{alcoholic} potassium cyanide gives a \cle{nitrile}:

\[ \ce{CH3CH2Br(l) + KCN(alc) ->[ethanol, reflux] CH3CH2C#N(l) + KBr(s)} \]

Notice carefully: the carbon of the $\ce{C#N}$ group becomes part of the chain, so bromoethane (2 carbons) gives propanenitrile (3 carbons).

\begin{propriete}[Chain extension via nitriles]
The reaction $\ce{R-X ->[KCN] R-C#N}$ is the standard method of \textbf{lengthening a carbon chain by exactly one carbon atom}. The nitrile can then be hydrolysed to a carboxylic acid or reduced to a primary amine, making this step a cornerstone of organic synthesis in the GCE examination.
\end{propriete}

\begin{methode}[Choosing reagents: aqueous vs alcoholic KOH; KCN vs AgCN]
\begin{enumerate}
\item \textbf{Want an alcohol?} Use \emph{aqueous} $\ce{KOH}$, warm gently --- substitution.
\item \textbf{Want an alkene?} Use \emph{alcoholic} $\ce{KOH}$, heat under reflux --- elimination.
\item \textbf{Want a nitrile, $\ce{R-C#N}$?} Use $\ce{KCN}$ in ethanol under reflux: the \emph{carbon} end of $\ce{^-C#N}$ attacks.
\item \textbf{Exam trap:} with $\ce{AgCN}$, the \emph{nitrogen} end attacks and the product is the \emph{isocyanide} (isonitrile), $\ce{R-N^+#C^-}$, not the nitrile. The question "why do KCN and AgCN give different products?" is a classic: KCN is ionic (free $\ce{CN^-}$, attack through C), whereas AgCN is largely covalent, so only the nitrogen lone pair is available.
\end{enumerate}
\end{methode}

\textbf{Reaction with ammonia: formation of amines.} Heating a haloalkane with a large \emph{excess} of ammonia in ethanol, in a sealed tube, gives a primary amine:

\[ \ce{CH3CH2Br(l) + NH3(alc) ->[ethanol, sealed tube, heat] CH3CH2NH2(l) + HBr(g)} \]

Ammonia is a nucleophile thanks to its lone pair. The danger is \cle{multiple alkylation}: the product $\ce{CH3CH2NH2}$ is itself nucleophilic and can attack a second molecule of haloalkane to give a secondary amine $\ce{(CH3CH2)2NH}$, then a tertiary amine, and finally a quaternary ammonium salt $\ce{(CH3CH2)4N^+Br^-}$. Using a large excess of ammonia keeps the primary amine as the major product.

\begin{exoresolu}[From 1-bromopropane to butanoic acid]
Statement: Show, with equations and conditions, how 1-bromopropane can be converted into butanoic acid, $\ce{CH3CH2CH2COOH}$.
\tcblower
\textbf{Solution.} The target acid has \emph{four} carbons while the starting haloalkane has three: we must lengthen the chain by one carbon, so we go through the nitrile.
\begin{enumerate}
\item \textbf{Nitrile formation} (alcoholic KCN, reflux):
\[ \ce{CH3CH2CH2Br(l) + KCN(alc) ->[ethanol, reflux] CH3CH2CH2C#N(l) + KBr(s)} \]
\item \textbf{Acid hydrolysis} of the nitrile (dilute $\ce{HCl}$, reflux):
\[ \ce{CH3CH2CH2C#N(l) + 2H2O(l) + HCl(aq) ->[reflux] CH3CH2CH2COOH(aq) + NH4Cl(aq)} \]
\end{enumerate}
Each step is a standard reaction; the key insight is recognising that $\ce{KCN}$ provides the extra carbon.
\end{exoresolu}

\courssub{Elimination Reactions of Haloalkanes}

\textbf{Observation.} Treat 2-bromopropane with \emph{aqueous} KOH and you obtain propan-2-ol (substitution). Treat the same compound with \emph{alcoholic} KOH and heat strongly, and the gas evolved decolourises bromine water --- an \emph{alkene} has formed. Same substrate, different reagent and solvent, different reaction.

\begin{definition}[Elimination reaction]
An \cle{elimination} is a reaction in which two atoms or groups are removed from adjacent carbon atoms of a molecule, with formation of a carbon--carbon double bond. In haloalkanes, loss of H and X from neighbouring carbons is called \cle{dehydrohalogenation}:
\[ \ce{CH3-CHBr-CH3(l) + KOH(alc) ->[ethanol, heat] CH3-CH=CH2(g) + KBr(s) + H2O(l)} \]
\end{definition}

In alcoholic solution, hydroxide behaves more as a \emph{base} than as a nucleophile: it removes a $\beta$-hydrogen while the halide leaves and the C=C bond forms. Heat favours elimination over substitution.

\begin{popfigure}
\begin{center}
\begin{tikzpicture}[scale=0.9, >=stealth]
% Substrate
\node[draw=popblue, rounded corners, fill=popblueL, text width=3.2cm, align=center] (start) at (0,0) {\textbf{2-bromopropane}\\ $\ce{CH3CHBrCH3}$};
% Substitution product
\node[draw=popgreen, rounded corners, fill=popgreenL, text width=3.4cm, align=center] (sub) at (-4.6,-3) {\textbf{Substitution}\\ propan-2-ol\\ $\ce{CH3CH(OH)CH3}$};
% Elimination product
\node[draw=poporange, rounded corners, fill=popgreenL, text width=3.4cm, align=center] (elim) at (4.6,-3) {\textbf{Elimination}\\ propene\\ $\ce{CH3CH=CH2}$};
% Arrows
\draw[->, very thick, popgreen] (start) -- (sub) node[midway, left, text width=2.6cm, align=center, font=\small] {aqueous $\ce{KOH}$\\ warm};
\draw[->, very thick, poporange] (start) -- (elim) node[midway, right, text width=2.6cm, align=center, font=\small] {alcoholic $\ce{KOH}$\\ heat, reflux};
\end{tikzpicture}
\end{center}
\legende{One substrate, two pathways: the solvent and temperature decide whether 2-bromopropane undergoes substitution (aqueous $\ce{KOH}$) or elimination (alcoholic $\ce{KOH}$, heat).}
\end{popfigure}

\textbf{Mechanism of elimination (E2).} The reaction is concerted: the base ($\ce{OH^-}$) abstracts a $\beta$-hydrogen, the C--H and C--X bonds break simultaneously, and the $\pi$-bond forms.

\begin{popfigure}
\begin{center}
\begin{tikzpicture}[scale=1.1, >=stealth]
% Reactants
\node[align=center] at (-3,0) {$\ce{CH3-CH-CH3}$\\ $\phantom{CH3-}\mid\phantom{CH3}$\\ $\phantom{CH3-}Br$};
\node[align=center] at (-3,-1.5) {$\ce{^-OH}$};
% Arrow for base attacking H
\draw[->, thick, poporange] (-2.2,-1.2) .. controls (-1.5,-0.5) .. (-1.2,0.2) node[above right, font=\small] {$\ce{OH^-}$ abstracts $\beta$-H};
% Arrow for C-H bond breaking to form C=C
\draw[->, thick, poporange] (-1.5,0.2) .. controls (-0.8,0.5) .. (-0.2,0.2);
% Arrow for C-Br bond breaking
\draw[->, thick, poporange] (-0.5,-0.3) .. controls (0.2,-0.8) .. (0.8,-0.3) node[right, font=\small] {$\ce{Br^-}$ leaves};
% Products
\node[align=center] at (3,0) {$\ce{CH3-CH=CH2}$};
\node[align=center] at (3,-1.5) {$\ce{H2O + Br^-}$};
\end{tikzpicture}
\end{center}
\legende{E2 elimination mechanism: concerted removal of $\beta$-H by base and departure of halide.}
\end{popfigure}

When the haloalkane has two different $\beta$-carbons, two alkenes are possible. For 2-bromobutane:

\[ \ensuremath{\mathrm{CH_{3}CH_{2}CHBrCH_{3}(l) \rightarrow [alc. KOH, heat] CH_{3}CH=CHCH_{3}(g) \ (major) + CH_{2}=CHCH_{2}CH_{3}(g) \ (minor)}} \]

\begin{definition}[Zaitsev's rule]
In the dehydrohalogenation of a haloalkane, the \textbf{major product is the more substituted alkene} --- the one in which the double bond carries the greater number of alkyl groups. Equivalently, "the hydrogen is removed preferentially from the $\beta$-carbon that has the \emph{fewer} hydrogens". The more substituted alkene is more stable because alkyl groups stabilise the double bond by electron donation (hyperconjugation).
\end{definition}

\begin{propriete}[Order of reactivity of haloalkanes]
For both substitution and elimination, the reactivity order is
\[ \ce{R-I} > \ce{R-Br} > \ce{R-Cl} \gg \ce{R-F} \]
because the C--X bond strength decreases down the group (C--I weakest, C--Cl strongest), so the halide leaves more easily from the heavier halogen compound. (C--F is so strong that fluoroalkanes are essentially unreactive under normal conditions.)
\end{propriete}

\begin{attention}[Substitution vs elimination: the classic confusion]
\begin{itemize}
\item \emph{Aqueous} $\ce{KOH}$ $\Rightarrow$ substitution $\Rightarrow$ \textbf{alcohol}.
\item \emph{Alcoholic} $\ce{KOH}$ $+$ heat $\Rightarrow$ elimination $\Rightarrow$ \textbf{alkene}.
\item Tertiary haloalkanes eliminate easily even with mild base; primary haloalkanes substitute easily. Do not write "alcoholic KOH gives an alcohol" --- this is the most common error in this chapter.
\end{itemize}
\end{attention}

\courssub{Other Reactions of Haloalkanes: Reduction and the Wurtz Reaction}

\textbf{Reduction to alkanes.} The C--X bond can be replaced by C--H using a reducing agent such as zinc and dilute hydrochloric acid, or lithium tetrahydridoaluminate, $\ce{LiAlH4}$, in dry ether:

\[ \ce{CH3CH2Br(l) + 2[H] ->[Zn(s)/HCl(aq)] CH3CH3(g) + HBr(aq)} \qquad \ce{CH3CH2Br(l) ->[LiAlH4, dry ether] CH3CH3(g)} \]

This is useful when a halogen has served as a "handle" and must be removed at the end of a synthesis.

\textbf{Wurtz reaction (brief).} Heating a haloalkane with sodium in dry ether couples two alkyl groups:

\[ \ce{2CH3CH2Br(l) + 2Na(s) ->[dry ether] CH3CH2CH2CH3(g) + 2NaBr(s)} \]

The carbon skeleton doubles: bromoethane gives butane. The reaction is mainly of historical and exam interest as a way of building symmetrical alkanes.

\courssub{Reactions of Aryl Halides}

\textbf{Observation.} Chlorobenzene, $\ce{C6H5Cl}$, can be boiled for hours with aqueous NaOH without any phenol forming --- conditions under which chloroethane would be hydrolysed in minutes. Industrially, phenol is obtained from chlorobenzene only under drastic conditions (NaOH at about $\SI{350}{\ensuremath{{}^{\circ}\mathrm{C}}}$ and high pressure). Why is the aryl C--Cl bond so unreactive?

\begin{popfigure}
\begin{center}
\begin{tikzpicture}[scale=1.0, >=stealth]
% chlorobenzene ring
\draw[thick] (0:1) -- (60:1) -- (120:1) -- (180:1) -- (240:1) -- (300:1) -- cycle;
% Circle for aromaticity
\draw (30:0.8) arc (30:90:0.8);
\draw (150:0.8) arc (150:210:0.8);
\draw (270:0.8) arc (270:330:0.8);
% Cl atom
\draw[thick] (0:1) -- (0:1.9) node[right] {$\ddot{\mathrm{Cl}}$};
% lone pair delocalisation arrow
\draw[->, very thick, poporange, bend left=35] (0:2.05) to node[above, font=\small, text width=2.8cm, align=center] {lone pair of Cl delocalises into the ring} (35:1.25);
\node[font=\small, align=center, text width=4.6cm] at (0,-1.9) {C--Cl acquires \textbf{partial double-bond character}:\\ shorter, stronger, harder to break};
\end{tikzpicture}
\end{center}
\legende{In chlorobenzene a lone pair on chlorine overlaps with the $\pi$-system of the benzene ring (resonance). The C--Cl bond gains partial double-bond character and resists cleavage.}
\end{popfigure}

\begin{attention}[Aryl halides resist nucleophilic substitution]
Unlike haloalkanes, \cle{aryl halides} do not undergo $S_{\mathrm{N}}1$ or $S_{\mathrm{N}}2$ under ordinary conditions, for two reinforcing reasons:
\begin{enumerate}
\item \textbf{Resonance:} a lone pair on the halogen is delocalised into the benzene ring, giving the C--X bond \emph{partial double-bond character}. It is shorter and stronger than in a haloalkane, so the halide cannot leave easily.
\item \textbf{Geometry:} the $\pi$-electron cloud of the ring repels approaching nucleophiles, and back-side attack ($S_{\mathrm{N}}2$) is blocked by the ring itself.
\end{enumerate}
Substitution on aryl halides therefore requires \textbf{harsh conditions} (high temperature, high pressure, or a copper catalyst). Never propose "warm aqueous NaOH" to convert chlorobenzene into phenol in an exam.
\end{attention}

\begin{savaistu}[Coupling reactions of aryl halides (enrichment)]
Aryl halides do react with metals. With sodium in dry ether, bromobenzene couples to give biphenyl (the \emph{Fittig} reaction, the aromatic analogue of Wurtz); with copper powder, two aryl halides couple in the \emph{Ullmann} reaction. These are beyond the core syllabus but show that aryl halides are not inert --- they simply follow different pathways from haloalkanes.
\end{savaistu}

\begin{exoresolu}[Predicting the major product]
Statement: 2-bromobutane is heated with KOH in ethanol. (a) Give the type of reaction. (b) Write the structures of the two possible organic products and state which is major, justifying your answer.
\tcblower
\textbf{Solution.}
\begin{enumerate}
\item[(a)] Alcoholic KOH $+$ heat $\Rightarrow$ \textbf{elimination} (dehydrohalogenation).
\item[(b)] Removal of H from C-1 gives but-1-ene, $\ce{CH2=CHCH2CH3}$; removal of H from C-3 gives but-2-ene, $\ce{CH3CH=CHCH3}$. By \textbf{Zaitsev's rule}, the major product is \textbf{but-2-ene}, the more substituted (and therefore more stable) alkene; but-1-ene is the minor product.
\end{enumerate}
\end{exoresolu}

\begin{aretenir}[Key points of this section]
\begin{itemize}
\item Haloalkanes give \textbf{ethers} with alkoxides (Williamson, $S_{\mathrm{N}}2$), \textbf{esters} with silver carboxylates, \textbf{nitriles} with alcoholic $\ce{KCN}$ (chain $+1$ carbon), and \textbf{amines} with excess $\ce{NH3}$ in ethanol (beware multiple alkylation).
\item \textbf{Aqueous} KOH $\Rightarrow$ substitution (alcohol); \textbf{alcoholic} KOH $+$ heat $\Rightarrow$ elimination (alkene); \textbf{Zaitsev}: the major alkene is the more substituted one.
\item Reactivity: $\ce{RI} > \ce{RBr} > \ce{RCl}$; haloalkanes are reduced to alkanes by $\ce{Zn/HCl}$ or $\ce{LiAlH4}$; the Wurtz reaction couples two alkyl groups with Na.
\item \textbf{Aryl halides} resist nucleophilic substitution because resonance gives the C--X bond partial double-bond character; they react only under harsh conditions.
\end{itemize}
\end{aretenir}

\begin{exercice}[Practice]
\begin{enumerate}
\item Write the equation and conditions for the conversion of 1-bromobutane into (a) pentanenitrile, (b) butan-1-ol, (c) but-1-ene.
\item Explain why the reaction of bromoethane with KCN gives propanenitrile whereas AgCN gives ethyl isocyanide.
\item 2-bromo-2-methylpropane is warmed with ethanolic KOH. Name and draw the organic product, and state the rule you used.
\item Explain, with the help of resonance, why chlorobenzene does not react with warm aqueous NaOH while chloroethane does.
\end{enumerate}
\end{exercice}

\begin{corrige}[Exercise 4]
\begin{enumerate}
\item (a) $\ce{CH3CH2CH2CH2Br(l) + KCN(alc) ->[ethanol, reflux] CH3CH2CH2CH2C#N(l) + KBr(s)}$; (b) aqueous KOH, warm: $\ce{CH3CH2CH2CH2Br(l) + KOH(aq) -> CH3CH2CH2CH2OH(l) + KBr(aq)}$; (c) alcoholic KOH, heat: $\ce{CH3CH2CH2CH2Br(l) + KOH(alc) -> CH3CH2CH=CH2(g) + KBr(s) + H2O(l)}$.
\item KCN is ionic: free $\ce{CN^-}$ attacks through carbon, giving the nitrile $\ce{R-C#N}$. AgCN is covalent: only the nitrogen lone pair is free to attack, giving the isocyanide $\ce{R-NC}$.
\item 2-methylpropene, $\ce{(CH3)2C=CH2}$, by elimination; it is the only possible alkene (all $\beta$-carbons are equivalent methyl groups), consistent with Zaitsev's rule (most substituted alkene).
\item In chlorobenzene, a lone pair on Cl is delocalised into the ring, giving C--Cl partial double-bond character: the bond is stronger and the ring $\pi$-cloud repels nucleophiles. In chloroethane the C--Cl bond is a simple, weaker single bond open to back-side attack, so hydrolysis by warm aqueous NaOH is easy.
\end{enumerate}
\end{corrige}

\coursec{Comparative Reactivity and Uses of Halogeno Compounds}

In the previous sections we saw that haloalkanes undergo nucleophilic substitution and elimination, while aryl halides are far more reluctant to react. It is now time to step back and answer a question examiners love: \emph{which halogeno compound reacts fastest, and why?} Understanding the trends in reactivity allows us to predict the outcome of experiments such as the famous \ce{AgNO3}/ethanol test, and also explains why certain halogeno compounds were chosen as solvents, refrigerants, anaesthetics and pesticides --- and why some of them had to be banned.

\courssub{Reactivity of Alkyl Halides versus Aryl Halides}

\begin{experience}[A simple observation]
Two test tubes each contain \SI{2}{cm^3} of ethanol and a few drops of aqueous silver nitrate. To the first we add two drops of 1-bromobutane, \ce{CH3CH2CH2CH2Br}; to the second, two drops of bromobenzene, \ce{C6H5Br}. Within a minute or two, the first tube develops a cream precipitate of \ce{AgBr}. The second tube remains clear even after warming for a long time. Same halogen, same reagent --- completely different behaviour. Why?
\end{experience}

The answer lies in the nature of the \ce{C-X} bond in each case.

\begin{itemize}
\item In a \cle{haloalkane} (alkyl halide), the halogen is attached to an $sp^3$ hybridised carbon. The \ce{C-X} bond is a simple $\sigma$ bond, polarised $\ensuremath{\mathrm{C^{\delta^{+}}-X^{\delta^{-}}}}$, and the carbon is open to attack by nucleophiles.
\item In a \cle{halobenzene} (aryl halide), the halogen is attached to an $sp^2$ hybridised carbon of the ring. A lone pair on the halogen overlaps with the delocalised $\pi$ system of the benzene ring, giving the \ce{C-X} bond \textbf{partial double-bond character}. The bond is therefore \textbf{shorter and stronger}, and the electron-rich ring repels approaching nucleophiles.
\end{itemize}

\begin{propriete}[Relative reactivity of alkyl and aryl halides]
Aryl halides are \textbf{much less reactive} than alkyl halides towards nucleophilic substitution. Under ordinary laboratory conditions, halobenzenes do not react with aqueous \ce{NaOH}, \ce{KCN} or \ce{NH3}; drastic conditions (high temperature, high pressure, or a copper catalyst) are required. Alkyl halides react under mild conditions.
\end{propriete}

\begin{attention}[Do not confuse the two families]
Never write that bromobenzene "reacts slowly with \ce{NaOH} to give phenol" at room temperature --- at GCE level, the conversion \ce{C6H5Cl -> C6H5OH} requires \ce{NaOH} at about \SI{350}{\ensuremath{{}^{\circ}\mathrm{C}}} and high pressure (the Dow process). In contrast, \ce{CH3CH2Br} is hydrolysed by warm aqueous \ce{NaOH} in minutes.
\end{attention}

\courssub{Effect of the Halogen: Bond Strength versus Bond Polarity}

Here is a beautiful trap that catches many candidates. Consider the \ce{C-X} bonds:

\begin{center}
\begin{tabularx}{\linewidth}{|l|X|X|X|X|}
\hline
\rowcolor{popblueL} Bond & \ce{C-F} & \ce{C-Cl} & \ce{C-Br} & \ce{C-I} \\
\hline
Bond enthalpy / $\mathrm{kJ\,mol^{-1}}$ & 467 & 346 & 290 & 228 \\
\hline
Electronegativity of X & 4.0 & 3.0 & 2.8 & 2.5 \\
\hline
Bond polarity & \multicolumn{4}{l|}{decreases \ce{C-F} $>$ \ce{C-Cl} $>$ \ce{C-Br} $>$ \ce{C-I}} \\
\hline
\end{tabularx}
\end{center}

If \emph{polarity} controlled reactivity, fluoroalkanes would react fastest (most polar bond, largest $\delta+$ on carbon). In fact, the \textbf{opposite} is observed: iodoalkanes react fastest.

\begin{propriete}[Reactivity order of haloalkanes]
For a given alkyl group, the rate of nucleophilic substitution (e.g.\ hydrolysis) increases in the order
\[ \ce{R-F} \ll \ce{R-Cl} < \ce{R-Br} < \ce{R-I} \]
because the \cle{bond strength} of \ce{C-X} decreases down the group. The rate-determining step of both $S_N1$ and $S_N2$ mechanisms involves \textbf{breaking} the \ce{C-X} bond; a weaker bond breaks more easily, so the reaction is faster. Bond \emph{strength}, not bond \emph{polarity}, controls the rate.
\end{propriete}

\begin{attention}[The classic exam trap]
\textbf{Bond polarity alone does not predict reactivity.} A common wrong answer is: "\ce{R-F} is most reactive because F is the most electronegative halogen, so the bond is most polar." This is false. The \ce{C-F} bond is so \emph{strong} (\SI{467}{kJ.mol^{-1}}) that it is very hard to break, making fluoroalkanes almost inert. Always argue in terms of \textbf{bond enthalpy} when explaining rates of substitution.
\end{attention}

\courssub{Effect of the Substrate: Primary, Secondary and Tertiary Haloalkanes}

For a given halogen, the structure of the carbon skeleton also matters, because it decides \emph{which mechanism} operates.

\begin{itemize}
\item \textbf{Tertiary haloalkanes} (e.g.\ 2-bromo-2-methylpropane) form stable carbocations and are sterically blocked to back-side attack: they react by \cle{$S_N1$} and eliminate easily ($E1$). They hydrolyse \textbf{fastest} with aqueous/ethanolic \ce{AgNO3}.
\item \textbf{Primary haloalkanes} (e.g.\ 1-bromobutane) cannot form stable carbocations but are unhindered: they react by \cle{$S_N2$}. Their hydrolysis with \ce{AgNO3}/ethanol is slower than that of tertiary compounds.
\item \textbf{Secondary haloalkanes} are intermediate and can follow either pathway depending on conditions.
\end{itemize}

\begin{propriete}[Overall trend in rate of hydrolysis]
For the \ce{AgNO3}/ethanol hydrolysis, the rate order is
\[ \text{tertiary} > \text{secondary} > \text{primary} \quad\text{and}\quad \ce{R-I} > \ce{R-Br} > \ce{R-Cl} \]
because the rate-determining step is \ce{C-X} bond breaking, which is favoured by a weak \ce{C-X} bond and by a stable carbocation (or an easy $S_N2$ displacement).
\end{propriete}

\begin{popfigure}
\begin{tikzpicture}[scale=1.0, every node/.style={font=\small}]
% Axes of the summary map
\draw[->, thick, popdark] (0,0) -- (10.4,0) node[below left, popdark] {\textbf{halogen:} \ce{Cl} $\to$ \ce{Br} $\to$ \ce{I}};
\draw[->, thick, popdark] (0,0) -- (0,4.6) node[rotate=90, above left, popdark, align=center, text width=2.6cm] {\textbf{substrate:} primary $\to$ tertiary};
% Grid cells coloured by speed
\fill[popgreenL]  (0.2,0.2) rectangle (3.4,1.4);
\fill[popgreenL]  (3.6,0.2) rectangle (6.8,1.4);
\fill[popblueL]   (7.0,0.2) rectangle (10.2,1.4);
\fill[popgreenL]  (0.2,1.6) rectangle (3.4,2.8);
\fill[popblueL]   (3.6,1.6) rectangle (6.8,2.8);
\fill[popblueL]   (7.0,1.6) rectangle (10.2,2.8);
\fill[popblueL]   (0.2,3.0) rectangle (3.4,4.2);
\fill[popblueL]   (3.6,3.0) rectangle (6.8,4.2);
\fill[poporange!40] (7.0,3.0) rectangle (10.2,4.2);
% Labels
\node[align=center] at (1.8,0.8) {primary \ce{R-Cl}\\ slowest};
\node[align=center] at (5.2,0.8) {primary \ce{R-Br}\\ slow};
\node[align=center] at (8.6,0.8) {primary \ce{R-I}\\ moderate};
\node[align=center] at (1.8,2.2) {secondary \ce{R-Cl}\\ slow};
\node[align=center] at (5.2,2.2) {secondary \ce{R-Br}\\ moderate};
\node[align=center] at (8.6,2.2) {secondary \ce{R-I}\\ fast};
\node[align=center] at (1.8,3.6) {tertiary \ce{R-Cl}\\ moderate};
\node[align=center] at (5.2,3.6) {tertiary \ce{R-Br}\\ fast};
\node[align=center] at (8.6,3.6) {tertiary \ce{R-I}\\ \textbf{fastest}};
% Grid lines
\foreach \x in {3.4,6.8} \draw[popline] (\x,0.2) -- (\x,4.2);
\foreach \y in {1.4,2.8} \draw[popline] (0.2,\y) -- (10.2,\y);
\draw[popline] (0.2,0.2) rectangle (10.2,4.2);
\end{tikzpicture}
\legende{Reactivity map for hydrolysis of haloalkanes: rate increases down the halogen group (\ce{Cl} $\to$ \ce{I}, weaker \ce{C-X} bond) and from primary to tertiary substrate (easier carbocation formation, $S_N1$). Aryl halides lie far below this map --- essentially unreactive under these conditions.}
\end{popfigure}

\courssub{The \ce{AgNO3}/Ethanol Test: Comparing Rates and Identifying the Halogen}

\begin{methode}[Comparing haloalkane reactivity with \ce{AgNO3}/ethanol]
\begin{enumerate}
\item Place \SI{2}{cm^3} of ethanol in each of three test tubes standing in a water bath at about \SI{50}{\ensuremath{{}^{\circ}\mathrm{C}}}. (Ethanol is the solvent because haloalkanes are insoluble in water; the water bath keeps temperature constant for a fair test.)
\item Add a few drops of aqueous silver nitrate solution to each tube.
\item Simultaneously add two drops of 1-chlorobutane to tube 1, 1-bromobutane to tube 2 and 1-iodobutane to tube 3. Start a stopwatch.
\item Record the time for a precipitate to appear in each tube.
\item Identify the halide ion from the precipitate colour: \ce{AgCl} \textbf{white}, \ce{AgBr} \textbf{cream}, \ce{AgI} \textbf{pale yellow}.
\end{enumerate}
The hydrolysis \ce{R-X + H2O -> R-OH + H+ + X-} releases halide ions, which precipitate with silver ions: \ce{Ag+(aq) + X-(aq) -> AgX(s)}. The faster the \ce{C-X} bond breaks, the sooner the precipitate appears.
\end{methode}

\begin{popfigure}
\begin{tikzpicture}
\begin{axis}[
    ybar, bar width=22pt,
    width=11cm, height=6.5cm,
    ymin=0, ymax=14,
    ylabel={Time for precipitate / min (typical)},
    symbolic x coords={1-chlorobutane, 1-bromobutane, 1-iodobutane},
    xtick=data,
    x tick label style={font=\small},
    ytick={0,2,4,6,8,10,12,14},
    nodes near coords,
    every node near coord/.append style={font=\small},
    axis lines=left,
    ymajorgrids=true, grid style={popline!50},
]
\addplot[fill=popgreen, draw=popdark] coordinates {(1-chlorobutane,12)};
\addplot[fill=popgold, draw=popdark] coordinates {(1-bromobutane,4)};
\addplot[fill=poporange, draw=popdark] coordinates {(1-iodobutane,1)};
\end{axis}
\end{tikzpicture}
\legende{Typical precipitate appearance times in the \ce{AgNO3}/ethanol test at about \SI{50}{\ensuremath{{}^{\circ}\mathrm{C}}}: \ce{AgCl} (white) appears slowly, \ce{AgBr} (cream) faster, \ce{AgI} (pale yellow) almost immediately --- direct evidence that \ce{C-I} breaks most easily.}
\end{popfigure}

\begin{exoresolu}[Interpreting hydrolysis data]
Three isomeric bromoalkanes A, B and C of formula \ce{C4H9Br} were each warmed with ethanolic silver nitrate. Precipitates of \ce{AgBr} appeared after 30 s (A), 3 min (B) and 12 min (C). Identify A, B and C and justify your answer.
\tcblower
\textbf{Solution.} The three isomers are 1-bromobutane (primary), 2-bromobutane (secondary) and 2-bromo-2-methylpropane (tertiary). The rate of hydrolysis follows tertiary $>$ secondary $>$ primary, because tertiary forms the most stable carbocation ($S_N1$), while the primary compound must undergo the slower $S_N2$ displacement.
\begin{itemize}
\item \textbf{A} (fastest, 30 s) = 2-bromo-2-methylpropane, \ce{(CH3)3CBr} (tertiary).
\item \textbf{B} (3 min) = 2-bromobutane, \ce{CH3CH2CHBrCH3} (secondary).
\item \textbf{C} (slowest, 12 min) = 1-bromobutane, \ce{CH3CH2CH2CH2Br} (primary).
\end{itemize}
In each case the precipitate is cream \ce{AgBr}: \ce{Ag+(aq) + Br-(aq) -> AgBr(s)}.
\end{exoresolu}

\courssub{Uses of Halogeno Compounds}

Halogeno compounds are rarely found in nature, but chemists have put their special properties (low flammability, controlled reactivity, volatility, stability) to many uses.

\begin{aretenir}[One example per use --- learn this list]
\begin{itemize}
\item \textbf{Solvents:} trichloromethane (chloroform, \ce{CHCl3}) and tetrachloromethane (\ce{CCl4}) dissolve greases and many organic compounds; \ce{CCl4} was used in dry-cleaning but is now restricted because of its toxicity.
\item \textbf{Refrigerants:} chlorofluorocarbons (CFCs) such as \ce{CCl2F2} (Freon-12) were ideal --- non-toxic, non-flammable, easily liquefied --- until their effect on ozone was discovered. They are being replaced by hydrofluorocarbons (HFCs) and hydrochlorofluorocarbons (HCFCs).
\item \textbf{Anaesthetics:} halothane, \ce{CF3CHBrCl}, is a safe, non-flammable general anaesthetic that replaced explosive ether in operating theatres.
\item \textbf{Pesticides:} DDT, \ce{(ClC6H4)2CHCCl3}, was a powerful insecticide used against malaria-carrying mosquitoes.
\item \textbf{Plastics:} chloroethene (vinyl chloride, \ce{CH2=CHCl}) is polymerised to give poly(chloroethene), \cle{PVC}, used for pipes, electrical insulation and raincoats.
\end{itemize}
\end{aretenir}

\begin{attention}[Exam tip: quote real uses]
When a question asks for "uses of halogeno compounds", give \textbf{one named compound per use} with its function, e.g.\ "\ce{CHCl3} is used as a solvent", "halothane \ce{CF3CHBrCl} is an anaesthetic", "chloroethene is the monomer for PVC". Vague answers such as "used in industry" earn no marks.
\end{attention}

\courssub{Environmental Impact of Halogeno Compounds}

The very stability that made CFCs and DDT so useful is also their downfall: they persist in the environment.

\textbf{CFCs and the ozone layer.} CFCs are unreactive in the lower atmosphere, so they survive for decades and drift up into the stratosphere. There, intense ultraviolet radiation breaks the \ce{C-Cl} bond, releasing chlorine free radicals:
\[ \ensuremath{\mathrm{CCl_{2}F_{2} \rightarrow [UV] CClF_{2}^{\bullet} + Cl^{\bullet}}} \]
Each chlorine radical destroys an ozone molecule and is regenerated, so a single radical can destroy thousands of \ce{O3} molecules in a chain:
\[ \ensuremath{\mathrm{Cl^{\bullet} + O_{3} \rightarrow  ClO^{\bullet} + O_{2}}} \qquad \ensuremath{\mathrm{ClO^{\bullet} + O^{\bullet} \rightarrow  Cl^{\bullet} + O_{2}}} \]
The thinning of the ozone layer increases harmful UV radiation reaching the Earth's surface (skin cancers, eye damage, harm to crops and plankton).

\textbf{DDT and persistence.} DDT is very stable and fat-soluble. It accumulates in the fatty tissues of organisms and becomes concentrated along food chains (\cle{bioaccumulation}), where it caused, for example, thinning of birds' eggshells. Its agricultural use is now banned or severely restricted in most countries, though limited use against malaria mosquitoes continues where no alternative exists.

\begin{savaistu}[The Montreal Protocol]
In 1987, the nations of the world signed the \textbf{Montreal Protocol}, agreeing to phase out CFCs and other ozone-depleting substances. It is often described as the most successful international environmental agreement ever: the ozone layer is now slowly recovering. This is enrichment beyond the strict syllabus, but a short, well-stated reference to it can earn credit in "discuss" questions on the uses and dangers of halogeno compounds.
\end{savaistu}

\begin{exercice}[Consolidation]
\begin{enumerate}
\item Explain, in terms of bond enthalpy, why 1-iodopropane is hydrolysed faster than 1-chloropropane by warm aqueous \ce{NaOH}.
\item A student claims that chlorobenzene must react faster than chloroethane with \ce{AgNO3}/ethanol because the ring is electron-rich. Correct the student's reasoning and state what is actually observed.
\item Give one halogeno compound used as (a) a refrigerant, (b) an anaesthetic, (c) a monomer for a plastic, and state one environmental problem associated with halogeno compounds.
\end{enumerate}
\end{exercice}

\begin{corrige}[Exercise 1]
\begin{enumerate}
\item The \ce{C-I} bond (\SI{228}{kJ.mol^{-1}}) is much weaker than the \ce{C-Cl} bond (\SI{346}{kJ.mol^{-1}}). Since the rate-determining step of hydrolysis involves breaking the \ce{C-X} bond, the weaker \ce{C-I} bond breaks more easily, so 1-iodopropane reacts faster.
\item The reasoning is wrong: the electron-rich ring actually \emph{repels} nucleophiles, and the \ce{C-Cl} bond in chlorobenzene has partial double-bond character (lone-pair delocalisation into the ring), making it stronger and harder to break. In practice, chloroethane gives a white precipitate of \ce{AgCl} on warming, while chlorobenzene gives no precipitate even after prolonged heating.
\item (a) \ce{CCl2F2} (a CFC, now replaced by HFCs); (b) halothane \ce{CF3CHBrCl}; (c) chloroethene \ce{CH2=CHCl} for PVC. Environmental problem: CFCs release chlorine radicals in the stratosphere which catalytically destroy ozone (or: DDT persists and bioaccumulates in food chains).
\end{enumerate}
\end{corrige}

\coursec{Alcohols: Classification, Nomenclature, Isomerism and Physical Properties}

In the previous section we saw how versatile the halogeno compounds are: a haloalkane such as bromoethane can be converted into many other families simply by replacing the halogen atom. One of the most important of these conversions is the nucleophilic substitution with aqueous hydroxide ions, \ce{C2H5Br + OH- -> C2H5OH + Br-}, which replaces the halogen by an \ce{-OH} group. The product, ethanol, belongs to a new and enormous family: the \cle{alcohols}. Alcohols are everywhere in Cameroonian daily life — ethanol in beverages and in methylated spirit used in laboratories and hospitals, glycerol in soaps and cosmetics, menthol in balms. In this section we learn to recognise, classify and name them, to count and draw their isomers (including a new type, optical isomerism), and to explain their characteristic physical properties.

\courssub{The Alcohol Family: General Formula and Functional Group}

\begin{definition}[Alcohol]
An \cle{alcohol} is an organic compound in which a hydroxyl group, \ce{-OH}, is bonded to a saturated ($sp^3$) carbon atom of an alkyl group. The saturated aliphatic alcohols form a homologous series of general molecular formula \ensuremath{\mathrm{C_{n}H_{2n+1}OH}} (equivalently \ensuremath{\mathrm{C_{n}H_{2n+2}O}}), with $n \geq 1$.
\end{definition}

The \cle{functional group} is the hydroxyl group \ce{-OH}. It is this group, not the hydrocarbon chain, that dictates the chemistry of the family. Two points must be stressed from the start:

\begin{itemize}
\item The oxygen of \ce{-OH} must be attached to a \emph{saturated} carbon. If \ce{-OH} is attached directly to a benzene ring, the compound is a \cle{phenol}, with quite different properties (studied later in this chapter).
\item The \ce{-OH} group is covalent: alcohols do \emph{not} give \ce{OH-} ions in water and are \emph{not} alkaline. Ethanol is neutral to litmus.
\end{itemize}

\courssub{Classification of Alcohols}

Just as we classified haloalkanes, we classify alcohols according to the number of alkyl groups attached to the carbon bearing the \ce{-OH} (the \emph{carbinol carbon}).

\begin{definition}[Primary, secondary and tertiary alcohols]
Let $C^*$ be the carbon atom carrying the \ce{-OH} group.
\begin{itemize}
\item A \cle{primary alcohol} ($1^\circ$): $C^*$ is bonded to \textbf{one} other carbon (or none). General pattern \ce{R-CH2-OH}. Example: propan-1-ol, \ce{CH3CH2CH2OH}.
\item A \cle{secondary alcohol} ($2^\circ$): $C^*$ is bonded to \textbf{two} other carbons. General pattern \ce{R2CH-OH}. Example: propan-2-ol, \ce{CH3CH(OH)CH3}.
\item A \cle{tertiary alcohol} ($3^\circ$): $C^*$ is bonded to \textbf{three} other carbons. General pattern \ce{R3C-OH}. Example: 2-methylpropan-2-ol, \ce{(CH3)3COH}.
\end{itemize}
\end{definition}

\begin{attention}[Why the class matters]
The class of an alcohol determines its behaviour on oxidation (primary $\rightarrow$ aldehyde $\rightarrow$ acid; secondary $\rightarrow$ ketone; tertiary $\rightarrow$ resistant), its speed in the Lucas test, and its ease of dehydration. Always state the class before predicting a reaction.
\end{attention}

\courssub{IUPAC Nomenclature of Alcohols}

\begin{methode}[Naming an alcohol]
\begin{enumerate}
\item Identify the \textbf{longest carbon chain containing the \ce{-OH} group}: this gives the parent alkane name. (For cyclic alcohols, the ring is the parent and the carbon bearing the \ce{-OH} receives number 1, e.g.\ cyclohexanol.)
\item Replace the final \emph{-e} of the alkane by the suffix \cle{-ol}: methan\textbf{ol}, ethan\textbf{ol}, propan\textbf{ol}.
\item \textbf{Number the chain from the end that gives the \ce{-OH} the lowest locant}; place the locant before the suffix: propan-1-ol, butan-2-ol.
\item Name and locate substituents as usual (alphabetical order, lowest set of locants).
\item For several \ce{-OH} groups (polyols), keep the \emph{-e} and use \emph{-diol}, \emph{-triol}: ethane-1,2-diol, propane-1,2,3-triol.
\end{enumerate}
\end{methode}

\begin{exemplebox}[Applying the rules]
\begin{itemize}
\item \ce{CH3-CH2-CH2-OH}: propan-1-ol (primary).
\item \ce{CH3-CH(OH)-CH3}: propan-2-ol (secondary); numbering from either end gives locant 2.
\item \ce{(CH3)2CH-CH2-OH}: 2-methylpropan-1-ol (primary — note the class depends on the carbinol carbon, not on the chain branching).
\item \ce{HO-CH2-CH2-OH}: ethane-1,2-diol (ethylene glycol, used in antifreeze).
\item \ce{HOCH2-CH(OH)-CH2OH}: propane-1,2,3-triol (glycerol, obtained from saponification of fats and oils such as palm oil).
\end{itemize}
\end{exemplebox}

\begin{attention}[Common naming mistakes]
Never write ``propan-3-ol'': numbering must give the \ce{-OH} the lowest locant, so it is propan-1-ol. Do not drop the locant when ambiguity is possible: ``butanol'' could be butan-1-ol or butan-2-ol. The \ce{-OH} takes priority over alkyl substituents and multiple bonds when numbering.
\end{attention}

\courssub{Structural Isomerism in Alcohols}

Alcohols with four or more carbons exhibit all three types of structural isomerism.

\begin{itemize}
\item \textbf{Chain isomerism}: the carbon skeleton differs (straight vs branched), e.g.\ butan-1-ol and 2-methylpropan-1-ol.
\item \textbf{Position isomerism}: the \ce{-OH} occupies a different position on the same skeleton, e.g.\ butan-1-ol and butan-2-ol.
\item \textbf{Functional group isomerism}: the same molecular formula corresponds to a different functional group. An alcohol \ensuremath{\mathrm{C_{n}H_{2n+2}O}} is isomeric with an \cle{ether} \ensuremath{\mathrm{R-O-R^{\prime}}}, e.g.\ ethanol \ce{C2H5OH} and methoxymethane \ce{CH3OCH3} are both \ce{C2H6O}.
\end{itemize}

\begin{methode}[Finding all structural isomers of a given formula, e.g. \ce{C4H10O}]
\begin{enumerate}
\item Write every possible \textbf{carbon skeleton} for \ce{C4}: the straight chain \ce{C-C-C-C} and the branched chain \ce{C-C(C)-C}.
\item On each skeleton, place \ce{-OH} on every \textbf{non-equivalent carbon}: these are the alcohols (chain and position isomers).
\item Then insert an oxygen \textbf{between two carbons} (\ce{C-O-C}) in all distinct ways: these are the ethers (functional isomers).
\item Check that no structure is counted twice and that the total matches the expected count (7 for \ce{C4H10O}).
\end{enumerate}
\end{methode}

Applying this to \ce{C4H10O} gives four alcohols (butan-1-ol, butan-2-ol, 2-methylpropan-1-ol, 2-methylpropan-2-ol) and three ethers (methoxypropane, 2-methoxypropane, ethoxyethane).

\begin{popfigure}
\begin{center}
\begin{tikzpicture}[scale=0.62, every node/.style={font=\small}]
% Alcohols
\node[font=\small\bfseries, popblue] at (1.5,2.6) {Alcohols (chain \& position isomers)};
% butan-1-ol
\draw[popline] (0,0)--(1,0.55)--(2,0)--(3,0.55);
\draw[popline] (3,0.55)--(3.75,0.55);
\node[popink] at (4.1,0.55) {OH};
\node[popmuted, font=\scriptsize] at (1.7,-0.7) {butan-1-ol ($1^\circ$)};
% butan-2-ol
\draw[popline] (5,0)--(6,0.55)--(7,0)--(8,0.55);
\draw[popline] (6,0.55)--(6,1.45);
\node[popink] at (6,1.85) {OH};
\node[popmuted, font=\scriptsize] at (6.5,-0.7) {butan-2-ol ($2^\circ$)};
% 2-methylpropan-1-ol
\draw[popline] (10,0)--(11,0.55)--(12,0);
\draw[popline] (11,0.55)--(11,1.45);
\draw[popline] (12,0)--(12.75,0);
\node[popink] at (13.1,0) {OH};
\node[popmuted, font=\scriptsize] at (11.3,-0.7) {2-methylpropan-1-ol ($1^\circ$)};
% 2-methylpropan-2-ol
\draw[popline] (14.5,0)--(15.5,0.55)--(16.5,0);
\draw[popline] (15.5,0.55)--(15.5,1.45);
\node[popink] at (15.5,1.85) {OH};
\node[popmuted, font=\scriptsize] at (15.5,-0.7) {2-methylpropan-2-ol ($3^\circ$)};
% Ethers
\node[font=\small\bfseries, poporange] at (1.5,-2.4) {Ethers (functional isomers)};
\draw[popline] (0,-4)--(1,-3.45)--(2,-4);
\node[popink] at (2.6,-4) {O};
\draw[popline] (3.2,-4)--(4.2,-3.45);
\node[popmuted, font=\scriptsize] at (2.1,-4.9) {methoxypropane};
\draw[popline] (6,-4)--(7,-3.45);
\node[popink] at (7.6,-3.45) {O};
\draw[popline] (8.2,-3.45)--(9.2,-4)--(10.2,-3.45);
\node[popmuted, font=\scriptsize] at (8.1,-4.9) {ethoxyethane};
\draw[popline] (12,-4)--(13,-3.45);
\node[popink] at (13.6,-3.45) {O};
\draw[popline] (14.2,-3.45)--(15.2,-4);
\draw[popline] (14.2,-3.45)--(14.2,-2.55);
\node[popmuted, font=\scriptsize] at (13.6,-4.9) {2-methoxypropane};
\end{tikzpicture}
\end{center}
\legende{The seven structural isomers of \ce{C4H10O}: four alcohols and three ethers.}
\end{popfigure}

\begin{attention}[Ethers are functional isomers of alcohols]
Same molecular formula, completely different chemistry! Ethoxyethane, \ce{C2H5OC2H5}, has the same formula \ce{C4H10O} as butan-1-ol, yet it has no \ce{O-H} bond: it does not react with sodium, cannot be oxidised by \ce{Cr2O7^{2-}}/\ce{H+}, and does not form hydrogen bonds with itself. In an examination, never forget the ethers when asked for ``all isomers of \ensuremath{\mathrm{C_{n}H_{2n+2}O}}''.
\end{attention}

\courssub{Optical Isomerism: Chirality and Enantiomers}

Look again at butan-2-ol, \ce{CH3-CH(OH)-CH2-CH3}. Its second carbon carries four \emph{different} groups: \ce{-OH}, \ce{-H}, \ce{-CH3} and \ce{-CH2CH3}. Such a carbon is called an \cle{asymmetric} or \cle{chiral carbon} (marked $C^*$). A molecule containing one chiral carbon cannot be superposed on its mirror image, exactly as your left hand cannot be superposed on your right hand.

\begin{definition}[Chirality and enantiomers]
A molecule is \cle{chiral} if it is not superposable on its mirror image. The two non-superposable mirror-image forms are called \cle{enantiomers} (optical isomers). Enantiomers have identical physical properties (melting point, boiling point, density) and identical chemical properties towards ordinary reagents, but they rotate the plane of \cle{plane-polarised light} by equal angles in \emph{opposite} directions: one is dextrorotatory ($+$), the other laevorotatory ($-$). An equimolar mixture of the two, a \cle{racemic mixture}, is optically inactive.
\end{definition}

\begin{popfigure}
\begin{center}
\begin{tikzpicture}[scale=0.9, every node/.style={font=\small}]
% Left enantiomer
\node[popink] at (0,0) {C};
\node at (-1.1,0.6) {\ce{CH3}};
\node at (1.1,0.6) {\ce{C2H5}};
\node at (-0.9,-0.8) {OH};
\node at (0.9,-0.8) {H};
\draw[popline] (0,0)--(-0.9,0.45);
\draw[popline] (0,0)--(0.9,0.45);
\draw[popline, line width=2.2pt] (0,0)--(-0.7,-0.6);
\draw[popline, dashed] (0,0)--(0.7,-0.6);
\node[popmuted, font=\scriptsize] at (0,-1.5) {($+$)-butan-2-ol};
% Mirror
\draw[popmuted, line width=1.5pt] (2.2,-1.6)--(2.2,1.1);
\node[popmuted, font=\scriptsize] at (2.2,1.35) {mirror};
% Right enantiomer
\node[popink] at (4.4,0) {C};
\node at (3.3,0.6) {\ce{C2H5}};
\node at (5.5,0.6) {\ce{CH3}};
\node at (3.5,-0.8) {H};
\node at (5.3,-0.8) {OH};
\draw[popline] (4.4,0)--(3.5,0.45);
\draw[popline] (4.4,0)--(5.3,0.45);
\draw[popline, dashed] (4.4,0)--(3.7,-0.6);
\draw[popline, line width=2.2pt] (4.4,0)--(5.1,-0.6);
\node[popmuted, font=\scriptsize] at (4.4,-1.5) {($-$)-butan-2-ol};
\end{tikzpicture}
\end{center}
\legende{The two enantiomers of butan-2-ol. The thick wedge bond projects out of the page, the dashed bond goes behind it. The two forms are mirror images and cannot be superposed.}
\end{popfigure}

\begin{exemplebox}[Spotting chiral carbons]
In 2-methylpropan-1-ol, \ce{(CH3)2CHCH2OH}, no carbon bears four different groups (the middle carbon carries two identical \ce{CH3} groups), so the molecule is \emph{achiral} — no optical isomerism. Among the four \ce{C4H10O} alcohols, only butan-2-ol is chiral. Rule: look for a carbon with \textbf{four different substituents}.
\end{exemplebox}

\courssub{Physical Properties of Alcohols}

Why is ethanol a liquid boiling at \SI{78}{\ensuremath{{}^{\circ}\mathrm{C}}} while ethane, of almost the same molar mass, is a gas boiling at \SI{-89}{\ensuremath{{}^{\circ}\mathrm{C}}}? The answer lies in the \ce{O-H} bond. Oxygen is much more electronegative than hydrogen, so the \ce{O-H} bond is strongly polar ($\ensuremath{\mathrm{O^{\delta^{-}}-H^{\delta^{+}}}}$). The hydrogen of one molecule is attracted to the lone pairs of the oxygen of a neighbouring molecule: this is a \cle{hydrogen bond}.

\begin{propriete}[Consequences of hydrogen bonding]
\begin{itemize}
\item Alcohols have \textbf{much higher boiling points} than alkanes (and ethers) of comparable molar mass, because extra energy is needed to break the hydrogen bonds before the molecules can escape into the vapour.
\item Within the homologous series, boiling point \textbf{increases with chain length} as van der Waals forces grow.
\item For a given formula, branching lowers the boiling point: butan-1-ol (\SI{118}{\ensuremath{{}^{\circ}\mathrm{C}}}) $>$ 2-methylpropan-1-ol (\SI{108}{\ensuremath{{}^{\circ}\mathrm{C}}}) $>$ 2-methylpropan-2-ol (\SI{83}{\ensuremath{{}^{\circ}\mathrm{C}}}).
\end{itemize}
(The explanation in terms of hydrogen bonding is examinable; learn to state it in one or two clear sentences.)
\end{propriete}

\begin{popfigure}
\begin{center}
\begin{tikzpicture}
\begin{axis}[
width=0.85\linewidth, height=6.2cm,
xlabel={Number of carbon atoms $n$},
ylabel={Boiling point / \si{\ensuremath{{}^{\circ}\mathrm{C}}}},
xmin=0.5, xmax=5.5, ymin=-120, ymax=130,
xtick={1,2,3,4,5},
legend pos=south east,
legend style={font=\scriptsize},
grid=both, grid style={popline!40},
]
\addplot[color=popblue, mark=*, thick] coordinates {(1,-161.5) (2,-89) (3,-42) (4,-0.5) (5,36)};
\addlegendentry{Alkanes \ensuremath{\mathrm{C_{n}H_{2n+2}}}}
\addplot[color=poporange, mark=square*, thick] coordinates {(1,64.7) (2,78.4) (3,97.2) (4,117.7) (5,137.8)};
\addlegendentry{Alcohols \ensuremath{\mathrm{C_{n}H_{2n+1}OH}}}
\draw[popmuted, dashed] (axis cs:2,-89) -- (axis cs:2,78.4);
\node[popmuted, font=\scriptsize, align=center] at (axis cs:3.1,-10) {gap due to\\ hydrogen bonding};
\end{axis}
\end{tikzpicture}
\end{center}
\legende{Boiling points of straight-chain alkanes and primary alcohols. The large gap (about \SI{170}{\ensuremath{{}^{\circ}\mathrm{C}}} for $n=2$) is the energetic signature of hydrogen bonding in alcohols.}
\end{popfigure}

\begin{propriete}[Solubility in water]
Small alcohols (methanol, ethanol, propan-1-ol) are \textbf{miscible with water in all proportions}, because their \ce{-OH} group forms hydrogen bonds with water molecules. As the hydrocarbon chain lengthens, the hydrophobic alkyl part dominates and \textbf{solubility decreases}: butan-1-ol is only moderately soluble, and alcohols beyond about \ce{C6} are practically insoluble. Polyols (ethane-1,2-diol, glycerol), with several \ce{-OH} groups, are very soluble and hygroscopic.
\end{propriete}

\begin{savaistu}[Glycerol from palm oil]
When palm oil — abundant in the South-West and Littoral regions — is boiled with sodium hydroxide (saponification), the products are soap (sodium salts of fatty acids) and propane-1,2,3-triol (glycerol). Glycerol's three \ce{-OH} groups make it an excellent humectant: it retains water, which is why it appears in skin creams and medicinal syrups.
\end{savaistu}

\begin{exoresolu}[Isomers, classes and properties of \ce{C4H10O}]
\textbf{Statement.} (a) Give the displayed or structural formulae and names of all alcohols of molecular formula \ce{C4H10O}, and state the class of each. (b) Which of them shows optical isomerism? Justify. (c) Name the functional isomers of these alcohols. (d) Arrange the four alcohols in order of increasing boiling point and explain.
\tcblower
\textbf{Solution.} (a) The four alcohols are:
\begin{itemize}
\item \ce{CH3CH2CH2CH2OH}, butan-1-ol, primary;
\item \ce{CH3CH2CH(OH)CH3}, butan-2-ol, secondary;
\item \ce{(CH3)2CHCH2OH}, 2-methylpropan-1-ol, primary;
\item \ce{(CH3)3COH}, 2-methylpropan-2-ol, tertiary.
\end{itemize}
(b) Butan-2-ol is chiral: $C_2$ bears four different groups (\ce{-OH}, \ce{-H}, \ce{-CH3}, \ce{-C2H5}), so it exists as two enantiomers that rotate plane-polarised light in opposite directions.
(c) The functional isomers are the ethers: methoxypropane \ce{CH3OCH2CH2CH3}, 2-methoxypropane \ce{CH3OCH(CH3)2} and ethoxyethane \ce{C2H5OC2H5}.
(d) Increasing boiling point: 2-methylpropan-2-ol (\SI{83}{\ensuremath{{}^{\circ}\mathrm{C}}}) $<$ butan-2-ol (\SI{99.5}{\ensuremath{{}^{\circ}\mathrm{C}}}) $<$ 2-methylpropan-1-ol (\SI{108}{\ensuremath{{}^{\circ}\mathrm{C}}}) $<$ butan-1-ol (\SI{118}{\ensuremath{{}^{\circ}\mathrm{C}}}). All four form hydrogen bonds, but branching makes the molecules more compact, reducing surface contact and hence van der Waals forces; the most branched isomer therefore boils lowest.
\end{exoresolu}

\begin{exercice}[Check your understanding]
\begin{enumerate}
\item Name the following and give the class of each: (a) \ce{CH3CH2CH2CH2CH2OH}; (b) \ce{CH3CH2CH(OH)CH2CH3}; (c) \ce{(CH3)2C(OH)CH2CH3}.
\item Write the structural formula of: (a) 3-methylbutan-2-ol; (b) propane-1,3-diol.
\item How many structural isomers (alcohols and ethers) correspond to \ce{C3H8O}? Name them all.
\item Explain why ethanol is miscible with water whereas hexan-1-ol is almost insoluble, and why both are far more soluble than the corresponding alkanes.
\item Which of these molecules is chiral: propan-2-ol, butan-2-ol, 2-methylbutan-2-ol, 3-methylpentan-3-ol? Justify each answer.
\end{enumerate}
\end{exercice}

\begin{corrige}[Exercise 1]
\begin{enumerate}
\item (a) Pentan-1-ol, primary. (b) Pentan-3-ol, secondary. (c) 2-methylbutan-2-ol, tertiary.
\item (a) \ce{(CH3)2CH-CH(OH)-CH3}. (b) \ce{HOCH2CH2CH2OH}.
\item Three: propan-1-ol \ce{CH3CH2CH2OH}, propan-2-ol \ce{CH3CH(OH)CH3}, and the ether methoxyethane \ce{CH3OC2H5}.
\item Ethanol's small hydrocarbon part allows its \ce{-OH} to hydrogen-bond extensively with water, giving miscibility. In hexan-1-ol the long hydrophobic chain disrupts water's hydrogen-bond network, so solubility is very low. Alkanes have no \ce{-OH} at all and cannot hydrogen-bond with water, so they are essentially insoluble.
\item Only butan-2-ol is chiral ($C_2$ bears \ce{H}, \ce{OH}, \ce{CH3}, \ce{C2H5}). Propan-2-ol has two identical \ce{CH3} groups on $C_2$; 2-methylbutan-2-ol has two \ce{CH3} groups on $C_2$; 3-methylpentan-3-ol has two identical \ce{C2H5} groups on $C_3$.
\end{enumerate}
\end{corrige}

\begin{aretenir}[Key points of this section]
\begin{itemize}
\item Alcohols: \ensuremath{\mathrm{C_{n}H_{2n+1}OH}}; functional group \ce{-OH} on a saturated carbon; classified primary, secondary or tertiary according to the carbinol carbon.
\item IUPAC: suffix \emph{-ol}, chain numbered to give \ce{-OH} the lowest locant; polyols: \emph{-diol}, \emph{-triol}.
\item Structural isomerism: chain, position, and functional (alcohols vs ethers); \ce{C4H10O} has 7 isomers.
\item Optical isomerism arises from a chiral carbon (four different groups); enantiomers rotate plane-polarised light in opposite directions.
\item Hydrogen bonding explains high boiling points and the water solubility of small alcohols; solubility falls as the chain lengthens.
\end{itemize}
\end{aretenir}

\coursec{Preparation and Reactions of Alcohols}

In the previous section we learnt to classify alcohols as primary, secondary or tertiary, to name them, and to explain their physical properties (high boiling points, water solubility) in terms of hydrogen bonding through the \ce{-OH} group. We now ask the two questions that make a functional group truly useful: \emph{how do we make alcohols?} and \emph{what do they do?} The answer to the second question is organised around a simple structural idea: an alcohol molecule \ce{R-O-H} can react by breaking either the \cle{\ce{C-O} bond} (the whole \ce{-OH} group leaves) or the \cle{\ce{O-H} bond} (only the hydrogen leaves). Keeping this dichotomy in mind will make every reaction in this section easy to remember.

\courssub{Preparation of Alcohols}

\textbf{Hydration of alkenes.}\par
Alkenes react with water to give alcohols. Two industrial variants exist:

\begin{itemize}
\item \emph{Direct (catalytic) hydration}: steam is passed with the alkene over a catalyst of phosphoric acid adsorbed on silica, at about $300^\circ\mathrm{C}$ and high pressure:
\[\ensuremath{\mathrm{CH_{2}=CH_{2}(g) + H_{2}O(g) \rightarrow [H_{3}PO_{4},\ 300^\circ\mathrm{C},\ 60\ atm] CH_{3}CH_{2}OH(g)}}\]
\item \emph{Indirect hydration}: the alkene is first absorbed in cold concentrated sulphuric acid to give an alkyl hydrogensulphate, which is then hydrolysed by dilution with water:
\[\ce{CH2=CH2 + H2SO4 -> CH3CH2OSO3H ->[H2O] CH3CH2OH + H2SO4}\]
\end{itemize}

With unsymmetrical alkenes, Markovnikov's rule applies: the \ce{-OH} attaches to the carbon bearing fewer hydrogens, so propene gives propan-2-ol, not propan-1-ol.

\textbf{Hydrolysis of haloalkanes.}\par
Heating a haloalkane with aqueous alkali (e.g.\ \ce{NaOH(aq)}) replaces the halogen by \ce{-OH} by nucleophilic substitution:
\[\ensuremath{\mathrm{CH_{3}CH_{2}Br + NaOH(aq) \rightarrow [\text{heat}] CH_{3}CH_{2}OH + NaBr}}\]
This is an excellent laboratory route when the haloalkane is available, and it is the reverse of the alcohol-to-haloalkane conversions we shall meet below. The mechanism is \ce{S_N2} for primary haloalkanes and \ce{S_N1} for tertiary ones.

\textbf{Fermentation of sugars.}\par
The oldest preparation of ethanol --- and one every Cameroonian student knows --- is the anaerobic fermentation of sugars by enzymes (zymase) present in yeast:
\[\ensuremath{\mathrm{C_{6}H_{12}O_{6}(aq) \rightarrow [\text{yeast enzymes, warm}] 2C_{2}H_{5}OH(aq) + 2CO_{2}(g)}}\]
The natural sugars of palm sap ferment spontaneously within hours, which is why fresh palm wine is sweet and mildly alcoholic in the morning but increasingly strong and sour by evening. The same biochemistry underlies traditional brews such as \emph{odontol} (from maize or millet) and \emph{bil-bil}. Distillation of the fermented liquor concentrates the ethanol.

\begin{attention}[Exam tip: fermentation]
Quote the fermentation equation \emph{balanced}: \ce{C6H12O6 -> 2C2H5OH + 2CO2}. State the conditions: yeast (source of enzymes), warm (about $25$--$35^\circ\mathrm{C}$; heat kills the yeast), and absence of air (anaerobic). Ethanol concentrations above roughly 15\% kill the yeast, which is why distillation is needed for stronger drinks.
\end{attention}

\textbf{Reduction of carbonyl compounds and carboxylic acids.}\par
Aldehydes, ketones and carboxylic acids can be reduced to alcohols. The choice of reagent matters:

\begin{center}
\begin{tabularx}{\linewidth}{|l|X|X|}
\hline
\rowcolor{popblueL} \textbf{Reagent} & \textbf{Reduces} & \textbf{Product} \\
\hline
\ce{NaBH4} (in water/ethanol) & aldehydes, ketones & $1^\circ$ / $2^\circ$ alcohol \\
\hline
\ce{LiAlH4} (in dry ether) & aldehydes, ketones, \textbf{carboxylic acids}, esters & $1^\circ$ / $2^\circ$ / $1^\circ$ alcohol \\
\hline
\ce{H2}/Ni or Pt (heat) & aldehydes, ketones & $1^\circ$ / $2^\circ$ alcohol \\
\hline
\end{tabularx}
\end{center}

\[\ensuremath{\mathrm{CH_{3}CHO \rightarrow [NaBH_{4},\ then\ H_{2}O] CH_{3}CH_{2}OH}} \qquad \ensuremath{\mathrm{CH_{3}COCH_{3} \rightarrow [LiAlH_{4},\ then\ H_{2}O] CH_{3}CH(OH)CH_{3}}}\]
\[\ensuremath{\mathrm{CH_{3}COOH \rightarrow [1.\ LiAlH_{4}\ 2.\ H_{2}O] CH_{3}CH_{2}OH}}\]

\begin{attention}
\ce{NaBH4} is \emph{not} strong enough to reduce carboxylic acids or esters; only \ce{LiAlH4} does that. \ce{LiAlH4} reacts violently with water, so it is used in dry ether and the product is released by a separate aqueous work-up. A very common exam slip is to write \ce{NaBH4} for the reduction of \ce{CH3COOH} to ethanol --- it scores zero.
\end{attention}

\courssub{Reactions Involving Cleavage of the C--OH Bond}

\textbf{Dehydration to alkenes (intramolecular).}\par
When an alcohol is heated with excess concentrated sulphuric or phosphoric acid (about $170^\circ\mathrm{C}$ for ethanol), a molecule of water is eliminated and an alkene forms:
\[\ensuremath{\mathrm{CH_{3}CH_{2}OH \rightarrow [conc.\ H_{2}SO_{4},\ 170^\circ\mathrm{C}] CH_{2}=CH_{2} + H_{2}O}}\]
Where several alkenes are possible, the major product is the \emph{more substituted} alkene (Zaitsev's rule). Thus butan-2-ol gives mainly but-2-ene, with some but-1-ene. The ease of dehydration follows the stability of the carbocation intermediate:
\[3^\circ > 2^\circ > 1^\circ\]

\begin{popfigure}
\begin{center}
\begin{tikzpicture}[scale=1.0, every node/.style={font=\small}]
% Mechanism of acid-catalysed dehydration of propan-2-ol
\node[font=\bfseries] at (0,4.2) {Step 1: protonation of --OH};
\draw (0,3.4) node {\ce{CH3-CH(OH)-CH3}};
\draw[->, thick, popblue] (1.8,3.4) -- (3.0,3.4) node[midway, above] {\ce{H+}};
\draw (4.6,3.4) node {\ce{CH3-CH(OH2^+)-CH3}};
\node[font=\bfseries] at (0,1.9) {Step 2: loss of water (carbocation)};
\draw[->, thick, popblue] (1.8,1.4) -- (3.0,1.4);
\draw (0,1.4) node {\ce{CH3-CH(OH2^+)-CH3}};
\draw (4.6,1.4) node {\ce{CH3-CH^+-CH3} + \ce{H2O}};
\node[font=\bfseries] at (0,-0.6) {Step 3: loss of \ce{H+} from adjacent C};
\draw[->, thick, popblue] (1.8,-1.1) -- (3.0,-1.1);
\draw (0,-1.1) node {\ce{CH3-CH^+-CH3}};
\draw (4.6,-1.1) node {\ce{CH3-CH=CH2} + \ce{H+}};
% curly arrows
\draw[->, poporange, thick] (4.2,3.15) to[bend right=40] (4.75,3.15);
\draw[->, poporange, thick] (4.3,1.15) to[bend left=40] (4.9,1.15);
\draw[->, poporange, thick] (3.9,-1.35) to[bend right=45] (4.5,-1.35);
\end{tikzpicture}
\end{center}
\legende{Acid-catalysed dehydration of propan-2-ol: protonation of the \ce{-OH} group converts it into the good leaving group \ce{H2O}; loss of water gives a secondary carbocation; loss of a proton from an adjacent carbon gives propene. The acid is regenerated --- it is a true catalyst.}
\end{popfigure}

\textbf{Formation of haloalkanes.}\par
The \ce{-OH} group can be replaced by a halogen:
\begin{itemize}
\item With hydrogen halides: \ce{CH3CH2OH + HBr -> CH3CH2Br + H2O} (reactivity \ce{HI} > \ce{HBr} > \ce{HCl}; \ce{HCl} needs anhydrous \ce{ZnCl2} as catalyst --- the basis of the Lucas test).
\item With phosphorus pentachloride: \ce{CH3CH2OH + PCl5 -> CH3CH2Cl + POCl3 + HCl} (misty fumes of \ce{HCl} --- a useful test for \ce{-OH}).
\item With thionyl chloride: \ce{CH3CH2OH + SOCl2 -> CH3CH2Cl + SO2 + HCl} (clean method: the by-products are gases).
\end{itemize}
The mechanism with \ce{HX} is nucleophilic substitution: \ce{S_N2} for primary alcohols, \ce{S_N1} (via a carbocation) for tertiary alcohols.

\textbf{Intermolecular dehydration to ethers.}\par
Under milder conditions --- excess alcohol with concentrated \ce{H2SO4} at $140^\circ\mathrm{C}$ --- two alcohol molecules condense to an ether:
\[\ensuremath{\mathrm{2CH_{3}CH_{2}OH \rightarrow [conc.\ H_{2}SO_{4},\ 140^\circ\mathrm{C}] CH_{3}CH_{2}-O-CH_{2}CH_{3} + H_{2}O}}\]

\begin{methode}[Controlling dehydration: alkene or ether?]
\begin{enumerate}
\item Want the \textbf{alkene}? Use \emph{excess concentrated acid} and heat strongly (about $170^\circ\mathrm{C}$); the alkene distils off as it forms.
\item Want the \textbf{ether}? Use the alcohol in \emph{excess} and keep the temperature at about $140^\circ\mathrm{C}$.
\item Remember: same reagents, different conditions --- temperature and which reagent is in excess decide the product. Examiners love this contrast.
\end{enumerate}
\end{methode}

\courssub{Reactions Involving Cleavage of the O--H Bond}

\textbf{Reaction with sodium or potassium.}\par
Alcohols behave as very weak acids towards active metals, giving an \cle{alkoxide} and hydrogen gas:
\[\ce{2CH3CH2OH + 2Na -> 2CH3CH2ONa + H2 ^}\]
The reaction is gentler than that of sodium with water: steady effervescence, and the sodium does not usually melt. The product, sodium ethoxide, is a strong base and a good nucleophile.

\begin{propriete}[Acidity of alcohols]
Alcohols liberate \ce{H2} from sodium (effervescence), yet they are \emph{neutral to litmus}: the \ce{O-H} bond is far less ionised than in carboxylic acids. (Proof not examinable; the observation is.) Never confuse ``reacts with Na'' with ``acidic to litmus''.
\end{propriete}

\textbf{Esterification.}\par
With a carboxylic acid (conc.\ \ce{H2SO4} catalyst, reflux), an alcohol gives an ester --- a slow, reversible reaction:
\[\ensuremath{\mathrm{CH_{3}COOH + CH_{3}CH_{2}OH \rightleftharpoons [conc.\ H_{2}SO_{4}] CH_{3}COOCH_{2}CH_{3} + H_{2}O}}\]
With an acyl chloride the reaction is fast and irreversible at room temperature:
\[\ce{CH3COCl + CH3CH2OH -> CH3COOCH2CH3 + HCl}\]

\courssub{Other Reactions and the Reaction Map}

\textbf{Oxidation (overview).}\par
With acidified \ce{K2Cr2O7} or \ce{KMnO4}, primary alcohols oxidise to aldehydes then carboxylic acids, secondary alcohols to ketones, while tertiary alcohols resist oxidation under normal conditions. Full details and the distinguishing tests follow in the next section.

\textbf{Combustion.}\par
Alcohols burn with a clean blue flame: \ce{C2H5OH + 3O2 -> 2CO2 + 3H2O}. This is why ethanol is used as a fuel and in spirit lamps.

\begin{attention}[Two opposite trends --- do not mix them up]
Ease of \textbf{dehydration}: $3^\circ > 2^\circ > 1^\circ$ (carbocation stability). Ease of \textbf{oxidation}: $1^\circ > 2^\circ \gg 3^\circ$ (tertiary alcohols resist oxidation because they have no H on the carbinol carbon). A tertiary alcohol dehydrates most easily but does not oxidise --- stating the reverse is a classic error.
\end{attention}

\begin{popfigure}
\begin{center}
\begin{tikzpicture}[scale=1.0, every node/.style={font=\small}]
\node[draw, popblue, thick, rounded corners, fill=popblueL, minimum width=2.4cm, minimum height=0.9cm] (eth) at (0,0) {ethanol \ce{C2H5OH}};
\node[draw, popgreen, thick, rounded corners] (alk) at (-5,2.2) {ethene \ce{CH2=CH2}};
\node[draw, popgreen, thick, rounded corners] (ether) at (-5,-2.2) {ethoxyethane \ce{C2H5OC2H5}};
\node[draw, popgreen, thick, rounded corners] (halo) at (0,3.0) {bromoethane \ce{C2H5Br}};
\node[draw, popgreen, thick, rounded corners] (ester) at (5,2.2) {ethyl ethanoate \ce{CH3COOC2H5}};
\node[draw, popgreen, thick, rounded corners] (alkox) at (5,-2.2) {sodium ethoxide \ce{C2H5ONa}};
\node[draw, popgreen, thick, rounded corners] (ald) at (0,-3.0) {ethanal \ce{CH3CHO}};
\draw[->, thick, popdark] (eth) -- (alk) node[midway, above, sloped, font=\scriptsize] {conc.\ \ce{H2SO4}, $170^\circ\mathrm{C}$};
\draw[->, thick, popdark] (eth) -- (ether) node[midway, below, sloped, font=\scriptsize] {conc.\ \ce{H2SO4}, $140^\circ\mathrm{C}$};
\draw[->, thick, popdark] (eth) -- (halo) node[midway, right, font=\scriptsize] {\ce{HBr} / \ce{PCl5} / \ce{SOCl2}};
\draw[->, thick, popdark] (eth) -- (ester) node[midway, above, sloped, font=\scriptsize] {\ce{CH3COOH}, \ce{H+}};
\draw[->, thick, popdark] (eth) -- (alkox) node[midway, below, sloped, font=\scriptsize] {\ce{Na}};
\draw[->, thick, popdark] (eth) -- (ald) node[midway, right, font=\scriptsize] {\ce{K2Cr2O7}/\ce{H+}, distil};
\end{tikzpicture}
\end{center}
\legende{Reaction map of ethanol. Reactions above the horizontal involve \ce{C-OH} cleavage (alkene, ether, haloalkane); those below involve \ce{O-H} cleavage (alkoxide) or oxidation (ethanal).}
\end{popfigure}

\begin{exoresolu}[Choosing the right preparation]
\textbf{Statement.} A student needs pure propan-2-ol. Suggest two laboratory/industrial routes, giving reagents, conditions and equations. \tcblower
\textbf{Solution.}
\begin{enumerate}
\item \emph{Hydration of propene} (Markovnikov): \ensuremath{\mathrm{CH_{3}CH=CH_{2} + H_{2}O \rightarrow [H_{3}PO_{4},\ 300^\circ\mathrm{C}] CH_{3}CH(OH)CH_{3}}}. The \ce{-OH} adds to the central carbon, which bears fewer hydrogens.
\item \emph{Reduction of propanone}: \ensuremath{\mathrm{CH_{3}COCH_{3} \rightarrow [NaBH_{4},\ then\ H_{2}O] CH_{3}CH(OH)CH_{3}}}. \ce{NaBH4} suffices because propanone is a ketone; \ce{LiAlH4} would also work.
\end{enumerate}
Hydrolysis of 2-bromopropane with aqueous \ce{NaOH} would be a valid third route.
\end{exoresolu}

\begin{exercice}[Products and conditions]
\begin{enumerate}
\item Give the organic product(s) when butan-2-ol is heated with excess concentrated \ce{H2SO4} at $170^\circ\mathrm{C}$, naming the major product and the rule used.
\item Write the equation for the reaction of ethanol with (a) sodium, (b) \ce{PCl5}, (c) ethanoyl chloride.
\item Explain why 2-methylpropan-2-ol dehydrates readily but gives no colour change with acidified \ce{K2Cr2O7}.
\end{enumerate}
\end{exercice}

\begin{corrige}[Exercise 1]
\begin{enumerate}
\item But-2-ene (major, Zaitsev's rule --- the more substituted alkene) and but-1-ene (minor); water is eliminated.
\item (a) \ce{2C2H5OH + 2Na -> 2C2H5ONa + H2}; (b) \ce{C2H5OH + PCl5 -> C2H5Cl + POCl3 + HCl}; (c) \ce{CH3COCl + C2H5OH -> CH3COOC2H5 + HCl}.
\item It is a tertiary alcohol: it dehydrates easily because the tertiary carbocation is very stable, but it has no hydrogen on the carbinol carbon, so it resists oxidation and the orange dichromate stays orange.
\end{enumerate}
\end{corrige}

\begin{aretenir}[Key points]
\begin{itemize}
\item Prepare alcohols by: hydration of alkenes (Markovnikov), hydrolysis of haloalkanes (\ce{NaOH(aq)}, heat), fermentation (\ce{C6H12O6 -> 2C2H5OH + 2CO2}), and reduction (\ce{NaBH4} for aldehydes/ketones; \ce{LiAlH4} also for acids and esters).
\item \ce{C-OH} cleavage: dehydration to alkenes ($170^\circ\mathrm{C}$, excess acid, Zaitsev, ease $3^\circ>2^\circ>1^\circ$); ethers at $140^\circ\mathrm{C}$ with excess alcohol; haloalkanes with \ce{HX}, \ce{PCl5} or \ce{SOCl2}.
\item \ce{O-H} cleavage: alkoxides with \ce{Na}/\ce{K} (\ce{H2} evolved, but neutral to litmus); esters with carboxylic acids or acyl chlorides.
\item Oxidation: $1^\circ \to$ aldehyde $\to$ acid; $2^\circ \to$ ketone; $3^\circ$ resists --- the opposite trend to dehydration.
\end{itemize}
\end{aretenir}

\coursec{Distinguishing Tests for Alcohols and Introduction to Phenols}

In the previous section we prepared alcohols and studied their reactions. A natural question now arises in the laboratory: \emph{given an unknown colourless liquid that is known to be an alcohol, how can we tell whether it is primary, secondary or tertiary?} This is a classic GCE Advanced Level practical question, and it is answered by three powerful tests: \cle{oxidation}, the \cle{Lucas test} and the \cle{iodoform test}. We then open the door to a new family of compounds, the \cle{phenols}, in which the hydroxyl group is attached directly to a benzene ring.

\courssub{Oxidation as a Distinguishing Test for Alcohols}

\begin{experience}[A first observation]
Three unlabelled bottles contain propan-1-ol, propan-2-ol and 2-methylpropan-2-ol. A few drops of acidified potassium dichromate(VI) solution are warmed with each liquid. The first two bottles turn the orange solution green; the third remains orange even after prolonged heating. Oxidation has already separated the tertiary alcohol from the other two.
\end{experience}

\begin{definition}[Oxidation of an alcohol]
In organic chemistry, \cle{oxidation} of an alcohol corresponds to the loss of hydrogen (removal of the O--H hydrogen and of one hydrogen from the carbon bearing the --OH group) or, equivalently, the gain of oxygen. The usual oxidising agents are \cle{acidified potassium dichromate(VI)}, \ce{K2Cr2O7/H+}, or \cle{acidified potassium manganate(VII)}, \ce{KMnO4/H+}. The oxidant is conventionally written $[\mathrm{O}]$.
\end{definition}

The outcome depends entirely on the \textbf{class} of the alcohol:

\begin{itemize}
\item A \textbf{primary alcohol} is oxidised first to an \textbf{aldehyde}, then to a \textbf{carboxylic acid}:
\[ \ce{CH3CH2OH + [O] -> CH3CHO + H2O} \qquad \ce{CH3CHO + [O] -> CH3COOH} \]
\item A \textbf{secondary alcohol} is oxidised to a \textbf{ketone}, which resists further oxidation under these conditions:
\[ \ce{CH3CH(OH)CH3 + [O] -> CH3COCH3 + H2O} \]
\item A \textbf{tertiary alcohol} has \textbf{no hydrogen on the carbon bearing the --OH group}; it therefore \textbf{resists oxidation} under normal conditions. (Vigorous oxidation with hot concentrated \ce{KMnO4} cleaves C--C bonds, but this is not a standard distinguishing test.)
\end{itemize}

\begin{propriete}[Colour changes during oxidation]
When a primary or secondary alcohol is oxidised by acidified dichromate(VI), the solution changes from \cle{orange} (\ce{Cr2O7^{2-}}) to \cle{green} (\ce{Cr^{3+}}). With acidified \ce{KMnO4}, the purple solution is decolourised. No colour change occurs with a tertiary alcohol. (The mechanism is not examinable, but the observations must be quoted precisely.)
\end{propriete}

\begin{attention}[Exam tip: state the colour change]
In the GCE, writing ``the alcohol is oxidised'' earns little credit unless you add the \textbf{observation}: \emph{orange acidified dichromate turns green} (or \emph{purple permanganate decolourises}). This colour change is the experimental \textbf{proof} that oxidation of a primary or secondary alcohol has occurred.
\end{attention}

\textbf{Controlled oxidation: distillation versus reflux.}\par
Because a primary alcohol can be oxidised in two stages, the technique used decides the product:

\begin{itemize}
\item To obtain the \textbf{aldehyde}, the alcohol is added to \emph{excess} acidified dichromate and the mixture is heated so that the aldehyde \textbf{distils off as it forms}. Since the aldehyde (e.g. ethanal, b.p. \SI{21}{\ensuremath{{}^{\circ}\mathrm{C}}}) is more volatile than the alcohol, it leaves the oxidising medium before it can be oxidised further.
\item To obtain the \textbf{carboxylic acid}, the mixture is heated under \cle{reflux} with excess oxidant, so that nothing escapes and oxidation proceeds to completion.
\end{itemize}

\begin{exemplebox}[Oxidation of butan-2-ol]
Butan-2-ol is a secondary alcohol. Warmed with acidified \ce{K2Cr2O7}, the orange solution turns green and butanone is formed:
\[ \ce{CH3CH(OH)CH2CH3 + [O] -> CH3COCH2CH3 + H2O} \]
No carboxylic acid can result, because further oxidation of a ketone would require breaking a C--C bond.
\end{exemplebox}

\courssub{The Lucas Test}

Oxidation distinguishes tertiary alcohols from the other two classes, but it cannot separate a primary from a secondary alcohol, since both are oxidised. The \cle{Lucas test} completes the identification.

\begin{definition}[Lucas reagent]
The \cle{Lucas reagent} is a solution of anhydrous zinc chloride, \ce{ZnCl2}, in concentrated hydrochloric acid. It converts an alcohol into the corresponding chloroalkane:
\[ \ce{R-OH + HCl ->[ZnCl2] R-Cl + H2O} \]
The chloroalkane is \textbf{insoluble} in the aqueous medium, so its appearance is revealed by \cle{turbidity} (cloudiness) in the test tube.
\end{definition}

The reaction is a nucleophilic substitution whose rate depends on the stability of the carbocation-like transition state: tertiary alcohols react fastest (via \cle{$S_N1$}), primary alcohols slowest (via \cle{$S_N2$}). Hence the \emph{time taken for cloudiness to appear} reveals the class of the alcohol.

\begin{methode}[Lucas test: procedure and interpretation]
\begin{enumerate}
\item Place about \SI{1}{cm^3} of the alcohol in a dry test tube.
\item Add about \SI{5}{cm^3} of Lucas reagent, stopper and shake.
\item Observe the time taken for the mixture to turn cloudy.
\end{enumerate}
\begin{center}
\begin{tabularx}{\linewidth}{|l|X|}
\hline
\rowcolor{popblueL} \textbf{Class of alcohol} & \textbf{Observation} \\
\hline
Tertiary ($3^{\circ}$) & Turbidity appears almost \textbf{immediately} \\
\hline
Secondary ($2^{\circ}$) & Turbidity appears within about \textbf{5--10 minutes} \\
\hline
Primary ($1^{\circ}$) & Turbidity appears only \textbf{very slowly} (long heating; often no visible change at room temperature) \\
\hline
\end{tabularx}
\end{center}
\end{methode}

\begin{popfigure}
\begin{tikzpicture}[scale=0.9, every node/.style={transform shape}]
% Lucas test tubes
\draw[thick, popblue] (0,0) -- (0,2.6) arc (180:360:0.45) -- (0.9,0);
\draw[thick, popblue] (0,2.6) -- (0,2.9); \draw[thick, popblue] (0.9,2.6) -- (0.9,2.9);
\fill[popblueL, opacity=0.8] (0.02,0.05) -- (0.02,1.6) arc (180:360:0.43) -- (0.88,0.05) -- cycle;
\foreach \x in {0.15,0.35,0.55,0.75} {\fill[gray, opacity=0.55] (\x,0.9) circle (0.05);}
\node[align=center, font=\small] at (0.45,-0.5) {immediate\\ ($3^{\circ}$)};
% second tube
\begin{scope}[xshift=2.6cm]
\draw[thick, popblue] (0,0) -- (0,2.6) arc (180:360:0.45) -- (0.9,0);
\draw[thick, popblue] (0,2.6) -- (0,2.9); \draw[thick, popblue] (0.9,2.6) -- (0.9,2.9);
\fill[popblueL, opacity=0.55] (0.02,0.05) -- (0.02,1.6) arc (180:360:0.43) -- (0.88,0.05) -- cycle;
\fill[gray, opacity=0.4] (0.3,0.9) circle (0.04); \fill[gray, opacity=0.4] (0.6,1.1) circle (0.04);
\node[align=center, font=\small] at (0.45,-0.5) {5--10 min\\ ($2^{\circ}$)};
\end{scope}
% third tube
\begin{scope}[xshift=5.2cm]
\draw[thick, popblue] (0,0) -- (0,2.6) arc (180:360:0.45) -- (0.9,0);
\draw[thick, popblue] (0,2.6) -- (0,2.9); \draw[thick, popblue] (0.9,2.6) -- (0.9,2.9);
\fill[popblueL, opacity=0.25] (0.02,0.05) -- (0.02,1.6) arc (180:360:0.43) -- (0.88,0.05) -- cycle;
\node[align=center, font=\small] at (0.45,-0.5) {very slow\\ ($1^{\circ}$)};
\end{scope}
% iodoform tube
\begin{scope}[xshift=8.6cm]
\draw[thick, popgold] (0,0) -- (0,2.6) arc (180:360:0.45) -- (0.9,0);
\draw[thick, popgold] (0,2.6) -- (0,2.9); \draw[thick, popgold] (0.9,2.6) -- (0.9,2.9);
\fill[popblueL, opacity=0.4] (0.02,0.4) -- (0.02,1.6) arc (180:360:0.43) -- (0.88,0.4) -- cycle;
\fill[popgold] (0.02,0.05) -- (0.02,0.4) arc (180:360:0.43) -- (0.88,0.05) -- cycle;
\node[align=center, font=\small] at (0.45,-0.5) {yellow \ce{CHI3}\\ precipitate};
\end{scope}
\end{tikzpicture}
\legende{Schematic observations: Lucas turbidity for $3^{\circ}$, $2^{\circ}$ and $1^{\circ}$ alcohols, and the yellow iodoform precipitate.}
\end{popfigure}

\courssub{The Iodoform Test}

\begin{definition}[Iodoform test]
An alcohol containing the group \ce{CH3-CH(OH)-} (i.e. ethanol or any \textbf{methyl secondary alcohol}) reacts with iodine in the presence of sodium hydroxide to give a \textbf{yellow precipitate of tri-iodomethane (iodoform)}, \ce{CHI3}, with a characteristic antiseptic smell.
\[ \ce{CH3CH(OH)R + 4I2 + 6NaOH -> CHI3 v + RCOONa + 5NaI + 5H2O} \]
For ethanol ($\mathrm{R = H}$): \quad $\ce{CH3CH2OH + 4I2 + 6NaOH -> CHI3 v + HCOONa + 5NaI + 5H2O}$
\end{definition}

The reagent first oxidises the alcohol to a carbonyl compound bearing a \ce{CH3-CO-} group, which is then tri-iodinated at the methyl group and cleaved by hydroxide. Consequently the test is also positive for \textbf{aldehydes and ketones} containing the \ce{CH3CO-} group (e.g. ethanal, propanone).

\begin{attention}[Warning: which alcohols respond?]
Only alcohols bearing the \ce{CH3-CH(OH)-} skeleton give a positive iodoform test: \textbf{ethanol} and \textbf{secondary alcohols of the type \ce{CH3CH(OH)R}} (e.g. propan-2-ol, butan-2-ol, pentan-2-ol). Methanol, propan-1-ol, butan-1-ol, 2-methylpropan-1-ol, 3-methylbutan-2-ol (no \ce{CH3CH(OH)-}) and 2-methylpropan-2-ol give a \textbf{negative} test. Do not write that ``all secondary alcohols'' respond — pentan-3-ol (\ce{CH3CH2CH(OH)CH2CH3}) gives no iodoform!
\end{attention}

\begin{exoresolu}[Identifying an unknown alcohol]
An alcohol $X$, \ce{C4H10O}, gives a negative iodoform test. Warmed with acidified \ce{K2Cr2O7}, the orange solution turns green, and the organic product does not react with Fehling's solution. Identify $X$.
\tcblower
\textbf{Step 1: Molecular formula \ce{C4H10O}.} The four structural isomers are: butan-1-ol, butan-2-ol, 2-methylpropan-1-ol, 2-methylpropan-2-ol.

\textbf{Step 2: Negative iodoform test.} This excludes ethanol (not \ce{C4}) and any alcohol with the \ce{CH3CH(OH)-} group. Among the isomers, \textbf{butan-2-ol} (\ce{CH3CH(OH)CH2CH3}) has this group and would give a positive test. So $X$ is \textbf{not butan-2-ol}.

\textbf{Step 3: Acidified dichromate turns green.} Oxidation occurs $\Rightarrow$ $X$ is primary or secondary. This excludes \textbf{2-methylpropan-2-ol} (tertiary).

\textbf{Step 4: Oxidation product does not reduce Fehling's solution.} The product is not an aldehyde $\Rightarrow$ $X$ is not a primary alcohol. This excludes \textbf{butan-1-ol} and \textbf{2-methylpropan-1-ol}.

\textbf{Conclusion:} The only remaining isomer is \textbf{butan-2-ol} — but Step 2 excluded it! The data are inconsistent. In a well-set GCE question, the three tests give a unique fingerprint for each isomer:
\begin{itemize}
\item Butan-1-ol: green dichromate + Fehling's positive (aldehyde) + no iodoform.
\item Butan-2-ol: green dichromate + Fehling's negative (ketone) + \textbf{yellow iodoform}.
\item 2-Methylpropan-1-ol: green dichromate + Fehling's positive (aldehyde) + no iodoform.
\item 2-Methylpropan-2-ol: no colour change (tertiary) + no Fehling's test needed + no iodoform.
\end{itemize}
The question as stated contains a contradiction; the correct interpretation is that \emph{each isomer gives a distinct combination of results}.
\end{exoresolu}

\begin{popfigure}
\begin{tikzpicture}[
box/.style={draw=popblue, thick, rounded corners, fill=popblueL, align=center, font=\small, text width=3.1cm, inner sep=4pt},
res/.style={draw=popgreen, thick, rounded corners, fill=popgreenL, align=center, font=\small, text width=2.9cm, inner sep=4pt},
arr/.style={-{Stealth[length=2.5mm]}, thick, popdark}]
\node[box, align=center] (start) at (0,0){Unknown alcohol\\ Warm with acidified \ce{K2Cr2O7}};
\node[box, align=center] (ox) at (-3.4,-2.4){Orange $\rightarrow$ green\\ Oxidation product tested with Fehling's};
\node[res, align=center] (tert) at (3.6,-2.4){No colour change\\ $\Rightarrow$ \textbf{tertiary alcohol}};
\node[res, align=center] (prim) at (-5.6,-5.0){Fehling's: brick-red ppt (aldehyde)\\ $\Rightarrow$ \textbf{primary}};
\node[res, align=center] (sec) at (-1.4,-5.0){Fehling's: no change (ketone)\\ $\Rightarrow$ \textbf{secondary}};
\draw[arr] (start) -- (ox) node[midway, left, font=\scriptsize]{oxidised};
\draw[arr] (start) -- (tert) node[midway, right, font=\scriptsize]{not oxidised};
\draw[arr] (ox) -- (prim);
\draw[arr] (ox) -- (sec);
\node[align=center, font=\scriptsize, text width=8.5cm] at (-0.8,-6.6) {Confirmation: Lucas test --- turbidity immediate ($3^{\circ}$), in minutes ($2^{\circ}$), very slow ($1^{\circ}$). Iodoform test detects the \ce{CH3CH(OH)-} group.};
\end{tikzpicture}
\legende{Identification tree for primary, secondary and tertiary alcohols.}
\end{popfigure}

\begin{exercice}[Distinguishing three isomers]
You are given three liquids: butan-1-ol, butan-2-ol and 2-methylpropan-2-ol. Using only (a) acidified dichromate, (b) Lucas reagent and (c) \ce{I2/NaOH}, describe how you would identify each liquid, giving all observations.
\end{exercice}

\begin{corrige}[Distinguishing three isomers]
\textbf{Acidified dichromate:} butan-1-ol and butan-2-ol turn the orange solution green; 2-methylpropan-2-ol gives no change (tertiary, resists oxidation). \textbf{Lucas test:} 2-methylpropan-2-ol clouds immediately, butan-2-ol within 5--10 minutes, butan-1-ol only very slowly (or on heating). \textbf{Iodoform test:} only butan-2-ol (\ce{CH3CH(OH)CH2CH3}) gives a yellow \ce{CHI3} precipitate. Each test alone separates the three alcohols.
\end{corrige}

\courssub{Introduction to Phenols}

\begin{definition}[Phenol versus aromatic alcohol]
A \cle{phenol} is a compound in which a hydroxyl group is attached \textbf{directly to a carbon atom of a benzene ring}: \ce{C6H5OH}. If the --OH group is carried by a side chain (e.g. \ce{C6H5CH2OH}, phenylmethanol), the compound is an \textbf{aromatic alcohol}, \emph{not} a phenol, and it behaves essentially like an aliphatic alcohol.
\end{definition}

\begin{attention}[Common mistake]
\ce{C6H5CH2OH} is \textbf{not} a phenol. The decisive criterion is that the --OH must be bonded to a ring carbon. Examiners frequently set this trap.
\end{attention}

\textbf{Nomenclature.} The simplest member, \ce{C6H5OH}, is \cle{phenol} itself. Methyl derivatives are the \cle{methylphenols}, traditionally called \textbf{cresols}: 2-methylphenol (\emph{ortho}-cresol), 3-methylphenol (\emph{meta}) and 4-methylphenol (\emph{para}). Compounds with two --OH groups on the ring are the \textbf{dihydric phenols}: benzene-1,2-diol (catechol), benzene-1,3-diol (resorcinol) and benzene-1,4-diol (hydroquinone).

\textbf{Physical properties.} Phenol is a colourless crystalline solid (m.p. about \SI{41}{\ensuremath{{}^{\circ}\mathrm{C}}}) that turns pink on exposure to air and light owing to oxidation. Like alcohols, phenols form \cle{hydrogen bonds}, so their boiling points are high (phenol b.p. \SI{182}{\ensuremath{{}^{\circ}\mathrm{C}}}) and the lower members are moderately soluble in water (phenol itself: about \SI{9}{g} per \SI{100}{g} of water at room temperature, fully miscible above about \SI{66}{\ensuremath{{}^{\circ}\mathrm{C}}}). Phenol is toxic and causes white burns on the skin; handle with gloves.

\courssub{Preparation of Phenol}

\textbf{1. From benzene sulphonic acid (laboratory/classical route).} Benzene is sulphonated with concentrated sulphuric acid; the sodium salt of benzene sulphonic acid is then fused with solid sodium hydroxide to give sodium phenoxide, which is acidified:
\[ \ce{C6H6 + H2SO4(conc) -> C6H5SO3H + H2O} \]
\[ \ensuremath{\mathrm{C_{6}H_{5}SO_{3}Na + 2NaOH(s) \rightarrow [\text{fuse, 300 \ensuremath{{}^{\circ}\mathrm{C}}}] C_{6}H_{5}ONa + Na_{2}SO_{3} + H_{2}O}} \qquad \ce{C6H5ONa + HCl -> C6H5OH + NaCl} \]

\textbf{2. From benzene diazonium salts.} Phenylamine (aniline) is diazotised with \ce{NaNO2/HCl} at $0$--$\SI{5}{\ensuremath{{}^{\circ}\mathrm{C}}}$; warming the diazonium salt with water then gives phenol (this route is studied in detail with the nitrogen compounds):
\[ \ensuremath{\mathrm{C_{6}H_{5}N_{2}+Cl^{-} + H_{2}O \rightarrow [\text{warm}] C_{6}H_{5}OH + N_{2} ^ + HCl}} \]

\textbf{3. The cumene process (industrial).} Benzene is alkylated with propene to give cumene (2-phenylpropane), which is oxidised by air to a hydroperoxide and then cleaved by dilute acid to \textbf{phenol and propanone} — both valuable products:
\[ \ce{C6H5CH(CH3)2 ->[O2] C6H5C(CH3)2OOH ->[H+] C6H5OH + CH3COCH3} \]

\courssub{Reactions of the --OH Group of Phenol}

\begin{propriete}[Acidity of phenol]
Phenol is a \cle{weak acid}, stronger than water and alcohols but \textbf{weaker than carbonic acid}:
\[ \ce{C6H5OH + H2O <=> C6H5O- + H3O+} \qquad pK_a \approx 10 \]
The phenoxide ion \ce{C6H5O-} is stabilised by \textbf{delocalisation of the negative charge into the benzene ring}, which is why phenol is far more acidic than ethanol ($pK_a \approx 16$), whose ethoxide ion has no such stabilisation.
\end{propriete}

Three consequences of this acidity are examinable:

\textbf{(a) Reaction with sodium hydroxide.} Phenol dissolves in aqueous \ce{NaOH}, forming sodium phenoxide:
\[ \ce{C6H5OH + NaOH -> C6H5ONa + H2O} \]

\textbf{(b) No reaction with sodium carbonate.} Because phenol is weaker than carbonic acid ($pK_{a1}(\ce{H2CO3}) \approx 6.3$), it \textbf{cannot displace \ce{CO2}} from carbonates or hydrogencarbonates — unlike carboxylic acids, which effervesce with \ce{Na2CO3}. This difference is used to distinguish phenol from a carboxylic acid.

\textbf{(c) Reaction with sodium metal.} Like alcohols, phenol liberates hydrogen:
\[ \ce{2C6H5OH + 2Na -> 2C6H5ONa + H2 ^} \]

\textbf{Ester formation.} Phenol does \textbf{not} esterify directly with carboxylic acids (the equilibrium is unfavourable), but it reacts readily with the more reactive \textbf{acyl chlorides} to give esters:
\[ \ce{C6H5OH + CH3COCl -> CH3COOC6H5 + HCl} \]
The product is phenyl ethanoate. This reaction also illustrates that phenol retains the nucleophilic character of the --OH group even though its acidity is enhanced.

\begin{aretenir}[Key points]
\begin{itemize}
\item Acidified \ce{K2Cr2O7}: $1^{\circ}$ alcohol $\rightarrow$ aldehyde (distil off) $\rightarrow$ acid (reflux); $2^{\circ}$ $\rightarrow$ ketone; $3^{\circ}$ resists. Observation: orange $\rightarrow$ green.
\item Lucas test (\ce{ZnCl2/HCl}): turbidity immediate ($3^{\circ}$), within 5--10 min ($2^{\circ}$), very slow ($1^{\circ}$).
\item Iodoform test (\ce{I2/NaOH}): yellow \ce{CHI3} precipitate only for \ce{CH3CH(OH)-} alcohols (and \ce{CH3CO-} carbonyls).
\item Phenol = --OH directly on the benzene ring; \ce{C6H5CH2OH} is an aromatic alcohol.
\item Phenol is acidic enough to react with \ce{NaOH} and \ce{Na}, but too weak to liberate \ce{CO2} from \ce{Na2CO3}; it is esterified by acyl chlorides, not by carboxylic acids.
\end{itemize}
\end{aretenir}

\begin{exercice}[Phenol or not?]
For each compound, state whether it is a phenol and whether it reacts with (i) \ce{NaOH}, (ii) \ce{Na2CO3}, (iii) \ce{Na}: (a) \ce{C6H5OH}; (b) \ce{C6H5CH2OH}; (c) 4-methylphenol.
\end{exercice}

\begin{corrige}[Phenol or not?]
(a) \ce{C6H5OH}: a phenol; reacts with \ce{NaOH} (phenoxide formed), with \ce{Na} (\ce{H2} evolved), but \textbf{not} with \ce{Na2CO3}. (b) \ce{C6H5CH2OH}: \textbf{not} a phenol (OH on the side chain); reacts with \ce{Na} only. (c) 4-methylphenol: a phenol; behaves like phenol — reacts with \ce{NaOH} and \ce{Na}, not with \ce{Na2CO3}.
\end{corrige}

\begin{savaistu}[Phenol and antiseptics]
Dilute solutions of phenol were among the first surgical antiseptics used in the 19th century, and phenol is still a reference standard for rating disinfectants. Many antiseptic liquids sold in Cameroonian pharmacies contain chloroxylenol, a chlorinated methylphenol — a direct application of this chemistry.
\end{savaistu}

\coursec{Ring Reactions of Phenol and Comparison with Alcohols}

In the previous section we learnt how the Lucas test, the iodoform test and controlled oxidation allow us to tell primary, secondary and tertiary alcohols apart, and we met phenol, \ce{C6H5OH}, for the first time. We saw that although phenol carries an \ce{-OH} group, it behaves in many ways \emph{unlike} an ordinary alcohol. In this section we complete the picture: we study the reactions that take place \emph{on the benzene ring} of phenol, and then we draw a rigorous, examination-ready comparison between the hydroxyl group in alcohols and in phenols.

\courssub{Naming and Physical Properties of Phenol}

\textbf{Naming.} The parent compound is \cle{phenol} (IUPAC: benzenol). Substituted phenols are named by numbering the ring so that the carbon bearing the \ce{-OH} group is C-1. The positions \emph{ortho} (2,6), \emph{meta} (3,5) and \emph{para} (4) are used in common names (e.g., \emph{o}-cresol = 2-methylphenol, \emph{p}-nitrophenol = 4-nitrophenol). For di- and tri-substituted phenols, numbering gives the lowest set of locants to the substituents.

\textbf{Physical properties.} Phenol is a colourless crystalline solid at room temperature (m.p. $41^\circ\mathrm{C}$) but is often encountered as a liquid due to its hygroscopic nature. It has a characteristic sweet, tarry odour. It is moderately soluble in water (\SI{8.3}{g/100 mL} at $20^\circ\mathrm{C}$) but infinitely miscible above $66^\circ\mathrm{C}$. Its boiling point ($182^\circ\mathrm{C}$) is much higher than that of ethanol ($78^\circ\mathrm{C}$) due to stronger hydrogen bonding and a larger electron cloud. Phenol is neutral to litmus (unlike carboxylic acids) but turns blue litmus red very slightly, hinting at its weak acidity.

\courssub{Preparation of Phenol}

Three routes are required for the GCE examination.

\begin{enumerate}
\item \textbf{From benzene diazonium chloride (laboratory method).} The diazonium salt, prepared from aniline, is hydrolysed by warming with water:
\[
\ce{C6H5N2+Cl- (aq) + H2O(l) ->[warm] C6H5OH(aq) + N2(g) + HCl(aq)}
\]
This is a key step in aromatic synthesis because it allows the introduction of an \ce{-OH} group onto the ring via the amino group.

\item \textbf{From chlorobenzene (Dow process, industrial).} Chlorobenzene is fused with concentrated \ce{NaOH} at high temperature and pressure:
\[
\ensuremath{\mathrm{C_{6}H_{5}Cl(l) + 2NaOH(aq) \rightarrow [350^\circ\mathrm{C}, 300 atm] C_{6}H_{5}ONa(aq) + NaCl(aq) + H_{2}O(l)}}
\]
The sodium phenoxide is then acidified to liberate phenol.

\item \textbf{From cumene (cumene-phenol process, major industrial route).} Cumene (isopropylbenzene) is oxidised to cumene hydroperoxide, which is cleaved by acid:
\[
\ce{C6H5CH(CH3)2 + O2 -> C6H5C(OOH)CH3 ->[H+] C6H5OH + CH3COCH3}
\]
Acetone is a valuable co-product.
\end{enumerate}

\begin{aretenir}[Exam focus]
Be able to write the equation for the hydrolysis of benzene diazonium chloride with state symbols and conditions. The industrial routes are often asked as ``outline the process'' or ``name the co-product'' (acetone).
\end{aretenir}

\courssub{Activation of the Benzene Ring by the \ce{-OH} Group}

\textbf{Observation.} Benzene itself reacts with bromine only in the presence of an iron(III) bromide catalyst, and only one hydrogen is replaced. Phenol, by contrast, decolourises bromine water \emph{instantly, in the cold, with no catalyst at all}, and \emph{three} hydrogens are replaced. Something about the \ce{-OH} group has made the ring dramatically more reactive.

\textbf{Explanation.} The oxygen atom of the \ce{-OH} group possesses lone pairs of electrons. One of these lone pairs can be delocalised into the $\pi$-system of the benzene ring. The resonance structures of phenol (and more strongly of the phenoxide ion) place a \cle{negative charge} on the ring carbons at the \emph{ortho} (2 and 6) and \emph{para} (4) positions:

\begin{popfigure}
\begin{center}
\begin{tikzpicture}[scale=0.85]
% Structure 1: phenoxide
\draw[thick] (0:1) -- (60:1) -- (120:1) -- (180:1) -- (240:1) -- (300:1) -- cycle;
\draw (30:0.78) arc (30:150:0.78);
\draw (150:0.78) arc (150:270:0.78);
\draw (270:0.78) arc (270:390:0.78);
\node at (90:1.45) {$\ce{O^{-}}$};
\draw[thick] (90:1) -- (90:1.25);
\node[below] at (0,-1.35) {phenoxide ion};
% arrow
\draw[<->,thick,poporange] (1.7,0) -- (2.9,0);
% Structure 2: charge at ortho
\begin{scope}[xshift=4.6cm]
\draw[thick] (0:1) -- (60:1) -- (120:1) -- (180:1) -- (240:1) -- (300:1) -- cycle;
\draw[thick] (90:1) -- (90:1.25);
\node at (90:1.5) {$\ce{O}$};
\draw (60:0.8) arc (60:120:0.8);
\draw (180:0.8) arc (180:240:0.8);
\draw (300:0.8) arc (300:360:0.8);
\node[popdark] at (30:1.5) {$\delta^{-}$};
\node[below] at (0,-1.35) {charge at C-2 (ortho)};
\end{scope}
\draw[<->,thick,poporange] (6.3,0) -- (7.5,0);
% Structure 3: charge at para
\begin{scope}[xshift=9.2cm]
\draw[thick] (0:1) -- (60:1) -- (120:1) -- (180:1) -- (240:1) -- (300:1) -- cycle;
\draw[thick] (90:1) -- (90:1.25);
\node at (90:1.5) {$\ce{O}$};
\draw (0:0.8) arc (0:60:0.8);
\draw (120:0.8) arc (120:240:0.8);
\draw (300:0.8) arc (300:360:0.8);
\node[popdark] at (270:1.55) {$\delta^{-}$};
\node[below] at (0,-1.35) {charge at C-4 (para)};
\end{scope}
\end{tikzpicture}
\end{center}
\legende{Delocalisation of the oxygen lone pair into the ring (shown here for the phenoxide ion, where the effect is strongest). Negative charge appears at the 2-, 4- and 6-positions.}
\end{popfigure}

\begin{propriete}[Activating and 2,4-directing effect of \ce{-OH}]
The \ce{-OH} group is an \cle{electron-donating} group by resonance (+M effect). It (i) \emph{activates} the benzene ring towards electrophilic substitution, making phenol far more reactive than benzene, and (ii) directs incoming electrophiles to the \cle{ortho} (2, 6) and \cle{para} (4) positions, because these positions carry increased electron density in the resonance hybrid. A proof is not examinable, but you must be able to \emph{explain} the effect using lone-pair delocalisation.
\end{propriete}

\begin{attention}[Do not confuse the two effects of oxygen]
Oxygen is electronegative, so by the \emph{inductive} effect (-I) it withdraws electron density. In phenol, however, the \emph{resonance (mesomeric) donation} of the lone pair (+M) dominates for ring reactivity. Never write that \ce{-OH} deactivates the ring: experimentally, phenol reacts much faster than benzene.
\end{attention}

\courssub{Halogenation of Phenol}

\begin{experience}[Reaction of phenol with bromine water]
Add a few drops of bromine water to a dilute aqueous solution of phenol at room temperature. The orange colour of bromine disappears \emph{immediately} and a \cle{white precipitate} of 2,4,6-tribromophenol forms. No catalyst, no heat, no special conditions are required.
\end{experience}

\[
\ce{C6H5OH(aq) + 3Br2(aq) -> C6H2Br3OH(s) + 3HBr(aq)}
\]

The product is \textbf{2,4,6-tribromophenol}, a white solid. Contrast this with benzene, which reacts with bromine only as liquid bromine in the presence of \ce{FeBr3}, giving monobromobenzene:

\[
\ce{C6H6(l) + Br2(l) ->[FeBr3] C6H5Br(l) + HBr(g)}
\]

\begin{aretenir}[Exam tip: the bromine water test]
A white precipitate with bromine water is characteristic of phenol. Quote the equation in full, with the \emph{three} molecules of \ce{Br2} and the name of the product:
\ce{C6H5OH + 3Br2 -> C6H2Br3OH(s) + 3HBr}. Forgetting the ``3'' or naming the product ``bromophenol'' are the two most common ways of losing marks here.
\end{aretenir}

\courssub{Nitration of Phenol}

Because the ring is activated, phenol is nitrated under much milder conditions than benzene.

\textbf{With dilute nitric acid} at room temperature, a mixture of \textbf{2-nitrophenol} (ortho) and \textbf{4-nitrophenol} (para) is obtained:

\[
\ce{C6H5OH(aq) + HNO3(aq, dil) -> 2-NO2-C6H4OH + 4-NO2-C6H4OH + H2O(l)}
\]

The two isomers can be separated by steam distillation: the 2-isomer forms intramolecular hydrogen bonds and is more volatile, while the 4-isomer forms intermolecular hydrogen bonds and is less volatile.

\textbf{With concentrated nitric acid} (in the presence of concentrated sulphuric acid), all three activated positions are substituted, giving \textbf{2,4,6-trinitrophenol}, commonly called \cle{picric acid}:

\[
\ensuremath{\mathrm{C_{6}H_{5}OH + 3HNO_{3}(conc) \rightarrow [conc\ H_{2}SO_{4}] C_{6}H_{2}(NO_{2})_{3}OH + 3H_{2}O}}
\]

\begin{savaistu}[Picric acid]
2,4,6-Trinitrophenol is a bright yellow solid once used as a dye and as an explosive. It is also a much stronger acid than phenol itself ($pK_a \approx 0.4$), because the three nitro groups stabilise the phenoxide ion even further by withdrawing electron density (-I, -M). It is a beautiful illustration of how substituents tune acidity.
\end{savaistu}

\courssub{Sulphonation and Friedel--Crafts Alkylation of Phenol}

\textbf{Sulphonation.} Phenol reacts with concentrated sulphuric acid. At low temperature (about $15^\circ\mathrm{C}$ to $20^\circ\mathrm{C}$) the main product is \textbf{2-hydroxybenzenesulphonic acid} (kinetic product, ortho); on heating to about $100^\circ\mathrm{C}$ the \textbf{4-hydroxybenzenesulphonic acid} (thermodynamic product, para) predominates:

\[
\ce{C6H5OH + H2SO4(conc) -> HO-C6H4-SO3H + H2O}
\]

\textbf{Friedel--Crafts alkylation.} Phenol reacts with a haloalkane in the presence of anhydrous aluminium chloride, \ce{AlCl3}, to introduce an alkyl group at the 2- and 4-positions. For example, with chloromethane:

\[
\ensuremath{\mathrm{C_{6}H_{5}OH + CH_{3}Cl \rightarrow [anhyd.\ AlCl_{3}] 2-CH_{3}-C_{6}H_{4}OH + 4-CH_{3}-C_{6}H_{4}OH + HCl}}
\]

The products (2-methylphenol and 4-methylphenol) are the \emph{cresols}, used in disinfectants. Note that, as always in Friedel--Crafts reactions, the catalyst must be \emph{anhydrous} \ce{AlCl3}.

\begin{popfigure}
\begin{center}
\begin{tikzpicture}[scale=1.0]
\draw[thick] (0:1) -- (60:1) -- (120:1) -- (180:1) -- (240:1) -- (300:1) -- cycle;
\draw (30:0.78) arc (30:150:0.78);
\draw (150:0.78) arc (150:270:0.78);
\draw (270:0.78) arc (270:390:0.78);
\draw[thick] (90:1) -- (90:1.4);
\node at (90:1.7) {$\ce{OH}$};
\node[popdark] at (30:1.6) {\textbf{2}};
\node[popdark] at (330:1.6) {\textbf{6}};
\node[popdark] at (270:1.6) {\textbf{4}};
\node[popmuted] at (150:1.55) {3};
\node[popmuted] at (210:1.55) {5};
\node[align=center] at (0,-2.2) {Electrophiles attack positions \textbf{2, 4, 6}};
\end{tikzpicture}
\end{center}
\legende{Numbering of the phenol ring. The \ce{-OH} group is at C-1; substitution occurs at the ortho (2, 6) and para (4) positions.}
\end{popfigure}

\courssub{Distinguishing Test for Phenols: Neutral Iron(III) Chloride}

\begin{experience}[The \ce{FeCl3} test]
Add a few drops of \emph{neutral} iron(III) chloride solution to a solution of the suspected phenol. A \cle{violet (purple) colouration} appears immediately, due to the formation of a coloured complex between \ce{Fe^{3+}} and the phenoxide ion. Alcohols give \emph{no} colouration.
\end{experience}

This test is given by virtually all compounds containing the \ce{-OH} group attached directly to a benzene ring (the \emph{enol}-type arrangement), and it is the standard GCE test for phenol.

\courssub{Comparison of the \ce{-OH} Group in Alcohols and Phenols}

We now bring together everything learnt in Lessons 41--49 into a systematic comparison. The central idea is the \emph{acidity} of the O--H bond.

\textbf{Why is phenol acidic but ethanol not?} When phenol loses \ce{H+}, the phenoxide ion \ce{C6H5O-} is formed. Its negative charge is delocalised over the ring (positions 2, 4 and 6), as shown in the resonance figure above. When ethanol loses \ce{H+}, the ethoxide ion \ce{C2H5O-} has its charge \emph{localised} on oxygen, and the ethyl group even pushes electron density onto it (+I effect), destabilising it further.

\begin{propriete}[Phenol is more acidic than ethanol]
The \cle{phenoxide ion} is stabilised by \cle{resonance}: the lone pair on oxygen is delocalised into the benzene ring, spreading the negative charge over the 2-, 4- and 6-positions. The ethoxide ion has no such stabilisation. Phenol therefore ionises appreciably ($K_a \approx 1.3 \times 10^{-10}$, $pK_a \approx 10$) whereas ethanol is essentially neutral ($K_a \approx 10^{-16}$, $pK_a \approx 16$). You must be able to state this reasoning in words for the examination.
\end{propriete}

\textbf{Reaction with sodium hydroxide.} Phenol is acidic enough to react with \ce{NaOH}, forming soluble sodium phenoxide:

\[
\ce{C6H5OH(aq) + NaOH(aq) -> C6H5O^{-}Na^{+}(aq) + H2O(l)}
\]

Ethanol does \emph{not} react with \ce{NaOH}. This is a practical distinguishing test: phenol dissolves in \ce{NaOH} solution, ethanol (already miscible with water) shows no reaction.

\begin{attention}[Phenol does NOT react with carbonates or hydrogencarbonates]
Phenol is a \emph{weaker} acid than carbonic acid ($H_2CO_3$), so it does \textbf{not} liberate \ce{CO2} from \ce{Na2CO3} or \ce{NaHCO3} solutions. Carboxylic acids, being stronger, do:
\ce{RCOOH + NaHCO3 -> RCOONa + H2O + CO2 ^}.
Hence the hydrogencarbonate test distinguishes a carboxylic acid (effervescence) from a phenol (no reaction). This is a favourite GCE trap.
\end{attention}

\textbf{Reaction with sodium metal.} Both phenol and ethanol react with sodium to liberate hydrogen gas, though phenol reacts more vigorously due to its greater acidity:
\[
\ce{2C6H5OH + 2Na -> 2C6H5ONa + H2 ^}
\quad
\ce{2C2H5OH + 2Na -> 2C2H5ONa + H2 ^}
\]

\textbf{Esterification.} Alcohols react readily with carboxylic acids (acid catalyst) or, better, with acyl chlorides, to give esters. Phenol reacts \emph{very slowly} with carboxylic acids because the lone pair on its oxygen is partly tied up in the ring, making it a poorer nucleophile; in practice phenyl esters are prepared using \emph{acyl chlorides} or \emph{acid anhydrides}:

\[
\ce{C6H5OH + CH3COCl -> CH3COOC6H5 + HCl}
\]

\textbf{Oxidation.} Primary and secondary alcohols are easily oxidised by acidified \ce{K2Cr2O7} or \ce{KMnO4} (to aldehydes/acids and ketones respectively). Phenol resists oxidation under these conditions because oxidation would require breaking into the stable aromatic ring; forcing conditions destroy the ring and give complex mixtures (e.g., benzoquinone, then tar).

\textbf{Hydrogen bonding.} Both alcohols and phenols form \cle{hydrogen bonds}, hence both have higher boiling points than alkanes of similar mass and the smaller members are water-soluble. Phenol (b.p.\ $182^\circ\mathrm{C}$) boils higher than ethanol (b.p.\ $78^\circ\mathrm{C}$) because of its larger electron cloud and stronger intermolecular forces; it is only moderately soluble in water because the large hydrocarbon ring is hydrophobic.

\begin{center}
\begin{tabularx}{\linewidth}{|l|X|X|}
\hline
\rowcolor{popblueL} \textbf{Property} & \textbf{Ethanol (alcohol)} & \textbf{Phenol} \\
\hline
Acidity & Neutral; $K_a \approx 10^{-16}$ & Weakly acidic; $K_a \approx 1.3 \times 10^{-10}$ \\
\hline
With \ce{NaOH} & No reaction & Reacts; forms sodium phenoxide \\
\hline
With \ce{Na2CO3}/\ce{NaHCO3} & No reaction & No reaction (no \ce{CO2}) \\
\hline
With neutral \ce{FeCl3} & No colouration & Violet colouration \\
\hline
With bromine water & No reaction & White precipitate of 2,4,6-tribromophenol \\
\hline
Esterification & Easy with \ce{RCOOH}/\ce{H+} & Needs acyl chloride or anhydride \\
\hline
Oxidation by \ce{K2Cr2O7}/\ce{H+} & Readily oxidised & Resists oxidation \\
\hline
Reaction with Na & \ce{H2} evolved & \ce{H2} evolved (both have \ce{-OH}) \\
\hline
\end{tabularx}
\end{center}

\begin{methode}[Three tests to distinguish phenol from ethanol]
\begin{enumerate}
\item \textbf{Neutral \ce{FeCl3} solution:} phenol gives a violet colouration; ethanol gives none.
\item \textbf{Bromine water:} phenol instantly decolourises it with a white precipitate of 2,4,6-tribromophenol; ethanol leaves the orange colour unchanged.
\item \textbf{\ce{NaOH} solution:} phenol reacts and dissolves to form sodium phenoxide; ethanol shows no acid--base reaction (confirm with universal indicator: phenol solution is weakly acidic, ethanol is neutral).
\end{enumerate}
Always state the \emph{reagent}, the \emph{observation with phenol} and the \emph{observation with ethanol} to earn full marks.
\end{methode}

\begin{exoresolu}[Distinguishing three unknown liquids]
A student is given three unlabelled liquids: ethanol, phenol solution and ethanoic acid. Using only bromine water, neutral iron(III) chloride solution and sodium hydrogencarbonate solution, show how each liquid can be identified, giving observations and equations where appropriate.
\tcblower
\textbf{Step 1: sodium hydrogencarbonate.} Add \ce{NaHCO3} solution to a sample of each. Only ethanoic acid, a carboxylic acid strong enough to displace carbonic acid, gives effervescence of \ce{CO2}:
\[
\ce{CH3COOH(aq) + NaHCO3(aq) -> CH3COONa(aq) + H2O(l) + CO2(g) ^}
\]
Ethanol and phenol give no reaction (phenol is too weak an acid).

\textbf{Step 2: bromine water.} To the two remaining liquids add bromine water. Phenol instantly decolourises it and gives a white precipitate of 2,4,6-tribromophenol:
\[
\ce{C6H5OH(aq) + 3Br2(aq) -> C6H2Br3OH(s) + 3HBr(aq)}
\]
Ethanol leaves the orange colour of bromine unchanged.

\textbf{Step 3 (confirmation): neutral \ce{FeCl3}.} Phenol gives a violet colouration; ethanol gives none. Hence: effervescence with \ce{NaHCO3} $\Rightarrow$ ethanoic acid; white precipitate with bromine water and violet with \ce{FeCl3} $\Rightarrow$ phenol; the remaining liquid, inert to all three reagents, is ethanol.
\end{exoresolu}

\begin{exercice}[Consolidation]
\begin{enumerate}
\item Explain, with the aid of resonance structures, why phenol reacts with bromine water without a catalyst whereas benzene requires \ce{FeBr3}.
\item Write equations, with conditions, for the nitration of phenol using (a) dilute and (b) concentrated nitric acid, naming all organic products.
\item A colourless crystalline solid dissolves in \ce{NaOH} solution but does not liberate \ce{CO2} from \ce{NaHCO3} solution, and gives a violet colour with neutral \ce{FeCl3}. Identify the functional class of the compound and justify each deduction.
\item Outline the industrial preparation of phenol from cumene, naming the co-product.
\end{enumerate}
\end{exercice}

\begin{corrige}[Exercise 1]
\textbf{1.} The lone pair on the oxygen of \ce{-OH} delocalises into the ring, placing negative charge at the 2-, 4- and 6-positions. The ring is therefore electron-rich and attacks the electrophile \ce{Br2} directly; benzene has no such activation and needs \ce{FeBr3} to polarise \ce{Br2} into a stronger electrophile.

\textbf{2.} (a) \ce{C6H5OH + HNO3(dil) -> 2-/4-NO2-C6H4OH + H2O} at room temperature; (b) \ensuremath{\mathrm{C_{6}H_{5}OH + 3HNO_{3}(conc) \rightarrow [conc\ H_{2}SO_{4}] C_{6}H_{2}(NO_{2})_{3}OH}} — 2,4,6-trinitrophenol (picric acid) — \ce{+ 3H2O}.

\textbf{3.} The solid is a \emph{phenol}: it is acidic enough to react with \ce{NaOH} but too weak to decompose \ce{NaHCO3}, and the violet \ce{FeCl3} colouration is characteristic of an \ce{-OH} group bonded to a benzene ring.

\textbf{4.} Cumene (isopropylbenzene) is oxidised by air to cumene hydroperoxide, which on acid cleavage yields phenol and acetone (the co-product).
\end{corrige}

\begin{attention}[Frequent errors at the GCE]
\begin{itemize}
\item Writing that phenol gives effervescence with \ce{NaHCO3} — it does \emph{not}; only carboxylic acids do.
\item Forgetting the coefficient 3 in the bromine water equation, or naming the product ``tribromobenzene'' instead of 2,4,6-tribromophenol.
\item Claiming phenol is easily oxidised like a primary alcohol; the aromatic ring protects it.
\item Saying the \ce{-OH} group directs \emph{meta}: it directs \emph{ortho/para} (2, 4, 6).
\item Confusing the conditions for nitration: dilute \ce{HNO3} at room temperature gives mono-nitrophenols; concentrated \ce{HNO3}/\ce{H2SO4} gives trinitrophenol.
\end{itemize}
\end{attention}

With the chemistry of the hydroxyl group in alcohols and phenols now fully compared, we move on in the next section to compounds in which the oxygen of the carbonyl group, \ce{C=O}, dominates the reactivity: the aldehydes and ketones.

\coursec{Aldehydes and Ketones: Structure, Properties and Preparation}

In the previous section we compared the behaviour of the hydroxyl group in alcohols and phenols, and we saw how strongly the carbon skeleton modifies the reactivity of \ce{-OH}. We now leave the \ce{O-H} bond and turn to a functional group in which oxygen is joined to carbon by a \emph{double} bond: the \cle{carbonyl group}, \ce{C=O}. This single group gives rise to two of the most important families in organic chemistry — the \cle{aldehydes} and the \cle{ketones} — and it controls the chemistry of the carboxylic acids and their derivatives that we shall meet later in this chapter.

\courssub{The Carbonyl Group: Structure and Polarity}

\begin{experience}[A first observation]
Warm a spiral of copper wire in a Bunsen flame until it is black (coated with copper(II) oxide), then dip it into ethanol held in a test tube. The black coating disappears, the copper shines again, and the sharp, fruity smell of \emph{ethanal} replaces the smell of ethanol. The alcohol has lost hydrogen and gained a \ce{C=O} double bond. Something new — a carbonyl compound — has been born.
\end{experience}

\begin{definition}[Aldehydes and ketones]
The \cle{carbonyl group} is a carbon atom joined to an oxygen atom by a double bond, \ce{C=O}.
\begin{itemize}
\item An \cle{aldehyde} has the carbonyl group at the \emph{end} of a carbon chain: the carbonyl carbon carries at least one hydrogen atom. General structure \ce{R-CHO}, where R may be H (methanal) or an alkyl/aryl group.
\item A \cle{ketone} has the carbonyl group \emph{within} the chain: the carbonyl carbon is bonded to two carbon groups. General structure \ensuremath{\mathrm{R-CO-R^{\prime}}}, where R and R' are alkyl or aryl groups.
\end{itemize}
Both families share the same general molecular formula \ensuremath{\mathrm{C_nH_{2n}O}} for saturated, acyclic members.
\end{definition}

\textbf{Electronic structure of \ce{C=O}.} The carbonyl carbon is \cle{$sp^2$ hybridised}: it forms three $\sigma$ bonds in a plane (to oxygen and to its two other substituents) with bond angles close to $120^\circ$, and the remaining p orbitals on carbon and oxygen overlap sideways to give the $\pi$ bond. Oxygen is far more electronegative than carbon (3.44 vs 2.55 on the Pauling scale), so the $\pi$ electrons are pulled strongly towards oxygen. The bond is therefore strongly \cle{polar}. This polarity is reinforced by resonance: the canonical form \ce{C^+-O^-} makes a significant contribution.

\begin{center}
\begin{tikzpicture}[scale=1.1]
% trigonal planar carbonyl carbon
\coordinate (C) at (0,0);
\coordinate (O) at (1.5,0);
\coordinate (R1) at (-1.1,0.95);
\coordinate (R2) at (-1.1,-0.95);
\draw[double, double distance=2.5pt, thick, popblue] (C) -- (O);
\draw[thick] (C) -- (R1);
\draw[thick] (C) -- (R2);
\node[popink, font=\bfseries] at (C) {C};
\node[popink, font=\bfseries] at (O) {O};
\node[popmuted] at (R1) {R};
\node[popmuted] at (R2) {R'};
% partial charges
\node[poporange, font=\small\bfseries] at (0.15,0.42) {$\delta+$};
\node[popblue, font=\small\bfseries] at (1.5,0.42) {$\delta-$};
% dipole arrow
\draw[->, very thick, poporange] (0.25,-0.55) -- (1.35,-0.55);
\node[poporange, font=\small] at (0.8,-0.85) {dipole $\mu$};
% lone pairs on oxygen
\fill[popblue] (1.75,0.18) circle (1.2pt);
\fill[popblue] (1.75,-0.18) circle (1.2pt);
% angle marker
\draw[popmuted] (-0.55,0.48) arc[start angle=139, end angle=221, radius=0.72];
\node[popmuted, font=\small] at (-1.35,0) {$\approx 120^\circ$};
\end{tikzpicture}
\legende{The polar carbonyl group: trigonal planar ($sp^2$) carbon, $\delta+$ on carbon and $\delta-$ on oxygen.}
\end{center}

\begin{aretenir}[Consequences of polarity]
The $\delta+$ carbonyl carbon is \emph{electron-poor} and is therefore attacked by \cle{nucleophiles}; the $\delta-$ oxygen is electron-rich and attracts \emph{electrophiles} (and protons). Nearly all the reactions of aldehydes and ketones — nucleophilic addition, reduction, condensation — follow from this simple picture.
\end{aretenir}

\courssub{Nomenclature of Aldehydes and Ketones}

\begin{methode}[Naming aldehydes and ketones (IUPAC)]
\begin{enumerate}
\item Identify the \textbf{longest continuous carbon chain containing the carbonyl carbon}; replace the final \emph{-e} of the alkane name by \textbf{-al} (aldehyde) or \textbf{-one} (ketone).
\item \textbf{Aldehydes:} the carbonyl carbon is always C-1, so no locant is needed for the \ce{CHO} group. Number the chain from the \ce{CHO} end.
\item \textbf{Ketones:} number the chain to give the carbonyl carbon the \textbf{lowest possible locant}, and state it: e.g. pentan-2-one.
\item Name and locate substituents as usual, in alphabetical order. Prefixes: \emph{hydroxy}, \emph{chloro}, \emph{nitro}, etc.
\item For cyclic ketones, the carbonyl carbon is C-1; the ring name takes the suffix \emph{-one} (e.g. cyclohexanone).
\end{enumerate}
\end{methode}

\begin{exemplebox}[Worked examples of naming]
\begin{itemize}
\item \ce{CH3-CH2-CH2-CO-CH3}: longest chain of 5 carbons containing \ce{C=O}; numbering from the right gives the carbonyl carbon locant 2 $\Rightarrow$ \textbf{pentan-2-one}.
\item \ce{CH3-CH(CH3)-CH2-CHO}: aldehyde carbon is C-1; a methyl group on C-3 $\Rightarrow$ \textbf{3-methylbutanal}.
\item \ce{CH3CH2CHO} is \textbf{propanal}; \ce{CH3COCH3} is \textbf{propanone}.
\item \ce{CH3CH2CHClCOCH3}: 5 carbons, carbonyl at C-2, chloro at C-4 $\Rightarrow$ \textbf{4-chloropentan-2-one}.
\item \ce{HOCH2CH2CHO}: 3 carbons, aldehyde at C-1, hydroxy at C-3 $\Rightarrow$ \textbf{3-hydroxypropanal}.
\end{itemize}
\end{exemplebox}

\begin{savaistu}[Common names you will still meet]
Methanal (\ce{HCHO}) is universally called \emph{formaldehyde} — its aqueous solution, formalin, is used to preserve biological specimens in school laboratories. Propanone (\ce{CH3COCH3}) is \emph{acetone}, the familiar solvent in nail-varnish remover. Ethanal (\ce{CH3CHO}) is \emph{acetaldehyde}. In examinations, IUPAC names are required, but you must recognise the common ones.
\end{savaistu}

\courssub{Isomerism in Aldehydes and Ketones}

Because aldehydes and ketones share the formula \ensuremath{\mathrm{C_nH_{2n}O}} (for saturated, acyclic compounds), they show rich isomerism:

\begin{itemize}
\item \textbf{Chain isomerism:} different carbon skeletons. Example: \ce{CH3CH2CH2CHO} (butanal) and \ce{(CH3)2CHCHO} (2-methylpropanal), both \ce{C4H8O}.
\item \textbf{Position isomerism (ketones only):} the carbonyl group occupies different positions on the same carbon skeleton. Example: \ce{CH3COCH2CH2CH3} (pentan-2-one) and \ce{CH3CH2COCH2CH3} (pentan-3-one), both \ce{C5H10O}. Aldehydes \emph{cannot} show position isomerism because \ce{CHO} must be at C-1.
\item \textbf{Functional group isomerism:} an aldehyde and a ketone with the same molecular formula. Example: propanal (\ce{CH3CH2CHO}) and propanone (\ce{CH3COCH3}), both \ce{C3H6O}.
\item \textbf{Metamerism (a type of structural isomerism):} ketones with different alkyl groups on either side of the carbonyl, e.g. \ce{CH3COCH2CH3} (butanone) and \ce{CH3CH2COCH3} (same compound, but for higher homologues like \ce{C5H10O}: pentan-2-one vs 3-methylbutan-2-one).
\end{itemize}

\begin{attention}[Aldehydes and ketones of the same formula are functional isomers]
Propanal, \ce{CH3CH2CHO}, and propanone, \ce{CH3COCH3}, both have the molecular formula \ce{C3H6O}, yet they belong to \emph{different families} with different chemistry. Do not write them as the same compound, and remember that from \ce{C3H6O} upwards, every aldehyde has at least one ketone isomer. A frequent exam error is to claim that \ce{C3H6O} shows position isomerism — it does not; it shows functional group isomerism. Another error: forgetting that ketones start at \ce{C3} (propanone), so \ce{C2H4O} can only be ethanal.
\end{attention}

\courssub{Physical Properties}

\textbf{Boiling points.} Carbonyl compounds are polar molecules (dipole moments typically 2.3–2.7 D), so \cle{dipole–dipole attractions} exist between them. However, the oxygen atom is bonded only to carbon, never to hydrogen, so a pure aldehyde or ketone \emph{cannot hydrogen-bond to itself}. Consequently, for comparable molar mass:

\[ T_b(\text{alcohol}) > T_b(\text{aldehyde/ketone}) > T_b(\text{alkane}) \]

\begin{center}
\begin{tikzpicture}
\begin{axis}[
    width=0.85\linewidth, height=6.2cm,
    xlabel={Number of carbon atoms}, ylabel={Boiling point / $^\circ$C},
    xmin=1.6, xmax=4.4, ymin=-60, ymax=130,
    xtick={2,3,4},
    legend style={at={(0.02,0.98)}, anchor=north west, font=\small},
    grid=both, grid style={popline!40},
]
\addplot[color=popmuted, mark=*, thick] coordinates {(2,-88.6) (3,-42.1) (4,-0.5)};
\addlegendentry{Alkanes}
\addplot[color=poporange, mark=square*, thick] coordinates {(2,20.2) (3,48.8) (4,74.8)};
\addlegendentry{Aldehydes (\ce{RCHO})}
\addplot[color=popblue, mark=triangle*, thick] coordinates {(2,78.4) (3,97.2) (4,117.7)};
\addlegendentry{Alcohols (\ce{ROH})}
\end{axis}
\end{tikzpicture}
\legende{Boiling points for two-, three- and four-carbon compounds. The polar carbonyl compounds boil well above the alkanes (dipole–dipole forces) but below the alcohols, which hydrogen-bond to themselves.}
\end{center}

\textbf{Melting points and physical state.} Methanal is a gas at room temperature (b.p. $-19^\circ$C); ethanal is a volatile liquid (b.p. $20.2^\circ$C); propanone and higher members are liquids. All have characteristic, often pungent or fruity odours.

\textbf{Solubility.} Although they cannot hydrogen-bond to each other, the carbonyl oxygen \emph{can accept hydrogen bonds from water}. Hence the small members — methanal, ethanal, propanone, propanal, butanone — are completely miscible with water. Solubility falls rapidly as the hydrocarbon chain lengthens; by about five carbons (pentanal, pentan-2-one) the compounds are only sparingly soluble ($< 10\ \mathrm{g/100\ mL}$).

\textbf{Density.} All common aldehydes and ketones are less dense than water (typical $0.78–0.85\ \mathrm{g/cm^3}$), so they form a separate upper layer when mixed with water in a separating funnel.

\begin{aretenir}[Physical properties summary]
Volatile, often pleasant- or sharp-smelling liquids (methanal is a gas); boiling points between those of alkanes and alcohols of similar $M_r$; small members water-soluble through hydrogen bonding \emph{to water}, not to themselves; less dense than water.
\end{aretenir}

\courssub{Preparation of Aldehydes and Ketones}

The syllabus requires mastery of the following laboratory and industrial routes.

\textbf{1. Oxidation of alcohols (the standard laboratory route).} Acidified potassium dichromate(VI), \ce{K2Cr2O7/H2SO4}, oxidises primary alcohols to aldehydes and secondary alcohols to ketones:

\begin{align*}
\ce{R-CH2OH + [O] &-> R-CHO + H2O} \quad \text{(primary alcohol $\to$ aldehyde)}\\
\ensuremath{\mathrm{R-CH(OH)-R^\prime + [O] }}&\ensuremath{\mathrm{\rightarrow  R-CO-R^\prime + H_{2}O}} \quad \text{(secondary alcohol $\to$ ketone)}
\end{align*}

The orange dichromate(VI) ion (\ce{Cr2O7^2-}) is reduced to green \ce{Cr^3+}, giving a visible colour change. The reaction is carried out by adding the alcohol slowly to the acidified oxidant.

\begin{attention}[Exam tip — distil the aldehyde off immediately]
An aldehyde is oxidised \emph{further} to a carboxylic acid if it stays in contact with the oxidising agent:
\[ \ce{R-CHO + [O] -> R-COOH} \]
To stop at the aldehyde, use \emph{limited} oxidising agent and \textbf{distil the aldehyde off as it forms} — possible because the aldehyde (e.g. ethanal, b.p. $20.2\ ^\circ$C) boils well below the alcohol (ethanol, $78.4\ ^\circ$C). For a ketone, simply \emph{reflux}: ketones resist further oxidation under these conditions, so no special precaution is needed. Confusing "distil" (aldehyde) with "reflux" (ketone, or full oxidation to acid) is one of the most common errors at GCE.
\end{attention}

\textbf{2. Catalytic dehydrogenation of alcohols (industrial).} Passing alcohol vapour over hot copper (about $300\ ^\circ$C) removes hydrogen:

\[ \ensuremath{\mathrm{CH_{3}CH_{2}OH \rightarrow [Cu,\ 300\ ^\circ C] CH_{3}CHO + H_{2}}} \qquad \ensuremath{\mathrm{CH_{3}CH(OH)CH_{3} \rightarrow [Cu,\ 300\ ^\circ C] CH_{3}COCH_{3} + H_{2}}} \]

This is the reverse of the catalytic hydrogenation (reduction) of carbonyl compounds, and is how much industrial ethanal and propanone are produced.

\textbf{3. Hydration of alkynes.} Ethyne adds water in the presence of dilute \ce{H2SO4} and mercury(II) sulphate catalyst (\ce{HgSO4}). The unstable enol formed first tautomerises to ethanal — the only aldehyde accessible this way, since higher alkynes give ketones (Markovnikov addition):

\[ \ce{HC#CH + H2O ->[HgSO4/H2SO4] [CH2=CH-OH] -> CH3CHO} \]

For a terminal alkyne \ce{R-C#CH}, Markovnikov hydration gives a methyl ketone \ce{R-CO-CH3}.

\textbf{4. Ozonolysis of alkenes.} Ozone cleaves the \ce{C=C} bond, and reductive work-up (e.g. \ce{Zn/H2O} or \ce{(CH3)2S}) converts each alkene carbon into a carbonyl carbon. A tetrasubstituted carbon gives a ketone; a carbon bearing hydrogen gives an aldehyde:

\[ \ensuremath{\mathrm{(CH_{3})_{2}C=CHCH_{3} \rightarrow [1.\ O_{3}][2.\ Zn/H_{2}O] (CH_{3})_{2}C=O + CH_{3}CHO}} \]

Ozonolysis is valuable both as a preparative route and as a means of locating double bonds in unknown alkenes.

\textbf{5. Friedel–Crafts acylation (aromatic ketones).} Benzene reacts with an acyl chloride in the presence of anhydrous \ce{AlCl3} to give an aromatic ketone — the route to compounds such as phenylethanone:

\[ \ce{C6H6 + CH3COCl ->[AlCl3] C6H5COCH3 + HCl} \]

The acyl cation \ce{CH3CO+} is the electrophile. This reaction does not work for aldehydes (formyl chloride is unstable).

\begin{exoresolu}[Identifying and preparing carbonyl compounds]
\textbf{Statement.} Compound A, \ce{C4H10O}, is oxidised by acidified \ce{K2Cr2O7} to compound B, \ce{C4H8O}. B does not react with Fehling's solution. (a) Identify the family of B. (b) Give the possible structures of A and B. (c) Describe how you would carry out the oxidation in the laboratory.
\tcblower
\textbf{Solution.} (a) Loss of \ce{H2} on oxidation shows A is an alcohol. Since B resists further oxidation (no reaction with Fehling's), B is a \textbf{ketone} and A is a \textbf{secondary alcohol}.
(b) The only secondary alcohol of formula \ce{C4H10O} is \textbf{butan-2-ol}, \ce{CH3CH(OH)CH2CH3}; B is therefore \textbf{butanone}, \ce{CH3COCH2CH3}.
(c) Heat butan-2-ol under \emph{reflux} with acidified potassium dichromate(VI); the mixture turns from orange to green. No immediate distillation is required because the ketone cannot be oxidised further under these conditions.
\[ \ce{CH3CH(OH)CH2CH3 + [O] -> CH3COCH2CH3 + H2O} \]
\end{exoresolu}

\begin{exercice}[Check your understanding]
\begin{enumerate}
\item Name the following: (a) \ce{CH3CH2COCH2CH3}; (b) \ce{(CH3)2CHCH2CHO}; (c) \ce{CH3CH2CH2CHO}; (d) \ce{CH3COCH2CH2Cl}.
\item Draw and classify all carbonyl isomers of \ce{C4H8O}, stating the type of isomerism between each pair.
\item Explain why propanone (b.p. $56\ ^\circ$C) boils much higher than butane (b.p. $-0.5\ ^\circ$C) but lower than propan-1-ol (b.p. $97\ ^\circ$C), although all three have similar $M_r$.
\item Outline, with equations and conditions, two different preparations of ethanal and one preparation of propanone.
\item An unknown compound \ce{C5H10O} gives a positive 2,4-DNP test but does not react with Tollens' reagent. It gives a yellow precipitate with \ce{I2/NaOH}. Deduce its structure and name it.
\end{enumerate}
\end{exercice}

\begin{corrige}[Exercise 1]
\textbf{1.} (a) pentan-3-one; (b) 3-methylbutanal; (c) butanal; (d) 4-chloropentan-2-one.

\textbf{2.} \ce{CH3CH2CH2CHO} butanal and \ce{(CH3)2CHCHO} 2-methylpropanal (chain isomers of each other); \ce{CH3COCH2CH3} butanone. Butanone is a \emph{functional group isomer} of both aldehydes. The two aldehydes are chain isomers.

\textbf{3.} Propanone is polar (\ce{C=O}), so dipole–dipole attractions hold its molecules together, unlike non-polar butane (only van der Waals forces) — hence a much higher boiling point. Propan-1-ol forms intermolecular \emph{hydrogen bonds} through its \ce{O-H} group, which are stronger than dipole–dipole forces, so it boils highest.

\textbf{4.} Ethanal: (i) oxidise ethanol with limited acidified \ce{K2Cr2O7}, distilling off the ethanal as it forms: \ce{CH3CH2OH + [O] -> CH3CHO + H2O}; (ii) hydrate ethyne with \ce{HgSO4/H2SO4}: \ce{HC#CH + H2O -> CH3CHO}. Propanone: oxidise propan-2-ol under reflux with acidified \ce{K2Cr2O7}: \ce{CH3CH(OH)CH3 + [O] -> CH3COCH3 + H2O}.

\textbf{5.} Positive 2,4-DNP $\Rightarrow$ carbonyl compound. Negative Tollens' $\Rightarrow$ not an aldehyde $\Rightarrow$ ketone. Positive iodoform test (\ce{I2/NaOH}) $\Rightarrow$ methyl ketone (\ce{CH3CO-}). Formula \ce{C5H10O} $\Rightarrow$ \ce{CH3COC3H7}. Structure: \ce{CH3COCH2CH2CH3}, \textbf{pentan-2-one}.
\end{corrige}

\begin{aretenir}[Key points of this section]
\begin{itemize}
\item Aldehydes (\ce{RCHO}) and ketones (\ensuremath{\mathrm{RCOR^{\prime}}}) contain the polar carbonyl group \ce{C=O}; the carbonyl carbon is $sp^2$ hybridised and $\delta+$, hence attacked by nucleophiles.
\item Shared formula \ensuremath{\mathrm{C_nH_{2n}O}} $\Rightarrow$ chain, position (ketones) and functional group isomerism.
\item Boiling points lie between alkanes and alcohols; small members water-soluble; less dense than water.
\item Prepare by oxidation of primary alcohols (\emph{distil off the aldehyde!}) or secondary alcohols (reflux); also by dehydrogenation of alcohols, hydration of alkynes, ozonolysis, and Friedel–Crafts acylation for aromatic ketones.
\end{itemize}
\end{aretenir}

\coursec{Reactions of Aldehydes and Ketones}

In the previous section we saw that aldehydes and ketones both contain the \cle{carbonyl group} $\ce{C=O}$, and that this group is strongly \cle{polar}: because oxygen is more electronegative than carbon, the bonding electrons are pulled towards oxygen, leaving the carbonyl carbon electron-poor ($\delta^{+}$) and the oxygen electron-rich ($\delta^{-}$). Almost everything in this section follows from that single fact. A reagent carrying a lone pair or a negative charge — a \cle{nucleophile} — is attracted to the $\delta^{+}$ carbon and attacks it. This is why the characteristic reaction of aldehydes and ketones is \cle{nucleophilic addition}, in sharp contrast to alkenes, whose electron-rich $\ce{C=C}$ bond undergoes \emph{electrophilic} addition.

\courssub{Nucleophilic Addition: the General Mechanism}

\begin{definition}[Nucleophilic addition on a carbonyl]
A nucleophilic addition is a reaction in which a nucleophile $\ce{Nu^{-}}$ (or $\ce{Nu{:}}$) attacks the electron-deficient carbonyl carbon, the $\pi$ electrons of $\ce{C=O}$ move onto oxygen, and the resulting alkoxide ion is then protonated. Overall a molecule $\ce{H-Nu}$ adds across the double bond:
\[ \ensuremath{\mathrm{R-CO-R^\prime + H-Nu \rightarrow  R-C(OH)(Nu)-R^\prime}} \]
\end{definition}

The mechanism has two steps, and you must be able to draw the curly arrows:

\begin{enumerate}
\item \textbf{Step 1 (slow):} the nucleophile donates its lone pair to the carbonyl carbon; the $\ce{C=O}$ $\pi$-bond breaks and its electron pair moves onto oxygen, giving a negatively charged \cle{alkoxide intermediate}.
\item \textbf{Step 2 (fast):} the alkoxide oxygen abstracts a proton (from $\ce{H-Nu}$, $\ce{H2O}$ or $\ce{H3O+}$) to give the neutral addition product.
\end{enumerate}

\begin{attention}[Aldehydes react faster than ketones]
Ketones are generally \emph{less reactive} towards nucleophilic addition than aldehydes, for two reasons: (i) the two alkyl groups push electrons onto the carbonyl carbon (+I effect), making it less $\delta^{+}$; (ii) the two bulky alkyl groups \emph{sterically hinder} the approach of the nucleophile. Thus ethanal reacts more readily than propanone with HCN.
\end{attention}

\courssub{Addition of Hydrogen Cyanide: Formation of Hydroxynitriles}

\begin{experience}[Observation]
When ethanal is mixed with potassium cyanide solution and a little dilute sulphuric acid (which generates HCN \emph{in situ}), the two substances combine to give a single product, \cle{2-hydroxypropanenitrile}, $\ce{CH3-CH(OH)-CN}$. The molecule has grown: a new C--C bond has been created.
\end{experience}

Hydrogen cyanide adds across the carbonyl group to give a \cle{hydroxynitrile} (historically called a \emph{cyanohydrin}):

\[ \ce{CH3CHO + HCN ->[KCN/H^{+}] CH3-CH(OH)-CN} \]

Because HCN is a very weak acid, the concentration of the true nucleophile $\ce{CN^{-}}$ in pure HCN is tiny; the $\ce{KCN/H^{+}}$ mixture supplies enough $\ce{CN^{-}}$ for step 1 and enough $\ce{H^{+}}$ for step 2.

\begin{propriete}[Mechanism of cyanohydrin formation — examinable]
Step 1: $\ce{CN^{-}}$ attacks the carbonyl carbon; the $\pi$ electrons of $\ce{C=O}$ move onto oxygen, giving an alkoxide. Step 2: the alkoxide is protonated by HCN (regenerating $\ce{CN^{-}}$, which therefore acts catalytically). The curly arrows are shown in the figure below. \emph{This mechanism is examinable at GCE Advanced Level.}
\end{propriete}

\begin{popfigure}
\begin{tikzpicture}[scale=1.0, every node/.style={font=\small}]
% Ethanal
\node at (0,2.6) {\textbf{Step 1}};
\draw[thick] (0,1.6) -- (0.8,2.1);
\draw[thick] (0.8,2.1) -- (1.6,1.6);
\draw[thick] (0.8,2.1) -- (0.8,3.0);
\draw[thick] (0.9,2.1) -- (0.9,3.0);
\node at (0.85,3.3) {O};
\node at (-0.35,1.6) {\ce{CH3}};
\node at (2.0,1.55) {H};
\node at (0.85,2.35) [right=2pt] {\scriptsize $\delta^{+}$};
\node at (1.15,3.25) {\scriptsize $\delta^{-}$};
% CN-
\node at (0.8,0.6) {$\ce{^{-}CN}$};
% curly arrow CN -> C
\draw[->, thick, poporange] (0.8,0.95) .. controls (0.2,1.3) .. (0.65,1.95);
% curly arrow pi -> O
\draw[->, thick, poporange] (0.95,2.75) .. controls (1.35,2.9) .. (1.15,3.2);
% arrow to intermediate
\draw[->, very thick, popblue] (2.6,2.1) -- (3.6,2.1);
% alkoxide intermediate
\draw[thick] (4.2,1.6) -- (5.0,2.1);
\draw[thick] (5.0,2.1) -- (5.8,1.6);
\draw[thick] (5.0,2.1) -- (5.0,3.0);
\node at (5.0,3.3) {$\ce{O^{-}}$};
\node at (3.85,1.6) {\ce{CH3}};
\node at (6.2,1.55) {H};
\draw[thick] (5.0,2.1) -- (5.0,1.2);
\node at (5.0,0.95) {CN};
% Step 2
\node at (7.6,2.6) {\textbf{Step 2}};
\node at (6.9,3.3) {$\ce{H-CN}$};
\draw[->, thick, poporange] (5.25,3.25) .. controls (6.0,3.7) .. (6.75,3.45);
\draw[->, very thick, popblue] (7.6,2.1) -- (8.6,2.1);
% product
\draw[thick] (9.2,1.6) -- (10.0,2.1);
\draw[thick] (10.0,2.1) -- (10.8,1.6);
\draw[thick] (10.0,2.1) -- (10.0,3.0);
\node at (10.0,3.3) {OH};
\node at (8.85,1.6) {\ce{CH3}};
\node at (11.2,1.55) {H};
\draw[thick] (10.0,2.1) -- (10.0,1.2);
\node at (10.0,0.95) {CN};
\node at (10.0,0.35) {\scriptsize 2-hydroxypropanenitrile};
\end{tikzpicture}
\legende{Curly-arrow mechanism for the addition of $\ce{CN^{-}}$ to ethanal, followed by protonation of the alkoxide by HCN.}
\end{popfigure}

\begin{savaistu}[Why cyanohydrins matter]
Hydroxynitriles are valuable intermediates because the $\ce{-CN}$ group can be hydrolysed to $\ce{-COOH}$. Addition of HCN to ethanal, followed by acid hydrolysis, gives \emph{lactic acid} ($\ce{CH3-CH(OH)-COOH}$), the acid formed in sour milk and in our muscles during hard exercise. Notice also that HCN addition \emph{lengthens the carbon chain by one carbon} — a favourite exam trick in synthesis questions.
\end{savaistu}

\courssub{Reduction of Aldehydes and Ketones}

The carbonyl group can be reduced to an alcohol. Two families of reducing agents are used:

\begin{itemize}
\item \cle{Complex metal hydrides}: lithium tetrahydridoaluminate(III), $\ce{LiAlH4}$ (in dry ether), or sodium tetrahydridoborate(III), $\ce{NaBH4}$ (in aqueous or alcoholic solution). Both deliver the hydride ion $\ce{H^{-}}$, a powerful nucleophile, to the carbonyl carbon.
\item \cle{Catalytic hydrogenation}: $\ce{H2}$ over a nickel catalyst, on heating.
\end{itemize}

\begin{propriete}[Products of reduction]
Aldehydes are reduced to \textbf{primary alcohols}; ketones are reduced to \textbf{secondary alcohols}:
\[ \ce{CH3CHO + 2[H] -> CH3CH2OH} \qquad \ce{CH3COCH3 + 2[H] -> CH3CH(OH)CH3} \]
\end{propriete}

The notation $\ce{[H]}$ represents hydrogen supplied by the reducing agent. Mechanistically, with $\ce{NaBH4}$, $\ce{H^{-}}$ attacks the carbonyl carbon exactly as $\ce{CN^{-}}$ did, and the alkoxide formed is then protonated by the solvent.

\begin{exemplebox}[Choosing the product]
Reduction of butanone, $\ce{CH3COCH2CH3}$, with $\ce{LiAlH4}$ gives butan-2-ol, $\ce{CH3CH(OH)CH2CH3}$ — a secondary alcohol. Reduction of propanal gives propan-1-ol — a primary alcohol. In exams, always state the \emph{class} of alcohol formed; it is a quick way to earn marks.
\end{exemplebox}

\courssub{Condensation with 2,4-Dinitrophenylhydrazine: a Test for the Carbonyl Group}

Aldehydes and ketones react with \cle{2,4-dinitrophenylhydrazine} (2,4-DNPH, Brady's reagent) in a \cle{condensation} reaction: the two molecules join with the \emph{elimination of a small molecule}, water.

\[ \ce{R2C=O + H2N-NH-C6H3(NO2)2 -> R2C=N-NH-C6H3(NO2)2 + H2O} \]

The product, a \cle{2,4-dinitrophenylhydrazone}, is an \textbf{orange (or yellow-orange) crystalline precipitate}.

\begin{propriete}[The 2,4-DNPH test]
Adding a few drops of a carbonyl compound to Brady's reagent produces an \textbf{orange precipitate} if an aldehyde or a ketone is present. The test detects the $\ce{C=O}$ group but does \emph{not} by itself distinguish an aldehyde from a ketone. (The sharp melting point of the purified derivative can, however, be used to identify a specific carbonyl compound.)
\end{propriete}

\begin{attention}[Condensation is not addition]
Students often write the DNPH reaction as a simple addition. It is not: water is eliminated, so the product contains $\ce{C=N}$, not $\ce{C(OH)(NHNHAr)}$. Always balance the equation with $\ce{+ H2O}$ on the right.
\end{attention}

\courssub{Oxidation: the Key Difference Between Aldehydes and Ketones}

Here lies the most examined difference between the two families. The aldehyde group $\ce{-CHO}$ carries a hydrogen atom directly attached to the carbonyl carbon; this C--H bond is easily oxidised to $\ce{-OH}$, converting the aldehyde into a \cle{carboxylic acid}:

\[ \ce{RCHO + [O] -> RCOOH} \]

A ketone has \emph{no} hydrogen on the carbonyl carbon, so it \textbf{resists mild oxidising agents}. (Very strong oxidation breaks C--C bonds, giving a mixture of smaller acids — not required in detail here.)

\textbf{Oxidising agents and observations:}

\begin{itemize}
\item \textbf{Acidified potassium dichromate(VI)}, $\ce{K2Cr2O7/H2SO4}$, on warming: an aldehyde turns the solution from \textbf{orange to green} (formation of $\ce{Cr^{3+}}$) as it is oxidised to the acid. A ketone gives \emph{no change}.
\item \textbf{Tollens' reagent} (ammoniacal silver nitrate), warmed gently: an aldehyde reduces $\ce{Ag^{+}}$ to metallic silver, producing a \textbf{silver mirror} on the wall of the test tube (or a grey-black precipitate). A ketone gives \emph{no change}.
\item \textbf{Fehling's solution} (blue, containing $\ce{Cu^{2+}}$ in alkaline tartrate solution), warmed: an aldehyde reduces $\ce{Cu^{2+}}$ to $\ce{Cu2O}$, giving a \textbf{brick-red precipitate}. A ketone gives \emph{no change}.
\end{itemize}

\begin{propriete}[Tollens' and Fehling's tests]
Tollens' reagent gives a \textbf{silver mirror} with aldehydes; Fehling's solution gives a \textbf{brick-red precipitate of $\ce{Cu2O}$} with aldehydes. Ketones give a negative result with both reagents. These are the standard distinguishing tests at GCE Advanced Level.
\end{propriete}

\begin{exemplebox}[Half-equations for the dichromate oxidation]
When balancing oxidation equations with acidified dichromate, write the two half-equations:
\begin{align*}
\text{Reduction:}\quad & \ce{Cr2O7^{2-} + 14H+ + 6e- -> 2Cr^{3+} + 7H2O}\\
\text{Oxidation:}\quad & \ce{CH3CHO + H2O -> CH3COOH + 2H+ + 2e-}
\end{align*}
Multiplying the second equation by 3 and adding:
\[ \ce{3CH3CHO + Cr2O7^{2-} + 8H+ -> 3CH3COOH + 2Cr^{3+} + 4H2O} \]
\end{exemplebox}

\begin{methode}[Test sequence: identify an unknown carbonyl compound]
\begin{enumerate}
\item Add Brady's reagent (2,4-DNPH). \textbf{Orange precipitate} $\Rightarrow$ the compound contains $\ce{C=O}$ (aldehyde or ketone). No precipitate $\Rightarrow$ it is neither.
\item Warm a fresh sample with \textbf{Tollens' reagent} (or Fehling's solution).
\item \textbf{Silver mirror} (or brick-red precipitate) $\Rightarrow$ the compound is an \textbf{aldehyde}.
\item \textbf{No change} $\Rightarrow$ the compound is a \textbf{ketone}.
\end{enumerate}
\end{methode}

\begin{popfigure}
\begin{tikzpicture}[
  grow=down,
  every node/.style={font=\small, align=center},
  box/.style={draw=popblue, thick, rounded corners, fill=popblueL, text width=3.4cm, inner sep=4pt},
  res/.style={draw=popgreen, thick, rounded corners, fill=popgreenL, text width=3.0cm, inner sep=4pt},
  lbl/.style={font=\scriptsize\itshape, fill=white, inner sep=1pt},
  level distance=1.7cm, sibling distance=5.2cm]
\node[box] {Unknown organic liquid};
\node[box, below=1.2cm] (dnph) {Add 2,4-DNPH (Brady's reagent)};
\draw[->, thick, popdark] (0,-0.55) -- (dnph);
\node[res, below left=1.2cm and -1.4cm of dnph, align=center] (no){No precipitate \\ $\Rightarrow$ \textbf{not} a carbonyl};
\node[box, below right=1.2cm and -1.4cm of dnph] (toll) {Orange precipitate: $\ce{C=O}$ present. Warm with Tollens'};
\draw[->, thick, popdark] (dnph) -- (no);
\draw[->, thick, popdark] (dnph) -- (toll);
\node[res, below left=1.2cm and -1.2cm of toll, align=center] (ald){Silver mirror \\ $\Rightarrow$ \textbf{aldehyde}};
\node[res, below right=1.2cm and -1.2cm of toll, align=center] (ket){No change \\ $\Rightarrow$ \textbf{ketone}};
\draw[->, thick, popdark] (toll) -- (ald);
\draw[->, thick, popdark] (toll) -- (ket);
\end{tikzpicture}
\legende{Decision tree: 2,4-DNPH confirms the carbonyl group; Tollens' (or Fehling's) then distinguishes an aldehyde from a ketone.}
\end{popfigure}

\courssub{Chlorination of the Alpha-Carbon and the Iodoform Reaction}

The hydrogen atoms on the carbon \emph{adjacent} to the carbonyl group — the \cle{$\alpha$-hydrogens} — are weakly acidic, because the carbanion left behind when one is removed is stabilised by the neighbouring $\ce{C=O}$. Halogens therefore replace $\alpha$-hydrogens easily. With chlorine, an $\alpha$-hydrogen is substituted:

\[ \ce{CH3COCH3 + Cl2 -> CH3COCH2Cl + HCl} \]

The most important application is the \cle{iodoform (tri-iodomethane) reaction}. When a compound containing the group $\ce{CH3-CO-}$ (a \emph{methyl ketone}) — or ethanal, $\ce{CH3CHO}$ — is warmed with iodine and sodium hydroxide, all three $\alpha$-hydrogens are replaced by iodine, and the molecule then splits to give \textbf{iodoform}, $\ce{CHI3}$, a \textbf{yellow crystalline precipitate} with a characteristic antiseptic smell:

\[ \ce{CH3COCH3 + 3I2 + 4NaOH -> CHI3 + CH3COONa + 3NaI + 3H2O} \]

For ethanal:
\[ \ce{CH3CHO + 3I2 + 4NaOH -> CHI3 + HCOONa + 3NaI + 3H2O} \]

\begin{attention}[Who gives a positive iodoform test?]
\emph{Only} \textbf{methyl ketones} ($\ce{CH3CO-R}$) and \textbf{ethanal} (the only aldehyde with a $\ce{CH3-CO-}$ unit) give the yellow $\ce{CHI3}$ precipitate. Propanal, benzaldehyde and pentan-3-one, for example, give a \emph{negative} test. Note also that \emph{ethanol} and \emph{secondary alcohols of the type $\ce{CH3-CH(OH)-R}$} give a positive test too, because $\ce{I2/NaOH}$ first oxidises them to the corresponding carbonyl. Do not claim that every aldehyde or every ketone gives iodoform — that is the classic exam trap.
\end{attention}

\courssub{Uses of Aldehydes and Ketones}

\begin{itemize}
\item \textbf{Methanal (formaldehyde):} its aqueous solution (\emph{formalin}) is used as a \textbf{preservative} for biological specimens and as a disinfectant; it is also a raw material for \textbf{plastics} such as Bakelite and urea-formaldehyde resins.
\item \textbf{Propanone (acetone):} an excellent \textbf{solvent} for varnishes, paints, nail-polish removers and many organic reactions.
\item \textbf{Flavouring aldehydes:} benzaldehyde gives the smell of \textbf{almonds}; cinnamaldehyde flavours \textbf{cinnamon}; vanillin gives \textbf{vanilla} — widely used in Cameroonian bakeries and confectionery.
\end{itemize}

\courssub{Worked Example and Consolidation}

\begin{exoresolu}[Identifying three unknown liquids]
Three colourless liquids A, B and C are propanal, propanone and propan-1-ol, in some order. A gives an orange precipitate with 2,4-DNPH and a silver mirror with Tollens' reagent. B gives an orange precipitate with 2,4-DNPH but no reaction with Tollens'. C gives no precipitate with 2,4-DNPH. Identify A, B and C, and write the equation for the reaction of A with Tollens' reagent.
\tcblower
\textbf{Solution.}
\begin{itemize}
\item A reacts with 2,4-DNPH $\Rightarrow$ carbonyl compound; the silver mirror with Tollens' $\Rightarrow$ \textbf{aldehyde}. Hence A is \textbf{propanal}, $\ce{CH3CH2CHO}$.
\item B reacts with 2,4-DNPH but not with Tollens' $\Rightarrow$ \textbf{ketone}. Hence B is \textbf{propanone}, $\ce{CH3COCH3}$.
\item C does not react with 2,4-DNPH $\Rightarrow$ no $\ce{C=O}$. Hence C is \textbf{propan-1-ol}, $\ce{CH3CH2CH2OH}$.
\end{itemize}
Equation for A with Tollens' reagent:
\[ \ce{CH3CH2CHO + 2[Ag(NH3)2]+ + 3OH- -> CH3CH2COO- + 2Ag + 4NH3 + 2H2O} \]
The silver deposited forms the mirror; propanal is oxidised to the propanoate ion (propanoic acid in acid work-up).
\end{exoresolu}

\begin{exercice}[Reactions of butanone]
Butanone, $\ce{CH3COCH2CH3}$, is treated successively with: (a) $\ce{KCN/H+}$; (b) $\ce{NaBH4}$ in ethanol; (c) $\ce{I2/NaOH}$ on warming; (d) Fehling's solution on warming. For each, state whether a reaction occurs, name the organic product where relevant, and give any observation.
\end{exercice}

\begin{corrige}[Exercise 1]
\begin{itemize}
\item (a) Reaction occurs: nucleophilic addition of HCN gives the hydroxynitrile $\ce{CH3C(OH)(CN)CH2CH3}$ (2-hydroxy-2-methylbutanenitrile).
\item (b) Reduction to the secondary alcohol \textbf{butan-2-ol}, $\ce{CH3CH(OH)CH2CH3}$.
\item (c) Butanone is a \emph{methyl ketone} ($\ce{CH3CO-}$), so the iodoform test is positive: a \textbf{yellow precipitate of $\ce{CHI3}$} forms, together with $\ce{CH3CH2COONa}$.
\item (d) No reaction: ketones resist mild oxidising agents, so the blue Fehling's solution remains unchanged.
\end{itemize}
\end{corrige}

\begin{aretenir}[Key points]
\begin{itemize}
\item The polar $\ce{C=O}$ group undergoes \textbf{nucleophilic addition}: $\ce{Nu^{-}}$ attacks the $\delta^{+}$ carbon, then the alkoxide is protonated.
\item $\ce{HCN}$ ($\ce{KCN/H+}$) gives \textbf{hydroxynitriles}; the chain grows by one carbon.
\item $\ce{NaBH4}$, $\ce{LiAlH4}$ or $\ce{H2/Ni}$ reduce aldehydes to \textbf{primary} and ketones to \textbf{secondary} alcohols.
\item \textbf{2,4-DNPH}: orange precipitate $\Rightarrow$ carbonyl present (condensation, $\ce{+ H2O}$).
\item \textbf{Tollens'} (silver mirror) and \textbf{Fehling's} (brick-red $\ce{Cu2O}$) oxidise aldehydes only — the distinguishing tests.
\item \textbf{Iodoform}: yellow $\ce{CHI3}$ only from methyl ketones, ethanal (and $\ce{CH3-CH(OH)-R}$ alcohols).
\end{itemize}
\end{aretenir}

\coursec{Carboxylic Acids: Structure, Properties, Preparation and Reactions}

In the previous section we saw that aldehydes are easily oxidised: Fehling's solution, Tollens' reagent and even acidified dichromate all convert them into a new family of compounds. That family is the \cle{carboxylic acids} --- the first truly \emph{acidic} organic compounds we meet, and the parents of a whole dynasty of derivatives (esters, amides, acyl chlorides, anhydrides) studied in the next section. The sour taste of the vinegar used to season \emph{ndol\'e} or to preserve vegetables in Cameroonian kitchens is due to the most familiar member of the family, ethanoic acid.

\courssub{Structure, General Formula and Nomenclature}

\begin{definition}[The carboxyl group]
A \cle{carboxylic acid} is an organic compound containing the \textbf{carboxyl group} \ce{-COOH}, in which a \textbf{carbonyl group} \ce{C=O} and a \textbf{hydroxyl group} \ce{-OH} are attached to the \emph{same} carbon atom. The two groups do not behave independently: the carbonyl withdraws electron density from the \ce{O-H} bond, making the proton far easier to lose than in an alcohol. Saturated aliphatic monocarboxylic acids have the general molecular formula \ensuremath{\mathrm{C_{n}H_{2n}O2}} ($n \geq 1$), or \ce{R-COOH} where \ce{R} is \ce{H} or an alkyl group.
\end{definition}

\textbf{Nomenclature.} The longest chain is chosen to include the carboxyl carbon, which is always carbon number 1. The ending \emph{-e} of the corresponding alkane is replaced by \cle{-oic acid}:

\begin{center}
\begin{tabularx}{\linewidth}{|l|X|l|}
\hline
\rowcolor{popblueL} \textbf{Formula} & \textbf{Name} & \textbf{Common name} \\
\hline
\ce{HCOOH} & methanoic acid & formic acid \\
\hline
\ce{CH3COOH} & ethanoic acid & acetic acid \\
\hline
\ce{CH3CH2COOH} & propanoic acid & propionic acid \\
\hline
\ce{CH3CH2CH2COOH} & butanoic acid & butyric acid \\
\hline
\ce{C6H5COOH} & benzoic acid & benzoic acid \\
\hline
\end{tabularx}
\end{center}

Substituents are located by numbering from the \ce{COOH} carbon: \ce{CH3CH(Cl)COOH} is 2-chloropropanoic acid. Acids with two carboxyl groups use the suffix \emph{-dioic acid}: \ce{HOOC-COOH} is ethanedioic (oxalic) acid.

\textbf{Isomerism.} Carboxylic acids show \cle{chain isomerism} (e.g.\ \ce{CH3CH2CH2COOH} butanoic acid and \ce{(CH3)2CHCOOH} 2-methylpropanoic acid, both \ce{C4H8O2}) and \cle{functional group isomerism} with \textbf{esters}: \ce{C2H4O2} is ethanoic acid \ce{CH3COOH} or methyl methanoate \ce{HCOOCH3}. This is a favourite GCE question: an ester and an acid with the same molecular formula are easily distinguished by the effervescence test (see below).

\courssub{Physical Properties and Uses}

Carboxylic acids have \textbf{anomalously high boiling points} --- higher even than alcohols of similar molar mass. Ethanoic acid ($M_r = 60$) boils at \SI{118}{\ensuremath{{}^{\circ}\mathrm{C}}}, while propan-1-ol ($M_r = 60$) boils at \SI{97}{\ensuremath{{}^{\circ}\mathrm{C}}}. The reason is \cle{hydrogen bonding} taken one step further: two acid molecules pair up through \emph{two} hydrogen bonds to form a stable \textbf{dimer}, which persists even in the vapour state, so the effective molar mass is nearly doubled.

\begin{popfigure}
\begin{center}
\begin{tikzpicture}[scale=1.0, every node/.style={font=\small}]
% left molecule
\node at (0,1.0) {\ce{CH3}};
\node at (0,-1.0) {\ce{CH3}};
\draw[thick] (0.55,1.0) -- (1.25,1.0);
\draw[thick] (0.55,-1.0) -- (1.25,-1.0);
\node at (1.5,1.0) {\ce{C}};
\node at (1.5,-1.0) {\ce{C}};
\draw[thick,double] (1.75,1.15) -- (2.45,1.6);
\draw[thick,double] (1.75,-1.15) -- (2.45,-1.6);
\node at (2.75,1.7) {\ce{O}};
\node at (2.75,-1.7) {\ce{O}};
\draw[thick] (1.75,0.85) -- (2.45,0.45);
\draw[thick] (1.75,-0.85) -- (2.45,-0.45);
\node at (2.7,0.35) {\ce{O}};
\node at (2.7,-0.35) {\ce{O}};
\draw[thick] (2.95,0.35) -- (3.55,0.35);
\draw[thick] (2.95,-0.35) -- (3.55,-0.35);
\node at (3.8,0.35) {\ce{H}};
\node at (3.8,-0.35) {\ce{H}};
% hydrogen bonds
\draw[dashed,thick,poporange] (4.05,0.35) -- (4.75,0.35);
\draw[dashed,thick,poporange] (4.05,-0.35) -- (4.75,-0.35);
% right molecule
\node at (5.0,0.35) {\ce{O}};
\node at (5.0,-0.35) {\ce{O}};
\draw[thick] (5.25,0.45) -- (5.95,0.85);
\draw[thick] (5.25,-0.45) -- (5.95,-0.85);
\node at (6.2,1.0) {\ce{C}};
\node at (6.2,-1.0) {\ce{C}};
\draw[thick,double] (5.95,1.15) -- (5.25,1.6);
\draw[thick,double] (5.95,-1.15) -- (5.25,-1.6);
\node at (4.95,1.7) {\ce{O}};
\node at (4.95,-1.7) {\ce{O}};
\draw[thick] (6.45,1.0) -- (7.15,1.0);
\draw[thick] (6.45,-1.0) -- (7.15,-1.0);
\node at (7.7,1.0) {\ce{CH3}};
\node at (7.7,-1.0) {\ce{CH3}};
\node[popmuted] at (3.85,-2.3) {two hydrogen bonds \ce{O-H...O}};
\end{tikzpicture}
\end{center}
\legende{The hydrogen-bonded dimer of ethanoic acid: two \ce{O-H...O} bridges (dashed) hold the pair of molecules together.}
\end{popfigure}

The \textbf{small acids} (methanoic to butanoic) are \textbf{completely miscible with water}, because the \ce{-COOH} group hydrogen-bonds strongly with water molecules. Solubility falls as the hydrocarbon chain lengthens; benzoic acid is only sparingly soluble in cold water but dissolves in hot water. The first members are colourless liquids with sharp, vinegar-like smells; higher members are waxy solids.

\begin{savaistu}[Carboxylic acids around us]
Vinegar is a solution of about 5--8\% ethanoic acid in water, used in cooking and food preservation. Citric acid (2-hydroxypropane-1,2,3-tricarboxylic acid) gives lemons and oranges their sourness and is used as a food acidulant. Methanoic acid is the weapon of stinging ants and nettles --- its old name, formic acid, comes from the Latin \emph{formica}, ant. Ethanoic acid is produced industrially on a large scale for the manufacture of plastics, solvents and dyes.
\end{savaistu}

\courssub{Preparation of Carboxylic Acids}

\begin{methode}[Main laboratory and industrial preparations]
\begin{enumerate}
\item \textbf{Oxidation of a primary alcohol or an aldehyde} with acidified potassium dichromate (\ce{K2Cr2O7/H+}), heating under reflux (for the alcohol, excess oxidant and reflux ensure oxidation goes past the aldehyde stage):
\[ \ce{CH3CH2OH + 2[O] ->[{\ce{Cr2O7^{2-}/H+},\ \Delta}] CH3COOH + H2O} \]
\[ \ce{CH3CHO + [O] ->[{\ce{Cr2O7^{2-}/H+}}] CH3COOH} \]
\item \textbf{Hydrolysis of a nitrile}, by boiling with dilute acid (or dilute alkali followed by acidification):
\[ \ensuremath{\mathrm{CH_{3}CN + 2H_{2}O \rightarrow [H+,\ \Delta] CH_{3}COOH + NH_{4}+}} \]
This route \emph{adds one carbon} to the chain, since the nitrile itself was made from a haloalkane by substitution with \ce{CN-}.
\item \textbf{Hydrolysis of an ester} with dilute acid, or with alkali (saponification) followed by acidification:
\[ \ce{CH3COOCH2CH3 + H2O <=>[H+] CH3COOH + CH3CH2OH} \]
\item \textbf{Oxidation of an alkylbenzene} with hot alkaline \ce{KMnO4}, then acidification: whatever the length of the side chain, it is cut down to \ce{-COOH}:
\[ \ensuremath{\mathrm{C_{6}H_{5}CH_{2}CH_{3} \rightarrow [{\ce{KMnO4/OH^{-}},\ \Delta}][\text{then } \ce{H^{+}}] C_{6}H_{5}COOH}} \quad\text{(benzoic acid)} \]
\item \textbf{Carbonation of a Grignard reagent}: bubbling \ce{CO2} through the reagent, then acidic hydrolysis --- another useful chain-lengthening route:
\[ \ce{RMgX + CO2 -> RCOOMgX ->[H3O+] RCOOH} \]
\end{enumerate}
\end{methode}

\courssub{Acidity of Carboxylic Acids}

Carboxylic acids are \textbf{weak acids}: they ionise only partially in water,
\[ \ce{CH3COOH(aq) + H2O(l) <=> CH3COO-(aq) + H3O+(aq)} \qquad pK_a(\ce{CH3COOH}) = 4.8 \]

Why does a carboxylic acid lose \ce{H+} so much more readily than an alcohol ($pK_a \approx 16$)? The answer lies in the stability of the \cle{carboxylate anion}.

\begin{propriete}[Resonance stabilisation of the carboxylate anion]
In the carboxylate ion \ce{RCOO-}, the negative charge is \textbf{delocalised equally over the two oxygen atoms}: the two \ce{C-O} bonds are identical, intermediate in length between single and double. This resonance stabilisation makes the anion much more stable than an alkoxide ion \ce{RO-} (where the charge is localised on one oxygen), so the equilibrium of ionisation lies further to the right and the acid is stronger. (The explanation, not a formal proof, is examinable.)
\end{propriete}

\begin{popfigure}
\begin{center}
\begin{tikzpicture}[scale=1.1]
% left structure
\node at (0,0) {\ce{R}};
\draw[thick] (0.35,0) -- (1.05,0) node[right] {\ce{C}};
\draw[thick,double] (1.32,0.12) -- (2.0,0.6) node[right] {\ce{O}};
\draw[thick] (1.32,-0.12) -- (2.0,-0.6) node[right] {\ce{O^{-}}};
% resonance arrow
\draw[<->,thick,poporange] (3.1,0) -- (4.1,0);
% right structure
\node at (4.6,0) {\ce{R}};
\draw[thick] (4.95,0) -- (5.65,0) node[right] {\ce{C}};
\draw[thick] (5.92,0.12) -- (6.6,0.6) node[right] {\ce{O^{-}}};
\draw[thick,double] (5.92,-0.12) -- (6.6,-0.6) node[right] {\ce{O}};
% hybrid
\draw[->,thick,popdark] (3.6,-1.1) -- (3.6,-1.7);
\node at (1.9,-2.2) {\ce{R}};
\draw[thick] (2.25,-2.2) -- (2.95,-2.2) node[right] {\ce{C}};
\draw[thick,dash pattern=on 2pt off 2pt] (3.22,-2.08) -- (3.9,-1.6);
\draw[thick,dash pattern=on 2pt off 2pt] (3.22,-2.32) -- (3.9,-2.8);
\node at (4.3,-1.55) {\ensuremath{\mathrm{O^{\delta^{-}}}}};
\node at (4.3,-2.85) {\ensuremath{\mathrm{O^{\delta^{-}}}}};
\node[popmuted] at (3.6,-3.4) {resonance hybrid};
\end{tikzpicture}
\end{center}
\legende{Resonance in the carboxylate anion: the charge is shared equally by the two oxygens.}
\end{popfigure}

\textbf{Effect of substituents.} An \cle{electron-withdrawing group} (e.g.\ \ce{Cl}) attached near the \ce{COOH} pulls electron density away through the $\sigma$-bonds (the \textbf{inductive effect}, $-I$), spreading the negative charge of the anion still further and \emph{increasing} the acidity. The effect grows with the number of halogens and decreases with distance from the \ce{COOH}:
\[ \ce{Cl3CCOOH} > \ce{Cl2CHCOOH} > \ce{ClCH2COOH} > \ce{CH3COOH} \quad\text{(decreasing acidity)} \]
Conversely, an electron-\emph{donating} alkyl group ($+I$) pushes electron density onto the anion and destabilises it slightly, so methanoic acid ($pK_a = 3.8$) is stronger than ethanoic acid ($pK_a = 4.8$).

\begin{popfigure}
\begin{center}
\begin{tikzpicture}
\begin{axis}[
ybar, bar width=16pt, width=11.5cm, height=6.2cm,
ymin=0, ymax=5.5,
ylabel={$pK_a$}, symbolic x coords={CCl3, CHCl2, CH2Cl, H, CH3},
xtick=data, xticklabels={\ce{Cl3CCOOH},\ce{Cl2CHCOOH},\ce{ClCH2COOH},\ce{HCOOH},\ce{CH3COOH}},
x tick label style={font=\small, rotate=15, anchor=east},
nodes near coords, nodes near coords style={font=\small},
axis line style={draw=popline}, tick style={draw=popline},
every axis plot/.append style={fill=popblue, draw=popdark}
]
\addplot coordinates {(CCl3,0.7) (CHCl2,1.3) (CH2Cl,2.9) (H,3.8) (CH3,4.8)};
\end{axis}
\end{tikzpicture}
\end{center}
\legende{$pK_a$ values of substituted ethanoic acids: more electron-withdrawing chlorine atoms $\Rightarrow$ lower $pK_a$ $\Rightarrow$ stronger acid.}
\end{popfigure}

\begin{attention}[Reading a $pK_a$ chart]
A \textbf{lower} $pK_a$ means a \textbf{stronger} acid. Do not say ``\ce{Cl3CCOOH} has the smallest $pK_a$ so it is the weakest'' --- it is the \emph{strongest} of the series.
\end{attention}

\courssub{Reactions of Carboxylic Acids}

\textbf{1. Salt formation (reactions as an acid).} Carboxylic acids neutralise bases and displace carbon dioxide from carbonates and hydrogencarbonates:
\[ \ce{CH3COOH + NaOH -> CH3COONa + H2O} \]
\[ \ce{2CH3COOH + Na2CO3 -> 2CH3COONa + H2O + CO2 ^} \]
\[ \ce{CH3COOH + NaHCO3 -> CH3COONa + H2O + CO2 ^} \]

\begin{propriete}[Exam tip: the effervescence test]
Adding a carboxylic acid to solid \ce{NaHCO3} (or a solution of \ce{Na2CO3}) produces \textbf{effervescence of \ce{CO2}} (which turns limewater milky). \textbf{Phenols do not react} with \ce{NaHCO3} because they are weaker acids than carbonic acid. This test therefore \emph{distinguishes a carboxylic acid from a phenol} --- a classic GCE practical question.
\end{propriete}

\textbf{2. Esterification.} With an alcohol, in the presence of a few drops of \textbf{concentrated \ce{H2SO4}} and heat under reflux, an acid forms an \textbf{ester} in a slow, \textbf{reversible} reaction:
\[ \ensuremath{\mathrm{CH_{3}COOH + CH_{3}CH_{2}OH \rightleftharpoons [{\text{conc. } \ce{H2SO4},\ \Delta}] CH_{3}COOCH_{2}CH_{3} + H_{2}O}} \]

\begin{methode}[Esterification: conditions and role of the catalyst]
\begin{enumerate}
\item Mix the acid and alcohol, add a few drops of concentrated \ce{H2SO4}.
\item Heat under reflux (the reaction is slow at room temperature).
\item Concentrated \ce{H2SO4} acts as a \textbf{catalyst} (provides \ce{H+}, speeds up attainment of equilibrium) \emph{and} as a \textbf{dehydrating agent} (absorbs the water formed, shifting the equilibrium towards the ester by Le Chatelier's principle).
\item Remember: it is the \textbf{acid that loses \ce{OH}} and the \textbf{alcohol that loses \ce{H}}.
\end{enumerate}
\end{methode}

\textbf{3. Formation of acyl chlorides.} The \ce{-OH} of the carboxyl group is replaced by \ce{Cl} using \ce{PCl5} or, better, \textbf{thionyl chloride} \ce{SOCl2} (gaseous by-products make purification easy):
\[ \ce{CH3COOH + PCl5 -> CH3COCl + POCl3 + HCl} \qquad \ce{CH3COOH + SOCl2 -> CH3COCl + SO2 ^ + HCl ^} \]

\textbf{4. Reduction.} Unlike aldehydes, carboxylic acids resist mild reducing agents; only the powerful \textbf{lithium tetrahydridoaluminate} \ce{LiAlH4} (in dry ether, then hydrolysis) reduces them to \textbf{primary alcohols}:
\[ \ensuremath{\mathrm{CH_{3}COOH + 4[H] \rightarrow [{\ce{LiAlH4}}] CH_{3}CH_{2}OH + H_{2}O}} \]

\textbf{5. Amide formation.} The acid first reacts with ammonia to give the \textbf{ammonium salt}, which on strong heating loses water to give the \textbf{amide}:
\[ \ensuremath{\mathrm{CH_{3}COOH + NH_{3} \rightarrow  CH_{3}COONH_{4} \rightarrow [\Delta] CH_{3}CONH_{2} + H_{2}O}} \]

\textbf{6. Dehydration to anhydrides.} Heating two molecules of acid with a dehydrating agent (\ce{P2O5}) removes water to give an \textbf{acid anhydride}:
\[ \ensuremath{\mathrm{2CH_{3}COOH \rightarrow [{\ce{P2O5},\ \Delta}] (CH_{3}CO)_{2}O + H_{2}O}} \]

\textbf{7. Alpha-halogenation (HVZ reaction).} With \ce{Cl2} or \ce{Br2} in the presence of red phosphorus, a hydrogen on the \textbf{$\alpha$-carbon} (the carbon next to \ce{COOH}) is replaced (Hell--Volhard--Zelinsky reaction):
\[ \ensuremath{\mathrm{CH_{3}COOH + Cl_{2} \rightarrow [{\text{red } P}] ClCH_{2}COOH + HCl}} \]

\textbf{8. Special properties of methanoic acid.} \ce{HCOOH} is unique: its molecule contains an \textbf{aldehyde-like} \ce{-CHO} fragment (\ce{H-C(=O)-OH}), so it is a \textbf{reducing agent}.

\begin{attention}[Methanoic acid is exceptional!]
\ce{HCOOH} (and its salts) \textbf{reduces Tollens' reagent} (silver mirror) and \textbf{Fehling's solution} (brick-red \ce{Cu2O}), exactly like an aldehyde. No other carboxylic acid does this. On warming with concentrated \ce{H2SO4} it also dehydrates to carbon monoxide: $\ensuremath{\mathrm{HCOOH \rightarrow [{\text{conc. } \ce{H2SO4}}] CO ^ + H_{2}O}}$.
\end{attention}

\begin{popfigure}
\begin{center}
\begin{tikzpicture}[scale=1.0, every node/.style={font=\small}]
\node[draw=popdark, fill=popblueL, rounded corners, thick, minimum width=2.6cm, align=center] (acid) at (0,0) {\textbf{Ethanoic acid}\\ \ce{CH3COOH}};
\node[align=center] (salt) at (-4.6,1.8) {\ce{CH3COONa}\\ sodium ethanoate};
\node[align=center] (ester) at (4.6,1.8) {\ce{CH3COOC2H5}\\ ethyl ethanoate};
\node[align=center] (acyl) at (4.6,-1.8) {\ce{CH3COCl}\\ ethanoyl chloride};
\node[align=center] (alcohol) at (-4.6,-1.8) {\ce{CH3CH2OH}\\ ethanol};
\node[align=center] (amide) at (0,2.6) {\ce{CH3CONH2}\\ ethanamide};
\node[align=center] (anh) at (0,-2.6) {\ce{(CH3CO)2O}\\ ethanoic anhydride};
\draw[->,thick,popblue] (acid) -- (salt) node[midway, above, sloped, popmuted]{\scriptsize \ce{NaOH}};
\draw[->,thick,popblue] (acid) -- (ester) node[midway, above, sloped, popmuted]{\scriptsize \ce{C2H5OH}/\ce{H+}};
\draw[->,thick,popblue] (acid) -- (acyl) node[midway, below, sloped, popmuted]{\scriptsize \ce{SOCl2}};
\draw[->,thick,popblue] (acid) -- (alcohol) node[midway, below, sloped, popmuted]{\scriptsize \ce{LiAlH4}};
\draw[->,thick,popblue] (acid) -- (amide) node[midway, right, popmuted]{\scriptsize \ce{NH3} then $\Delta$};
\draw[->,thick,popblue] (acid) -- (anh) node[midway, right, popmuted]{\scriptsize \ce{P2O5}, $\Delta$};
\end{tikzpicture}
\end{center}
\legende{Reaction map centred on ethanoic acid.}
\end{popfigure}

\begin{exoresolu}[Worked exercise: identifying and comparing acids]
(a) An organic compound \textbf{A}, \ce{C3H6O2}, gives effervescence with \ce{NaHCO3}. Its isomer \textbf{B} does not, but has a fruity smell. Identify A and B. (b) Arrange in order of increasing acidity: \ce{CH3COOH}, \ce{ClCH2COOH}, \ce{HCOOH}, explaining briefly.
\tcblower
\textbf{(a)} Effervescence with \ce{NaHCO3} is the test for a \textbf{carboxylic acid}, so \textbf{A} is propanoic acid, \ce{CH3CH2COOH}. The isomer \textbf{B} with a fruity smell and no reaction with \ce{NaHCO3} is an \textbf{ester}: methyl ethanoate, \ce{CH3COOCH3} (or ethyl methanoate, \ce{HCOOCH2CH3}). This illustrates functional group isomerism in \ce{C3H6O2}.

\textbf{(b)} Increasing acidity: $\ce{CH3COOH} < \ce{HCOOH} < \ce{ClCH2COOH}$. The \ce{CH3} group is electron-donating ($+I$) and destabilises the carboxylate anion, so ethanoic acid is weaker than methanoic acid; the \ce{Cl} atom is electron-withdrawing ($-I$) and stabilises the anion, making chloroethanoic acid the strongest of the three ($pK_a$ 4.8, 3.8, 2.9 respectively).
\end{exoresolu}

\begin{exercice}[Practice]
\begin{enumerate}
\item Write the equation for the reaction of ethanoic acid with (i) solid sodium carbonate, (ii) propan-1-ol under esterification conditions, (iii) thionyl chloride. Name each organic product.
\item Explain, using the structures of the anions, why ethanoic acid is a much stronger acid than ethanol.
\item A colourless liquid gives a silver mirror with Tollens' reagent and effervesces with \ce{NaHCO3}. Identify it and justify both observations.
\end{enumerate}
\end{exercice}

\begin{corrige}[Exercise]
\begin{enumerate}
\item (i) $\ce{2CH3COOH + Na2CO3 -> 2CH3COONa + H2O + CO2 ^}$, sodium ethanoate. (ii) $\ce{CH3COOH + CH3CH2CH2OH <=>[H+] CH3COOCH2CH2CH3 + H2O}$, propyl ethanoate. (iii) $\ce{CH3COOH + SOCl2 -> CH3COCl + SO2 + HCl}$, ethanoyl chloride.
\item On losing \ce{H+}, ethanoic acid gives the carboxylate anion \ce{CH3COO-} in which the negative charge is delocalised over \emph{two} oxygens by resonance; ethanol gives the ethoxide ion \ce{CH3CH2O-} with the charge localised on one oxygen. The stabilised anion forms more readily, so ethanoic acid ionises more and is the stronger acid.
\item Methanoic acid, \ce{HCOOH}. It effervesces with \ce{NaHCO3} because it is a carboxylic acid, and it reduces Tollens' reagent because its structure contains an aldehyde-type \ce{-CHO} group, making it a reducing agent --- unique among carboxylic acids.
\end{enumerate}
\end{corrige}

\begin{aretenir}[Key points]
\begin{itemize}
\item Carboxylic acids \ce{R-COOH} ($\ensuremath{\mathrm{C_nH_{2n}O2}}$) contain the carboxyl group: carbonyl $+$ hydroxyl acting together; named with the suffix \emph{-oic acid}.
\item Hydrogen-bonded \textbf{dimers} explain high boiling points; small acids are water-soluble.
\item Prepared by oxidation of primary alcohols/aldehydes, hydrolysis of nitriles or esters, oxidation of alkylbenzenes, or carbonation of Grignard reagents.
\item Acidity arises from \textbf{resonance stabilisation} of \ce{RCOO-}; electron-withdrawing substituents strengthen the acid: $\ce{Cl3CCOOH} > \ce{ClCH2COOH} > \ce{CH3COOH}$.
\item Effervescence with \ce{NaHCO3} distinguishes carboxylic acids from phenols.
\item Key reactions: salt formation, esterification (conc.\ \ce{H2SO4}, reflux, reversible), acyl chloride formation (\ce{SOCl2}), reduction by \ce{LiAlH4}, amide and anhydride formation, $\alpha$-halogenation (HVZ).
\item \textbf{Methanoic acid is exceptional}: it reduces Tollens' and Fehling's reagents.
\end{itemize}
\end{aretenir}

\coursec{Carboxylic Acid Derivatives: Acyl Chlorides, Anhydrides, Esters and Amides}

In the previous section we studied carboxylic acids, \ce{RCOOH}, and saw that their chemistry is dominated by the carbonyl group \ce{C=O} and the hydroxyl group \ce{-OH}. We also saw that acids can be converted into more reactive species: acyl chlorides with \ce{SOCl2}, esters with alcohols, amides with ammonia. These compounds all share a common skeleton: an \cle{acyl group} \ce{RCO-} attached to a replaceable group \ce{Z}. They are called the \cle{carboxylic acid derivatives}, and they form one of the most coherent families in organic chemistry: nearly all their reactions follow a single pattern, \cle{nucleophilic acyl substitution}, and a single reactivity ladder governs which derivative can be converted into which.

\courssub{General Structure and the Reactivity Ladder}

\begin{definition}[Carboxylic acid derivative]
A carboxylic acid derivative is a compound of general formula \ce{RCOZ}, where \ce{RCO-} is an acyl group and \ce{Z} is an electronegative group that can be replaced by a nucleophile:
\begin{itemize}
\item \textbf{Acyl chloride}: \ce{RCOCl} (\ce{Z} = \ce{Cl})
\item \textbf{Acid anhydride}: \ensuremath{\mathrm{RCO-O-COR^{\prime}}} (\ce{Z} = \ensuremath{\mathrm{OCOR^{\prime}}})
\item \textbf{Ester}: \ensuremath{\mathrm{RCOOR^{\prime}}} (\ce{Z} = \ensuremath{\mathrm{OR^{\prime}}})
\item \textbf{Amide}: \ce{RCONH2} (\ce{Z} = \ce{NH2})
\end{itemize}
\end{definition}

\begin{methode}[The general mechanism: nucleophilic acyl substitution (addition--elimination)]
Every reaction of these derivatives with a nucleophile \ce{Nu^-} (or \ce{NuH}) proceeds in two stages:
\begin{enumerate}
\item \textbf{Addition.} The nucleophile attacks the electron-poor carbonyl carbon; the \ce{C=O} $\pi$ electrons move onto oxygen, giving a tetrahedral intermediate with \ce{O^-}.
\item \textbf{Elimination.} The \ce{O^-} lone pair collapses back down to re-form \ce{C=O}, expelling the leaving group \ce{Z^-}.
\end{enumerate}
Overall: \ce{RCOZ + Nu^- -> RCONu + Z^-}. The carbonyl is \emph{restored} — this is substitution, not the permanent addition seen with aldehydes and ketones, because here a good leaving group is present.
\end{methode}

\begin{propriete}[Reactivity order of acid derivatives]
Towards nucleophiles, the reactivity decreases in the order:
\[ \text{acyl chloride} > \text{anhydride} > \text{ester} > \text{amide} \]
\emph{Justification (examinable).} Reactivity depends on the \cle{leaving-group ability} of \ce{Z^-}: the weaker the base \ce{Z^-}, the better it leaves. \ce{Cl^-} is the conjugate base of the strong acid \ce{HCl}, so it is an excellent leaving group; \ce{RCOO^-} is moderate; \ce{RO^-} and especially \ce{NH2^-} are strong bases and leave poorly. In addition, in esters and amides the lone pair on \ce{O} or \ce{N} is delocalised into the carbonyl (\ce{C=O} acquires partial single-bond character), making the carbonyl carbon less electrophilic. This delocalisation is strongest in amides (nitrogen donates best), which is why amides are the least reactive of all.
\end{propriete}

\begin{attention}[A practical consequence]
A more reactive derivative can always be converted into a less reactive one (acyl chloride $\rightarrow$ anhydride $\rightarrow$ ester $\rightarrow$ amide), but \emph{never} the reverse directly. To go "up" the ladder you must first hydrolyse back to the carboxylic acid, then reactivate it (e.g.\ with \ce{SOCl2}).
\end{attention}

\begin{popfigure}
\begin{center}
\begin{tikzpicture}[scale=1.0]
% ladder
\draw[thick, popmuted] (0,0) -- (0,4.6);
\draw[->, thick, poporange] (-0.9,0.3) -- (-0.9,4.3) node[midway, left, align=center, text width=2.2cm]{\small decreasing reactivity};
\node[draw=popblue, fill=popblueL, rounded corners, minimum width=5.2cm, minimum height=0.8cm] at (3.2,4.1) {\ce{RCOCl} \quad acyl chloride};
\node[draw=popgreen, fill=popgreenL, rounded corners, minimum width=5.2cm, minimum height=0.8cm] at (3.2,2.9) {\ce{(RCO)2O} \quad anhydride};
\node[draw=poppurple, fill=white, rounded corners, minimum width=5.2cm, minimum height=0.8cm] at (3.2,1.7) {\ensuremath{\mathrm{RCOOR^{\prime}}} \quad ester};
\node[draw=popgold, fill=white, rounded corners, minimum width=5.2cm, minimum height=0.8cm] at (3.2,0.5) {\ce{RCONH2} \quad amide};
\node[right, align=left, text width=4.6cm] at (6.1,4.1) {\small leaving group \ce{Cl^-}: excellent};
\node[right, align=left, text width=4.6cm] at (6.1,2.9) {\small leaving group \ce{RCOO^-}: good};
\node[right, align=left, text width=4.6cm] at (6.1,1.7) {\small leaving group \ce{RO^-}: poor};
\node[right, align=left, text width=4.6cm] at (6.1,0.5) {\small leaving group \ce{NH2^-}: very poor};
\end{tikzpicture}
\end{center}
\legende{The reactivity ladder of carboxylic acid derivatives, governed by leaving-group ability.}
\end{popfigure}

\courssub{Acyl Chlorides}

\textbf{Nomenclature.}\par
Acyl chlorides are named by replacing the \emph{-oic acid} ending of the parent acid with \emph{-oyl chloride}. The carbonyl carbon is always C-1. Examples: \ce{CH3COCl} is \textbf{ethanoyl chloride}; \ce{C6H5COCl} is \textbf{benzoyl chloride}.

\textbf{Physical properties.}\par
Acyl chlorides are colourless liquids (low molecular mass) with sharp, irritating odours. They are not hydrogen-bonded, so boiling points are lower than the parent acids of similar molar mass. They fume in moist air due to reaction with water vapour.

\textbf{Preparation.}\par
Acyl chlorides are prepared from carboxylic acids using \cle{thionyl chloride}, \ce{SOCl2}:
\[ \ce{CH3COOH(l) + SOCl2(l) -> CH3COCl(l) + SO2(g) + HCl(g)} \]
This reagent is preferred over \ce{PCl5} or \ce{PCl3} because both by-products are gases, leaving pure acyl chloride after distillation. Alternative: \ce{CH3COOH + PCl5 -> CH3COCl + POCl3 + HCl}.

\textbf{Reactions.}\par
Acyl chlorides are the most reactive derivatives; all the reactions below are rapid, often vigorous, and occur at room temperature.

\begin{itemize}
\item \textbf{Hydrolysis} (water): \ce{CH3COCl(l) + H2O(l) -> CH3COOH(aq) + HCl(aq)} — vigorous, with steamy fumes of \ce{HCl}.
\item \textbf{Alcoholysis} (alcohols): \ce{CH3COCl(l) + C2H5OH(l) -> CH3COOC2H5(l) + HCl(g)} — a fast, irreversible route to esters, superior to direct esterification. Pyridine is often added to neutralise \ce{HCl}.
\item \textbf{Aminolysis} (ammonia or amines): \ce{CH3COCl(l) + 2NH3(aq) -> CH3CONH2(s) + NH4Cl(aq)}; with a primary amine, an $N$-substituted amide forms: \ce{CH3COCl + 2CH3NH2 -> CH3CONHCH3 + CH3NH3+Cl^-}. Two moles of amine are needed: one as nucleophile, one to neutralise \ce{HCl}.
\item \textbf{Friedel--Crafts acylation of benzene} (catalyst: anhydrous \ce{AlCl3}): \ce{C6H6(l) + CH3COCl(l) ->[AlCl3] C6H5COCH3(l) + HCl(g)}, giving the ketone phenylethanone (acetophenone).
\item \textbf{Reduction} (with \ce{LiAlH4} in dry ether): acyl chlorides are reduced to primary alcohols: \ce{CH3COCl ->[LiAlH4] CH3CH2OH}. (Partial reduction to aldehyde using \ce{LiAlH(OtBu)3} is beyond GCE scope.)
\end{itemize}

\begin{attention}[Exam tip: the fuming test]
Acyl chlorides \cle{fume in moist air}: they react with atmospheric water vapour to release \ce{HCl} gas, which appears as white steamy fumes and turns damp blue litmus red. This is a classic observation question — and a quick way to distinguish an acyl chloride from an ester or amide.
\end{attention}

\courssub{Acid Anhydrides}

\textbf{Nomenclature.}\par
Symmetrical anhydrides are named by replacing \emph{acid} with \emph{anhydride}: \ce{(CH3CO)2O} is \textbf{ethanoic anhydride}. Unsymmetrical anhydrides name both acids alphabetically: \ce{CH3CO-O-COCH2CH3} is \textbf{ethanoic propanoic anhydride}.

\textbf{Physical properties.}\par
Anhydrides are colourless liquids (low mass) or solids with pungent odours. They hydrolyse slowly in moist air. Boiling points are higher than acyl chlorides but lower than acids of similar mass due to lack of hydrogen bonding.

\textbf{Preparation.}\par
\begin{itemize}
\item From an acyl chloride and a carboxylate salt (laboratory): \ce{CH3COCl(l) + CH3COO^-Na+(s) -> (CH3CO)2O(l) + NaCl(s)}.
\item Industrial: dehydration of ethanoic acid via ketene (\ce{CH2=C=O}) at high temperature.
\end{itemize}

\textbf{Reactions.}\par
Anhydrides undergo the same nucleophilic acyl substitutions as acyl chlorides, but \emph{less violently} — they are the chemist's choice when a gentler, more controllable acylating agent is needed:
\begin{itemize}
\item \textbf{Hydrolysis:} \ce{(CH3CO)2O(l) + H2O(l) -> 2CH3COOH(aq)} (slow when cold, fast on warming).
\item \textbf{Alcoholysis:} \ce{(CH3CO)2O(l) + C2H5OH(l) -> CH3COOC2H5(l) + CH3COOH(l)}.
\item \textbf{Aminolysis:} \ce{(CH3CO)2O(l) + 2NH3(aq) -> CH3CONH2(s) + CH3COO^-NH4+(aq)}.
\item \textbf{Reduction} (\ce{LiAlH4}): gives two molecules of primary alcohol: \ce{(CH3CO)2O ->[LiAlH4] 2CH3CH2OH}.
\end{itemize}

\begin{savaistu}[Enrichment: aspirin from ethanoic anhydride]
Aspirin (acetylsalicylic acid) is manufactured by reacting 2-hydroxybenzoic acid (salicylic acid) with \ce{ethanoic anhydride}: the phenolic \ce{-OH} is acetylated to an ester. The anhydride is used rather than ethanoyl chloride because it is safer and easier to handle on an industrial scale — and no corrosive \ce{HCl} is produced.
\end{savaistu}

\courssub{Esters: Nomenclature, Properties and Preparation}

\begin{methode}[Naming esters]
An ester \ensuremath{\mathrm{RCOOR^{\prime}}} is named as an \textbf{alkyl alkanoate}:
\begin{enumerate}
\item Name the \textbf{alcohol part} (\ensuremath{\mathrm{R^{\prime}}}) first, as an alkyl group: methyl, ethyl, propyl\ldots
\item Name the \textbf{acid part} (\ce{RCO}) second, replacing \emph{-oic acid} by \emph{-oate}.
\end{enumerate}
Example: \ce{CH3COOC2H5} is \textbf{ethyl ethanoate} (ethyl from ethanol, ethanoate from ethanoic acid); \ce{HCOOCH3} is \textbf{methyl methanoate}; \ce{C6H5COOCH3} is \textbf{methyl benzoate}.
\end{methode}

\textbf{Physical properties.}\par
Esters are polar but cannot hydrogen-bond to each other (no \ce{O-H}), so their boiling points are \emph{lower} than those of the parent acids of similar mass. Small esters are pleasantly fragrant liquids — many natural fruit flavours are esters: ethyl ethanoate smells of pear drops, pentyl ethanoate of banana, octyl ethanoate of orange. They are only slightly soluble in water but are excellent organic solvents (e.g., ethyl ethanoate for chromatography, nail varnish remover).

\textbf{Preparation.}\par
\begin{itemize}
\item \textbf{Esterification} (reversible, acid-catalysed): \ensuremath{\mathrm{CH_{3}COOH(l) + C_{2}H_{5}OH(l) \rightleftharpoons [conc. H_{2}SO_{4}, \Delta] CH_{3}COOC_{2}H_{5}(l) + H_{2}O(l)}}. Concentrated \ce{H2SO4} acts as catalyst and dehydrating agent to shift equilibrium.
\item From \textbf{acyl chlorides} or \textbf{anhydrides} + alcohol: fast and irreversible (see above).
\item From carboxylic acid + diazomethane (specialised, not GCE core).
\end{itemize}

\courssub{Reactions of Esters}

\begin{propriete}[Hydrolysis of esters]
\begin{itemize}
\item \textbf{Acid hydrolysis} (dilute \ce{HCl} or \ce{H2SO4}, heat): \ce{CH3COOC2H5(l) + H2O(l) <=>[H+] CH3COOH(aq) + C2H5OH(aq)} — \emph{reversible}, the exact reverse of esterification.
\item \textbf{Alkaline hydrolysis / saponification} (\ce{NaOH(aq)}, heat): \ce{CH3COOC2H5(l) + NaOH(aq) -> CH3COO^-Na+(aq) + C2H5OH(aq)} — \emph{irreversible}, giving the sodium carboxylate (a soap if the chain is long) and the alcohol.
\end{itemize}
\end{propriete}

\begin{attention}[Reversible or irreversible?]
Acid hydrolysis of an ester is \cle{reversible} because the carboxylic acid produced can re-esterify. Alkaline hydrolysis is \cle{irreversible}: the carboxylic acid is deprotonated by \ce{OH^-} to the carboxylate ion \ce{RCOO^-}, whose negatively charged, resonance-stabilised carbonyl is no longer attacked by the alcohol. Do not write an equilibrium arrow for saponification — a frequent exam error.
\end{attention}

\begin{experience}[Saponification: soap from palm oil]
Palm oil, abundant in the South-West and Littoral regions, consists mainly of triglycerides — triesters of glycerol (propane-1,2,3-triol) with long-chain fatty acids such as palmitic acid, \ce{C15H31COOH}. Boiling palm oil with concentrated \ce{NaOH} hydrolyses the three ester bonds: glycerol is released and the sodium salts of the fatty acids precipitate as \cle{soap}. This is exactly the traditional village process in which palm oil is boiled with alkaline ash extract.
\end{experience}

\begin{popfigure}
\begin{center}
\begin{tikzpicture}[scale=0.95]
\node[align=center] (oil) at (0,0) {\ce{CH2-OOC-C15H31}\\ \ce{CH-OOC-C15H31}\\ \ce{CH2-OOC-C15H31}};
\node[below=0.15cm of oil, align=center]{\small triglyceride (palm oil)};
\node[right=1.3cm of oil] (plus) {$+$};
\node[right=1.5cm of plus] (naoh) {\ce{3NaOH}};
\node[right=1.1cm of naoh] (arr) {};
\draw[->, very thick, poporange] (arr.west) -- ++(1.2,0) node[midway, above]{\small heat};
\node[right=2.9cm of arr, align=center] (gly) {\ce{CH2-OH}\\ \ce{CH-OH}\\ \ce{CH2-OH}};
\node[below=0.15cm of gly, align=center]{\small glycerol};
\node[right=1.1cm of gly] (plus2) {$+$};
\node[right=1.5cm of plus2, align=center] (soap) {\ce{3C15H31COO^-Na+}};
\node[below=0.35cm of soap, align=center]{\small sodium palmitate (soap)};
\end{tikzpicture}
\end{center}
\legende{Saponification of a triglyceride from palm oil: alkaline hydrolysis gives glycerol and soap.}
\end{popfigure}

\textbf{Other reactions of esters.}\par
\begin{itemize}
\item \textbf{Reduction} (with \ce{LiAlH4} in dry ether, then \ce{H2O} work-up): esters are reduced all the way to primary alcohols (two alcohols formed): \ce{CH3COOC2H5 ->[LiAlH4] CH3CH2OH + C2H5OH}. \ce{NaBH4} does \emph{not} reduce esters.
\item \textbf{Aminolysis} (concentrated ammonia, heat, pressure): \ce{CH3COOC2H5 + NH3 -> CH3CONH2 + C2H5OH} — another route to amides, slower than from acyl chlorides.
\item \textbf{Transesterification} (alcohol + acid/base catalyst): $\mathrm{RCOOR' + R''OH \rightleftharpoons RCOOR'' + R'OH}$ — equilibrium reaction.
\item \textbf{Claisen condensation} (ester + strong base, e.g., \ce{NaOEt}): forms \ensuremath{\mathrm{\beta}}-ketoesters. (Mentioned for completeness; detailed mechanism not required at GCE.)
\end{itemize}

\courssub{Amides}

\textbf{Nomenclature and structure.}\par
Amides, \ce{RCONH2}, are named by replacing \emph{-oic acid} with \emph{-amide}: \ce{CH3CONH2} is \textbf{ethanamide}, \ce{HCONH2} is \textbf{methanamide}, \ce{C6H5CONH2} is \textbf{benzamide}. $N$-substituted amides (from amines) carry the prefix $N$-: \ce{CH3CONHCH3} is \textbf{$N$-methylethanamide}, \ce{CH3CON(CH3)2} is \textbf{$N,N$-dimethylethanamide}. The \ce{C-N} bond has partial double-bond character due to delocalisation of the nitrogen lone pair into the carbonyl — hence their low reactivity, planar geometry, and lack of basicity (the lone pair is unavailable).

\textbf{Physical properties.}\par
Primary and secondary amides have \ce{N-H} bonds and form strong hydrogen bonds, giving high melting/boiling points (e.g., ethanamide m.p. 82 \ensuremath{{}^{\circ}} C). Tertiary amides have no \ce{N-H} and lower boiling points. Small amides are soluble in water.

\textbf{Preparation.}\par
\begin{itemize}
\item Acyl chloride or anhydride + ammonia/amine (fast, standard laboratory route).
\item From an ester + concentrated ammonia (slow, heat).
\item By \textbf{heating an ammonium salt} of the acid: \ensuremath{\mathrm{CH_{3}COO^{-}NH_{4}+(s) \rightarrow [\Delta] CH_{3}CONH_{2}(s) + H_{2}O(g)}}.
\item Industrial: direct reaction of acid with ammonia at high temperature over a catalyst.
\end{itemize}

\textbf{Reactions.}\par
\begin{itemize}
\item \textbf{Hydrolysis} (needs heating with acid or alkali): with \ce{NaOH}, \ce{CH3CONH2(s) + NaOH(aq) -> CH3COO^-Na+(aq) + NH3(g)} — the ammonia evolved (turns damp red litmus blue) is a useful test for amides. Acid hydrolysis gives the acid + \ce{NH4+}: \ce{CH3CONH2 + H2O + H+ -> CH3COOH + NH4+}.
\item \textbf{Reduction} (\ce{LiAlH4} in dry ether): amides are reduced to \textbf{amines} (carbonyl oxygen removed entirely): \ce{CH3CONH2 ->[LiAlH4] CH3CH2NH2}. $N$-substituted amides give $N$-substituted amines.
\item \textbf{Hofmann degradation} (brief): a primary amide treated with \ce{Br2}/\ce{NaOH} gives a primary amine with \emph{one carbon fewer}: \ce{CH3CONH2 ->[Br2/NaOH] CH3NH2}. This is a useful way to shorten a carbon chain by one carbon.
\item \textbf{Dehydration} (with \ce{P2O5} or \ce{SOCl2}): primary amides dehydrate to nitriles: \ce{CH3CONH2 ->[P2O5] CH3CN + H2O}.
\item \textbf{Polyamides} (brief): diamines + dicarboxylic acids condense to give polyamides such as \textbf{nylon-6,6}, the polymer of ropes, fishing lines and clothing fibres.
\end{itemize}

\begin{exoresolu}[Worked example: choosing the right derivative]
\textbf{Statement.} A student must convert ethanoic acid into ethanamide in good yield. She considers two routes: (a) heating the acid directly with ammonia; (b) converting the acid to ethanoyl chloride first, then adding concentrated ammonia. Which route is better, and why? Write all equations.
\tcblower
\textbf{Solution.} Route (b) is better. Direct heating of \ce{CH3COOH} with \ce{NH3} first gives the ammonium salt \ce{CH3COO^-NH4+}, which must be strongly heated to dehydrate it to the amide — conditions under which yields are modest. The acyl chloride route exploits the reactivity ladder:
\begin{align*}
\ce{CH3COOH(l) + SOCl2(l) &-> CH3COCl(l) + SO2(g) + HCl(g)}\\
\ce{CH3COCl(l) + 2NH3(aq) &-> CH3CONH2(s) + NH4Cl(aq)}
\end{align*}
Ethanoyl chloride is far more reactive than the acid (excellent leaving group \ce{Cl^-}), so aminolysis is fast and complete at room temperature. Note that \textbf{two} moles of ammonia are needed: one acts as the nucleophile, the second neutralises the \ce{HCl} produced.
\end{exoresolu}

\begin{exoresolu}[Worked example: mechanism of alkaline hydrolysis of an ester]
\textbf{Statement.} Using curly arrows, outline the mechanism for the alkaline hydrolysis of ethyl ethanoate. State why the reaction is irreversible.
\tcblower
\textbf{Solution.}
\begin{enumerate}
\item \ce{OH^-} attacks the carbonyl carbon of \ce{CH3COOC2H5}; \ce{C=O} $\pi$ bond breaks to \ce{O^-} (tetrahedral intermediate).
\item The \ce{O^-} reforms \ce{C=O}, expelling \ce{C2H5O^-} (ethoxide).
\item \ce{C2H5O^-} is a stronger base than \ce{OH^-}, so it immediately abstracts a proton from the carboxylic acid formed (\ce{CH3COOH}), giving \ce{CH3COO^-} and \ce{C2H5OH}.
\end{enumerate}
The carboxylate ion \ce{CH3COO^-} is resonance-stabilised and its carbonyl carbon is not electrophilic enough to be attacked by ethanol. Hence no reverse reaction.
\end{exoresolu}

\begin{exercice}
(a) Arrange \ce{CH3COCl}, \ce{CH3CONH2}, \ce{(CH3CO)2O} and \ce{CH3COOCH3} in order of decreasing reactivity towards water, and justify in two sentences. (b) Write the equation for the reaction of propanoyl chloride with methanol and name the product. (c) Explain why saponification goes to completion while acid-catalysed ester hydrolysis does not. (d) Write the equation for the reduction of $N$-methylethanamide with \ce{LiAlH4} and name the organic product. (e) How would you distinguish experimentally between ethanoyl chloride and ethyl ethanoate?
\end{exercice}

\begin{corrige}[Exercise 1]
(a) \ce{CH3COCl} $>$ \ce{(CH3CO)2O} $>$ \ce{CH3COOCH3} $>$ \ce{CH3CONH2}. Reactivity follows leaving-group ability: \ce{Cl^-} is the weakest base and leaves best; \ce{NH2^-} is the strongest base and leaves worst, and delocalisation of the N lone pair further stabilises the amide. (b) \ce{C2H5COCl(l) + CH3OH(l) -> C2H5COOCH3(l) + HCl(g)}; the product is \textbf{methyl propanoate}. (c) In alkaline hydrolysis the acid product is deprotonated to the resonance-stabilised carboxylate \ce{RCOO^-}, which cannot be re-attacked by the alcohol, so the reverse reaction is impossible and hydrolysis is complete. In acid hydrolysis the free carboxylic acid formed can react again with the alcohol, so an equilibrium is established. (d) \ce{CH3CONHCH3 ->[LiAlH4] CH3CH2NHCH3}; product is \textbf{$N$-methylethanamine}. (e) Expose to moist air: ethanoyl chloride fumes white (\ce{HCl}) and turns damp blue litmus red; ethyl ethanoate has a fruity smell and does not fume. Alternatively, add water: ethanoyl chloride reacts violently with steamy fumes; ethyl ethanoate reacts slowly on heating.
\end{corrige}

\begin{aretenir}[Key points]
\begin{itemize}
\item All four derivatives \ce{RCOZ} react by \cle{nucleophilic acyl substitution}: addition to \ce{C=O}, then elimination of \ce{Z^-}.
\item Reactivity: \textbf{acyl chloride $>$ anhydride $>$ ester $>$ amide}, set by leaving-group ability; you can descend the ladder, never climb it directly.
\item Acyl chlorides fume in moist air (\ce{HCl}); prepared from acids + \ce{SOCl2}; best acylating agents.
\item Anhydrides are milder acylating agents (aspirin manufacture); hydrolysis gives two acid molecules.
\item Esters are alkyl alkanoates with fruity odours; \textbf{saponification} (\ce{NaOH}) is irreversible and turns palm oil into soap; acid hydrolysis is reversible.
\item Amides are the least reactive; planar \ce{C-N} with partial double bond; hydrolyse to acids (releasing \ce{NH3} with alkali), reduce to amines with \ce{LiAlH4}, dehydrate to nitriles, Hofmann degradation gives amine with one carbon fewer.
\end{itemize}
\end{aretenir}

\coursec{Nitrogen Compounds: Nitro Compounds, Nitriles, Amines, Diazonium Salts and Amino Acids}

In the previous section we met the carboxylic acid derivatives, and among them the \cle{amides}, $\ce{RCONH2}$, in which a nitrogen atom is attached to an acyl group. Nitrogen is the element that gives organic chemistry its biological dimension: it is found in proteins, in DNA, in most drugs, in dyes and in many fertilisers produced in Cameroon. In this final section we study the main families of nitrogen-containing organic compounds: \cle{nitro compounds}, \cle{nitriles}, \cle{amines}, \cle{diazonium salts} and \cle{amino acids}. A beautiful thread runs through them all: starting from benzene, we can travel from nitrobenzene to phenylamine, then to a diazonium salt, and finally to brightly coloured azo dyes.

\courssub{Nitro Compounds and Nitriles}

\textbf{Nitro compounds}\par

A \cle{nitro compound} contains the nitro group $\ce{-NO2}$ bonded to a carbon atom: general formula $\ce{R-NO2}$. They are classified as \emph{aliphatic nitro compounds} ($\ce{R}$ = alkyl) and \emph{aromatic nitro compounds} ($\ce{R}$ = aryl). The most important example is \cle{nitrobenzene}, $\ce{C6H5NO2}$, a pale yellow oily liquid with a smell of bitter almonds.

\begin{experience}[Preparation of nitrobenzene — Nitration of benzene]
Benzene is warmed at $50$--$\SI{60}{\ensuremath{{}^{\circ}\mathrm{C}}}$ with a nitrating mixture of concentrated nitric acid and concentrated sulphuric acid. The sulphuric acid generates the electrophile $\ce{NO2+}$ (nitronium ion):
\[ \ce{HNO3 + 2H2SO4 -> NO2+ + H3O+ + 2HSO4-} \]
\[ \ensuremath{\mathrm{C_{6}H_{6}(l) + HNO_{3}(conc) \rightarrow [conc.\ H_{2}SO_{4},\ 50-60\ ^{\circ}C] C_{6}H_{5}NO_{2}(l) + H_{2}O(l)}} \]
The mixture is then poured into cold water; nitrobenzene separates as the lower oily layer.
\end{experience}

\begin{propriete}[Chemical properties of nitro compounds]
\begin{itemize}
\item \textbf{Reduction to primary amines} (the key reaction): using tin and concentrated hydrochloric acid, followed by alkali to liberate the free amine:
\[ \ensuremath{\mathrm{C_{6}H_{5}NO_{2}(l) + 6[H] \rightarrow [Sn/conc.\ HCl] C_{6}H_{5}NH_{2}(aq) + 2H_{2}O(l)}} \]
This is the standard laboratory route from benzene to phenylamine, since benzene cannot be aminated directly.
\item \textbf{Uses}: Nitrobenzene is a precursor to phenylamine (aniline) for dyes, pharmaceuticals, and rubber chemicals. Aliphatic nitro compounds (e.g., nitromethane) are used as solvents and fuel additives.
\end{itemize}
\end{propriete}

\textbf{Nitriles}\par

\cle{Nitriles} contain the cyano group $\ce{-C#N}$; general formula $\ce{R-C#N}$. They are classified as \emph{aliphatic nitriles} ($\ce{R}$ = alkyl) and \emph{aromatic nitriles} ($\ce{R}$ = aryl). They are named with the suffix \emph{-nitrile}, counting the carbon of $\ce{C#N}$ in the chain: $\ce{CH3CN}$ is ethanenitrile, $\ce{CH3CH2CN}$ is propanenitrile, $\ce{C6H5CN}$ is benzonitrile.

\begin{propriete}[Preparation and reactions of nitriles]
\begin{itemize}
\item \textbf{Preparation from haloalkanes} (nucleophilic substitution, chain extension by one carbon): refluxing a haloalkane with alcoholic potassium cyanide:
\[ \ensuremath{\mathrm{CH_{3}CH_{2}Br(l) + KCN(alc) \rightarrow [ethanol,\ reflux] CH_{3}CH_{2}CN(l) + KBr(s)}} \]
\item \textbf{Preparation by dehydration of primary amides}: heating a primary amide with $\ce{P2O5}$ (or $\ce{SOCl2}$, $\ce{POCl3}$):
\[ \ensuremath{\mathrm{CH_{3}CONH_{2}(s) \rightarrow [P_{2}O_{5},\ heat] CH_{3}CN(l) + H_{2}O(g)}} \]
\item \textbf{Hydrolysis to carboxylic acids}: reflux with dilute acid (or alkali, giving the carboxylate salt):
\[ \ensuremath{\mathrm{CH_{3}CN(l) + 2H_{2}O(l) \rightarrow [H+,\ reflux] CH_{3}COOH(aq) + NH_{3}(aq)}} \]
\[ \ce{CH3CN(l) + H2O(l) + NaOH(aq) ->[reflux] CH3COONa(aq) + NH3(aq)} \]
\item \textbf{Reduction to primary amines}: with $\ce{LiAlH4}$ in dry ether, or catalytic hydrogenation ($\ce{H2}$/Ni, Pt):
\[ \ce{CH3CN(l) + 4[H] -> CH3CH2NH2(l)} \]
\item \textbf{Uses}: Nitriles are precious in synthesis because they allow a \emph{chain extension by one carbon}. Acrylonitrile ($\ce{CH2=CHCN}$) is polymerised to polyacrylonitrile for synthetic fibres (Orlon, Courtelle). Many pharmaceuticals contain nitrile groups.
\end{itemize}
\end{propriete}

\courssub{Amines: Classification, Nomenclature and Physical Properties}

Amines are organic derivatives of ammonia in which one, two or three hydrogen atoms have been replaced by alkyl or aryl groups.

\begin{definition}[Classification of amines]
\begin{itemize}
\item \textbf{Primary ($1^{\circ}$)}: one organic group attached to nitrogen, $\ce{RNH2}$, e.g. ethylamine $\ce{CH3CH2NH2}$.
\item \textbf{Secondary ($2^{\circ}$)}: two organic groups attached to nitrogen, $\ce{R2NH}$, e.g. dimethylamine $\ce{(CH3)2NH}$.
\item \textbf{Tertiary ($3^{\circ}$)}: three organic groups attached to nitrogen, $\ce{R3N}$, e.g. trimethylamine $\ce{(CH3)3N}$.
\item \textbf{Quaternary ammonium salts}: four organic groups, $\ce{R4N+ X-}$ (ionic, not amines proper).
\end{itemize}
Amines are \textbf{aliphatic} when the groups are alkyl (e.g. ethylamine) and \textbf{aromatic} when nitrogen is directly attached to a benzene ring (e.g. phenylamine, $\ce{C6H5NH2}$). In arylamines the lone pair is delocalised into the ring.
\end{definition}

\begin{attention}[A classic trap]
Do not confuse the classification of amines with that of alcohols! In amines we count the number of carbon groups attached to \emph{nitrogen}; in alcohols we count the carbon groups attached to the carbon bearing $\ce{-OH}$. Thus $\ce{(CH3)3CNH2}$ is a \emph{primary} amine even though the carbon skeleton is tertiary.
\end{attention}

\textbf{Nomenclature.} Simple amines are named as \emph{alkylamines}: $\ce{CH3NH2}$ methylamine, $\ce{C2H5NH2}$ ethylamine. In the systematic IUPAC style the suffix \emph{-amine} or the prefix \emph{amino-} is used: $\ce{CH3CH(NH2)CH3}$ is propan-2-amine; $\ce{C6H5NH2}$ is phenylamine (or aminobenzene). For secondary and tertiary amines, the largest alkyl group is the parent and the others are N-alkyl substituents: $\ce{(CH3)2NH}$ is N-methylethanamine (or dimethylamine).

\textbf{Physical properties.} Primary and secondary amines form \cle{hydrogen bonds} ($\ensuremath{\mathrm{N-H\cdots N}}$), so their boiling points are higher than those of alkanes of similar mass but lower than those of the corresponding alcohols (the $\ce{N-H}$ bond is less polar than $\ce{O-H}$). Tertiary amines have no $\ce{N-H}$ bond and cannot hydrogen-bond among themselves. Small amines are very soluble in water (hydrogen bonding with water molecules). Lower amines have a characteristic \emph{fishy}, ammoniacal smell. Amines are used in the manufacture of drugs (many antimalarials used in Cameroon contain amine groups), dyes, and agricultural chemicals.

\courssub{Preparation of Amines}

\begin{propriete}[Main preparations of amines]
\begin{enumerate}
\item \textbf{From haloalkanes and ammonia} (nucleophilic substitution): heating a haloalkane with excess concentrated ammonia in a sealed tube gives a primary amine:
\[ \ensuremath{\mathrm{CH_{3}CH_{2}Br(l) + NH_{3}(conc) \rightarrow [sealed\ tube,\ heat] CH_{3}CH_{2}NH_{2}(aq) + HBr(aq)}} \]
Excess ammonia limits further substitution to secondary and tertiary amines.
\item \textbf{Reduction of nitriles}: $\ce{RCN + 4[H] -> RCH2NH2}$ ($\ce{LiAlH4}$ in dry ether or $\ce{H2}$/Ni).
\item \textbf{Reduction of amides}: $\ce{RCONH2 + 4[H] ->[LiAlH4] RCH2NH2 + H2O}$.
\item \textbf{Reduction of nitro compounds} (route to aromatic amines): $\ce{C6H5NO2 + 6[H] ->[Sn/HCl] C6H5NH2 + 2H2O}$.
\item \textbf{Hofmann degradation of amides} (chain shortening by one carbon): a primary amide treated with bromine and concentrated alkali gives a primary amine with \emph{one carbon fewer}:
\[ \ce{CH3CH2CONH2(s) + Br2(aq) + 4NaOH(aq) -> CH3CH2NH2(aq) + 2NaBr(aq) + Na2CO3(aq) + 2H2O(l)} \]
\item \textbf{Gabriel phthalimide synthesis} (preparation of pure primary aliphatic amines): potassium phthalimide reacts with a haloalkane, then hydrolysis (or hydrazinolysis) releases the primary amine.
\end{enumerate}
\end{propriete}

\courssub{Reactions of Amines: Basicity, Salts, Acylation and Nitrous Acid}

The chemistry of amines is dominated by the \cle{lone pair} on nitrogen, which can accept a proton (Brønsted base) or attack electron-deficient centres (nucleophile).

\begin{propriete}[Basicity of amines]
Amines are weak bases: $\ce{RNH2 + H2O <=> RNH3+ + OH-}$. The order of base strength is:
\[ \text{aliphatic amines} > \text{ammonia} > \text{aromatic amines} \]
For example $\mathrm{p}K_b$ (at $\SI{25}{\ensuremath{{}^{\circ}\mathrm{C}}}$): ethylamine $\approx 3.3$, ammonia $\approx 4.7$, phenylamine $\approx 9.4$ (the \emph{smaller} the $\mathrm{p}K_b$, the \emph{stronger} the base).
\textbf{Explanation:}
\begin{itemize}
\item Alkyl groups are \emph{electron-donating} (positive inductive effect, $+I$): they increase the electron density on nitrogen, making the lone pair more available, and they stabilise the cation $\ce{RNH3+}$. Hence aliphatic amines are stronger bases than $\ce{NH3}$.
\item In phenylamine, the lone pair on nitrogen is \emph{delocalised into the benzene ring} by resonance; it is therefore much less available to accept a proton. Hence phenylamine is a far weaker base than ammonia.
\end{itemize}
\end{propriete}

\begin{popfigure}
\begin{center}
\begin{tikzpicture}
\begin{axis}[
    ybar, bar width=22pt, width=0.75\linewidth, height=5.5cm,
    symbolic x coords={Ethylamine, Ammonia, Phenylamine},
    xtick=data, ymin=0, ymax=11,
    ylabel={$\mathrm{p}K_b$},
    nodes near coords, every node near coord/.append style={font=\small},
    axis lines=left, tick style={draw=none},
    ymajorgrids, grid style={popline, dashed},
]
\addplot[fill=popblue, draw=popdark] coordinates {(Ethylamine,3.3) (Ammonia,4.7) (Phenylamine,9.4)};
\end{axis}
\end{tikzpicture}
\end{center}
\legende{Approximate $\mathrm{p}K_b$ values: smaller $\mathrm{p}K_b$ means a stronger base. Ethylamine is the strongest base, phenylamine the weakest.}
\end{popfigure}

\textbf{Salt formation.} Amines react with acids to form ammonium-type salts, soluble in water:
\[ \ce{CH3CH2NH2(aq) + HCl(aq) -> CH3CH2NH3+Cl-(aq)} \quad \text{(ethylammonium chloride)} \]
This reaction is used to distinguish and purify amines: the free amine is regenerated by adding alkali ($\ce{NaOH}$).

\textbf{Acylation.} Primary and secondary amines react vigorously with acyl chlorides at room temperature to give \emph{N-substituted amides}:
\[ \ce{CH3COCl(l) + C6H5NH2(aq) -> CH3CONHC6H5(s) + HCl(aq)} \]
The product here is \emph{N}-phenylethanamide. Tertiary amines, having no $\ce{N-H}$ hydrogen, cannot be acylated (they form salts with acyl chlorides).

\courssub{Reaction with Nitrous Acid and Benzene Diazonium Chloride}

Nitrous acid, $\ce{HNO2}$, is unstable, so it is generated \emph{in situ} from sodium nitrite and dilute hydrochloric acid: $\ce{NaNO2(aq) + HCl(aq) -> HNO2(aq) + NaCl(aq)}$. Its reaction with primary amines depends crucially on whether the amine is aliphatic or aromatic.

\begin{propriete}[Action of nitrous acid on primary amines]
\begin{itemize}
\item \textbf{Primary aliphatic amine} (room temperature): an alcohol is formed with effervescence of nitrogen gas:
\[ \ce{CH3CH2NH2(aq) + HNO2(aq) -> CH3CH2OH(aq) + N2(g) ^ + H2O(l)} \]
The rapid evolution of $\ce{N2}$ is a useful test for a primary aliphatic amine.
\item \textbf{Primary aromatic amine} (at $0$--$\SI{5}{\ensuremath{{}^{\circ}\mathrm{C}}}$): a \cle{diazonium salt} is formed, with \emph{no} nitrogen evolved:
\[ \ensuremath{\mathrm{C_{6}H_{5}NH_{2}(aq) + HNO_{2}(aq) + HCl(aq) \rightarrow [0-5\ ^{\circ}C] C_{6}H_{5}N_{2}+Cl^{-}(aq) + 2H_{2}O(l)}} \]
\end{itemize}
\end{propriete}

\begin{attention}[Phenylamine and nitrous acid]
Phenylamine does \emph{not} give $\ce{N2}$ with $\ce{HNO2}$ at low temperature: it forms the diazonium salt instead. Nitrogen gas is only evolved if the solution is warmed, when the diazonium salt decomposes to phenol. Examiners love this distinction!
\end{attention}

\begin{methode}[Diazotisation of phenylamine]
\begin{enumerate}
\item Dissolve phenylamine in excess dilute hydrochloric acid and cool the mixture in ice to below $\SI{5}{\ensuremath{{}^{\circ}\mathrm{C}}}$.
\item Prepare a cold solution of sodium nitrite, $\ce{NaNO2}$, in water.
\item Add the $\ce{NaNO2}$ solution slowly, keeping the temperature strictly between $0$ and $\SI{5}{\ensuremath{{}^{\circ}\mathrm{C}}}$.
\item Use the resulting benzene diazonium chloride solution immediately, still cold.
\end{enumerate}
The low temperature is essential: above about $\SI{10}{\ensuremath{{}^{\circ}\mathrm{C}}}$ the diazonium salt decomposes, liberating $\ce{N2}$ and forming phenol.
\end{methode}

\textbf{Reactions of benzene diazonium chloride.} The diazonium group $\ce{-N2+}$ is a superb leaving group (it leaves as stable $\ce{N2}$), which makes diazonium salts extraordinarily versatile.

\begin{propriete}[Replacement and coupling reactions of diazonium salts]
\textbf{Replacement reactions} (the $\ce{N2+}$ group is replaced):
\begin{itemize}
\item \textbf{Sandmeyer reactions}: with $\ce{CuCl}$/HCl $\rightarrow$ chlorobenzene; with $\ce{CuBr}$/HBr $\rightarrow$ bromobenzene; with $\ce{CuCN}$/KCN $\rightarrow$ benzonitrile $\ce{C6H5CN}$.
\item With warm \textbf{water}: phenol is formed, $\ce{C6H5N2+ + H2O -> C6H5OH + N2 + H+}$.
\item With \textbf{hypophosphorous acid} $\ce{H3PO2}$: the group is replaced by H, giving benzene (a way to remove an amino group from a ring).
\end{itemize}
\textbf{Coupling reactions} (the $\ce{N2+}$ group is retained): with alkaline phenol or with aromatic amines, the diazonium salt couples at the \emph{para} position (or ortho if para blocked) to give an \cle{azo compound} containing $\ce{-N=N-}$:
\[ \ensuremath{\mathrm{C_{6}H_{5}N_{2}+Cl^{-}(aq) + C_{6}H_{5}OH(aq) \rightarrow [NaOH,\ 0-5\ ^{\circ}C] p\text{-}HO-C_{6}H_{4}-N=N-C_{6}H_{5}(s) + HCl(aq)}} \]
Azo compounds are intensely coloured: they are the \emph{azo dyes}.
\end{propriete}

\begin{attention}[Exam tip: the orange dye]
Adding a cold diazonium salt solution to phenol in alkaline solution produces an \textbf{orange precipitate} of an azo dye (4-hydroxyazobenzene). This is a classic GCE observation question and also a confirmatory test for phenol.
\end{attention}

\begin{popfigure}
\begin{center}
\begin{tikzpicture}[
  box/.style={draw=popdark, fill=popblueL, rounded corners=2pt, align=center, font=\small, minimum height=0.9cm},
  dyebox/.style={draw=popdark, fill=popgreenL, rounded corners=2pt, align=center, font=\small, minimum height=0.9cm},
  arr/.style={-{Stealth[length=2.5mm]}, thick, popdark},
  lbl/.style={font=\scriptsize, align=center, text=popdark}]
\node[box, align=center] (bz) at (0,0){Benzene\\ $\ce{C6H6}$};
\node[box, align=center] (nb) at (3.6,0){Nitrobenzene\\ $\ce{C6H5NO2}$};
\node[box, align=center] (an) at (7.4,0){Phenylamine\\ $\ce{C6H5NH2}$};
\node[box, align=center] (dz) at (7.4,-2.6){Diazonium salt\\ $\ce{C6H5N2+Cl-}$};
\node[dyebox, align=center] (azo) at (3.4,-2.6){Azo dye\\ $\ce{C6H5N=NC6H4OH}$};
\node[dyebox, align=center] (ph) at (3.4,-5.0){Phenol\\ $\ce{C6H5OH}$};
\node[dyebox, align=center] (cb) at (7.4,-5.0){Chlorobenzene\\ $\ce{C6H5Cl}$};
\draw[arr] (bz) -- (nb) node[midway, above, lbl, align=center]{$\ce{HNO3}$/$\ce{H2SO4}$\\ $50$--$\SI{60}{\ensuremath{{}^{\circ}\mathrm{C}}}$};
\draw[arr] (nb) -- (an) node[midway, above, lbl, align=center]{Sn/conc.\ HCl\\ then $\ce{NaOH}$};
\draw[arr] (an) -- (dz) node[midway, right, lbl, align=center]{$\ce{NaNO2}$/HCl\\ $0$--$\SI{5}{\ensuremath{{}^{\circ}\mathrm{C}}}$};
\draw[arr] (dz) -- (azo) node[midway, above, lbl, align=center]{alkaline phenol\\ coupling};
\draw[arr] (dz) -- (ph) node[midway, right, lbl, xshift=-6mm] {warm $\ce{H2O}$};
\draw[arr] (dz) -- (cb) node[midway, right, lbl, align=center]{$\ce{CuCl}$/HCl\\ Sandmeyer};
\end{tikzpicture}
\end{center}
\legende{Synthesis web from benzene: nitration, reduction, diazotisation, then coupling or replacement reactions of the diazonium salt.}
\end{popfigure}

\courssub{Amino Acids: Structure, Classification, Isomerism and Preparation}

\cle{Amino acids} are the building blocks of proteins. Every protein in the beans, groundnuts and fish we eat is a long chain of amino acids.

\begin{definition}[$\alpha$-Amino acid]
An $\alpha$-amino acid has an amino group and a carboxyl group attached to the \emph{same} carbon atom (the $\alpha$-carbon):
\[ \ce{H2N-CHR-COOH} \]
where R is a side chain (H in glycine, $\ce{CH3}$ in alanine, etc.).
\end{definition}

\textbf{Classification.} By side chain: \emph{neutral} amino acids (one $\ce{-NH2}$, one $\ce{-COOH}$, e.g. glycine, alanine, valine), \emph{acidic} (an extra $\ce{-COOH}$, e.g. aspartic acid, glutamic acid), \emph{basic} (an extra $\ce{-NH2}$, e.g. lysine, arginine). \emph{Essential} amino acids cannot be synthesised by the human body and must come from food; \emph{non-essential} ones can be made by the body. They are usually known by their common names (glycine: aminoethanoic acid; alanine: 2-aminopropanoic acid).

\textbf{Optical isomerism.} The $\alpha$-carbon of $\ce{H2N-CHR-COOH}$ carries four different groups whenever $\ce{R} \neq \ce{H}$: it is a \cle{chiral centre}, so the amino acid exists as two non-superimposable mirror images (enantiomers). Naturally occurring amino acids are almost exclusively the \emph{L}-enantiomers.

\begin{attention}[Glycine, the exception]
Glycine, $\ce{H2N-CH2-COOH}$, is the \emph{only} common amino acid without optical isomerism, because its $\alpha$-carbon bears two hydrogen atoms and is therefore not chiral.
\end{attention}

\textbf{Preparation of amino acids.}
\begin{itemize}
\item \textbf{Strecker synthesis}: An aldehyde reacts with $\ce{HCN}$ and $\ce{NH3}$ to give an $\alpha$-aminonitrile, which is hydrolysed to the amino acid.
\[ \ce{RCHO + HCN + NH3 -> RCH(NH2)CN ->[H2O/H+] RCH(NH2)COOH} \]
\item \textbf{From $\alpha$-haloacids}: Nucleophilic substitution with ammonia.
\[ \ce{RCH(Cl)COOH + 2NH3 -> RCH(NH2)COOH + NH4Cl} \]
\item \textbf{Reduction of $\alpha$-oxoacids} (reductive amination): $\ce{RCOCOOH + NH3 + H2 ->[Ni] RCH(NH2)COOH + H2O}$.
\end{itemize}

\courssub{Reactions of Amino Acids: Zwitterions, Amphoterism and Peptide Bonds}

\begin{definition}[Zwitterion and amphoteric character]
In the solid state and in neutral aqueous solution, an amino acid does not exist as $\ce{H2N-CHR-COOH}$ but as a \cle{zwitterion} (dipolar ion), formed by internal proton transfer:
\[ \ce{H2N-CHR-COOH <=> H3N+-CHR-COO-} \]
Because it contains both a basic group ($\ce{-NH2}$/$\ce{-COO-}$) and an acidic group ($\ce{-COOH}$/$\ce{-NH3+}$), an amino acid is \cle{amphoteric}: it reacts with acids (accepting a proton on $\ce{-COO-}$) and with alkalis (losing a proton from $\ce{-NH3+}$). The pH at which the amino acid carries no net charge is its \emph{isoelectric point}.
\end{definition}

The zwitterion structure explains why amino acids are crystalline solids with high melting points and good solubility in water — properties of \emph{salts}, not of typical organic molecules.

\begin{popfigure}
\begin{center}
\begin{tikzpicture}[font=\small, every node/.style={align=center}]
\node (zw) at (0,0) {$\ensuremath{\mathrm{^{+}H_{3}N-\underset{\displaystyle\underset{\displaystyle R}{|}}{CH}-C(=O)-O^{-}}}$};
\node[below=2mm of zw, text=popmuted, font=\scriptsize] {zwitterion (neutral solution)};
\node (acid) at (-5.2,0) {$\ce{^{+}H3N-CHR-COOH}$};
\node[below=2mm of acid, text=popmuted, font=\scriptsize] {cation (acidic solution)};
\node (alk) at (5.2,0) {$\ce{H2N-CHR-COO^{-}}$};
\node[below=2mm of alk, text=popmuted, font=\scriptsize] {anion (alkaline solution)};
\draw[-{Stealth}, thick, popdark] (zw) -- (acid) node[midway, above, font=\scriptsize] {$\ce{H+}$};
\draw[-{Stealth}, thick, popdark] (zw) -- (alk) node[midway, above, font=\scriptsize] {$\ce{OH-}$};
\end{tikzpicture}
\end{center}
\legende{Amphoteric behaviour of an amino acid: the zwitterion accepts a proton in acid and donates one in alkali.}
\end{popfigure}

\textbf{Other reactions of amino acids.}
\begin{itemize}
\item \textbf{Reaction with nitrous acid}: $\alpha$-amino acids react with $\ce{HNO2}$ to give the corresponding $\alpha$-hydroxy acid with evolution of $\ce{N2}$:
\[ \ce{RCH(NH2)COOH + HNO2 -> RCH(OH)COOH + N2 ^ + H2O} \]
\item \textbf{Decarboxylation}: Heating with ninhydrin or enzymes yields amines and $\ce{CO2}$.
\item \textbf{Ninhydrin test}: Amino acids react with ninhydrin to give an intense purple-blue colour (Ruhemann's purple). This sensitive test is used in chromatography to locate amino acids — and by forensic scientists to reveal fingerprints, since sweat contains traces of amino acids!
\end{itemize}

\textbf{Peptide bond formation.} Two amino acids condense together: the $\ce{-COOH}$ of one reacts with the $\ce{-NH2}$ of the other, eliminating water and forming an amide linkage called a \cle{peptide bond}:
\[ \ensuremath{\mathrm{H_{2}N-CHR-COOH + H_{2}N-CHR^\prime-COOH \rightarrow  H_{2}N-CHR-CO-NH-CHR^\prime-COOH + H_{2}O}} \]
The product is a \emph{dipeptide}; it still has a free $\ce{-NH2}$ and a free $\ce{-COOH}$, so the process can continue to give a \emph{polypeptide}. \textbf{Proteins are polyamides} (polypeptides) built from $\alpha$-amino acids joined by peptide bonds. Hydrolysis of proteins (acid or enzyme) yields the constituent amino acids.

\begin{popfigure}
\begin{center}
\begin{tikzpicture}[font=\small, every node/.style={align=center}]
\node (a1) at (0,0) {$\mathrm{H_2N{-}CHR{-}COOH}$};
\node at (2.1,0) {$+$};
\node (a2) at (4.3,0) {$\mathrm{H_2N{-}CHR'{-}COOH}$};
\draw[-{Stealth[length=3mm]}, very thick, poporange] (5.9,0) -- (7.3,0) node[midway, above, font=\scriptsize, text=popdark] {$-\ce{H2O}$};

% Build dipeptide as chained nodes to mark peptide bond parts
\node[anchor=west] (dp1) at (8.5,0) {$\mathrm{H_2N{-}CHR{-}}$};
\node[anchor=west] (co) at (dp1.east) {$\mathrm{CO}$};
\node[anchor=west] (nh) at (co.east) {$\mathrm{{-}NH{-}}$};
\node[anchor=west] (dp2) at (nh.east) {$\mathrm{CHR'{-}COOH}$};

% Rectangle around peptide bond (CO-NH)
\draw[popblue, thick] ([xshift=-2pt,yshift=-3pt]co.south west) rectangle ([xshift=2pt,yshift=3pt]nh.north east);
\node[text=popblue, font=\scriptsize] at ($(co)!0.5!(nh)$) [yshift=-12pt] {peptide bond $\mathrm{{-}CO{-}NH{-}}$};
\end{tikzpicture}
\end{center}
\legende{Condensation of two amino acids: elimination of water creates the peptide (amide) bond of a dipeptide.}
\end{popfigure}

\begin{savaistu}[Ninhydrin and fingerprints (enrichment)]
Amino acids react with \emph{ninhydrin} to give an intense purple-blue colour. This sensitive test is used in chromatography to locate amino acids — and by forensic scientists to reveal fingerprints, since sweat contains traces of amino acids!
\end{savaistu}

\begin{exoresolu}[From benzene to an azo dye]
Describe, with equations and conditions, how benzene can be converted into an orange azo dye via phenylamine.
\tcblower
\textbf{Step 1 — Nitration.} Benzene is treated with concentrated $\ce{HNO3}$ and concentrated $\ce{H2SO4}$ at $50$--$\SI{60}{\ensuremath{{}^{\circ}\mathrm{C}}}$:
\[ \ensuremath{\mathrm{C_{6}H_{6}(l) + HNO_{3}(conc) \rightarrow [conc.\ H_{2}SO_{4}] C_{6}H_{5}NO_{2}(l) + H_{2}O(l)}} \]
\textbf{Step 2 — Reduction.} Nitrobenzene is reduced with tin and concentrated HCl, then the mixture is made alkaline to free the amine:
\[ \ce{C6H5NO2(l) + 6[H] ->[Sn/HCl] C6H5NH2(aq) + 2H2O(l)} \]
\textbf{Step 3 — Diazotisation.} Phenylamine in cold dilute HCl is treated with $\ce{NaNO2}$ at $0$--$\SI{5}{\ensuremath{{}^{\circ}\mathrm{C}}}$:
\[ \ensuremath{\mathrm{C_{6}H_{5}NH_{2}(aq) + HNO_{2}(aq) + HCl(aq) \rightarrow [0-5\ ^{\circ}C] C_{6}H_{5}N_{2}+Cl^{-}(aq) + 2H_{2}O(l)}} \]
\textbf{Step 4 — Coupling.} The cold diazonium solution is added to phenol in sodium hydroxide solution. An orange azo dye precipitates:
\[ \ensuremath{\mathrm{C_{6}H_{5}N_{2}+Cl^{-}(aq) + C_{6}H_{5}OH(aq) \rightarrow [NaOH] p\text{-}HO-C_{6}H_{4}-N=N-C_{6}H_{5}(s) + HCl(aq)}} \]
\end{exoresolu}

\begin{exercice}[Amines and amino acids]
\begin{enumerate}
\item Arrange ethylamine, ammonia and phenylamine in order of increasing basic strength and justify the order.
\item Write the equation for the reaction of propanamide with $\ce{Br2}$/NaOH (Hofmann degradation) and name the product.
\item Explain why alanine shows optical isomerism but glycine does not.
\item Write the equation for the formation of the dipeptide glycylalanine from glycine ($\ce{R = H}$) and alanine ($\ce{R = CH3}$).
\item Outline the Strecker synthesis of alanine from ethanal.
\end{enumerate}
\end{exercice}

\begin{corrige}[Exercise 1]
\begin{enumerate}
\item Phenylamine $<$ ammonia $<$ ethylamine. The ethyl group donates electrons (positive inductive effect), making the lone pair on N more available than in $\ce{NH3}$. In phenylamine the lone pair is delocalised into the benzene ring by resonance, so it is much less available: phenylamine is the weakest base.
\item $\ce{CH3CH2CONH2(s) + Br2(aq) + 4NaOH(aq) -> CH3CH2NH2(aq) + 2NaBr(aq) + Na2CO3(aq) + 2H2O(l)}$. The product is ethylamine (one carbon fewer than the amide).
\item In alanine, $\ce{H2N-CH(CH3)-COOH}$, the $\alpha$-carbon bears four different groups ($\ce{H}$, $\ce{CH3}$, $\ce{NH2}$, $\ce{COOH}$): it is chiral, so two enantiomers exist. In glycine, $\ce{H2N-CH2-COOH}$, the $\alpha$-carbon bears two identical H atoms: no chirality, no optical isomerism.
\item $\ce{H2N-CH2-COOH + H2N-CH(CH3)-COOH -> H2N-CH2-CO-NH-CH(CH3)-COOH + H2O}$.
\item Ethanal + $\ce{HCN}$ + $\ce{NH3}$ $\rightarrow$ 2-aminopropanenitrile $\rightarrow$ hydrolysis $\rightarrow$ alanine ($\ce{CH3CH(NH2)COOH}$).
\end{enumerate}
\end{corrige}

\begin{aretenir}[Key points of this section]
\begin{itemize}
\item Nitrobenzene (from nitration of benzene) is reduced by Sn/HCl to phenylamine; nitriles $\ce{RCN}$ come from haloalkanes + KCN or amide dehydration, and are hydrolysed to acids or reduced to amines.
\item Amines are classified $1^{\circ}$, $2^{\circ}$, $3^{\circ}$ by the number of carbon groups on nitrogen. Basicity: aliphatic $>$ $\ce{NH3}$ $>$ aromatic (inductive effect vs resonance delocalisation).
\item Primary aliphatic amines + $\ce{HNO2}$ give alcohols + $\ce{N2}$; phenylamine at $0$--$\SI{5}{\ensuremath{{}^{\circ}\mathrm{C}}}$ gives benzene diazonium chloride.
\item Diazonium salts undergo Sandmeyer replacements (Cl, Br, CN), hydrolysis to phenol, and azo coupling with alkaline phenol to give an \textbf{orange dye}.
\item Amino acids $\ce{H2N-CHR-COOH}$ are amphoteric zwitterions; all are optically active except glycine; they condense through peptide bonds to form proteins (polyamides). Preparation: Strecker synthesis, from $\alpha$-haloacids, reductive amination.
\end{itemize}
\end{aretenir}

\newpage

\coursec{Integration activity}

\courssub{Aims of the activity}
This integration activity is designed to consolidate your understanding of the major organic functional families studied in this chapter --- alcohols, aldehydes, ketones, carboxylic acids, phenols, esters and related derivatives --- by placing you in a realistic investigative context. You will:
\begin{itemize}
    \item Apply IUPAC nomenclature and structural knowledge to identify unknown organic liquids.
    \item Design an economical and safe sequence of qualitative distinguishing tests, predicting observations for each functional group.
    \item Write balanced chemical equations for the biochemical and chemical transformations involved in the spoilage of palm wine and the action of a paint solvent.
    \item Communicate your reasoning clearly in a poster format and critically evaluate the test strategies of other groups.
\end{itemize}
The activity draws on practical skills (test-tube reactions, observation recording), theoretical ideas (nucleophilic addition, oxidation, esterification, hydrolysis) and the Cameroonian context of traditional palm wine production and local workshop practices.

\courssub{Situation}
In the grassfields of the Northwest Region, a palm wine tapper from \emph{Ndop} delivers fresh palm sap to a local \emph{buvette} early on Monday morning. The sap is sweet, slightly fizzy and contains about $4\%$ ethanol by volume due to natural fermentation by wild yeasts (\emph{Saccharomyces} spp.) and lactic acid bacteria. By Wednesday, the same palm wine has turned distinctly sour, with a sharp vinegar-like smell and a pH around 3.5. The regular customers complain that the drink is ``spoilt''.

Meanwhile, a nearby furniture workshop uses a colourless, volatile liquid stored in an unlabelled \SI{250}{cm^3} bottle to clean paint brushes and remove dried paint splashes from wood. The workshop owner says the liquid ``smells like nail varnish remover'' and evaporates quickly.

During a routine inspection by the school's chemistry club, four unlabelled glass bottles (A, B, C, D) are discovered in a cupboard adjacent to the workshop. The club suspects they contain samples related to the palm wine spoilage and the paint solvent. The bottles are known to contain (in no particular order):
\begin{enumerate}
    \item Ethanol (\ce{CH3CH2OH})
    \item Ethanal (\ce{CH3CHO})
    \item Ethanoic acid (\ce{CH3COOH})
    \item Aqueous phenol solution (\ce{C6H5OH}, approximately $5\%$ w/v)
\end{enumerate}
No other chemicals are present. The club has access to standard Upper Sixth laboratory reagents: solid \ce{NaHCO3}, Fehling's solutions A and B, Tollens' reagent (ammoniacal \ce{AgNO3}), neutral \ce{FeCl3} solution, iodine in \ce{KI} solution (for the iodoform test), acidified potassium dichromate(VI) (\ce{K2Cr2O7}/\ce{H2SO4}), 2,4-dinitrophenylhydrazine (2,4-DNPH) reagent, dilute \ce{NaOH} and \ce{H2SO4}, and standard glassware. Safety equipment (fume cupboard, gloves, goggles) is available.

\courssub{Tasks}
Work in groups of three to four. Answer the following questions in order, recording your reasoning, equations and expected observations. Prepare a poster summarising your identification scheme for presentation and peer critique.

\begin{enumerate}
    \item \textbf{Biochemical pathway of palm wine spoilage.}
    \begin{enumerate}
        \item Write the balanced equation for the anaerobic conversion of glucose (\ce{C6H12O6}) to ethanol and carbon dioxide by yeast (zymase). State the conditions.
        \item Write the balanced equations for the aerobic oxidation of ethanol to ethanoic acid by acetic acid bacteria (\emph{Acetobacter}). Indicate the intermediate and the role of oxygen.
        \item Explain why the palm wine becomes sour over two days, linking the microbial processes to the observed pH change.
    \end{enumerate}

    \item \textbf{Design of a minimal distinguishing test sequence.}
    \begin{enumerate}
        \item Propose a logical order of \emph{at most five} tests from the available reagents that will unambiguously identify each of the four bottles (A--D). Justify the order (e.g., start with a test that splits the set into two groups).
        \item For each test, draw a table with columns: \textbf{Bottle}, \textbf{Reagent added}, \textbf{Observation}, \textbf{Inference}. Predict the exact observation for each bottle (e.g., ``effervescence, colourless gas'', ``brick-red precipitate'', ``violet colouration'', ``yellow precipitate'', ``green solution'').
        \item Explain how you would distinguish between ethanol and ethanal using the iodoform test and the acidified dichromate(VI) test, noting any overlap.
    \end{enumerate}

    \item \textbf{Identification of the paint solvent.}
    \begin{enumerate}
        \item The workshop owner mentions the solvent smells like nail varnish remover. Suggest two likely organic compounds (one ketone, one ester) that fit this description and are common paint solvents. Give their structural formulae and IUPAC names.
        \item Describe how you would use 2,4-DNPH to confirm the presence of a carbonyl group. Write the equation for the reaction of your chosen ketone with 2,4-DNPH.
        \item If the solvent is an ester, outline a hydrolysis procedure (reagents, conditions) to break it into its parent acid and alcohol. How would you then identify the hydrolysis products using tests from Task 2?
    \end{enumerate}

    \item \textbf{Poster presentation and peer critique.}
    \begin{enumerate}
        \item Prepare an A2 poster (or digital slide) that includes: the biochemical pathway (Task 1), the test sequence table (Task 2), the solvent identification logic (Task 3), and a safety note (hazards of Tollens' reagent, dichromate(VI), phenol).
        \item During the gallery walk, each group writes one constructive comment on another group's poster: either a suggestion to reduce the number of tests, a safety improvement, or a correction of an erroneous observation.
    \end{enumerate}
\end{enumerate}

\courssub{Model answer}
\textbf{Task 1 --- Biochemical pathway.}\par
\begin{enumerate}
    \item \textbf{Alcoholic fermentation (anaerobic):}
    \[\ce{C6H12O6(aq) ->[zymase (yeast)] 2 C2H5OH(aq) + 2 CO2(g)}\]
    Conditions: warm temperature (about $30$--$35\ ^\circ\mathrm{C}$), absence of oxygen, slightly acidic medium (pH 4--5).
    \item \textbf{Aerobic oxidation to vinegar (two steps):}
    \begin{align*}
    \ce{2 CH3CH2OH(aq) + O2(g) &->[Acetobacter] 2 CH3CHO(aq) + 2 H2O(l)} && \text{(ethanol to ethanal)} \\
    \ce{2 CH3CHO(aq) + O2(g) &->[Acetobacter] 2 CH3COOH(aq)} && \text{(ethanal to ethanoic acid)}
    \end{align*}
    Overall: \[\ce{CH3CH2OH(aq) + O2(g) ->[Acetobacter] CH3COOH(aq) + H2O(l)}\]
    Oxygen from the air is the oxidising agent (terminal electron acceptor); the bacteria use the energy released for their metabolism. Ethanal is the intermediate; it has the sharp smell noticed before the full sour taste develops.
    \item \textbf{Explanation:} Fresh palm sap contains sucrose and glucose. Yeasts ferment the sugars to ethanol (about $4\%$ v/v) and \ce{CO2}, which causes the fizz. When the wine is left exposed to air, \emph{Acetobacter} present on the gourd surface and in the air oxidise the ethanol first to ethanal, then to ethanoic acid. The accumulation of ethanoic acid, a weak acid that partially ionises as \ce{CH3COOH(aq) + H2O(l) <=> CH3COO-(aq) + H3O+(aq)}, lowers the pH to about 3.5 and gives the sour, vinegar-like taste. At ambient tropical temperatures ($25$--$30\ ^\circ\mathrm{C}$) the process takes 24--48 hours.
\end{enumerate}

\textbf{Task 2 --- Minimal distinguishing test sequence.}\par
\textbf{Proposed order (three decisive tests, plus one optional confirmation):}
\begin{enumerate}
    \item \textbf{Solid \ce{NaHCO3} --- tests for a carboxylic acid.} Only ethanoic acid effervesces, liberating \ce{CO2}. This splits the set: \fbox{C = ethanoic acid} versus \{A, B, D\}. Phenol does \emph{not} react: it is a weaker acid than carbonic acid, so it cannot liberate \ce{CO2} from a hydrogencarbonate.
    \item \textbf{Neutral \ce{FeCl3}(aq) --- tests for phenol.} Only phenol gives an intense violet colouration (formation of an iron(III)--phenoxide complex). \fbox{D = phenol}. Remaining \{A, B\} = ethanol and ethanal.
    \item \textbf{Fehling's test (or Tollens') --- tests for an aldehyde.} On warming, ethanal gives a brick-red precipitate of \ce{Cu2O} with Fehling's (or a silver mirror with Tollens'); ethanol does not react. \fbox{B = ethanal}, \fbox{A = ethanol}.
    \item \textbf{Confirmation (optional): iodoform test.} Both ethanol and ethanal give a yellow precipitate of \ce{CHI3}, because ethanol contains the \ce{CH3-CH(OH)-} unit (oxidised in situ to \ce{CH3CHO} by \ce{I2}/\ce{OH-}) and ethanal contains the \ce{CH3CO-} unit. This test therefore confirms the identities but cannot discriminate between A and B on its own.
\end{enumerate}

\textbf{Results table:}

\begin{center}
\begin{tabularx}{\linewidth}{|l|X|X|X|}
\hline
\rowcolor{popblueL}
\textbf{Bottle} & \textbf{Reagent added} & \textbf{Observation} & \textbf{Inference} \\
\hline
A (Ethanol) & \ce{NaHCO3} (s) & No visible reaction & Not a carboxylic acid \\
& \ce{FeCl3} (aq) & No violet colour (solution stays pale yellow) & Not phenol \\
& Fehling's (warm) & Blue solution remains, no precipitate & Not an aldehyde \\
& Iodoform (\ce{I2}/\ce{KI} + \ce{NaOH}, warm) & Yellow crystalline precipitate, antiseptic smell & \ce{CH3CH(OH)-} unit present \\
& Acidified \ce{K2Cr2O7} (warm) & Orange solution turns green & Primary alcohol oxidised \\
\hline
B (Ethanal) & \ce{NaHCO3} (s) & No visible reaction & Not a carboxylic acid \\
& \ce{FeCl3} (aq) & No violet colour & Not phenol \\
& Fehling's (warm) & Brick-red precipitate of \ce{Cu2O} forms & Aldehyde present \\
& Iodoform (\ce{I2}/\ce{KI} + \ce{NaOH}, warm) & Yellow crystalline precipitate & \ce{CH3CO-} group present \\
& Acidified \ce{K2Cr2O7} (warm) & Orange solution turns green & Aldehyde oxidised to acid \\
\hline
C (Ethanoic acid) & \ce{NaHCO3} (s) & \textbf{Effervescence; colourless gas (\ce{CO2}) turns limewater milky} & \textbf{Carboxylic acid confirmed} \\
& \ce{FeCl3} (aq) & No violet colour & Not phenol \\
& Fehling's (warm) & Blue solution remains & Not an aldehyde \\
& Iodoform & No precipitate & No \ce{CH3CO-} or \ce{CH3CH(OH)-} unit \\
& Acidified \ce{K2Cr2O7} (warm) & No change (solution stays orange) & Not oxidised \\
\hline
D (Phenol) & \ce{NaHCO3} (s) & No visible reaction (phenol is too weak an acid) & Not a carboxylic acid \\
& \ce{FeCl3} (aq) & \textbf{Intense violet colouration} & \textbf{Phenol confirmed} \\
& Fehling's (warm) & Blue solution remains & Not an aldehyde \\
& Iodoform & No precipitate & Negative \\
& Acidified \ce{K2Cr2O7} (warm) & No clean colour change under these conditions & Negative \\
\hline
\end{tabularx}
\end{center}

\textbf{Distinguishing ethanol from ethanal:} Both give a positive iodoform test (yellow \ce{CHI3}) and both turn acidified dichromate(VI) from orange to green, so neither of these tests can separate them. The decisive test is Fehling's (or Tollens'): ethanal, an aldehyde, reduces \ce{Cu^2+} to \ce{Cu2O} (brick-red precipitate) or \ce{Ag+} to Ag(s) (silver mirror), whereas ethanol, a primary alcohol, does not reduce these mild oxidising agents. Fehling's test is therefore the discriminator.

\textbf{Task 3 --- Paint solvent identification.}\par
\begin{enumerate}
    \item \textbf{Likely solvents:}
    \begin{itemize}
        \item \textbf{Propanone (acetone)} --- \ce{CH3COCH3}, the simplest ketone; the classic nail varnish remover and brush cleaner.
        \item \textbf{Ethyl ethanoate (ethyl acetate)} --- \ce{CH3COOCH2CH3}, a common ester solvent for paints, varnishes and adhesives, with a sweet fruity smell.
    \end{itemize}
    \item \textbf{2,4-DNPH test:} Add a few drops of the solvent to about \SI{2}{cm^3} of 2,4-DNPH reagent (Brady's reagent). A yellow-orange crystalline precipitate of the 2,4-dinitrophenylhydrazone forms within minutes, confirming a carbonyl group (aldehyde or ketone). Alcohols, acids, esters and phenol give no precipitate.
    \begin{center}
    \ce{CH3COCH3 + (NO2)2C6H3NHNH2 -> CH3C(=NNHC6H3(NO2)2)CH3 + H2O}
    \end{center}
    (Propanone 2,4-dinitrophenylhydrazone, yellow-orange precipitate.) This is a condensation (nucleophilic addition followed by elimination of water).
    \item \textbf{If the solvent is the ester (ethyl ethanoate):} hydrolyse it under reflux.
    \begin{itemize}
        \item \textbf{Alkaline hydrolysis (saponification):} reflux with \SI{1}{mol.dm^{-3}} \ce{NaOH}(aq) for about 30 minutes:
        \[\ce{CH3COOCH2CH3(l) + NaOH(aq) -> CH3COONa(aq) + CH3CH2OH(aq)}\]
        After cooling, acidify with dilute \ce{H2SO4} to liberate ethanoic acid from the sodium ethanoate. Test the mixture: effervescence with \ce{NaHCO3} confirms the carboxylic acid; distil and test the distillate --- positive iodoform, negative Fehling's, dichromate(VI) orange to green --- confirming ethanol.
        \item \textbf{Acid hydrolysis:} reflux with dilute \ce{H2SO4} (reversible reaction):
        \[\ce{CH3COOCH2CH3(l) + H2O(l) <=>[H+] CH3COOH(aq) + CH3CH2OH(aq)}\]
        The same distinguishing tests identify the two products.
    \end{itemize}
    For propanone: 2,4-DNPH positive (carbonyl present), iodoform positive (\ce{CH3CO-} group), Fehling's negative (ketones resist mild oxidation), and no hydrolysis products (a ketone is not an ester).
\end{enumerate}

\textbf{Task 4 --- Poster and critique.}\par
\begin{itemize}
    \item \textbf{Poster layout:} Title ``The Mystery of the Spoilt Palm Wine''. Sections: (1) biochemical pathway with balanced equations and conditions; (2) flowchart of the test sequence (decision tree: \ce{NaHCO3} $\to$ \ce{FeCl3} $\to$ Fehling's); (3) results table as above; (4) solvent identification logic; (5) safety box: ``Tollens' reagent --- prepare fresh, never store (forms explosive silver compounds on standing); acidified dichromate(VI) --- toxic and oxidising, wear gloves; phenol --- toxic and corrosive, absorbed through skin, use the fume cupboard; 2,4-DNPH reagent --- flammable and acidic, keep away from flames.''
    \item \textbf{Sample peer critique:} ``Group X used Tollens' test before Fehling's --- both detect aldehydes, so one is redundant; they could drop Tollens' and save silver nitrate. They also did not mention neutralising dichromate(VI) waste before disposal.''
\end{itemize}

\begin{aretenir}[Key take-aways for the GCE exam]
\begin{itemize}
\item A logical test sequence starts with the most discriminating reagent (e.g., \ce{NaHCO3} for carboxylic acids, \ce{FeCl3} for phenols) to reduce the number of unknowns quickly.
\item The iodoform test is positive for compounds containing the \ce{CH3CO-} group (ethanal and methyl ketones) and for alcohols with the \ce{CH3CH(OH)-} unit (ethanol and secondary alcohols of that type), which are oxidised in situ. Methanol, other primary alcohols, ordinary aldehydes and ketones without a methyl group next to \ce{C=O} give a negative test.
\item Acidified dichromate(VI) oxidises primary alcohols to carboxylic acids (via aldehydes), secondary alcohols to ketones, and aldehydes to carboxylic acids --- in each case the orange solution turns green. Tertiary alcohols and ketones resist oxidation under these conditions.
\item 2,4-DNPH identifies carbonyl compounds (aldehydes and ketones) through coloured 2,4-dinitrophenylhydrazone precipitates; it does not react with alcohols, carboxylic acids, esters or amides under normal conditions.
\item Alkaline ester hydrolysis (saponification) gives a carboxylate salt and an alcohol; acid hydrolysis gives the free carboxylic acid and the alcohol. The products are identified with the same distinguishing tests.
\end{itemize}
\end{aretenir}

\begin{attention}[Common pitfalls in this activity]
\begin{itemize}
\item Forgetting state symbols in the biochemical equations (e.g., \ce{CO2(g)}), or omitting the conditions (yeast/zymase, anaerobic; \emph{Acetobacter}, aerobic).
\item Confusing the iodoform test: ethanol and ethanal both give a yellow precipitate, but for different structural reasons (ethanol is first oxidised to ethanal in situ by \ce{I2}/\ce{OH-}).
\item Writing ``effervescence'' for phenol with \ce{NaHCO3}: phenol ($\mathrm{p}K_{\mathrm{a}} \approx 10$) is a weaker acid than carbonic acid and does not liberate \ce{CO2} from hydrogencarbonates. Only carboxylic acids do.
\item Saying dichromate(VI) is ``decolourised'': the correct observation is orange turning to green (reduction of \ce{Cr2O7^2-} to \ce{Cr^3+}).
\item Omitting the condition ``warm'' for Fehling's and dichromate(VI) tests; cold reagents react slowly or not at all.
\item Safety: never store Tollens' reagent after use (explosive silver compounds form on standing) and never heat it in a sealed tube; dispose of silver residues as instructed.
\end{itemize}
\end{attention}

\begin{savaistu}[Cameroon context]
\begin{itemize}
\item Palm wine (``matango'' or ``vin de palme''), tapped from raffia or oil palms, is a traditional beverage across the forest and grassfield zones of Cameroon. Its natural fermentation involves wild yeasts such as \emph{Saccharomyces cerevisiae} together with lactic acid bacteria. The souring on standing is due to \emph{Acetobacter} converting ethanol into ethanoic acid --- essentially natural vinegar production, the same chemistry used worldwide to make vinegar from wine.
\item In local workshops (menuiserie), ``white spirit'' (a petroleum distillate) and propanone are common paint thinners, while ethyl ethanoate appears in some industrial coatings and adhesives. Recognising such solvents by smell and simple chemical tests is a genuine practical skill for quality control and for laboratory safety.
\end{itemize}
\end{savaistu}

\newpage

\coursec{Methods and worked examples}

\courssub{Method 1 --- Naming Halogeno Compounds, Alcohols and Carbonyl Compounds}

\begin{aretenir}[Priority of functional groups for naming]
The principal functional group determines the suffix. Order of priority (highest first): carboxylic acid (\emph{-oic acid}) > aldehyde (\emph{-al}) > ketone (\emph{-one}) > alcohol (\emph{-ol}) > amine (\emph{-amine}) > halogen (prefix \emph{halo-}). The chain must contain the principal group; for haloalkanes (no higher group) choose the longest chain with the maximum number of halogen substituents.
\end{aretenir}

\begin{methode}[Naming an organic compound (haloalkane, alcohol, aldehyde, ketone, acid)]
\begin{enumerate}
\item Identify the \cle{principal functional group}: it fixes the suffix (\emph{-ol} for alcohols, \emph{-al} for aldehydes, \emph{-one} for ketones, \emph{-oic acid} for carboxylic acids). Halogens are always \emph{prefixes} (fluoro-, chloro-, bromo-, iodo-).
\item Select the \cle{longest carbon chain} containing the principal functional group (or the greatest number of halogen atoms for haloalkanes).
\item Number the chain so that the principal group gets the \cle{lowest locant}. For aldehydes and carboxylic acids the functional carbon is automatically C-1. For haloalkanes, give the lowest set of locants to the substituents (lowest sum rule).
\item Name substituents in \cle{alphabetical order} (bromo before chloro before fluoro before iodo before methyl), each with its locant; use di-, tri-, tetra- for repeats (ignored for alphabetical ordering).
\item Assemble: locant-prefix + parent chain + locant-suffix, e.g. \ce{2-bromo-3-methylbutane}; \ce{butan-2-ol}; \ce{3-methylpentanal}; \ce{pentan-2-one}.
\end{enumerate}
\textbf{Reflexes:} check that every carbon has four bonds; check that the suffix carbon is inside the chosen chain; for ketones the carbonyl cannot be on a terminal carbon.
\end{methode}

\begin{attention}[Common naming pitfalls]
\begin{itemize}
\item Do not write ``2-ethylpropanal'' --- the longest chain containing \ce{-CHO} is butane, so it is \ce{3-methylbutanal}.
\item For di-/tri-substituted halogens, the multiplying prefixes (di, tri) do \emph{not} affect alphabetical order: \ce{1,2-dibromo-1-chloroethane} (bromo before chloro).
\item The locant for the functional group suffix is placed immediately before the suffix (\ce{butan-2-ol}), not before the parent name (\ce{2-butanol} is old IUPAC).
\end{itemize}
\end{attention}

\begin{exoresolu}[Naming practice (GCE Paper 2 style)]
Give the IUPAC names of the following compounds:\\
(a) \ce{CH3-CH(Br)-CH(Cl)-CH3}\quad (b) \ce{CH3-CH2-CH(OH)-CH(CH3)-CH3}\quad (c) \ce{(CH3)2CH-CH2-CHO}\quad (d) \ce{CH3-CO-CH2-CH(CH3)2}
\tcblower
\textbf{(a)} Longest chain = 4 C (butane). Numbering from either end gives locants 2 and 3; tie broken alphabetically: bromo gets the lower number. Name: \textbf{2-bromo-3-chlorobutane}.\\[2pt]
\textbf{(b)} Longest chain containing \ce{-OH} = 5 C (pentane). Number from the right so that \ce{-OH} is on C-3 either way; then the methyl must get the lowest locant: methyl on C-2. Name: \textbf{2-methylpentan-3-ol}.\\[2pt]
\textbf{(c)} The \ce{-CHO} carbon is C-1; chain of 4 C (butanal) with a methyl on C-3. Name: \textbf{3-methylbutanal}.\\[2pt]
\textbf{(d)} Longest chain containing \ce{C=O} = 5 C (pentanone). Number from the left: carbonyl on C-2, methyl on C-4. Name: \textbf{4-methylpentan-2-one}.\\[2pt]
\textbf{Common error:} writing ``2-ethylpropanal'' type names --- always re-check that the chain chosen is the \emph{longest} one containing the functional group.
\end{exoresolu}

\courssub{Method 2 --- Predicting and Writing S$_N$1 / S$_N$2 Substitutions}

\begin{popfigure}
\begin{tikzpicture}[scale=0.9, baseline=(current bounding box.center)]
% SN2 backside attack
\node[align=center] at (0,2.5) {\textbf{S$_N$2: Concerted backside attack}};
\draw[thick,->] (-2,1.5) -- (-0.5,1.5) node[midway,above] {\ce{Nu^-}};
\draw[thick] (-0.5,1.5) -- (0,1.5);
\draw[thick,dashed] (0,1.5) -- (0.5,1.5);
\node[draw,circle,minimum size=1.2cm] (C) at (0,0) {};
\node at (0,0) {\ce{C}};
\draw[thick] (-0.6,-0.6) -- (0,0) node[midway,left] {\ce{R}};
\draw[thick] (0.6,-0.6) -- (0,0) node[midway,right] {\ensuremath{\mathrm{R'}}};
\draw[thick] (0,0.6) -- (0,1.2) node[midway,right] {\ensuremath{\mathrm{R''}}};
\draw[thick,->] (0,-0.6) -- (0,-1.5) node[midway,right] {\ce{X^-}};
\draw[thick,dashed] (0,-0.6) -- (0,-1.2);
\draw[<->,thick,popblue] (-1.2,0) arc (180:270:1.2) node[midway,left=2pt] {inversion};
% SN1
\node[align=center] at (5,2.5) {\textbf{S$_N$1: Carbocation intermediate}};
\draw[thick] (4,1) -- (4.5,1);
\draw[thick,->] (4.5,1) -- (5.5,1) node[midway,above] {slow};
\node[draw,circle,minimum size=1.2cm] (C2) at (5,0) {};
\node at (5,0) {\ce{C+}};
\draw[thick] (4.4,-0.6) -- (5,0) node[midway,left] {\ce{R}};
\draw[thick] (5.6,-0.6) -- (5,0) node[midway,right] {\ensuremath{\mathrm{R'}}};
\draw[thick] (5,0.6) -- (5,1.2) node[midway,right] {\ensuremath{\mathrm{R''}}};
\draw[thick,->] (6,0) -- (7,0) node[midway,above] {fast};
\draw[thick] (7,0) -- (7.5,0) node[midway,above] {\ce{Nu^-}};
\draw[thick] (7.5,0) -- (8,0);
\node at (8.5,0) {\ensuremath{\mathrm{RR'R''C-Nu}}};
\end{tikzpicture}
\legende{Key mechanistic differences: S$_N$2 is concerted with inversion; S$_N$1 proceeds via a planar carbocation (racemisation if chiral).}
\end{popfigure}

\begin{methode}[Choosing the mechanism and writing the products of nucleophilic substitution]
\begin{enumerate}
\item Classify the haloalkane: \cle{primary} ($1^\circ$), \cle{secondary} ($2^\circ$) or \cle{tertiary} ($3^\circ$) by counting carbons directly attached to the C--X carbon.
\item Apply the rule: $1^\circ$ $\rightarrow$ \cle{S$_N$2} (one step, backside attack, inversion, rate $= k[\ce{RX}][\ce{Nu^-}]$); $3^\circ$ $\rightarrow$ \cle{S$_N$1} (two steps, carbocation intermediate, racemisation, rate $= k[\ce{RX}]$); $2^\circ$ can go either way (polar protic solvents and weak nucleophiles favour S$_N$1; polar aprotic solvents and strong nucleophiles favour S$_N$2).
\item Identify the nucleophile and the product: \ce{OH^-} (aq) $\rightarrow$ alcohol; \ce{CN^-} (ethanolic) $\rightarrow$ nitrile (chain extended by one C!); \ce{NH3} (ethanolic, excess) $\rightarrow$ primary amine; \ce{RO^-} $\rightarrow$ ether; \ce{RCOO^-} $\rightarrow$ ester.
\item Write the balanced equation with state symbols, replacing X by the nucleophile residue; the displaced halide leaves as \ce{X^-} (or \ce{HX} with \ce{NH3}).
\item For S$_N$2 draw the curly arrow from the nucleophile's lone pair to C, and from the C--X bond to X, in \emph{one} concerted step; for S$_N$1 show the slow ionisation first, then fast attack on the planar carbocation.
\end{enumerate}
\textbf{Key idea:} steric hindrance slows S$_N$2; carbocation stability ($3^\circ > 2^\circ > 1^\circ$) speeds S$_N$1.
\end{methode}

\begin{attention}[GCE traps on mechanisms]
\begin{itemize}
\item Never draw a carbocation for a primary haloalkane --- it is too unstable.
\item In S$_N$1, the nucleophile attacks the \emph{planar} carbocation from both faces; if the carbon was chiral, a racemic mixture results.
\item \ce{KCN} in ethanol gives nitriles (\ce{R-CN}); \ce{AgCN} gives isocyanides (\ce{R-NC}) --- know the difference.
\item Rate equation: S$_N$1 is first order overall; S$_N$2 is second order overall.
\end{itemize}
\end{attention}

\begin{exoresolu}[Mechanism and products]
2-bromo-2-methylpropane reacts with warm aqueous potassium hydroxide. (a) Give the product and name the mechanism. (b) Write the rate equation and explain why the rate is unchanged if the \ce{KOH} concentration is doubled. (c) State what product forms if the same haloalkane is heated with \ce{KCN} in ethanol.
\tcblower
\textbf{(a)} 2-bromo-2-methylpropane, \ce{(CH3)3CBr}, is a \emph{tertiary} haloalkane: the bulky methyl groups block backside attack and the tertiary carbocation is stabilised by three electron-releasing methyl groups. Substitution therefore goes by \textbf{S$_N$1}:
\[
\ce{(CH3)3CBr(l) ->[slow] (CH3)3C+(solv) + Br^-(solv)} \qquad \ce{(CH3)3C+(solv) + OH^-(aq) ->[fast] (CH3)3COH(l)}
\]
Product: \textbf{2-methylpropan-2-ol}.\\[3pt]
\textbf{(b)} Rate $= k\,[\ce{(CH3)3CBr}]$. The slow, rate-determining step is the ionisation of the C--Br bond, which involves \emph{only} the haloalkane; \ce{OH^-} enters only in the fast second step. Doubling $[\ce{KOH}]$ therefore leaves the rate unchanged (first order overall).\\[3pt]
\textbf{(c)} With \ce{KCN}/ethanol, \ce{CN^-} attacks through carbon (S$_N$1 because tertiary):
\[
\ce{(CH3)3CBr(l) + CN^-(eth) -> (CH3)3C-CN(l) + Br^-(eth)}
\]
Product: \textbf{2,2-dimethylpropanenitrile}. Note the carbon chain grows by one carbon --- a classic way of ascending a homologous series.
\end{exoresolu}

\courssub{Method 3 --- Distinguishing Primary, Secondary and Tertiary Alcohols}

\begin{popfigure}
\begin{tikzpicture}[scale=0.8]
% Lucas test tubes
\foreach \i/\lbl/\color in {1/Tertiary/poporange, 2/Secondary/popyellow, 3/Primary/popgreen} {
  \begin{scope}[xshift=(\i-1)*5cm]
    \draw[thick] (0,0) rectangle (2,4);
    \fill[\color,opacity=0.5] (0,0) rectangle (2,3);
    \draw[thick,popblue] (0.2,3.2) -- (1.8,3.2) node[midway,above] {cloudy};
    \node at (1,3.7) {\lbl};
    \node at (1,-0.5) {Lucas reagent};
  \end{scope}
}
\node at (5,-1.5) {Time: immediate \hspace{1cm} 5--10 min \hspace{1cm} no change (RT)};
\end{tikzpicture}
\legende{Lucas test observations for the three classes of alcohols.}
\end{popfigure}

\begin{methode}[Using Lucas, oxidation and iodoform tests]
\begin{enumerate}
\item \cle{Lucas test} (\ce{ZnCl2}/conc.\ \ce{HCl}, room temperature): tertiary alcohol $\rightarrow$ cloudy \emph{immediately} (alkyl chloride insoluble); secondary $\rightarrow$ cloudy within 5--10 minutes; primary $\rightarrow$ no change at room temperature (may react on heating). The cloudiness is the insoluble chloroalkane.
\item \cle{Oxidation} (acidified \ce{K2Cr2O7}, \ce{H2SO4}, heat): orange $\rightarrow$ green (\ce{Cr2O7^{2-}} $\rightarrow$ \ce{Cr^{3+}}). Primary $\rightarrow$ aldehyde (distil immediately to isolate) $\rightarrow$ carboxylic acid (reflux); secondary $\rightarrow$ ketone (no further oxidation); tertiary $\rightarrow$ \emph{no reaction} (no $\alpha$-hydrogen on the carbinol carbon).
\item \cle{Iodoform test} (\ce{I2}/\ce{NaOH}, warm): positive (yellow \ce{CHI3} precipitate, antiseptic smell) for alcohols containing \ce{CH3-CH(OH)-} (e.g.\ ethanol, propan-2-ol) and for methyl ketones \ce{CH3-CO-R}. The reaction involves oxidation to the methyl carbonyl followed by halogenation and cleavage.
\item Combine the tests in a logical flow: Lucas first (speed), then oxidation to confirm, then iodoform if a methyl carbinol is suspected.
\end{enumerate}
\end{methode}

\begin{aretenir}[Summary of alcohol reactions]
\begin{center}
\begin{tabularx}{\linewidth}{|l|X|X|X|}\hline
\rowcolor{popblueL}
\textbf{Test} & \textbf{Primary} & \textbf{Secondary} & \textbf{Tertiary} \\ \hline
Lucas (\ce{ZnCl2}/\ce{HCl}) & No change (RT) & Cloudy in 5--10 min & Cloudy immediately \\ \hline
Oxidation (\ce{K2Cr2O7}/\ce{H+}) & \ce{->} aldehyde \ce{->} acid & \ce{->} ketone & No reaction \\ \hline
Iodoform (\ce{I2}/\ce{NaOH}) & Only ethanol & \ce{CH3CH(OH)R} positive & Negative \\ \hline
\end{tabularx}
\end{center}
\end{aretenir}

\begin{exoresolu}[Identifying three isomeric alcohols]
Three unlabelled bottles contain the isomers A, B and C of \ce{C4H10O}: butan-1-ol, butan-2-ol and 2-methylpropan-2-ol. With Lucas reagent, C turns cloudy instantly, B after a few minutes, A not at all. On warming with acidified \ce{K2Cr2O7}, A and B turn the solution green; C does not. Only B gives a yellow precipitate with \ce{I2}/\ce{NaOH}. Identify A, B and C, giving reasons and equations.
\tcblower
\textbf{C} reacts instantly with Lucas reagent and resists oxidation $\Rightarrow$ \emph{tertiary} alcohol: C is \textbf{2-methylpropan-2-ol}, \ce{(CH3)3COH}. Tertiary alcohols have no hydrogen on the carbinol carbon, so they cannot be oxidised without breaking C--C bonds.\\[3pt]
\textbf{B} reacts slowly with Lucas reagent $\Rightarrow$ \emph{secondary} alcohol: B is \textbf{butan-2-ol}, \ce{CH3CH2CH(OH)CH3}. Oxidation gives the ketone butanone:
\[
\ce{CH3CH2CH(OH)CH3(l) + [O] -> CH3CH2COCH3(l) + H2O(l)}
\]
Butan-2-ol contains the \ce{CH3-CH(OH)-} group, hence the positive iodoform test:
\[
\ce{CH3CH2CH(OH)CH3(l) + 4I2(aq) + 6NaOH(aq) -> CHI3(s) v + CH3CH2COONa(aq) + 5NaI(aq) + 5H2O(l)}
\]
\textbf{A} shows no Lucas reaction at room temperature but is oxidised $\Rightarrow$ \emph{primary} alcohol: A is \textbf{butan-1-ol}, \ce{CH3CH2CH2CH2OH}, oxidised via butanal to butanoic acid:
\[
\ce{CH3CH2CH2CH2OH(l) + [O] -> CH3CH2CH2CHO(l) + H2O(l) ->[O] CH3CH2CH2COOH(l)}
\]
Butan-1-ol lacks the \ce{CH3-CH(OH)-} skeleton, so its iodoform test is negative. All three identifications are consistent.
\end{exoresolu}

\courssub{Method 4 --- Nucleophilic Addition to Aldehydes and Ketones; Distinguishing Them}

\begin{popfigure}
\begin{tikzpicture}[scale=0.9]
% Nucleophilic addition to carbonyl
\node[align=center] at (0,2.5) {\textbf{Nucleophilic addition to \ce{C=O}}};
\draw[thick] (-2,1) -- (0,1) node[midway,above] {\ce{Nu^-}};
\draw[thick,->] (-1,1) -- (-0.2,1);
\node at (0,0) {\ce{C}};
\draw[thick] (0,0) -- (-0.6,-0.6) node[midway,left] {\ce{R}};
\draw[thick] (0,0) -- (0.6,-0.6) node[midway,right] {\ensuremath{\mathrm{R'}}};
\draw[thick,->] (0,0.6) -- (0,1.2) node[midway,right] {\ce{O^-}};
\draw[thick,->] (0,1.2) -- (0,1.8) node[midway,right] {\ce{H+}};
\node at (0,2.2) {\ensuremath{\mathrm{Nu-C-R-R'}}};
\node at (0,2.2) [above] {\ce{OH}};
\end{tikzpicture}
\legende{General mechanism: nucleophile attacks the $\delta^+$ carbonyl carbon; the resulting alkoxide is protonated.}
\end{popfigure}

\begin{methode}[Predicting carbonyl reactions]
\begin{enumerate}
\item Recognise the \cle{polar carbonyl}: \ensuremath{\mathrm{C^{\delta^{+}}=O^{\delta^{-}}}}; nucleophiles attack the carbon.
\item \cle{HCN} (generated in situ from \ce{KCN} + dilute acid) $\rightarrow$ hydroxynitrile (cyanohydrin): \ce{R2C=O + HCN -> R2C(OH)CN}. Draw the arrow from \ce{CN^-} to C, then protonation of \ce{O^-}. \emph{Chain lengthens by one carbon.}
\item \cle{Reduction}: \ce{NaBH4} (mild, reduces aldehydes/ketones only) or \ce{LiAlH4} (strong, reduces acids/esters too) or \ce{H2}/\ce{Ni} (catalytic). Aldehyde $\rightarrow$ primary alcohol; ketone $\rightarrow$ secondary alcohol.
\item \cle{2,4-DNPH} (Brady's reagent): orange/yellow precipitate with \emph{any} aldehyde or ketone --- a test for \ce{C=O}, not a distinction. Recrystallise and measure melting point to identify.
\item To \cle{distinguish aldehyde from ketone}: Fehling's or Benedict's (brick-red \ce{Cu2O} with aldehydes only) or Tollens' (silver mirror with aldehydes only). Ketones resist mild oxidation.
\item \cle{Iodoform}: only \ce{CH3CHO} (ethanal) and methyl ketones \ce{CH3CO-R} give yellow \ce{CHI3}.
\end{enumerate}
\end{methode}

\begin{attention}[Fehling's vs Tollens' vs Benedict's]
\begin{itemize}
\item Fehling's: \ce{Cu^{2+}} complexed with tartrate in alkali; brick-red \ce{Cu2O} precipitate.
\item Benedict's: \ce{Cu^{2+}} complexed with citrate; same observation, more stable reagent.
\item Tollens': \ce{[Ag(NH3)2]+} in ammonia; silver mirror or black \ce{Ag} precipitate.
\item \textbf{All three} are reduced by aldehydes (\ce{RCHO -> RCOO-}) but \emph{not} by ketones.
\item \ce{CH3CHO} gives positive iodoform \emph{and} Fehling's/Tollens'; other aldehydes give only Fehling's/Tollens'.
\end{itemize}
\end{attention}

\begin{exoresolu}[Reactions of propanone and propanal]
(a) Write the equation for the reaction of propanone with HCN and name the product. (b) State what is observed when propanal and propanone are separately warmed with Fehling's solution. (c) Give the organic products when propanone is (i) reduced by \ce{NaBH4}, (ii) warmed with \ce{I2}/\ce{NaOH}.
\tcblower
\textbf{(a)} The cyanide ion attacks the $\delta^+$ carbonyl carbon; the alkoxide formed is then protonated by HCN:
\[
\ce{CH3COCH3(l) + HCN(aq) -> (CH3)2C(OH)CN(l)}
\]
Product: \textbf{2-hydroxy-2-methylpropanenitrile} (propanone cyanohydrin). The carbon chain has grown by one carbon.\\[3pt]
\textbf{(b)} Propanal is an aldehyde, easily oxidised: the blue Fehling's solution gives a \textbf{brick-red precipitate of \ce{Cu2O}}:
\[
\ce{CH3CH2CHO(l) + 2Cu^{2+}(aq) + 5OH^-(aq) -> CH3CH2COO^-(aq) + Cu2O(s) v + 3H2O(l)}
\]
Propanone is a ketone: it is not oxidised by Fehling's, so the solution \textbf{remains blue}.\\[3pt]
\textbf{(c)} (i) Reduction of the ketone gives the secondary alcohol \textbf{propan-2-ol}:
\[
\ce{CH3COCH3(l) + 2[H] -> CH3CH(OH)CH3(l)}
\]
(ii) Propanone is a methyl ketone, so the iodoform test is positive --- a \textbf{yellow precipitate of tri-iodomethane}:
\[
\ce{CH3COCH3(l) + 3I2(aq) + 4NaOH(aq) -> CHI3(s) v + CH3COONa(aq) + 3NaI(aq) + 3H2O(l)}
\]
\end{exoresolu}

\courssub{Method 5 --- Esterification and Ester Hydrolysis (with Calculations)}

\begin{methode}[Handling esterification questions]
\begin{enumerate}
\item Write the equilibrium: $\ensuremath{\mathrm{RCOOH(l) + R'OH(l) \rightleftharpoons [conc.\ H_{2}SO_{4}][\Delta] RCOOR'(l) + H_{2}O(l)}}$. The acid contributes the \ce{RCO} part, the alcohol the \ensuremath{\mathrm{OR'}} part. Conc.\ \ce{H2SO4} is catalyst \emph{and} dehydrating agent.
\item Name the ester: alkyl (from alcohol) + alkanoate (from acid), e.g.\ ethanol + ethanoic acid $\rightarrow$ \emph{ethyl ethanoate}.
\item Remember the reaction is \cle{slow and reversible}; excess of the cheaper reagent shifts the equilibrium forward (Le Chatelier).
\item \cle{Hydrolysis}: acid hydrolysis reverses esterification (reflux with dilute \ce{H2SO4}); alkaline hydrolysis (\emph{saponification}) with \ce{NaOH} is complete (irreversible) and gives the \emph{carboxylate salt} + alcohol --- basis of soap making from fats.
\item For yield calculations: $n = m/M$, find the limiting reagent, apply the percentage yield to the theoretical amount of ester.
\end{enumerate}
\end{methode}

\begin{aretenir}[Esterification vs Saponification]
\begin{center}
\begin{tabularx}{\linewidth}{|l|X|X|}\hline
\rowcolor{popblueL}
\textbf{Aspect} & \textbf{Acid hydrolysis} & \textbf{Alkaline hydrolysis (saponification)} \\ \hline
Reagents & Dilute \ce{H2SO4} or \ce{HCl}, reflux & Aqueous \ce{NaOH} or \ce{KOH}, reflux \\ \hline
Reversibility & Reversible (equilibrium) & Irreversible (carboxylate resists nucleophilic attack) \\ \hline
Products & Carboxylic acid + alcohol & Carboxylate salt + alcohol \\ \hline
Use & Lab preparation of acids/alcohols & Soap making, quantitative analysis \\ \hline
\end{tabularx}
\end{center}
\end{aretenir}

\begin{exoresolu}[Esterification with a yield calculation]
A student heats \SI{6.0}{g} of ethanoic acid ($M = 60$) with \SI{4.6}{g} of ethanol ($M = 46$) and a few drops of concentrated \ce{H2SO4}. After purification she obtains \SI{5.3}{g} of ethyl ethanoate ($M = 88$). (a) Write the equation. (b) Calculate the percentage yield. (c) State two ways of increasing the yield. (d) Write the equation for the hydrolysis of the ester by hot aqueous \ce{NaOH}.
\tcblower
\textbf{(a)}
\[
\ensuremath{\mathrm{CH_{3}COOH(l) + CH_{3}CH_{2}OH(l) \rightleftharpoons [conc.\ H_{2}SO_{4}][\Delta] CH_{3}COOCH_{2}CH_{3}(l) + H_{2}O(l)}}
\]
\textbf{(b)} Moles of acid: $n = 6.0/60 = 0.10$ mol. Moles of ethanol: $n = 4.6/46 = 0.10$ mol. The equation is 1:1, so neither is in excess; theoretical moles of ester $= 0.10$ mol.
\[
m_{\text{theo}} = 0.10 \times 88 = 8.8\ \mathrm{g}; \qquad \text{yield} = \frac{5.3}{8.8}\times 100 \approx 60\ \%
\]
\textbf{(c)} Use an \emph{excess} of one reagent (e.g.\ cheap ethanol) to shift the equilibrium to the right; \emph{remove the ester} (or water) as it forms, e.g.\ by distillation --- Le Chatelier's principle.\\[3pt]
\textbf{(d)} Alkaline hydrolysis (saponification) is irreversible because the carboxylate ion formed resists attack by the alcohol:
\[
\ce{CH3COOCH2CH3(l) + NaOH(aq) -> CH3COONa(aq) + CH3CH2OH(l)}
\]
Products: sodium ethanoate and ethanol.
\end{exoresolu}

\courssub{Method 6 --- Explaining Acidity and Basicity Trends}

\begin{popfigure}
\begin{tikzpicture}[scale=0.8]
% Delocalisation in carboxylate
\node[align=center] at (-2,2) {\textbf{Carboxylate anion}};
\draw[thick] (-3,1) -- (-1,1);
\draw[thick] (-3,1) -- (-3.5,0.5) node[left] {\ce{O^-}};
\draw[thick] (-1,1) -- (-0.5,0.5) node[right] {\ce{O}};
\draw[thick,dashed] (-3,1) -- (-2,1.5);
\draw[thick,dashed] (-1,1) -- (-2,1.5);
\node at (-2,1.5) [above] {\ce{R}};
\draw[<->,popblue] (-2.5,0.8) -- (-1.5,0.8) node[midway,below] {delocalised};
% Delocalisation in phenoxide
\node[align=center] at (3,2) {\textbf{Phenoxide anion}};
\draw[thick] (2,1) -- (4,1);
\draw[thick] (2,1) -- (1.5,0.5) node[left] {\ce{O^-}};
\draw[thick] (4,1) -- (4.5,0.5);
\draw[thick] (2,1) -- (2,1.5);
\draw[thick] (4,1) -- (4,1.5);
\draw[thick,dashed] (2,1) -- (3,1.5);
\draw[thick,dashed] (4,1) -- (3,1.5);
\node at (3,1.5) [above] {ring};
\draw[<->,popblue] (2.5,0.8) -- (3.5,0.8) node[midway,below] {delocalised};
\end{tikzpicture}
\legende{Delocalisation stabilises carboxylate (over two O) better than phenoxide (into ring).}
\end{popfigure}

\begin{methode}[Comparing acid strengths (carboxylic acids, phenol, alcohols) and base strengths (amines)]
\begin{enumerate}
\item Acidity depends on the \cle{stability of the anion} formed: \ce{RCOOH <=> RCOO^- + H^+}. The carboxylate is stabilised by \emph{delocalisation} of the negative charge over two electronegative oxygens.
\item \cle{Electron-withdrawing groups} (e.g.\ \ce{Cl}, \ce{NO2}, \ce{CF3}) stabilise the anion further via $-I$ effect $\Rightarrow$ stronger acid, lower $pK_a$. \cle{Electron-donating groups} (alkyl, $+I$) destabilise it $\Rightarrow$ weaker acid. Order: \ce{ClCH2COOH} ($pK_a \approx 2.86$) $>$ \ce{HCOOH} ($3.75$) $>$ \ce{CH3COOH} ($4.76$).
\item Phenol is more acidic than alcohols because the \cle{phenoxide ion} is stabilised by delocalisation into the benzene ring; but it is weaker than carboxylic acids (charge delocalised over one O and ring carbons vs two O). Test: phenol does \emph{not} liberate \ce{CO2} from \ce{NaHCO3}, carboxylic acids do.
\item Amine basicity depends on the \cle{availability of the lone pair} on N: alkyl groups push electrons ($+I$) $\Rightarrow$ aliphatic amines are stronger bases than \ce{NH3}; in \ce{C6H5NH2} the lone pair is delocalised into the ring $\Rightarrow$ aniline is a much weaker base.
\end{enumerate}
\end{methode}

\begin{attention}[Acidity/basicity reasoning in exams]
\begin{itemize}
\item Always mention \textbf{stability of the conjugate base} for acidity, and \textbf{availability of the lone pair} for basicity.
\item Use the terms ``electron-withdrawing'' ($-I$) and ``electron-donating'' ($+I$) correctly.
\item Do \emph{not} say ``phenol is acidic because it has an \ce{-OH} group'' --- alcohols have \ce{-OH} too but are neutral.
\item For amines, steric hindrance can reduce basicity of tertiary amines in aqueous solution (solvation effects) --- but gas phase basicity follows $3^\circ > 2^\circ > 1^\circ > \ce{NH3}$.
\end{itemize}
\end{attention}

\begin{exoresolu}[Ranking and justifying]
(a) Arrange in increasing acidity: ethanol, phenol, ethanoic acid, chloroethanoic acid. Justify. (b) Arrange in increasing basicity: aniline, ammonia, methylamine. Justify. (c) Describe one chemical test that distinguishes phenol from ethanoic acid.
\tcblower
\textbf{(a)} Increasing acidity: \textbf{ethanol $<$ phenol $<$ ethanoic acid $<$ chloroethanoic acid}.\\
Ethoxide \ce{CH3CH2O^-} is destabilised by the electron-donating ethyl group ($+I$), so ethanol hardly ionises ($pK_a \approx 16$). The phenoxide ion \ce{C6H5O^-} is stabilised by delocalisation of the charge into the benzene ring, so phenol is more acidic ($pK_a \approx 10$). The ethanoate ion \ce{CH3COO^-} delocalises the charge over \emph{two electronegative oxygens} --- better stabilisation still ($pK_a \approx 4.76$). In \ce{ClCH2COO^-} the electronegative Cl withdraws electron density ($-I$ effect), stabilising the anion further: chloroethanoic acid is the strongest acid of the four ($pK_a \approx 2.86$).\\[3pt]
\textbf{(b)} Increasing basicity: \textbf{aniline $<$ ammonia $<$ methylamine}.\\
In aniline the nitrogen lone pair is delocalised into the benzene ring and is much less available to accept a proton ($pK_b \approx 9.4$). In methylamine the methyl group is electron-donating ($+I$ effect), increasing the electron density on N and stabilising \ce{CH3NH3+}, so methylamine is a stronger base than ammonia ($pK_b \approx 3.36$ vs $4.75$).\\[3pt]
\textbf{(c)} Add aqueous \ce{NaHCO3} to each: ethanoic acid gives \textbf{effervescence of \ce{CO2}} (it is a strong enough acid to decompose the hydrogencarbonate):
\[
\ce{CH3COOH(aq) + NaHCO3(aq) -> CH3COONa(aq) + CO2(g) ^ + H2O(l)}
\]
Phenol, being weaker than carbonic acid, gives \textbf{no effervescence}.
\end{exoresolu}

\courssub{Method 7 --- Amines, Diazonium Salts and Amino Acids: Synthetic Routes}

\begin{methode}[Planning short syntheses with nitrogen compounds]
\begin{enumerate}
\item \cle{Preparing amines}: reduce a nitrile (\ce{LiAlH4} in dry ether or \ce{H2}/\ce{Ni}) --- adds one carbon; or reduce a nitro compound \ce{C6H5NO2 ->[Sn/HCl] C6H5NH2} for aromatic amines (reduction of nitro group).
\item \cle{Diazotisation}: \ce{C6H5NH2 + NaNO2 + 2HCl} at $0$--$5\ ^\circ\mathrm{C}$ $\rightarrow$ benzene diazonium chloride \ce{C6H5N2+Cl^-}. Keep it \emph{cold} --- it decomposes above $10\ ^\circ\mathrm{C}$ giving phenol and \ce{N2}.
\item Use the diazonium salt (Sandmeyer reactions): warm with \ce{H2O} $\rightarrow$ phenol; \ce{CuCl} $\rightarrow$ chlorobenzene; \ce{CuBr} $\rightarrow$ bromobenzene; \ce{CuCN} $\rightarrow$ benzonitrile; \ce{HBF4} then heat $\rightarrow$ fluorobenzene. Couple with phenol (alkaline) or aromatic amine $\rightarrow$ azo dye.
\item \cle{Amino acids} \ce{H2N-CHR-COOH} are amphoteric: with acid they form \ce{^{+}H3N-CHR-COOH}; with base \ce{H2N-CHR-COO^-}; in neutral solution they exist as \cle{zwitterions} \ce{^{+}H3N-CHR-COO^-} (high melting points, water-soluble, crystalline).
\end{enumerate}
\end{methode}

\begin{savaistu}[Azo dyes and Cameroon]
Azo dyes (\ce{-N=N-} chromophore) are widely used in the textile industry. In Cameroon, traditional dyeing techniques (e.g., \emph{ndop} cloth in the Grassfields) use natural indigo, but modern fabrics rely on synthetic azo dyes. The coupling reaction must be done in cold, alkaline conditions to keep the diazonium ion electrophilic and the phenol nucleophilic (as phenoxide).
\end{savaistu}

\begin{exoresolu}[From benzene to a dye; behaviour of glycine]
(a) Outline, with reagents, conditions and equations, the conversion of nitrobenzene into phenylamine, then into benzene diazonium chloride, and finally into an orange azo dye. (b) Explain why glycine, \ce{H2NCH2COOH}, has a high melting point and reacts with both \ce{HCl} and \ce{NaOH}. Write the equations.
\tcblower
\textbf{(a)} \emph{Step 1 --- Reduction:} nitrobenzene is reduced by tin and concentrated \ce{HCl} (the \ce{Sn}/\ce{HCl} generates \ce{[H]} in situ):
\[
\ensuremath{\mathrm{C_{6}H_{5}NO_{2}(l) + 6[H] \rightarrow [Sn/conc.\ HCl] C_{6}H_{5}NH_{2}(l) + 2H_{2}O(l)}}
\]
\emph{Step 2 --- Diazotisation:} phenylamine is treated with sodium nitrite and dilute \ce{HCl} at $0$--$5\ ^\circ\mathrm{C}$ (the \ce{NaNO2} + \ce{HCl} generate nitrous acid \ce{HNO2} in situ):
\[
\ensuremath{\mathrm{C_{6}H_{5}NH_{2}(l) + NaNO_{2}(aq) + 2HCl(aq) \rightarrow [0-5\ ^\circ C] C_{6}H_{5}N_{2}+Cl^{-}(aq) + NaCl(aq) + 2H_{2}O(l)}}
\]
The low temperature is essential because benzene diazonium chloride decomposes on warming, giving phenol and nitrogen.\\
\emph{Step 3 --- Coupling:} the cold diazonium salt is added to phenol in alkaline solution; the diazonium ion acts as an electrophile and couples at the \emph{para} position (major) and \emph{ortho}:
\[
\ensuremath{\mathrm{C_{6}H_{5}N_{2}+Cl^{-}(aq) + C_{6}H_{5}OH(aq) \rightarrow [NaOH,\ cold] HO-C_{6}H_{4}-N=N-C_{6}H_{5}(s) + HCl(aq)}}
\]
The product, 4-hydroxyazobenzene, contains the \ce{-N=N-} chromophore and is an \textbf{orange azo dye}.\\[3pt]
\textbf{(b)} Glycine contains both a basic \ce{-NH2} group and an acidic \ce{-COOH} group. An internal proton transfer gives the \textbf{zwitterion} \ce{^{+}H3NCH2COO^-}; the strong ionic attractions between zwitterions in the crystal lattice explain the \emph{high melting point} (\SI{233}{\ensuremath{{}^{\circ}\mathrm{C}}}) and solubility in water.\\
With acid, the carboxylate end accepts a proton:
\[
\ce{^{+}H3NCH2COO^-(s) + HCl(aq) -> ^{+}H3NCH2COOH(aq) + Cl^-(aq)}
\]
With base, the ammonium end donates a proton:
\[
\ce{^{+}H3NCH2COO^-(s) + NaOH(aq) -> H2NCH2COO^-Na^+(aq) + H2O(l)}
\]
This dual behaviour is why amino acids are described as \emph{amphoteric}.
\end{exoresolu}

\newpage

\coursec{Exercises}

\courssub{I check my knowledge}

\begin{exercice}[Multiple choice questions]
For each question, choose the single correct answer.
\begin{enumerate}
\item The hydrolysis of 2-bromo-2-methylpropane by aqueous sodium hydroxide proceeds mainly by:
\begin{enumerate}
\item an $S_{\mathrm{N}}2$ mechanism because the carbon bearing bromine is unhindered;
\item an $S_{\mathrm{N}}1$ mechanism because a stable tertiary carbocation is formed;
\item an elimination mechanism because the substrate is tertiary;
\item a free-radical mechanism because the \ce{C-Br} bond is weak.
\end{enumerate}
\item Which reagent gives a silver mirror with propanal but not with propanone?
\begin{enumerate}
\item 2,4-dinitrophenylhydrazine
\item Acidified potassium dichromate(VI)
\item Tollens' reagent
\item Fehling's solution only
\end{enumerate}
\item The correct order of decreasing acidity is:
\begin{enumerate}
\item ethanol $>$ phenol $>$ ethanoic acid
\item ethanoic acid $>$ phenol $>$ ethanol
\item phenol $>$ ethanoic acid $>$ ethanol
\item ethanoic acid $>$ ethanol $>$ phenol
\end{enumerate}
\item Benzene diazonium chloride is prepared by reacting phenylamine with:
\begin{enumerate}
\item concentrated \ce{HNO3} at room temperature;
\item \ce{NaNO2} and dilute \ce{HCl} below $\SI{5}{\ensuremath{{}^{\circ}\mathrm{C}}}$;
\item \ce{NaNO3} and concentrated \ce{H2SO4} under reflux;
\item \ce{NH4Cl} and dilute \ce{HCl} at $\SI{50}{\ensuremath{{}^{\circ}\mathrm{C}}}$.
\end{enumerate}
\end{enumerate}
\end{exercice}

\begin{exercice}[True or false — justify your answer]
State whether each statement is true or false, and justify briefly.
\begin{enumerate}
\item Chlorobenzene reacts readily with aqueous sodium hydroxide under the same mild conditions as chloroethane.
\item Propan-2-ol gives a yellow precipitate of triiodomethane in the iodoform test.
\item All primary amines react with nitrous acid to give stable diazonium salts at room temperature.
\item An amide is formed when an acyl chloride reacts with ammonia.
\end{enumerate}
\end{exercice}

\begin{exercice}[Definitions]
Define each of the following terms, giving one example in each case:
\begin{enumerate}
\item nucleophilic substitution;
\item zwitterion;
\item esterification;
\item diazotisation.
\end{enumerate}
\end{exercice}

\begin{exercice}[Direct calculations]
\begin{enumerate}
\item Calculate the percentage by mass of bromine in 1-bromopropane, \ce{C3H7Br}. ($A_r$: C = 12, H = 1, Br = 80)
\item $\SI{6.0}{g}$ of ethanoic acid reacts with excess ethanol in the presence of concentrated \ce{H2SO4}. Calculate the theoretical mass of ethyl ethanoate formed. ($A_r$: C = 12, H = 1, O = 16)
\item A carboxylic acid has $\mathrm{p}K_a = 4.20$. Calculate its acid dissociation constant $K_a$.
\end{enumerate}
\end{exercice}

\courssub{I practise}

\begin{exercice}[Isomers of \ce{C4H10O}]
\begin{enumerate}
\item Write the structural formulae and names of all isomers of molecular formula \ce{C4H10O}.
\item Classify each alcohol isomer as primary, secondary or tertiary.
\item An isomer X of \ce{C4H10O} gives no reaction with sodium, while its isomer Y reacts with sodium and gives a positive iodoform test. Identify X and Y, giving reasons.
\end{enumerate}
\end{exercice}

\begin{exercice}[Distinguishing alcohols and phenol]
You are given four unlabelled bottles containing propan-1-ol, propan-2-ol, 2-methylpropan-2-ol and phenol. Using only Lucas reagent, neutral iron(III) chloride solution and acidified potassium dichromate(VI), devise a logical sequence of tests to identify each liquid. Present your expected observations in a table and write one equation for a positive test in each case.
\end{exercice}

\begin{exercice}[Identifying a carbonyl compound]
Compound Z, of molecular formula \ce{C3H6O}, gives an orange precipitate with 2,4-dinitrophenylhydrazine and a silver mirror with Tollens' reagent.
\begin{enumerate}
\item Identify Z, giving reasons for your answer.
\item Write equations for the two tests described.
\item State the structural formula of the isomer of Z and predict the result of each of the two tests on this isomer.
\end{enumerate}
\end{exercice}

\begin{exercice}[Acidity of substituted ethanoic acids]
The $\mathrm{p}K_a$ values of ethanoic, chloroethanoic and trichloroethanoic acids are 4.76, 2.86 and 0.65 respectively.
\begin{enumerate}
\item Explain the trend in these values, referring to the stability of the carboxylate anions.
\item Predict, with justification, an approximate $\mathrm{p}K_a$ value for dichloroethanoic acid.
\item State, with a reason, which of the four acids reacts most vigorously with sodium hydrogencarbonate.
\end{enumerate}
\end{exercice}

\courssub{I prepare for the GCE}

\begin{exercice}[Mechanisms of nucleophilic substitution — 20 marks]
\begin{enumerate}
\item Draw the $S_{\mathrm{N}}1$ mechanism for the hydrolysis of 2-bromo-2-methylpropane by aqueous sodium hydroxide, showing all steps, intermediates and curly arrows. \hfill (6 marks)
\item Draw the $S_{\mathrm{N}}2$ mechanism for the hydrolysis of 1-bromopropane under the same conditions, showing the transition state and curly arrows. \hfill (4 marks)
\item Explain, in terms of carbocation stability and steric hindrance, why each substrate follows its particular pathway. \hfill (4 marks)
\item Write the rate equation for each reaction and predict the effect of doubling the hydroxide ion concentration on the rate of each reaction. \hfill (4 marks)
\item State the optical consequence of $S_{\mathrm{N}}1$ hydrolysis of a single enantiomer of a chiral haloalkane, and explain why. \hfill (2 marks)
\end{enumerate}
\end{exercice}

\begin{exercice}[Organic synthesis from ethanol — 20 marks]
Starting from ethanol, show with reagents, conditions and balanced equations how you would obtain:
\begin{enumerate}
\item ethyl ethanoate; \hfill (6 marks)
\item propanoic acid, by a chain-extension route passing through a nitrile (state clearly the haloalkane intermediate and the role of \ce{KCN}); \hfill (8 marks)
\item ethylamine, using two different routes, one of which avoids the formation of secondary and tertiary amines. \hfill (6 marks)
\end{enumerate}
\end{exercice}

\begin{exercice}[The hydroxy group in ethanol, phenol and ethanoic acid — 20 marks]
``The $-\mathrm{OH}$ group behaves differently in ethanol, phenol and ethanoic acid.''

Discuss this statement with reference to the reactions of the three compounds with:
\begin{enumerate}
\item sodium metal; \hfill (6 marks)
\item aqueous sodium hydroxide; \hfill (6 marks)
\item aqueous sodium hydrogencarbonate. \hfill (4 marks)
\end{enumerate}
Support each point with a balanced equation where a reaction occurs, and explain clearly, in terms of the stability of the anions formed, the observed differences in acidity. \hfill (4 marks)
\end{exercice}

\newpage

\coursec{Solutions to the exercises}

\begin{corrige}[Exercise 1 — Multiple choice questions]
\textbf{Question 1: answer (b).} 2-Bromo-2-methylpropane, \ce{(CH3)3CBr}, is a \cle{tertiary} haloalkane. The three bulky methyl groups sterically block the back-side attack required for \cle{$S_{\mathrm{N}}2$}, while loss of \ce{Br-} gives the tertiary carbocation \ce{(CH3)3C+}, stabilised by the positive inductive effect (+I) of three methyl groups. Ionisation followed by fast nucleophilic capture is therefore the \cle{$S_{\mathrm{N}}1$} pathway.

\textbf{Question 2: answer (c).} Tollens' reagent (ammoniacal silver nitrate) is reduced to a \cle{silver mirror} by \emph{aldehydes} such as propanal, but not by ketones such as propanone. 2,4-DNPH (a) reacts with both aldehydes and ketones; acidified dichromate (b) also oxidises propanal but the question asks specifically for the silver mirror test; Fehling's solution (d) gives a red precipitate with aldehydes but is not the reagent named in the question.

\textbf{Question 3: answer (b).} Ethanoic acid ($\mathrm{p}K_{\mathrm{a}} \approx 4.8$) is a much stronger acid than phenol ($\mathrm{p}K_{\mathrm{a}} \approx 10$), which is itself more acidic than ethanol ($\mathrm{p}K_{\mathrm{a}} \approx 16$). The ethanoate ion is stabilised by \cle{resonance} over two equivalent oxygen atoms; the phenoxide ion is stabilised by \cle{delocalisation} of the negative charge into the benzene ring; the ethoxide ion has no stabilisation, and the electron-donating ethyl group slightly destabilises the anion.

\textbf{Question 4: answer (b).} Diazotisation of phenylamine uses \ce{NaNO2} and dilute \ce{HCl} (which generate nitrous acid, \ce{HNO2}, in situ) at $0$--$5^{\circ}\mathrm{C}$; above $5^{\circ}\mathrm{C}$ the diazonium salt decomposes to phenol with evolution of nitrogen.
\end{corrige}

\begin{corrige}[Exercise 2 — True or false]
\begin{enumerate}
\item \textbf{False.} In chlorobenzene the \ce{C-Cl} bond acquires partial double-bond character due to delocalisation of a lone pair on chlorine into the benzene ring (\cle{mesomeric effect}), strengthening the bond. Additionally, the electron-rich ring repels approaching nucleophiles. Hydrolysis requires drastic conditions (high temperature, high pressure, concentrated \ce{NaOH}), unlike chloroethane which undergoes hydrolysis on warming with aqueous \ce{NaOH}.
\item \textbf{True.} Propan-2-ol, \ce{CH3CH(OH)CH3}, contains the \ce{CH3CH(OH)-} group (\cle{iodoform reaction} structural requirement). It is oxidised by \ce{I2}/\ce{NaOH} to propanone, which then undergoes halogenation and cleavage to give a yellow precipitate of \ce{CHI3} (iodoform).
\item \textbf{False.} Only \emph{aromatic} primary amines give diazonium salts stable enough to be isolated (and only below $5^{\circ}\mathrm{C}$). Aliphatic primary amines react with nitrous acid to form unstable aliphatic diazonium salts which decompose immediately, losing \ce{N2} and forming a mixture of products (mainly alcohols via $S_{\mathrm{N}}1$).
\item \textbf{True.} Acyl chlorides react violently with ammonia by nucleophilic acyl substitution: \ce{CH3COCl + 2NH3 -> CH3CONH2 + NH4Cl}. The second mole of ammonia neutralises the \ce{HCl} formed, driving the reaction to completion.
\end{enumerate}
\end{corrige}

\begin{corrige}[Exercise 3 — Definitions]
\begin{enumerate}
\item \textbf{Nucleophilic substitution:} A reaction in which a \cle{nucleophile} (an electron-pair donor, e.g. \ce{OH-}, \ce{NH3}, \ce{CN-}) replaces a \cle{leaving group} (e.g. halide) on a saturated carbon atom. Example: \ce{CH3CH2Br + OH- -> CH3CH2OH + Br-}.
\item \textbf{Zwitterion:} A \cle{dipolar ion} formed by internal proton transfer in an amino acid, carrying both a positive charge (\ce{-NH3+}) and a negative charge (\ce{-COO-}) but electrically neutral overall. Example: \ce{^{+}H3N-CH2-COO-} (glycine zwitterion).
\item \textbf{Esterification:} The reversible reaction of a carboxylic acid with an alcohol in the presence of concentrated \ce{H2SO4} as catalyst (and dehydrating agent) under \cle{reflux} to give an ester and water. Example: \ensuremath{\mathrm{CH_{3}COOH(l) + C_{2}H_{5}OH(l) \rightleftharpoons [conc. H_{2}SO_{4}][\Delta] CH_{3}COOC_{2}H_{5}(l) + H_{2}O(l)}}.
\item \textbf{Diazotisation:} The reaction of a primary aromatic amine with nitrous acid (generated in situ from \ce{NaNO2} and dilute \ce{HCl}) at $0$--$5^{\circ}\mathrm{C}$ to form an \cle{aryl diazonium salt}. Example: \ensuremath{\mathrm{C_{6}H_{5}NH_{2} + HNO_{2} + HCl \rightarrow [0-5^{\circ}\mathrm{C}] C_{6}H_{5}N_{2}+Cl^{-} + 2H_{2}O}}.
\end{enumerate}
\end{corrige}

\begin{corrige}[Exercise 4 — Direct calculations]
\begin{enumerate}
\item Molar mass \ce{C3H7Br} = $3(12) + 7(1) + 80 = 123\ \mathrm{g\,mol^{-1}}$.
\[ \%\,\mathrm{Br} = \frac{80}{123}\times 100 = 65.0\ \% \]
\item Equation: \ce{CH3COOH + C2H5OH <=> CH3COOC2H5 + H2O} (1:1 mole ratio).
$M_{\mathrm{r}}(\ce{CH3COOH}) = 60\ \mathrm{g\,mol^{-1}}$; $n(\text{acid}) = \dfrac{6.0}{60} = 0.10\ \mathrm{mol}$.
$M_{\mathrm{r}}(\ce{CH3COOC2H5}) = 4(12) + 8(1) + 2(16) = 88\ \mathrm{g\,mol^{-1}}$.
Assuming 100\% yield: $m(\text{ester}) = 0.10 \times 88 = \SI{8.8}{g}$.
\item $K_{\mathrm{a}} = 10^{-\mathrm{p}K_{\mathrm{a}}} = 10^{-4.20} = 6.3\times 10^{-5}\ \mathrm{mol\,dm^{-3}}$.
\end{enumerate}
\end{corrige}

\begin{corrige}[Exercise 5 — Isomers of \ce{C4H10O}]
\textbf{(a)} Seven structural isomers exist: four alcohols and three ethers.
\begin{itemize}
\item \ce{CH3CH2CH2CH2OH}: butan-1-ol (primary);
\item \ce{CH3CH2CH(OH)CH3}: butan-2-ol (secondary);
\item \ce{(CH3)2CHCH2OH}: 2-methylpropan-1-ol (primary);
\item \ce{(CH3)3COH}: 2-methylpropan-2-ol (tertiary);
\item \ce{CH3OCH2CH2CH3}: 1-methoxypropane (methyl propyl ether);
\item \ce{CH3CH2OCH2CH3}: ethoxyethane (diethyl ether);
\item \ce{CH3OCH(CH3)2}: 2-methoxypropane (methyl isopropyl ether).
\end{itemize}

\textbf{(b)} Classification of the alcohols:
\begin{itemize}
\item Primary: butan-1-ol, 2-methylpropan-1-ol;
\item Secondary: butan-2-ol;
\item Tertiary: 2-methylpropan-2-ol.
\end{itemize}

\textbf{(c)} X gives no reaction with sodium $\Rightarrow$ X has no \ce{-OH} group $\Rightarrow$ X is an \emph{ether} (e.g. ethoxyethane). Y reacts with sodium (so Y is an alcohol) and gives a positive iodoform test $\Rightarrow$ Y must contain the \ce{CH3CH(OH)-} group. The only \ce{C4H10O} alcohol with this feature is \textbf{butan-2-ol}.
\end{corrige}

\begin{corrige}[Exercise 6 — Distinguishing alcohols and phenol]
\textbf{Step 1 — Neutral \ce{FeCl3} (aq):} Phenol gives an intense \emph{violet/purple coloration} due to formation of the complex ion \ce{[Fe(OC6H5)6]^{3-}}; the three alcohols show no colour change.
Equation (simplified): \ce{6C6H5OH + Fe^{3+}(aq) -> [Fe(OC6H5)6]^{3-}(aq) + 6H+(aq)}.

\textbf{Step 2 — Lucas reagent (\ce{ZnCl2}/conc.\ \ce{HCl})} on the three remaining liquids (tests for \ce{C-OH} cleavage via $S_{\mathrm{N}}1$/$S_{\mathrm{N}}2$):
\begin{itemize}
\item 2-methylpropan-2-ol (tertiary): \emph{immediate cloudiness} (two layers form at once) due to rapid formation of insoluble \ce{(CH3)3CCl};
\item Propan-2-ol (secondary): cloudiness appears after \emph{5--10 minutes};
\item Propan-1-ol (primary): \emph{no cloudiness} at room temperature (requires heating for $S_{\mathrm{N}}2$).
\end{itemize}
Equation (tertiary example): \ce{(CH3)3COH + HCl -> (CH3)3CCl + H2O}.

\textbf{Step 3 — Acidified \ce{K2Cr2O7}, warm} (confirms primary/secondary distinction via oxidation):
\begin{itemize}
\item Propan-1-ol (primary): orange $\to$ green (oxidised to propanal then propanoic acid);
\item Propan-2-ol (secondary): orange $\to$ green (oxidised to propanone);
\item 2-methylpropan-2-ol (tertiary): \emph{no colour change} (resists oxidation).
\end{itemize}
Equation (primary): \ce{CH3CH2CH2OH + 2[O] -> CH3CH2COOH + H2O}.

\begin{center}
\begin{tabularx}{\linewidth}{|l|X|X|X|}
\hline
\rowcolor{popblueL} Test & Phenol & 2-methylpropan-2-ol (3\textdegree) & Propan-2-ol (2\textdegree) / Propan-1-ol (1\textdegree) \\
\hline
Neutral \ce{FeCl3} & Violet coloration & No change & No change \\
\hline
Lucas reagent & — & Immediate cloudiness & 5--10 min (2\textdegree) / None at RT (1\textdegree) \\
\hline
\ce{K2Cr2O7}/\ce{H+} (warm) & — & No change (orange) & Orange $\to$ Green \\
\hline
\end{tabularx}
\end{center}
\end{corrige}

\begin{corrige}[Exercise 7 — Identifying a carbonyl compound]
\textbf{(a)} The orange/yellow precipitate with 2,4-DNPH (Brady's reagent) shows Z is a \cle{carbonyl compound} (aldehyde or ketone). The silver mirror with Tollens' reagent (\ce{[Ag(NH3)2]+}) shows Z is an \emph{aldehyde} (ketones are not oxidised by mild oxidising agents). With molecular formula \ce{C3H6O}, Z is \textbf{propanal}, \ce{CH3CH2CHO}.

\textbf{(b)} Condensation with 2,4-DNPH:
\ce{CH3CH2CHO + H2NNHC6H3(NO2)2 -> CH3CH2CH=NNHC6H3(NO2)2 + H2O} (orange precipitate, \cle{2,4-dinitrophenylhydrazone}).
Oxidation by Tollens' reagent:
\ce{CH3CH2CHO + 2[Ag(NH3)2]+ + 3OH- -> CH3CH2COO- + 2Ag(s) + 4NH3 + 2H2O} (silver mirror).

\textbf{(c)} The structural isomer is \textbf{propanone}, \ce{CH3COCH3}. With 2,4-DNPH it gives an orange precipitate (it is a carbonyl compound), but with Tollens' reagent it gives \emph{no reaction} (no silver mirror) because ketones resist mild oxidation.
\end{corrige}

\begin{corrige}[Exercise 8 — Acidity of substituted ethanoic acids]
\textbf{(a)} Chlorine is strongly electron-withdrawing (\cle{$-I$ inductive effect}). In the carboxylate anion \ensuremath{\mathrm{Cl_nCH_{3-n}COO-}}, each chlorine atom pulls electron density away from the \ce{COO-} group, dispersing the negative charge and stabilising the anion. The more stable the anion, the more the equilibrium \ce{RCOOH <=> RCOO- + H+} lies to the right, hence the stronger the acid (lower $\mathrm{p}K_{\mathrm{a}}$).
Trend: \ce{CH3COOH} ($\mathrm{p}K_{\mathrm{a}}=4.76$) $<$ \ce{ClCH2COOH} ($\mathrm{p}K_{\mathrm{a}}=2.86$) $<$ \ce{Cl2CHCOOH} ($\mathrm{p}K_{\mathrm{a}}\approx1.48$) $<$ \ce{Cl3CCOOH} ($\mathrm{p}K_{\mathrm{a}}=0.65$). Acid strength increases with number of Cl atoms.

\textbf{(b)} Dichloroethanoic acid, \ce{Cl2CHCOOH}, has two electron-withdrawing Cl atoms. Its $\mathrm{p}K_{\mathrm{a}}$ should lie between the mono- and trichloro acids, closer to the trichloro value. Expected $\mathrm{p}K_{\mathrm{a}} \approx 1.3$--$1.5$ (literature value $\approx 1.48$).

\textbf{(c)} Trichloroethanoic acid (\ce{Cl3CCOOH}) reacts most vigorously with \ce{NaHCO3} because it is the strongest acid (lowest $\mathrm{p}K_{\mathrm{a}}$), donating \ce{H+} most readily. The reaction produces effervescence of \ce{CO2}:
\ce{Cl3CCOOH + NaHCO3 -> Cl3CCOONa + H2O + CO2 ^}.
Phenol ($\mathrm{p}K_{\mathrm{a}}\approx10$) and ethanol ($\mathrm{p}K_{\mathrm{a}}\approx16$) are weaker acids than carbonic acid ($\mathrm{p}K_{\mathrm{a1}}\approx6.3$) and do not react.
\end{corrige}

\begin{corrige}[Exercise 9 — Mechanisms of nucleophilic substitution (20 marks)]
\textbf{(a) $S_{\mathrm{N}}1$ hydrolysis of \ce{(CH3)3CBr} (6 marks).}
Step 1 (slow, rate-determining): Heterolytic fission of the \ce{C-Br} bond. Curly arrow from the \ce{C-Br} bond to \ce{Br}.
\ce{(CH3)3C-Br -> (CH3)3C+ + Br-}
Step 2 (fast): Nucleophilic attack by \ce{OH-} on the planar carbocation. Curly arrow from lone pair on \ce{OH-} to the positively charged carbon.
\ce{(CH3)3C+ + OH- -> (CH3)3COH}
\textit{Marks:} Correct slow ionisation with curly arrow (1), carbocation intermediate shown (1), both curly arrows correct (2), \ce{Br-} and product identified (1), steps labelled slow/fast (1).

\textbf{(b) $S_{\mathrm{N}}2$ hydrolysis of \ce{CH3CH2CH2Br} (4 marks).}
One concerted step: Back-side attack by \ce{OH-}. Curly arrow from lone pair on \ce{OH-} to the $\delta+$ carbon; simultaneously curly arrow from \ce{C-Br} bond to \ce{Br}.
Transition state: $\ensuremath{\mathrm{[HO^{\delta^{-}}\cdots CH_{2} \cdots Br^{\delta^{-}}]^{-}}}$ (partial bonds shown dashed, negative charge distributed).
Products: \ce{CH3CH2CH2OH + Br-}. \cle{Inversion of configuration} occurs at the carbon centre (Walden inversion).
\textit{Marks:} Back-side attack shown (1), both curly arrows (1), transition state with partial bonds/charge (1), inversion noted (1).

\textbf{(c) Explanation of pathway preference (4 marks).}
Tertiary substrate: Forms a relatively stable tertiary carbocation (stabilised by three +I methyl groups) (1); the $\alpha$-carbon is sterically hindered, blocking back-side attack required for $S_{\mathrm{N}}2$ (1).
Primary substrate: Would form a highly unstable primary carbocation (1); the $\alpha$-carbon is unhindered, allowing easy back-side attack for the concerted $S_{\mathrm{N}}2$ path (1).

\textbf{(d) Kinetics (4 marks).}
$S_{\mathrm{N}}1$: Rate $= k[\ce{(CH3)3CBr}]$ (unimolecular) (1). Doubling $[\ce{OH-}]$ has \emph{no effect} on the rate (1).
$S_{\mathrm{N}}2$: Rate $= k[\ce{C3H7Br}][\ce{OH-}]$ (bimolecular) (1). Doubling $[\ce{OH-}]$ \emph{doubles} the rate (1).

\textbf{(e) Optical consequence (2 marks).}
$S_{\mathrm{N}}1$ on a chiral substrate leads to \cle{racemisation}: a racemic (optically inactive) mixture is obtained (1), because the planar \cle{sp$^{2}$} carbocation is attacked equally from either face by the nucleophile (1).
\end{corrige}

\begin{corrige}[Exercise 10 — Organic synthesis from ethanol (20 marks)]
\textbf{(a) Ethyl ethanoate (6 marks).}
\begin{enumerate}
\item Oxidation of ethanol to ethanoic acid: \ensuremath{\mathrm{CH_{3}CH_{2}OH(l) + 2[O] \rightarrow [K_{2}Cr_{2}O_{7}/H_{2}SO_{4}][\Delta] CH_{3}COOH(aq) + H_{2}O(l)}} (reflux to ensure complete oxidation to acid).
\item Esterification: \ensuremath{\mathrm{CH_{3}COOH(l) + CH_{3}CH_{2}OH(l) \rightleftharpoons [conc. H_{2}SO_{4}][\Delta] CH_{3}COOCH_{2}CH_{3}(l) + H_{2}O(l)}} (reflux with concentrated \ce{H2SO4} as catalyst and dehydrating agent; distil off ester).
\end{enumerate}
\textit{Marks:} Reagents/conditions for oxidation (2), reagents/conditions for esterification (2), both equations balanced with state symbols (2).

\textbf{(b) Propanoic acid via a nitrile (chain extension) (8 marks).}
\begin{enumerate}
\item \ce{CH3CH2OH + HBr ->[conc. H2SO4 / NaBr] CH3CH2Br + H2O} (or \ce{PBr3}). Haloalkane intermediate: \textbf{bromoethane}.
\item \ce{CH3CH2Br + KCN ->[ethanol, reflux] CH3CH2CN + KBr}. \ce{KCN} (in ethanolic solution) supplies the nucleophile \ce{CN-}; this step extends the carbon chain by one carbon (C2 $\to$ C3).
\item Hydrolysis of nitrile: \ce{CH3CH2CN + 2H2O + H+ ->[reflux] CH3CH2COOH + NH4+}. (Acid hydrolysis gives the acid directly; alkaline hydrolysis gives the salt which is then acidified).
\end{enumerate}
\textit{Marks:} Haloalkane identified (1), role of \ce{KCN}/chain extension explained (2), three balanced equations (3), conditions stated (2).

\textbf{(c) Ethylamine (\ce{CH3CH2NH2}), two routes from ethanol (6 marks).}
\textbf{Route 1: Ammonolysis of bromoethane (direct, but gives mixture).}
\ce{CH3CH2OH ->[HBr] CH3CH2Br}
\ce{CH3CH2Br + 2NH3(alc) ->[sealed tube, 373 K] CH3CH2NH2 + NH4Br}
Excess ammonia limits further substitution, but some diethylamine and triethylamine are usually formed as by-products.

\textbf{Route 2: Gabriel synthesis (gives pure primary amine).}
\begin{enumerate}
\item \ce{CH3CH2OH ->[HBr] CH3CH2Br}
\item \ce{Phthalimide + KOH -> Potassium phthalimide}
\item \ce{Potassium phthalimide + CH3CH2Br -> N-Ethylphthalimide}
\item \ce{N-Ethylphthalimide + NH2NH2*H2O (hydrazine) ->[reflux] CH3CH2NH2 + Phthalhydrazide}
\end{enumerate}
\textit{Alternative Route 2 (Reduction of nitrile derived from ethanal):}
\ce{CH3CH2OH ->[PCC/distill] CH3CHO ->[NH2OH] CH3CH=NOH ->[P2O5] CH3CN ->[LiAlH4 / H2, Ni] CH3CH2NH2}
\textit{Marks:} Two valid distinct routes (2), balanced equations/steps (2), explanation of selectivity/purity (2).
\end{corrige}

\begin{corrige}[Exercise 11 — The hydroxy group in ethanol, phenol and ethanoic acid (20 marks)]
\textbf{(a) Reaction with sodium metal (6 marks).} All three react, evolving hydrogen gas. Vigour increases with acidity.
\begin{align*}
\ce{2CH3CH2OH(l) + 2Na(s) &-> 2CH3CH2ONa(s) + H2(g)} &\text{(gentle effervescence)} \\
\ce{2C6H5OH(s) + 2Na(s) &-> 2C6H5ONa(s) + H2(g)} &\text{(steady reaction, phenol often melts)} \\
\ce{2CH3COOH(l) + 2Na(s) &-> 2CH3COONa(s) + H2(g)} &\text{(vigorous effervescence)}
\end{align*}
\textit{Marks:} Three correct equations with states (3), observations linked to acidity trend (3).

\textbf{(b) Reaction with aqueous \ce{NaOH} (6 marks).}
\begin{itemize}
\item \textbf{Ethanol:} \emph{No reaction}. Ethanol ($\mathrm{p}K_{\mathrm{a}}\approx16$) is a weaker acid than water ($\mathrm{p}K_{\mathrm{a}}\approx15.7$); equilibrium \ce{C2H5OH + OH- <=> C2H5O- + H2O} lies far to the left (2 marks).
\item \textbf{Phenol:} Reacts readily (neutralisation). \ce{C6H5OH(aq) + NaOH(aq) -> C6H5ONa(aq) + H2O(l)}. Phenol ($\mathrm{p}K_{\mathrm{a}}\approx10$) is a stronger acid than water (2 marks).
\item \textbf{Ethanoic acid:} Reacts readily (neutralisation). \ce{CH3COOH(aq) + NaOH(aq) -> CH3COONa(aq) + H2O(l)} (2 marks).
\end{itemize}

\textbf{(c) Reaction with \ce{NaHCO3} (aq) (4 marks).}
Only \textbf{ethanoic acid} reacts, with \emph{effervescence} (\ce{CO2} evolution).
\ce{CH3COOH(aq) + NaHCO3(aq) -> CH3COONa(aq) + H2O(l) + CO2(g)}
Phenol and ethanol do not react. Carbonic acid ($\mathrm{p}K_{\mathrm{a1}}\approx6.3$) is stronger than phenol ($\mathrm{p}K_{\mathrm{a}}\approx10$) and ethanol ($\mathrm{p}K_{\mathrm{a}}\approx16$), so the equilibrium lies to the left for the weaker acids (2 marks for equation/observation, 2 marks for explanation).

\textbf{(d) Explanation via anion stability (4 marks).}
\begin{itemize}
\item \textbf{Ethanoate ion (\ce{CH3COO-}):} Negative charge delocalised equally over two electronegative oxygen atoms by \cle{resonance} (\ensuremath{\mathrm{O=C-O^{-} <\rightarrow  -O-C=O}}). Highly stabilised (1 mark).
\item \textbf{Phenoxide ion (\ce{C6H5O-}):} Negative charge delocalised into the benzene ring (ortho/para positions) by resonance. Stabilised, but less effectively than in carboxylate because charge is dispersed onto less electronegative carbon atoms (1 mark).
\item \textbf{Ethoxide ion (\ce{C2H5O-}):} No resonance stabilisation. The electron-donating \cle{+I effect} of the ethyl group increases electron density on oxygen, \emph{destabilising} the anion (1 mark).
\end{itemize}
Order of anion stability: Ethanoate $>$ Phenoxide $>$ Ethoxide.
Hence order of acidity: \cle{Ethanoic acid $>$ Phenol $>$ Ethanol} (1 mark).
\end{corrige}

\newpage

\coursec{Summary sheet}

\begin{popfigure}
\begin{tikzpicture}[
  mindmap,
  grow cyclic,
  every node/.style={concept, text=white, font=\small\bfseries},
  root concept/.append style={concept color=popdark, font=\normalsize\bfseries},
  level 1/.append style={level distance=3.6cm, sibling angle=60, font=\scriptsize\bfseries},
  level 2/.append style={level distance=2.3cm, sibling angle=40, font=\tiny},
  concept connection/.append style={line width=1pt}
]
\node[root concept]{Organic Chemistry 2}
  child[concept color=popblue]{ node {Haloalkanes \ce{R-X}}
    child { node {SN1 / SN2} }
    child { node {Elimination $\to$ alkene} }
  }
  child[concept color=popgreen]{ node {Alcohols \ce{R-OH}}
    child { node {Oxidation: $1^\circ, 2^\circ, 3^\circ$} }
    child { node {Lucas / Iodoform} }
  }
  child[concept color=poporange]{ node {Phenol \ce{C6H5OH}}
    child { node {Ring: \ce{Br2}, \ce{HNO3}} }
    child { node {Weak acid} }
  }
  child[concept color=popgold]{ node {Carbonyls \ce{R-CHO}, \ensuremath{\mathrm{R-CO-R'}}}
    child { node {Nucleophilic addition} }
    child { node {2,4-DNPH, Tollens, Fehling} }
  }
  child[concept color=popred]{ node {Acids \ce{R-COOH}}
    child { node {Esterification} }
    child { node {Derivatives: \ce{RCOCl}, \ensuremath{\mathrm{RCOOR'}}, \ce{RCONH2}} }
  }
  child[concept color=popmuted]{ node {N-compounds}
    child { node {Amines: bases} }
    child { node {Diazonium salts} }
    child { node {Amino acids} }
  };
\end{tikzpicture}
\legende{Mind map of Chapter CHE02: families of organic compounds and their key reactions.}
\end{popfigure}

\begin{bilan}
\textbf{Key definitions}
\begin{itemize}
\item \cle{Haloalkane}: \ce{R-X} (X = F, Cl, Br, I); classified $1^\circ$, $2^\circ$, $3^\circ$ by the carbon bearing X.
\item \cle{Nucleophilic substitution}: a nucleophile (\ce{OH-}, \ce{CN-}, \ce{NH3}, \ce{RO-}) replaces X. \cle{SN2}: one step, backside attack, favoured by $1^\circ$; \cle{SN1}: two steps via a carbocation, favoured by $3^\circ$.
\item \cle{Elimination} (alcoholic \ce{KOH}, heat): loss of HX to give an alkene (Zaitsev: most substituted alkene major).
\item \cle{Alcohols}: \ce{R-OH}; $1^\circ \to$ aldehyde $\to$ acid; $2^\circ \to$ ketone; $3^\circ$ resist oxidation.
\item \cle{Phenol}: \ce{OH} on benzene; weak acid, stronger than alcohols; activates the ring (2,4,6-substitution).
\item \cle{Aldehyde} \ce{R-CHO} / \cle{ketone} \ensuremath{\mathrm{R-CO-R'}}: carbonyl \ce{C=O}, general formula \ensuremath{\mathrm{C_nH_{2n}O}}.
\item \cle{Carboxylic acid} \ce{R-COOH}: weak acid; derivatives: acyl chlorides \ce{RCOCl}, anhydrides \ce{(RCO)2O}, esters \ensuremath{\mathrm{RCOOR'}}, amides \ce{RCONH2}.
\item \cle{Amines}: \ce{RNH2}, \ce{R2NH}, \ce{R3N}; basic (lone pair on N). \cle{Amino acids}: \ce{H2N-CHR-COOH}, amphoteric, zwitterion.
\end{itemize}

\textbf{Essential reactions and formulas}
\begin{itemize}
\item \ce{R-X + OH- (aq) -> R-OH + X-};\quad \ce{R-X + CN- -> R-CN} (chain $+1$ C);\quad \ce{R-X + NH3 -> R-NH2}.
\item \ensuremath{\mathrm{R-CH_{2}-CHX-R' + KOH(alc) \rightarrow  R-CH=CH-R' + KX + H_{2}O}}.
\item \ensuremath{\mathrm{R-OH \rightarrow [conc.\ H_{2}SO_{4}, 170^\circ C] alkene + H_{2}O}};\quad \ce{R-OH + HX -> R-X + H2O}.
\item \ce{R-CHO + 2[H] -> R-CH2OH};\quad \ce{R2C=O + 2[H] -> R2CH-OH} (\ce{NaBH4} or \ce{LiAlH4}).
\item \ce{R-CHO + [O] -> R-COOH}; ketones resist mild oxidation.
\item \ensuremath{\mathrm{R-COOH + R'-OH \rightleftharpoons [conc.\ H_{2}SO_{4}] R-COO-R' + H_{2}O}} (esterification, reversible).
\item \ce{R-COOH + SOCl2 -> R-COCl + SO2 + HCl};\quad \ce{R-COCl + NH3 -> R-CONH2 + HCl}.
\item \ensuremath{\mathrm{C_{6}H_{5}NH_{2} + HNO_{2} + HCl \rightarrow [0-5^\circ C] C_{6}H_{5}N_{2}+Cl^{-}}} (diazonium salt $\to$ phenol, halides, azo dyes).
\item Acidity order: \ce{RCOOH} $>$ \ce{C6H5OH} $>$ \ce{H2O} $>$ \ce{R-OH}; electron-withdrawing groups increase acidity.
\end{itemize}

\textbf{Identification tests (practical)}
\begin{itemize}
\item \cle{2,4-DNPH}: orange precipitate for aldehydes \emph{and} ketones.
\item \cle{Tollens}: silver mirror with aldehydes only; \cle{Fehling/Benedict}: brick-red \ce{Cu2O} with aliphatic aldehydes.
\item \cle{Lucas test}: $3^\circ$ immediate cloudiness, $2^\circ$ in minutes, $1^\circ$ very slow.
\item \cle{Iodoform} (\ce{I2}/\ce{NaOH}): yellow \ce{CHI3} with \ce{CH3CO-} or \ce{CH3CH(OH)-} compounds.
\item \cle{Phenol}: violet colour with neutral \ce{FeCl3}; white precipitate of 2,4,6-tribromophenol with bromine water.
\item \cle{Carboxylic acids}: effervescence of \ce{CO2} with \ce{NaHCO3}.
\end{itemize}

\textbf{Methods}
\begin{itemize}
\item Distinguish SN1/SN2 by structure of substrate ($3^\circ$ vs $1^\circ$) and solvent polarity.
\item To lengthen a carbon chain by one C: \ce{R-X -> R-CN -> R-COOH} (hydrolysis) or \ce{R-CH2NH2} (reduction).
\item Always balance equations with state symbols and write conditions above the arrow.
\end{itemize}

\textbf{Common mistakes}
\begin{itemize}
\item Confusing aqueous \ce{KOH} (substitution) with alcoholic \ce{KOH} (elimination).
\item Oxidising a $3^\circ$ alcohol (no H on the carbinol carbon: no reaction).
\item Saying phenol behaves like an alcohol: phenol does \emph{not} undergo easy substitution of \ce{OH} nor dehydration.
\item Forgetting that ketones give a negative Tollens/Fehling test.
\item Diazotisation above \SI{5}{\ensuremath{{}^{\circ}\mathrm{C}}}: the diazonium salt decomposes.
\end{itemize}

\textbf{GCE tips}
\begin{itemize}
\item Draw mechanisms with curly arrows from the electron-rich centre (lone pair or $\pi$ bond) to the electron-poor centre.
\item Quote observations precisely (colour, precipitate, effervescence) in test questions.
\item For synthesis questions, work backwards from the target molecule and count the carbons.
\end{itemize}
\end{bilan}

\findecours

\end{document}
